\documentclass[11pt]{article}
\usepackage[T1]{fontenc}
\usepackage[utf8]{inputenc}
\usepackage{lmodern,amsmath,amssymb,amsthm,mathtools,microtype,booktabs,longtable,array}
\usepackage[a4paper,margin=27mm]{geometry}
\PassOptionsToPackage{hyphens}{url}
\usepackage[hidelinks]{hyperref}
\setlength{\emergencystretch}{3em}
\newtheorem{theorem}{Theorem}[section]
\newtheorem{proposition}[theorem]{Proposition}
\newtheorem{lemma}[theorem]{Lemma}
\newtheorem{corollary}[theorem]{Corollary}
\theoremstyle{definition}\newtheorem{definition}[theorem]{Definition}
\newtheorem{remark}[theorem]{Remark}
\newcommand{\Z}{\mathbb Z}\newcommand{\Q}{\mathbb Q}
\newcommand{\F}{\mathbb F}\newcommand{\Gal}{\operatorname{Gal}}
\newcommand{\Hom}{\operatorname{Hom}}\newcommand{\Cl}{\operatorname{Cl}}
\newcommand{\workbenchurl}{https://github.com/KokunoYumeto/k-theory-and-galois-groups-workbench}
\newcommand{\recordurl}{https://github.com/KokunoYumeto/mathematics-workbenches}
\title{Anderson and Ihara's question after Sharifi,\\and its CM elliptic analogue after Ishii}
\author{}
\date{25 November 2026}
\begin{document}\maketitle
\begin{abstract}
Sharifi proved Anderson and Ihara's kernel equality at odd regular primes from the Deligne--Ihara conjecture, subsequently proved by Brown. Starting from the class-two cup comparison in the recent extension under Greenberg's finiteness hypothesis, we construct canonical quotients of every finite nilpotency class at every odd prime. Their source-defined kernels intersect in the outer kernel, and their commutator modules contain every completed free Lie degree of the odd torsionfree quotient. Finiteness of a single even Iwasawa component ensures that its even abelian component survives in full, using Ohta's isomorphism and its completed CM-density input. The stated, unrefereed $K$-theory theorem of Mok and Zheng gives selected quotients of exact classes through $61$ at every odd prime. At $37$, every lift of the even abelian generator has a normal closure that fails to define a Galois subextension over the first cyclotomic field. McCallum--Sharifi's integral obstruction, with Brown's theorem, also disproves generation by the specified Soul\'e elements in degree $32$. After Ishii, we prove that the outer action detects the full elliptic-unit radical at every split prime in the class-number-one CM setting, and that the elliptic Soul\'e characters jointly separate it. Under finiteness of the conjugation-fixed class module, we construct a canonical detected torsion-free quotient of class exactly two with the full abelianization. Its commutator map is an explicit integral shear operator on Kummer measures. The source is free, and the abelian assertion also holds character by character with the exact Kummer twist. The full CM kernel equality remains open. The mathematics and text were produced by GPT-6 Astra (ChatGPT Pro mode), and the main theorems were independently re-proved by the same model at inherited, unexposed effort settings; no human has checked this work.
\end{abstract}

\section{The question and the results}
Let $p$ be an odd prime, let $F_r=\Q(\mu_{p^r})$, and put $F_\infty=\bigcup_{r\ge1}F_r$. In a fixed algebraic closure of $\Q$, let $\Omega_p$ be the maximal pro-$p$ extension of $F_1$ unramified outside $p$. Write
\[
 H=\Gal(\Omega_p/F_\infty),\qquad
 X=\varprojlim_r\Cl(F_r)\{p\},
\]
where the transition maps on class groups are norms. Let $\Pi_p$ be the maximal pro-$p$ quotient of the geometric \'etale fundamental group of $\mathbb P^1-\{0,1,\infty\}$, and let $\Omega_p^*$ be the fixed field of the kernel of $G_\Q\to\operatorname{Out}(\Pi_p)$.
Anderson and Ihara ask whether $\Omega_p^*=\Omega_p$ \cite[p.~272, question (a)]{AI}. Sharifi's Theorem~1 \cite[p.~2]{Sharifi} proves the equality for odd regular primes from the Deligne--Ihara conjecture; Brown's Theorem~1.1 \cite[p.~949]{Brown} proves the freeness input. Theorem~1.1(a), pp.~1--2, of \cite{Note} proves the equality under finiteness of $X^+$.
Ishii's Theorem~A \cite[p.~2; Theorem~2.14, p.~11]{Ishii} proves an elliptic analogue under his class-number, prime-decomposition and freeness hypotheses. Our Theorems~\ref{thm:canonical-all-stages} and~\ref{thm:cm-main}, pp.~\pageref{thm:canonical-all-stages} and~\pageref{thm:cm-main}, give the canonical tower in every finite class and the CM elliptic-unit and class-two fields; Theorem~\ref{thm:selected-all-stages}, p.~\pageref{thm:selected-all-stages}, gives the component theorem.

We specify the quotient in the component theorem. Let
$\Delta=\operatorname{Gal}(F_1/\mathbb Q)$, let $\omega$ be its
Teichm\"uller character, and write $V^{[j]}$ for the image of
$(p-1)^{-1}\sum_{\delta\in\Delta}\omega(\delta)^{-j}\delta$
on a compact $\mathbb Z_p[\Delta]$-module $V$. Exponents are
read modulo $p-1$. Put $\Gamma=\operatorname{Gal}(F_\infty/F_1)$,
$\Lambda=\mathbb Z_p[[\Gamma]]$, and
\[
 M=H^{\rm ab}=U\oplus T,\qquad U=M^-,\quad T=M^+,
 \quad T_e=T^{[e]},\quad D=U_{\Lambda\text{-tor}},\quad W=U/D.
\]
Here $+$ and $-$ are the eigenspaces of complex conjugation.
The abelianization has a canonical action of $\Delta\times\Gamma$.

In the class-two BCH Lie module of $H$, first kill the central
brackets with an input from $T$. Its negative part then maps
isomorphically to $U$. Kill the image of $D$ in this negative
part and its brackets. The resulting compact Lie module has the
form
\[
 \mathcal L=W\oplus\mathcal E,\qquad
 0\longrightarrow\widehat{\bigwedge}^{\,2}_{\mathbb Z_p}W
 \longrightarrow\mathcal E\longrightarrow T\longrightarrow0,
 \qquad [w,w']=w\wedge w',
\]
where $\mathcal E$ is central. Put $Z=\mathcal E_{\Lambda\text{-tor}}$, let
$Q=(W\oplus\mathcal E/Z,\star)$, where
$x\star y=x+y+\tfrac12[x,y]$, and write $q:H\twoheadrightarrow Q$.
The notation $\Lambda$-torsion means the elements annihilated
by some nonzero element of $\Lambda$.

Choose a topological generator \(\gamma\) of \(\Gamma\).
The canonical quotients are defined by an intrinsic recursion. Put
\(V=M_{\Lambda\text{-tor}}=D\oplus T\), and let
\(a(z)=\det(z-\gamma\mid V)\), with \(a=1\) when \(V=0\).
All commutator products and subgroups below are closed. Set
\[
 \mathcal F_1=H,\qquad
 J_{n+1}=\overline{\prod_{i+j=n+1}[\mathcal F_i,\mathcal F_j]},\qquad
 \mathcal R_n=\mathcal F_n/J_{n+1},\qquad
 \mathcal F_{n+1}=\operatorname{preim}_{\mathcal F_n}(\mathcal R_{n,\Lambda\text{-tor}}).
\]
Thus \(\mathcal F_2=\ker(H\to W)\). Define
\(Q_n=H/\mathcal F_{n+1}\) and \(L^{(n)}=\Omega_p^{\mathcal F_{n+1}}\).
For \(n\geq2\), put \(\mathcal S_n=J_n/J_{n+1}\).

\begin{theorem}[Canonical detection in every finite class]
\label{thm:canonical-all-stages}
For every odd prime \(p\), this recursion is well-defined at every
\(n\geq1\). There is an integer \(h_n\geq0\) such that
\[
 \mathcal R_{n,\Lambda\text{-tor}}=\ker\bigl(p^{h_n}a(\gamma):\mathcal R_n\to \mathcal R_n\bigr).
\]
In particular this torsion subgroup is closed. Every \(\mathcal F_n\) is stable
under the full arithmetic Galois group. The quotient \(Q_n\) is
torsionfree of nilpotency class exactly \(n\), and is the compact
image of \(H\) in the rational outer image through Brown-generator
length \(n\). For \(n\geq2\), the central extension
\[
 1\longrightarrow \mathcal R_n/\mathcal R_{n,\Lambda\text{-tor}}
 \longrightarrow Q_n\longrightarrow Q_{n-1}\longrightarrow1
\]
contains an injective commutator map from
\(\operatorname{Lie}^n(W)\) into its kernel. Here
\(\operatorname{Lie}^n(W)\) is the completed degree-\(n\) free
\(\mathbb Z_p\)-Lie module on \(W\). If \(K_{\rm out}\) denotes
the kernel of the pro-\(p\) outer action on \(H\), then
\[
 K_{\rm out}=\bigcap_{n\geq1}\mathcal F_n,\qquad
 L^{(1)}\subsetneq L^{(2)}\subsetneq L^{(3)}\subsetneq\cdots,
 \qquad \bigcup_{n\geq1}L^{(n)}=\Omega_p^*.
\]
\end{theorem}
This theorem reconstructs the detected field from the source group,
its cyclotomic action and successive torsion submodules. It does not
prove that the intersection is trivial. The first five quotients are
computed explicitly in the body; the general proof is in
Sections~\ref{sec:weighted-companion-saturation} and
\ref{sec:canonical-all-stages}.

We now specify a canonical component quotient in every class.
For \(n\geq2\), let \(B_n\) be the image of \(\mathcal F_2\) in \(Q_n\),
and put \(D_n=\operatorname{im}(\mathcal F_{n+1}\to U)\). Choose a tame
\(\Delta\) complement. For a set \(\Sigma\) of even exponents,
let \(J_{n,\Sigma}\) be the closed \(Q_n\)-normal subgroup generated
by every \(g\in B_n\) satisfying
\[
 \delta g\delta^{-1}=g^{\omega(\delta)^e}
 \quad\text{for all }\delta\in\Delta,
 \qquad e\text{ even},\quad e\notin\Sigma.
\]
Put \(Q_{n,\Sigma}=Q_n/J_{n,\Sigma}\),
\(T_\Sigma=\bigoplus_{e\in\Sigma}T_e\), and let \(L_{n,\Sigma}\)
be the fixed field of \(\ker(H\to Q_{n,\Sigma})\).

\begin{theorem}[Detection of selected even components]
\label{thm:selected-all-stages}
Let \(p\) be odd, \(n\geq2\), and let \(\Sigma\) be a set of
even exponents for which every \(X^{[e]}\), \(e\in\Sigma\), is finite.
The subgroup \(J_{n,\Sigma}\) is independent of the tame complement
and stable under the full arithmetic action. Its canonical field is
Galois over \(\mathbb Q\), and
\[
 L_{n,\Sigma}\subseteq\Omega_p^*,\qquad
 \operatorname{Gal}(L_{n,\Sigma}/F_\infty)=Q_{n,\Sigma},\qquad
 Q_{n,\Sigma}^{\rm ab}=(U/D_n)\oplus T_\Sigma.
\]
These groups form a compatible tower. For every \(n\) satisfying
\(-2,-4,\ldots,-2\lfloor n/2\rfloor\in\Sigma\), the group
\(Q_{n,\Sigma}\) has nilpotency class exactly \(n\).
For \(2\leq n\leq5\), it equals the component quotient defined
in the body by the integral class-two BCH Lie module of \(B_n\).
\end{theorem}
The proof lifts each tame eigenclass to a group eigen-element by
finite averaging in successive central subgroups. It then constructs
a homomorphism to a free nilpotent group on two odd letters to prove
the exact class. The arithmetic input for the selected abelian
components is Ohta's isomorphism and the Fukaya--Kato identity, with
their character maps in Subsection~\ref{subsec:ohta-component-gate}.

For \(\Sigma_{30}=\{-2m:1\leq m\leq30\}\),
Corollary~\ref{cor:k-canonical-detection} gives a strict selected tower
through class \(61\) at every odd prime, using Mok--Zheng's stated
unrefereed theorem and its specified archive dependencies.
For \(p<2^{31}\), the full equality is already known from
\cite{Note,HHO}; this application includes primes beyond that range.

The literal quotient by chosen lifts has a separate, explicit
obstruction. Theorem~\ref{thm:component-no-galois-p37} proves that at
\(37\) every primitive lift of the even abelian generator has a
normal closure which fails to be Galois-stable. Theorem~
\ref{thm:p37-cyclotomic-core} computes its cyclotomic core as the
identity, so even one such lift-wise containment is equivalent to
full faithfulness. The groups above use all eigen-elements in the
specified characters, which proves their invariance.
Proposition~\ref{prop:integral37} combines McCallum--Sharifi's
calculation with Brown's theorem to show that the specified Soul\'e
elements fail to generate degree \(32\) integrally at \(37\).

\subsection{The elliptic-unit field after Ishii}

Let $K$ be imaginary quadratic of class number one, let $p\geq5$
split in $K$, and let $E/K$ have complex multiplication by $O_K$.
Put $L_n=K(p^n)$, $L_\infty=\bigcup_nL_n$, and let $\Omega$
be the maximal pro-$p$ extension of $L_1$ unramified outside $p$.
Here $K(p^n)$ is the ray class field of modulus $p^nO_K$.
Let $\rho_E$ be the outer representation on the geometric pro-$p$
fundamental group of $E-\{O\}$, and define
$\Omega^*=\Omega\cap\bar K^{\ker\rho_E}$.
Ishii's proof of Lemma~2.13, p.~10, gives
\[
 \bar K^{\ker\rho_E}=K(E[p])\Omega^*,\qquad
 [K(E[p]):L_1]\text{ prime to }p.
\]
Let $C(L_n)$ be Rubin's elliptic-unit group, with the precise
convention in \cite[Definition~2.2, p.~8]{IshiiSoule}, and put
\[
 L_{\rm ell}=L_\infty(c^{1/p^r}:c\in C(L_n),\ n,r\geq1),
 \qquad A=\operatorname{Gal}(\Omega/L_\infty)^{\rm ab}.
\]
For $w=|O_K^\times|$, write
$I=\{(m_1,m_2)\in\mathbb Z_{\geq1}^2-\{(1,1)\}:
 m_1\equiv m_2\pmod w\}$ and let $\kappa_{\boldsymbol m}$
be the geometric elliptic Soul\'e character defined in
\cite[Definition~1.1, p.~3]{IshiiSoule}.

Write $H_E=\operatorname{Gal}(\Omega/L_\infty)$ and define
\[
 \mathcal X_2=\varprojlim_n\operatorname{Cl}(L_n)\{p\},\qquad
 \mathcal J_2=\varinjlim_n\operatorname{Cl}(L_n)\{p\},
\]
using norms and extension of ideals, respectively.
The superscript $+$ denotes the fixed subgroup of complex conjugation
$c$; it is considered as a $\mathbb Z_p$-module.
Let $\Delta_E=\operatorname{Gal}(L_1/K)$.
For a set $\mathcal O$ of its characters stable under
$\chi\mapsto\chi^c$, put
\[
 e_{\mathcal O}=\sum_{\chi\in\mathcal O}
       \frac1{|\Delta_E|}\sum_{\delta\in\Delta_E}\chi(\delta)^{-1}\delta,
 \qquad
 \mathcal O^\dagger=\{\omega_{\mathrm{cyc}}\chi^{-1}:\chi\in\mathcal O\},
\]
where $\omega_{\mathrm{cyc}}$ is the cyclotomic character on $\Delta_E$.
All these characters are $\mathbb Z_p$-valued.

For the nonabelian assertion, put \(V=T_p(E)\) and let
\(\alpha:A\to\mathbb Z_p[[V]](1)\) be Nakamura's canonical
metabelian homomorphism, as defined in
\cite[Definition~2.5, p.~8]{Ishii}. The completed group
ring is the module of bounded measures on the additive group \(V\).
On \(V\times V\), put
\[
 t_1(x,y)=(x,x+y),\qquad t_2(x,y)=(x+y,y),\qquad
 \mathcal J=(t_1)_*+(t_2)_*-1.
\]
Define the continuous homomorphism
\[
 J_A:\widehat{\bigwedge}^2 A\longrightarrow
                    \mathbb Z_p[[V\times V]](2),\qquad
 a\wedge b\longmapsto
 \mathcal J\bigl(\alpha(a)\otimes\alpha(b)-\alpha(b)\otimes\alpha(a)\bigr).
\]

\begin{theorem}[CM field detection after Ishii]\label{thm:cm-main}
For every $K,E,p$ just specified,
\[
 L_{\rm ell}\subseteq\Omega^*,\qquad
 \bigcap_{\boldsymbol m\in I}\ker(\kappa_{\boldsymbol m}:A
       \to\mathbb Z_p(\boldsymbol m))
 =\ker\bigl(A\to\operatorname{Gal}(L_{\rm ell}/L_\infty)\bigr).
\]
Thus the elliptic Soul\'e characters jointly separate the full
elliptic-unit radical quotient, which is detected by the outer action.

For every such $\mathcal O$ for which $(e_{\mathcal O}\mathcal X_2)^+$
is finite, $e_{\mathcal O}\mathcal J_2=0$ and the canonical field
\[
 L_{\mathcal O}=
  \Omega^{\ker(H_E\longrightarrow A
                 \xrightarrow{e_{\mathcal O^\dagger}}e_{\mathcal O^\dagger}A)}
\]
is contained in $\Omega^*$.
In particular, finiteness of $\mathcal X_2^+$ gives
\[
 \mathcal J_2=0,\qquad
 \Omega^{[H_E,H_E]}=L_{\rm ell}\subseteq\Omega^*,\qquad
 \bigcap_{\boldsymbol m\in I}\ker(\kappa_{\boldsymbol m}|_A)=0,
\]
and $H_E$ is a free pro-$p$ group. In this finite-plus case, let
\(N_2^{\rm CM}\) be the inverse image of \(\ker J_A\) under
\[
 \gamma_2H_E\longrightarrow\gamma_2H_E/\gamma_3H_E
                  \simeq\widehat{\bigwedge}^2 A.
\]
Then \(N_2^{\rm CM}\) is closed and arithmetically stable, and
\[
 \Omega^{N_2^{\rm CM}}\subseteq\Omega^*,\qquad
 Q_2^{\rm CM}=H_E/N_2^{\rm CM}
\]
is torsion-free of class exactly two. Its abelianization is \(A\)
and its central commutator module is \(\operatorname{im}J_A\). Every pair of elements whose abelian images are linearly independent over \(\mathbb Q_p\) has a nonzero commutator in this quotient.
\end{theorem}

Theorem~\ref{thm:cm-main} is proved in
Proposition~\ref{prop:cm-full-radical} and
Corollary~\ref{cor:cm-radical-joint-kernel}; its finite-component assertions are proved in Theorem~\ref{thm:cm-finite-plus-abelian}, and its class-two assertion in Theorem~\ref{thm:cm-measure-class-two}. The unconditional input is Ishii's
primitive-pair moment identity, Lemma~3.1, p.~11, and his theta
unit formulas. The radical containment imposes no ray-class-number, prime-decomposition,
or Lie-freeness hypothesis. Under the stated finiteness of the conjugation-fixed
class module, the theorem detects a nonabelian extension with the whole
abelianization and proves freeness of its source. The proof uses the exact full-class Kummer sequence,
cyclotomic weak Leopoldt, and the finite-index comparison for elliptic units. For $y^2=x^3-x$ over $\mathbb Q(i)$
at $5$ and $y^2=x^3+1$ over $\mathbb Q(\sqrt{-3})$ at $7$,
Proposition~\ref{prop:cm-unramified-zero} verifies the finite-plus
hypothesis, so the class-two field is detected in both examples. The passage to the
full nonabelian kernel equality remains open. Section~\ref{sec:cm-audit}
also proves the exact semilinear descent map and the nonzero
transgression in these examples, and gives the explicit derivation
relation used in Ishii's updated discussion of freeness.

Sharifi's current account of the Eisenstein-ideal programme
\cite[p.~5]{SharifiProgram} also reports characterwise order
comparisons of modular-symbol and class-group quotients for
levels at most $250$, computed in Sage through OpenAI's Codex.
These computations concern the same cup-product map whose
Iwasawa limit occurs in Subsection~\ref{subsec:ohta-component-gate}.
Our component argument uses the stated Ohta isomorphism and
Fukaya--Kato identity at that limit. The posted $K$-group
results and their shared Lean-archive inputs are described
separately in Section~\ref{sec:dependencies}.

Theorem~\ref{thm:cm-elliptic-euler-kernel} identifies the finite
class-two shear kernels with the duals of Goncharov's elliptic Euler
cokernels. Nakamura's theta coefficients give the exact comparison,
including the factor $144$, the origin coefficient and every transition
map. The completed kernel is the inverse limit of its integral
finite-level lattices. This relates the remaining class-two question
to the persistence of elliptic-unit triangle relations along the ray tower.

\section{Proof overview}

The source construction starts with the abelianization of the free pro-$p$ group $H$. Its $\Lambda$-torsion is finite over $\mathbb Z_p$; the quotient has regular Iwasawa coordinates after inverting $p$. Hain--Matsumoto's density theorem and Brown's freeness theorem identify the rational outer image with a completed free Lie image on odd generators. Their weight evaluations detect every bounded coefficient series in the regular coordinates. An even abelian class first enters this image through a commutator of odd generators, so its comparison requires the central extension above that class.

The class-two cup calculation supplies this comparison one character at a time. Finiteness of $X^{[e]}$ makes $T_e[1/p]$ cyclic through the local and global unit sequence. Ohta's isomorphism transports the resulting module to modular symbols. The Fukaya--Kato derivative identity makes the cup map invertible at every characteristic prime, including repeated roots. A finite integral cokernel gives one power of $p$ valid throughout the finite-level comparison. Norm-compatible unit classes vanish on the odd $\Lambda$-torsion. The explicit cochain calculation therefore gives the required nonzero companion numerator on each root divisor of the characteristic polynomial of $T_e$. It follows that the central torsion kernel has zero projection to the selected even component.

The construction then repeats in each finite nilpotency class. The first few classes display the commutator and Jacobi maps explicitly. The general proof uses a free Lie algebra with generators assigned the weights at which their abelian directions are first detected. A bounded Hall-coordinate calculation identifies its rational relation module and keeps the integral $p$-torsion with a fixed bound at each stage. The main coefficient argument proves division by a total-product divisor in this original weighted source. On that hypersurface, two interval poles can coincide at complementary boundary cuts. Expanding companion trees bounds the pole orders from words crossing the cut; finite coefficient projections force the exact source divisibility. This proves injectivity in every degree and supplies the next companion extension. The original abelian characteristic polynomial, with the finite-stage $p$-bound, defines a closed torsion subgroup and gives the exact truncated outer kernel.

The normalized Ihara action and a free-group centralizer calculation identify the intersection of these kernels with the full pro-$p$ outer kernel. Compactness then identifies its relation module with the inverse limit of the finite-stage torsion modules. Compact Nakayama reduces triviality of that kernel to the corresponding inverse limit modulo $(p,\gamma-1)$. For the component quotient, finite central averaging lifts tame eigenclasses integrally. Taking the normal closure of all omitted even eigen-elements gives a canonical arithmetic quotient in every class. Projection to two Brown generators proves its asserted nilpotency class. The calculation at $37$ also explains why choosing arbitrary lifts in $H$ does not give this construction: every pure even generator has a normal closure failing cyclotomic invariance. The $K$-theory application follows through the displayed Kummer and class-group norm maps, and the integral Soul\'e obstruction follows from McCallum--Sharifi with Brown's injectivity theorem.

For CM curves, Ishii's primitive-pair moments give a bounded Kummer measure including both coordinate axes. A uniform polynomial limit recovers the constant moment; an injective measure operator handles the additional auxiliary ideals on the axes. This proves detection of the full elliptic-unit radical. Under the stated finite-component hypothesis, class-group norm surjectivity and the norm--extension identity make the fixed direct-limit classes vanish. Their Kummer dual is then fixed by the anticyclotomic direction, whereas Hochschild--Serre and cyclotomic weak Leopoldt give zero invariants in the abelian source. The finite-index comparison of elliptic units and global units proves abelian detection with the exact character twist. For all components, the same Kummer sequence and the cohomological criterion also prove source freeness.

Nakamura's boundary-preserving lift permits the next integral calculation. On the second lower-central quotient of the geometric commutator subgroup, the displacement of an arithmetic commutator is diagonal translation minus identity applied to two shear pushforwards minus identity. The translation operator is injective on bounded measures, giving the canonical detected class-two source quotient. Differentiating its exponential transform proves that two independent abelian classes have a nonzero commutator; moment-isolating operators construct the free class-two source subgroups. The Euler-factor identity specifies the image on the mixed unit chart. The finite examples, transgression calculation, and geometric normalizer theorem give further exact maps at the remaining CM kernel question. The geometric proof uses the genus-ten cover, Nakamura's degree-six arithmetic quotient, and the verified Johnson contraction identity.

\section{Arithmetic and the class-two comparison}\label{sec:abelian}

This section supplies the arithmetic input and the cup comparison
for Corollary~\ref{thm:class-two-components}.
Write $\Delta=\Gal(F_1/\Q)$, choose the Teichm\"uller character
$\omega$, and read character exponents modulo $p-1$.
For a compact $\Z_p[\Delta]$-module $V$, put
\[
 V^{[j]}=\epsilon_jV,\qquad
 \epsilon_j=\frac1{p-1}\sum_{\delta\in\Delta}\omega(\delta)^{-j}\delta.
\]
Let $\Gamma=\Gal(F_\infty/F_1)$, choose $\gamma$ with
$u=\chi_p(\gamma)$ generating $1+p\Z_p$, and put
$\Lambda=\Z_p[[\Gamma]]$, $R=\Lambda[1/p]$.
The abelianization $M=H^{\rm ab}$ has its conjugation action of
$\Delta\times\Gamma$, independently of choices of lifts. Put
\[
 U=M^-,\quad T=M^+=\bigoplus_{e\ {\rm even}}T_e,\quad
 T_e=M^{[e]},\quad D=U_{\Lambda\text{-tor}}.
\]

Choose a lift of $\Delta$ by profinite Schur--Zassenhaus.
Coprime fixed-point lifting gives an element of
$\operatorname{Gal}(\Omega_p/F_1)^\Delta$ above a topological
generator of $\Gamma$. Its procyclic closure maps isomorphically
to $\Gamma$, giving commuting lifts and a splitting
$\operatorname{Gal}(\Omega_p/F_1)=H\rtimes\Gamma$.
The lifting assertion follows at each finite pro-$p$ quotient
from Schur--Zassenhaus applied to the preimage of the fixed
coset; compactness gives compatible lifts. These choices are
used for the cochain calculation. The full invariance calculation
in Section~\ref{sec:class-two-components} removes them from the
resulting quotient.

\subsection{Norms and the odd torsion}

The following computation removes the use of finiteness of the whole
plus module in the component theorem. It is a fixed-coefficient
form of the unit-radical phenomenon in
\cite[Theorem~3, p.~3]{McCallum}.
\begin{lemma}\label{lem:norms}
Put $P_r=H\rtimes\Gamma^{p^r}$ and $A_n=\Z/p^n(1)$.
Let $D_0\subseteq H^{\rm ab}$ be a $\Gamma$-stable finitely
generated $\Z_p$-submodule. A sequence
$f_r\in H^1(P_r,A_n)$ compatible under corestriction restricts
to zero on $D_0$ at every level where it is defined.
\end{lemma}
\begin{proof}
Restriction to $H$ factors through $H^{\rm ab}$, because $H$
acts trivially on the coefficients. Write
$\phi_r=(\operatorname{res}_Hf_r)|_{D_0}$ and
$V_n=\Hom_{\Z_p}(D_0,A_n)$. On this finite module $\gamma$ acts by
\[
                  (B\phi)(d)=u\phi(\gamma^{-1}d).
\]
The corestriction formula with coset representatives
$\gamma^{jp^r}$, $0\leq j<p$, gives
\[
 \phi_r=N_r(B)\phi_{r+1},\qquad
 N_r(B)=\sum_{j=0}^{p-1}B^{jp^r}.
\]
Continuity gives an $s$ with $B^{p^s}=1$ on $V_n$.
Hence $N_r(B)=p$ for $r\geq s$. Iterating $n$ times gives
$\phi_r=p^n\phi_{r+n}=0$ at all sufficiently high levels.
The transition formulas prove the assertion at earlier levels.
\end{proof}

For $p\nmid a$ the identity
\[
 N_{F_{r+1}/F_r}(1-\zeta_{p^{r+1}}^a)=1-\zeta_{p^r}^a
\]
provides such a sequence of Kummer classes at every fixed coefficient
exponent $n$. For $a=p^vb$, $p\nmid b$ and $v<r$, one has
\[
 1-\zeta_{p^r}^{p^vb}
   =\prod_{j=0}^{p^v-1}(1-\zeta_{p^r}^{b+jp^{r-v}}).
\]
All exponents on the right are prime to $p$. Products and the
$c$-even projector preserve corestriction compatibility.
Every cyclotomic unit class used below consequently annihilates $D$.

The group $H$ is free pro-$p$ by Ferrero--Washington and
Iwasawa \cite[Theorem~2, pp.~273--274]{Iwasawa}.
Its abelianization $M$ is finitely generated over $\Lambda$:
remove the $\Gamma$-coordinates from a finite generating set of
$\Gal(\Omega_p/F_1)$ and take their abelian classes in $H$.
Their $\Gamma$-conjugates generate $M$.
In particular $M$ is $\Z_p$-torsionfree.
The module $D$ is finite free over $\Z_p$. To see this, choose
an annihilator $p^af$ of its finitely generated $\Lambda$-torsion
module, with $p\nmid f$. Injectivity of multiplication by $p$
gives $fD=0$. Weierstrass preparation makes $\Lambda/(f)$ finite
over $\Z_p$. Moreover $D=\ker(f:U\to U)$, so it is closed.
It is saturated: $px\in D$ gives $fpx=0$, then $fx=0$, hence
$x\in D$.

Put $W=U/D$. Local and global unit duality gives $\Lambda$-rank
one in every odd character of $W$
\cite[Proposition~2.2.6 and its proof, pp.~3--4]{WWEOrd}.
The global minus units contribute roots of unity, and the
unramified quotient has rank zero. The ring $R$ is a PID:
$\Lambda$ is a two-dimensional regular local UFD, localization
at $p$ removes its height-two prime, and a one-dimensional
noetherian UFD is a PID. Therefore
\begin{equation}
                  W[1/p]=\bigoplus_{a\ {\rm odd}}Re_a .
                  \label{eq:Wfree}
\end{equation}
Finite generation clears the denominators of this isomorphism
and its inverse. Completed tensor powers thus have bounded
denominators. The exterior square has alternating coordinates
in $\Z_p[[T_1,T_2]][1/p]$.

\subsection{The two arithmetic characters}

The arithmetic inputs occur in two explicitly related components.
The local/global unit sequence in
\cite[Corollary~2.1.3, p.~8]{WWEPseudo} has the form
\begin{equation}
 0\longrightarrow\mathcal E_e/\mathcal C_e
 \longrightarrow\Lambda/(f_e)\longrightarrow T_e
 \longrightarrow X^{[e]}\longrightarrow0 ,
 \label{eq:unitseq}
\end{equation}
with equal characteristic ideals for $\mathcal E_e/\mathcal C_e$
and $X^{[e]}$. Here $\mathcal E_e,\mathcal C_e$ are the norm limits
of global and circular units in character $e$.
Finiteness of $X^{[e]}$ kills both end terms after inverting $p$;
hence $T_e[1/p]$ is cyclic.
The even $p$-ramified module $T$ is finite free over $\Z_p$ by
Kummer duality, Ferrero--Washington, and the torsionfreeness of $M$
\cite[Proposition~3.4.7 and Corollary~3.4.8, p.~92]{SharifiIwa}.
Let $q_e(z)$ be its characteristic polynomial on $T_e$ for
$\gamma$; cyclicity makes it the minimal polynomial as well.

The cup dual has action
\begin{equation}
 (g\lambda)(t)=\chi_p(g)^2\lambda(g^{-1}t).
 \label{eq:twodual}
\end{equation}
Thus $T_e^*(2)$ has character $\omega^{2-e}$; an eigenvalue
$\alpha$ becomes $u^2/\alpha$. Kummer adjunction identifies
$T_e[1/p]$ with the corresponding dual and twist of
$X^{[1-e]}[1/p]$, using the analogous formula with one twist.
The symbol theorem therefore uses
$\theta=\omega^{2-e}$ and $\chi=\omega^{1-e}$.
The excluded characters $\theta=1,\omega^2$ correspond to
$T_2,T_0$, which are zero by $X^{[-1]}=X^{[1]}=0$
\cite[proof of Lemma~5.2.7, p.~13]{WWEOrd}.

\subsection{The one-component cup input from Ohta's isomorphism}
\label{subsec:ohta-component-gate}

We specify the published isomorphism used here, including its recently
completed density input.  Ohta's Part~II, Theorem~(6.2.2), p.~409, in \cite{Ohta2020} proves the isomorphism given
by the upper-right coordinate of the Eisenstein Galois representation.
Its use of CM-point density in Part~I, p.~381, follows from Ohta's \cite[Theorem~(2.3.4), author pp.~15--16]{OhtaDensity}.  This match can be checked directly: the shifts
$\delta_{\mathfrak q}=\{\mathfrak q\}_\Gamma$ chosen on p.~378 of the
2020 paper lie in $\ker(\mathrm{Cl}_\infty\to\mathrm{Cl}_{n_1})$ and
are distinct modulo $\mathrm{Cl}^{\rm alg}$ by its
Proposition~(7.2.4)(iv).  They, together with the unbounded conductor
subsequence $n(\epsilon)$ and the classes in
$\ker(\mathrm{Cl}_{n(\epsilon)}\to\mathrm{Cl}_{n_1})$, are exactly the
data of the 2026 theorem with initial level $n_1$.  Forgetting the
finite level structures maps the former CM tuples onto the latter.
Finiteness and irreducibility of the modular cover then transfer the
density back to the source used in 2020.  Thus the required mod-$p$
nonvanishing input is a theorem.

Here is the action dictionary.  For an even nonexceptional $e$, set
\[
 \chi=\omega^{1-e},\qquad
 \theta_{\mathrm{WWE}}=\omega^{2-e},\qquad
 \theta_{\mathrm O}=\omega^{e-2}.
\]
We define the modular target separately from the arithmetic module
$M=H^{\mathrm{ab}}$.  In the modular-curve and dual-Hecke convention of
\cite[\S3.1]{WWEPseudo}, put $\theta=\theta_{\mathrm{WWE}}$ and
\[
 H'_{\mathrm{mod},\theta}
 =\varprojlim_r H^1_{\mathrm{et}}
   (X_1(p^r)_{\overline{\mathbb Q}},\mathbb Z_p)^{\mathrm{ord}}_\theta,
 \qquad
 H_{\mathrm{mod},\theta}
 =H'_{\mathrm{mod},\theta}\otimes_{\mathfrak h'_\theta}\mathfrak h_\theta.
\]
Here $X_1(p^r)$ is the compactified modular curve classifying
elliptic curves with a point of exact order $p^r$, and the ordinary
summand is the summand on which the dual $U_p$ operator is invertible.
The transitions are corestrictions, the subscript is the diamond
eigenspace, and $\mathfrak h_\theta$ is the localization of the algebra
$\mathfrak h'_\theta$ of dual Hecke operators at its Eisenstein maximal
ideal.  Let $I_\theta\subset\mathfrak h_\theta$ be the cuspidal Eisenstein
ideal, the image of the annihilator of the ordinary Eisenstein series
in the full Hecke algebra.  Set
\[
 M_{\mathrm{mod}}
 =H^-_{\mathrm{mod},\theta}/I_\theta H^-_{\mathrm{mod},\theta},
\]
where the minus sign denotes complex conjugation.  The coefficient
$\Lambda$ acts by $\gamma\mapsto\langle\kappa(\gamma)\rangle$.

Ohta's class character is
$(\theta_{\mathrm O}\omega)^{-1}=\chi$.  His tame level is $1$,
and the exponent $e-2$ is even, so that $-1$ belongs to its kernel,
as required in Theorem~(6.2.2).  The exceptional indices $0$ and $2$
have $T_0=T_2=0$ and need no cup-symbol argument.  Ohta's $X^-$ is a
splitting of the ordinary quotient $X/X^+$, with $X^+=X^{I_p}$
in his (6.1.14).  The comparison with the above minus space uses the
natural isomorphism $H^-_{\mathrm{mod},\theta}\to H_{\mathrm{quo},\theta}$
of \cite[Theorem~3.4.1 and Corollary~3.4.2]{WWEPseudo}, followed by the
modular-curve and diamond-action transport to that ordinary quotient.
In Ohta's (6.1.26), an element $\gamma$ with $\kappa(\gamma)=u$ acts on
$(X^-/I^*X^-)^\dagger$ as $u^{-1}\iota(u)^{-1}$.  After the Tate twist $(1)$,
this becomes the inverse diamond action $\iota(u)^{-1}$; its character
is $\theta_{\mathrm{WWE}}$.  The extension-coordinate map is
\[
 \sigma\longmapsto \sigma e^+-e^+\pmod{I^*X}.
\]
Sharifi's \cite[pp.~3 and 13--14]{SharifiProgram} explicitly uses the inverse-Galois convention and records the
Ohta isomorphism for this $\Upsilon$.  Equivalently, this is the
upper-right-coordinate definition in \cite[\S8.1]{WWEPseudo}:
evaluate the extension coordinate in
$\operatorname{Hom}_{\mathfrak h_\theta/I_\theta}
 (H^+_{\mathrm{mod},\theta}/I_\theta H^+_{\mathrm{mod},\theta},
 M_{\mathrm{mod}})$ on the distinguished element corresponding to
$1$ under $\mathfrak h_\theta/I_\theta\simeq
H^+_{\mathrm{mod},\theta}/I_\theta H^+_{\mathrm{mod},\theta}$.
Its restriction to the cyclotomic kernel is a homomorphism, and
ordinarity makes it factor through $X_\chi(1)$.
We therefore have the isomorphism
\[
 \Upsilon: V=X^{[1-e]}(1)\xrightarrow{\ \sim\ }M_{\mathrm{mod}}.
\]

Suppose now that $X^{[e]}$ is finite, so that the previously established
adjoint calculation makes $V[1/p]$ cyclic.  At a characteristic prime
$f$, let $R$ be the corresponding discrete valuation ring and write
$V_R\simeq R/(f^r)$.  Transport this generator to
$(M_{\mathrm{mod}})_R$ by $\Upsilon$.
The endomorphism $\Upsilon\varpi$ is multiplication by an element $a$,
and the Fukaya--Kato identity in
\cite[Theorem~5.21]{SharifiFK} gives
\[
 (a-1)\xi'=0\quad\hbox{in }R/(f^r).
\]
Writing $\xi=f^r g$ with $g\in R^\times$, separability in
characteristic zero gives $v_f(\xi')=r-1$.  Hence $a-1\in(f)$,
so $a$ is a unit.  Thus $\varpi_R$ is an isomorphism at every
characteristic prime, with no restriction on the multiplicity $r$.
It follows that the rational cup-symbol map is surjective on all of
$V[1/p]$. Both integral modules are finitely generated over $\Lambda$
with zero $\mu$-invariant, so the cokernel of the integral map
is a finite $\mathbb Z_p$-module. One power of $p$ kills it.
Reduction along the norm-compatible finite-level symbol maps
therefore gives the uniform integral bound used below.

Thus $X^{[e]}$ finite supplies cyclicity and the full rational
cup image in the twice-twisted companion.  The arithmetic character
transport is the displayed map $\theta_{\mathrm O}\mapsto
\theta_{\mathrm{WWE}}$ together with the Tate twist.

The cyclicity step is used in the direction displayed in the
Fukaya--Kato identity. For example, in a discrete valuation ring
with uniformizer $t$ and residue characteristic zero, take
$V=R/(t^2)$, $M_{\mathrm{mod}}=(R/t)^2$ and both maps zero. Then
$t(\Upsilon\varpi-1_{M_{\mathrm{mod}}})=0$, while $tV\ne0$.
In the argument above, Ohta's isomorphism identifies $M_{\mathrm{mod}}$ with the
cyclic module $V$ and turns the identity into multiplication by
$a$; that is the step which excludes this example.


For clarity we pass to group cohomology explicitly.
Put $S_n=\operatorname{Spec}\mathcal O_{F_{n+1}}[1/p]$.
The maximal pro-$p$ quotient of $\pi_1(S_n)$ is $P_n$:
normal closures of finite $p$-extensions of cyclotomic layers
remain $p$-extensions over $F_1$.
Inflation is an isomorphism on $H^1$ and an injection on $H^2$.
Indeed, an invariant homomorphism from the kernel of the
maximal pro-$p$ quotient to a finite $p$-module would give
a larger pro-$p$ quotient. The five-term sequence applies.
The universal finite \'etale tower has zero $H^1$ with finite
local coefficients, since every torsor becomes trivial in
that tower. Cartan--Leray consequently gives the same comparison
with \'etale cohomology \cite[Theorem~14.9 and
Proposition~14.13, pp.~96--98]{MilneEtale}.
These maps respect cups. The cup-symbol map just proved has
finite index in its integral target. Fix $p^{b_0}$ annihilating
that finite cokernel. The unique prime above $p$ is totally
ramified in every relative cyclotomic layer, so its intersection
with the lower Hilbert $p$-class field is trivial. The class norms
are therefore surjective. Consequently the projections
$V=X^{[1-e]}(1)\to H^2_{\mathrm{et}}(S_n,\Z/p^n(2))^{[2-e]}$
are onto, and $p^{b_0}$ times every such target lies in the
projection of the integral cup-symbol image. The comparison of
\cite[\S\S7.1.9--7.1.10, p.~86]{FK} then puts
$p^{b_0}$ times the asserted cohomology component in the image
of the cyclotomic-unit cups, uniformly at every level.

Reduction from coefficients $p^{n+1}$ to $p^n$ is surjective
by the class-group description in \cite[\S7.1.9]{FK}.
There is no Brauer contribution: the field is totally imaginary
and only one finite prime is omitted.
For a $p$-unit $v$, the quotient $v/c(v)$ is a global unit
with absolute value one at every embedding, hence a root of
unity. The odd unit classes are cyclic, and their alternating
self-cup is zero because $p$ is odd. Thus, for every $n\geq1$,
the subgroup $p^{b_0}H^2(P_n,\Z/p^n(2))^{[2-e]}$
is generated by cups of $c$-even classes annihilating $D$.
The exponent $b_0$ depends on $e$ and $p$, and is independent
of the level $n$.

\subsection{The projected class-two extension}

We give the comparison explicitly, because removing $D$ changes
the source of the cup calculation.
For a pro-$p$ group, brackets below mean closed commutator
subgroups. Since $p$ is odd, the quotient $H/[H,[H,H]]$
has its integral class-two Baker--Campbell--Hausdorff Lie
structure. Its central commutator module is
$\widehat{\bigwedge}^{\,2}_{\Z_p}M$. This follows first for a
finite free generating set by collecting commutators, and then
for $H$ by the inverse limit over finite subsets of a convergent
free generating sequence. Kill the central brackets involving
$T$. The remaining central module is
$\widehat{\bigwedge}^{\,2}U$, on which $c$ acts as $1$.
The minus part of this Lie module consequently maps
isomorphically to $U$; this gives a canonical odd section.
Kill its image of $D$, together with all brackets involving
$D$. The resulting Lie module has abelianization $W\oplus T$,
central module
\[
                   N=\widehat{\bigwedge}^{\,2}_{\Z_p}W,
\]
and plus part in the exact sequence
\begin{equation}
             0\longrightarrow N\longrightarrow \mathcal E
              \longrightarrow T\longrightarrow0 .
             \label{eq:projectedextension}
\end{equation}
All these quotients are stable under $\Gamma$ and $\Delta$.
The closedness assertions follow from compactness. The quotient
$W$ is $\Z_p$-torsionfree, so its completed exterior square
is torsionfree, as is seen in product $\Z_p$-bases.

Let $E$ denote the even exponents $e$ for which $X^{[e]}$
is finite. For $e\in E$ with $T_e\ne0$, choose a cyclic vector
$v_e\in T_e[1/p]$ and a lift $y_e\in\mathcal E^{[e]}[1/p]$.
Set
\begin{equation}
                C_e=q_e(\gamma)y_e\in N[1/p].
                \label{eq:projectedcompanion}
\end{equation}
The action of $\gamma$ on two odd coordinates is multiplication
by $z=(1+T_1)(1+T_2)$. Thus $C_e$ has alternating coefficient
series $C_{e;a,b}(T_1,T_2)$, with $a,b$ odd and
$a+b=e$ modulo $p-1$. Each series has a common bounded
denominator.

\begin{lemma}\label{lem:projectedpole}
For each root $\alpha$ of $q_e$, at least one matching coefficient
of $C_e$ has nonzero restriction to
$(1+T_1)(1+T_2)=\alpha$.
\end{lemma}
\begin{proof}
Let $L/\Q_p$ split $q_e$, with ring of integers $\mathcal O$.
Every eigenvalue satisfies $\alpha-1\in\mathfrak m_{\mathcal O}$,
since the action of $\Gamma$ on the finite lattice $T_e$ is
continuous. Suppose all the restrictions vanish. The coordinates
\[
 S=(1+T_1)(1+T_2)-\alpha,\qquad
 T_2=\frac{\alpha+S}{1+T_1}-1
\]
give an isomorphism of complete local power-series rings over
$\mathcal O$. Vanishing at $S=0$ is divisibility by $S$.
Division preserves bounded denominators and alternation.
It follows that $C_e=(\gamma-\alpha)A$ for some $A\in N_L$.

Write $r(z)=q_e(z)/(z-\alpha)$. In the cyclic module $T_{e,L}$,
the vector $w=r(\gamma)v_e$ is nonzero and spans an
$\alpha$-eigenline $J$. This remains true for a repeated root:
$r$ has degree one less than the annihilating polynomial.
The vector $r(\gamma)y_e-A$ lifts $w$ equivariantly.
Thus \eqref{eq:projectedextension} splits over $J$.

The lattice $J\cap T_{\mathcal O}$ is a direct summand of
$T_{\mathcal O}$: its quotient is torsionfree over the DVR
$\mathcal O$. Extend this equivariant lift to an
$\mathcal O$-linear section
$s:T_{\mathcal O}\to\mathcal E_{\mathcal O}[1/p]$.
Some fixed $p^b$ makes its image integral. For $h\in H$ write
its image in the class-two Lie module as
\[
        x(h)+s(t(h))+z(h),\quad
 x(h)\in W_{\mathcal O},\ t(h)\in T_{\mathcal O},\
 z(h)\in N_{\mathcal O}[1/p].
\]
The same $p^b$ makes every $z(h)$ integral. The section defect
\[
        \beta_\delta(t)=\delta s(t)-s(\delta t)
\]
vanishes on $J$ for $\delta\in\Gamma$.

Fix $n$ and take two of the even cyclotomic unit classes $f,g$
whose cups span the indicated finite-index subgroup. Their invariant cocycle
representatives are obtained by averaging with $(1+c)/2$.
They vanish on the $\Gamma^{p^n}$ section, since $c$ fixes
that section and negates $\mathcal O/p^n(1)$. On $H$ they
vanish on $T$ by parity and on $D$ by Lemma~\ref{lem:norms};
therefore they factor through $W$.
Let $\psi=f\wedge g:N_{\mathcal O}\to\mathcal O/p^n(2)$.
Define the continuous cochain
\[
 k(h)=-\frac{p^b}{2}f(h)g(h)-\psi(p^b z(h)),\qquad
 k(h\delta)=k(h).
\]
The class-two product adds $[x(h),x(h')]/2$ to its central
coordinate. Using $dk(h,h')=k(h')-k(hh')+k(h)$ on $H$
gives
\[
                       dk(h,h')=p^bf(h)g(h').
\]
Indeed the quadratic term contributes
$p^b(f(h)g(h')+f(h')g(h))/2$, and the central term contributes
$p^b(f(h)g(h')-f(h')g(h))/2$.
For $\delta,\delta'\in\Gamma^{p^n}$ the remaining cocycle is
\begin{equation}
 (p^bf\smile g-dk)(h\delta,h'\delta')
       =-\psi\bigl(p^b\beta_\delta(t(h'))\bigr).
 \label{eq:projectedcochain}
\end{equation}
To verify this formula, use equivariance of $f,g$ on $H$,
their vanishing on the section, and
$z(\delta h'\delta^{-1})=\delta z(h')+
\beta_\delta(t(h'))$. This also verifies the coefficient
twist $\chi_p(\delta)^2$ in the coboundary.
The edge class in \eqref{eq:projectedcochain} restricts to zero
on $J$.

We spell out why this contradicts the uniform finite-index assertion
without making compatible choices of cup expressions.
Both $H$ and $\Gamma^{p^n}$ have $p$-cohomological dimension
one. Hochschild--Serre gives
\[
 H^2(P_n,\mathcal O/p^n(2))^+
 \simeq H^1(\Gamma^{p^n},
             \Hom_{\Z_p}(T,\mathcal O/p^n)(2)).
\]
Put $V=\Hom_{\Z_p}(T,\mathcal O)(2)$ and let $B$ be its
$\gamma$-operator. Evaluation on the procyclic generator
identifies this group with
\begin{equation}
             V/(p^nV+(B^{p^n}-1)V).
             \label{eq:coinvariants}
\end{equation}
Corestriction sends a generator value to that same value
modulo $(B^{p^n}-1)V$: this follows by evaluating the
corestriction cocycle on $\gamma^{p^n}$, where exactly one
coset representative crosses a fundamental domain.
Thus these transitions are quotient maps.
Continuity makes their defining submodules cofinal with the
$p$-adic neighbourhoods of zero, so their inverse limit is $V$.

Restriction to $J\cap T_{\mathcal O}$ is a surjection on
integral dual lattices. Its target has character $2-e$ and
$\gamma$-eigenvalue $u^2/\alpha$ by \eqref{eq:twodual}.
It induces surjections on the finite quotients
\eqref{eq:coinvariants}. By the uniform finite-index assertion and
\eqref{eq:projectedcochain}, the image of the fixed submodule
$p^{b+b_0}V^{[2-e]}$ is zero in every such quotient of the dual
lattice of $J$. Their separated inverse limit is that
nonzero free $\mathcal O$-module. A surjection onto it cannot
annihilate $p^{b+b_0}V^{[2-e]}$. This contradiction proves the lemma.
\end{proof}

\subsection{The rational outer image and its odd coefficients}

We construct the rational target and prove that its linear
coefficients detect the quotient $W$.
Hain--Matsumoto's Theorem~8.4 and the surrounding construction,
together with Brown's Theorem~1.1, give the weight-completed
free Lie algebra on $s_m$ for odd $m\geq3$ as the Lie algebra
of the rational outer image \cite{HM,Brown}.
The image of $H$ is dense in every finite weight quotient.
The reason is that the cyclotomic kernel in
$G_{\Q,\{p\}}$ maps onto $H$, and the outer action factors
through $\Gal(\Omega_p/\Q)$.
A Levi splitting and successive conjugation put $\gamma$
in the torus: in weight $m>0$ the operator $u^m-1$ is
invertible over $\Q_p$. The commuting prime-to-$p$ action
then lies in the same torus, since no nonzero positive weight
centralizes $\gamma$. In these coordinates
\[
           \gamma s_m=u^ms_m,\qquad
           \delta s_m=\omega(\delta)^m s_m.
\]
Choose a nonzero annihilator of $D$ in $\Lambda$. It has
only finitely many roots $u^m$ with odd $m\geq3$.
Pass to the quotient of the completed free Lie algebra by
the closed ideal generated by the letters at those finitely
many weights. The quotient is the completed free Lie algebra
on the remaining letters, with the same torus action, and
the image of $H$ remains dense in every finite weight quotient.
Every odd residue class still contains infinitely many of
these weights. Let $\rho$ denote the resulting representation
and $K=\ker(\rho_p)\cap H$ the kernel of the original
pro-$p$ outer action. Functorial rational completion and
this quotient give $K\subseteq\ker\rho$.

For each remaining weight $m$, the coefficient of $s_m$
is an additive functional on $M$. It vanishes on $T$ by
parity and on $D$ because its annihilator specializes to
a nonzero scalar at $\gamma=u^m$. Thus it factors through $W$.
In the coordinates \eqref{eq:Wfree}, the coefficient has
the form
\[
              b_m\,F_a(u^m-1),\qquad a\equiv m\pmod{p-1}.
\]
Here $b_m\ne0$: otherwise density in the abelian weight-$m$
quotient would fail. The argument only uses these remaining weights.

For $h\in K$, every such coefficient is zero. A nonzero
bounded-denominator series in one variable has only finitely
many zeros in $p\Z_p$, by Weierstrass preparation. The
infinitely many weights in each odd residue class therefore
force every $F_a$ to vanish. Thus the image of $h$ in $W$
is zero.


% To follow abelian_components.tex.  Uses its notation and proved lemmas.
% No private macros beyond the main paper's existing conventions.

\section{A canonical nonabelian quotient of class two}
\label{sec:class-two-components}

We construct a canonical class-two quotient for every odd prime.
Its full odd commutator module is detected by the outer action;
the cup comparison determines which even abelian components survive.

Retain the notation $W=U/D$ and the class-two Lie module constructed
in Section~\ref{sec:abelian}.  Write it as
\[
 \mathcal L=W\oplus\mathcal E,\qquad
 0\longrightarrow\widehat{\bigwedge}^{\,2}W
 \longrightarrow\mathcal E\longrightarrow T\longrightarrow0.
\]
The first summand is its canonical $c$-negative part, the second is
central and $c$-positive, and $[w,w']=w\wedge w'$ in the displayed
central submodule.  The direct sum $W\oplus\mathcal E$ is an
identity of additive Lie modules; no section of
$\mathcal E\to T$ is chosen.

For a set $E$ of even characters define the central Lie ideal
and its quotient by
\[
 \mathcal I_E=\bigoplus_{b\notin E}\mathcal E^{[b]},\qquad
 \mathcal L_E=\mathcal L/\mathcal I_E,
 \qquad C_E=(\mathcal L_E,\star),\qquad
 x\star y=x+y+\tfrac12[x,y].
\]
Only even $b$ occur in the direct sum.  These are compact modules,
and $2$ is a unit in $\mathbb Z_p$, so $C_E$ is a pro-$p$ group.
Let $q_E:H\twoheadrightarrow C_E$ be the composite of the
class-two quotient, the quotient killing the canonical odd section
of $D$ and its brackets, and the quotient by $\mathcal I_E$.
In additive coordinates its bracket is
\[
 [w,w']_{C_E}=\operatorname{pr}_E(w\wedge w').
\]
Thus discarding $b$ removes both $T_b$ and the central brackets
in character $b$.  This formula specifies every map used in the
construction.

Put $Z=\mathcal E_{\Lambda\text{-tor}}$ and
$Q=(W\oplus\mathcal E/Z,\star)$. The proof below constructs
an injection of $Q$ into a quotient of the rational outer image.


\begin{theorem}[Canonical class-two detection]
\label{thm:canonical-torsion-quotient}
For every odd prime $p$, the group $Z$ is closed, finite free
 over $\mathbb Z_p$, and saturated in $\mathcal E$; it has
zero intersection with $N=\widehat{\bigwedge}^{\,2}W$.
The kernel of $q$ is normal in $\operatorname{Gal}(\Omega_p/\mathbb Q)$
and contains the pro-$p$ outer-action kernel in $H$. Thus
\[
 L^{(2)}=\Omega_p^{\ker q}\subseteq\Omega_p^*,\qquad
 \operatorname{Gal}(L^{(2)}/F_\infty)=Q.
\]
The group $Q$ is torsionfree and nonabelian of class two, with
\[
 [Q,Q]=N,\qquad Q^{\rm ab}=W\oplus T/\pi_T(Z),
\]
where $\pi_T:\mathcal E\to T$ is the displayed projection.
For every even $e$ such that $X^{[e]}$ is finite,
$Z^{[e]}=0$; consequently the full $T_e$ occurs in the even
abelianization of $Q$.
\end{theorem}

We first record the separation of the completed commutator module
used in the proof.  It is separation by generator length in the
full free Lie algebra of the rational outer image.

\begin{lemma}\label{lem:class-two-moments}
Let $\mathcal M$ be the odd weights retained after removing the
finite exceptional set in Section~\ref{sec:abelian}.
Write the linear coefficients of that rational outer
representation as
\[
 \kappa_m(w)=b_m F_a(u^m-1),\qquad
 a\equiv m\pmod{p-1},\quad b_m\ne0.
\]
The induced bracket map from
$\widehat{\bigwedge}^{\,2}_{\mathbb Z_p}W$
to the length-two part of the weight-completed free Lie algebra
on the $s_m$, $m\in\mathcal M$, is injective.
The assertion also holds after projecting source and target onto
any collection of $\Delta$-characters.
\end{lemma}

\begin{proof}
Choose the rational $R$-basis of $W[1/p]$ from
\eqref{eq:Wfree}.  Clearing denominators in this basis and its
inverse embeds $W$ in a finite sum of copies of $\Lambda$,
with cokernel killed by a fixed power of $p$.
The same construction embeds its completed tensor square
in a finite sum of $\mathbb Z_p[[T_1,T_2]][1/p]$.
For precision, a reverse map multiplied by $p^r$ gives a
composite equal to $p^r$ on $W$; tensoring twice gives
$p^{2r}$.  The completed tensor square is $p$-torsionfree,
as is seen from a product $\mathbb Z_p$-basis of $W$, so
the tensor map is injective.  Since $2$ is invertible,
the completed exterior square is its alternating summand.

Let $A_{a,b}(T_1,T_2)$ be its alternating coordinates.
The coefficient at $[s_m,s_n]$ is
\[
             b_mb_n A_{a,b}(u^m-1,u^n-1).
\]
Use antisymmetry to read a coefficient with $m>n$ as the
negative of that with the two indices reversed.  Thus these
coefficients are available for every distinct ordered pair,
not only for $m<n$.
If they all vanish, fix $n$ and vary $m$ through its residue
class in $\mathcal M$, excluding $m=n$.  This is an infinite
set of distinct points in $p\mathbb Z_p$.
The one-variable Weierstrass theorem makes
$A_{a,b}(T_1,u^n-1)$ identically zero.  Varying $n$ and
applying the same argument to each coefficient of $T_1$
gives $A_{a,b}=0$.  Every coordinate vanishes.
For a fixed character $e$, only pairs $a+b=e$ occur;
the same argument uses exactly those pairs and proves
the projected assertion.
\end{proof}

\begin{lemma}\label{lem:class-two-invariance}
The quotient $\mathcal L$ and each central character quotient $C_E$ have kernels stable under $\operatorname{Gal}(\Omega_p/\mathbb Q)$.
\end{lemma}
\begin{proof}
We first check the invariance and the canonical quotient map.
Before removing $D$, the first quotient can be written
intrinsically as
\[
 H/\overline{[\pi^{-1}(T),H]},\qquad \pi:H\longrightarrow M.
\]
Indeed $[H,H]\subseteq\pi^{-1}(T)$, so its third lower
central subgroup is killed, and the further killed central
brackets are precisely those involving $T$.
Its central subgroup above $T$ is the $c$-positive Lie
module $\mathcal E$ before removing $[D,U]$.
The ideal generated by the canonical negative lift of $D$
is that lift plus $[D,U]$.
It is preserved by $\Gamma$ and $\Delta$, and it is an ideal,
so it is preserved by conjugation from $H$ as well.
Here is the change-of-section calculation for full arithmetic
invariance. For $g\in\operatorname{Gal}(\Omega_p/\mathbb Q)$ write
$gcg^{-1}=\operatorname{Ad}(h)c$ and let $v$ be the odd abelian
part of $\log h$. If $s_c:U\to\mathcal L_0^-$ is the negative
section before removing $D$, write
$g s_c(u)=s_c(g_Uu)+\lambda_g(u)$ with central $\lambda_g(u)$.
The identity $gc=\operatorname{Ad}(h)cg$ gives
\[
 -s_c(g_Uu)-\lambda_g(u)
 =-s_c(g_Uu)+\lambda_g(u)-[v,g_Uu],
 \qquad \lambda_g(u)=\tfrac12[v,g_Uu].
\]
As $g_U(D)=D$, the difference belongs to $[D,U]$. Thus the
ideal is invariant under every $g$, independently of the section.
The action on the central positive part is unaffected by inner
conjugations.  Consequently
its character summands, and hence $\mathcal I_E$, are
independent of these lift choices.  The closed subgroup
$\ker q_E$ is normal in $\operatorname{Gal}(\Omega_p/\mathbb Q)$.

\end{proof}

\begin{proof}[Proof of Theorem~\ref{thm:canonical-torsion-quotient}]
We first establish the module assertions used to form the quotient.
In the bounded coordinates for $N$ constructed above, the action
of $\gamma$ is multiplication by $(1+T_1)(1+T_2)$.
For any nonzero $f\in\mathbb Z_p[[\gamma-1]]$, the series
\[
             f((1+T_1)(1+T_2)-1)
\]
is nonzero: setting $T_2=0$ recovers $f(T_1)$.
The coefficient rings $\mathbb Z_p[[T_1,T_2]][1/p]$ are
integral domains. Multiplication by every nonzero $f$ is
therefore injective on $N$, so $N$ is $\Lambda$-torsionfree.

The module $T$ is finite free over $\mathbb Z_p$.
Let $a(z)$ be the characteristic polynomial of $\gamma$ on
$T$; if $T=0$, put $a=1$. Cayley--Hamilton gives
$a(\gamma)\mathcal E\subseteq N$. If $x\in Z$, then
$a(\gamma)x$ is a $\Lambda$-torsion element of $N$, hence zero.
Conversely every element killed by $a(\gamma)$ is
$\Lambda$-torsion. Thus
\[
                     Z=\ker(a(\gamma):\mathcal E\to\mathcal E).
\]
This identifies $Z$ as a closed submodule. Its intersection
with $N$ is zero, so $\pi_T$ embeds $Z$ in the finite free
$\mathbb Z_p$-module $T$. In particular $Z$ is finite free
over $\mathbb Z_p$. The module $\mathcal E$ is
$\mathbb Z_p$-torsionfree, and the kernel of an endomorphism
of such a module is saturated: $px\in Z$ implies
$p\,a(\gamma)x=0$, hence $a(\gamma)x=0$ and $x\in Z$.
Consequently $\mathcal E/Z$ is also torsionfree.

The arithmetic action on the central module $\mathcal E$
factors through $\Delta\times\Gamma$. Its torsion submodule
$Z$ is invariant under that action. The preceding intrinsic
construction of $\mathcal L$, including the change-of-section
calculation, proves invariance under the entire arithmetic
group. Thus $q$ has Galois-stable kernel.

Take the rational free outer image constructed above. In
addition to the finite exceptional odd weights coming from
$D$, remove every odd letter occurring in a positive pair
$m,n\geq3$ with
\[
                              a(u^{m+n})=0.
\]
There are only finitely many such pairs: $a$ has finitely
many roots, $u$ has infinite order, and any fixed positive
sum has finitely many positive summands. The remaining
set of odd weights is cofinite. Density in each remaining
abelian weight quotient still makes its linear coefficient
$b_m$ nonzero.

Pass to generator length at most two in this free target.
The linear image of $T$ is zero by parity, so its brackets
have length at least three. The canonical odd section of $D$
has zero linear image by the choice of weights and zero
quadratic image by $c$-parity. Its brackets also vanish.
Thus the map factors through $\mathcal L$.
The image of $Z$ is central of length two. Its coefficient
at $[s_m,s_n]$ is killed by $a(u^{m+n})$, which is nonzero
by construction. Every such coefficient is therefore zero.
We obtain a factorization on the whole group
\[
 H\xrightarrow{q} Q\xrightarrow{j}
      \exp(\mathfrak f_{\leq2}),
\]
where $\mathfrak f$ is free on the remaining odd letters.

We prove that $j$ is injective. The degree-one moment
calculation detects $W$, so an element in its kernel lies
in $\mathcal E/Z$. Represent it by $x\in\mathcal E$.
Then $a(\gamma)x\in N$ has zero image in the length-two
target. The bracket moment lemma is still valid after
deleting any finite set of odd weights, and is injective
on all of $N$. Hence $a(\gamma)x=0$. The kernel identity
for $Z$ proved above gives $x\in Z$, as required.

The pro-$p$ outer kernel maps to zero in this rational
quotient. The factorization and injectivity therefore give
its containment in $\ker q$. Galois correspondence proves
the field assertions. Since $Z\cap N=0$, the commutator
module remains exactly $N$; the abelianization is
$W\oplus T/\pi_T(Z)$. The additive Lie module
$W\oplus\mathcal E/Z$ is torsionfree, and powers in a
class-two BCH group satisfy $x^{p^r}=p^rx$. Thus $Q$ has
no nontrivial finite-order element. Its commutator module
is nonzero, because $W[1/p]$ has infinite dimension.

Finally fix an even $e$ with $X^{[e]}$ finite.
Suppose $Z^{[e]}[1/p]$ is nonzero. Over a finite splitting
field $L/\mathbb Q_p$ it contains an eigenvector $x$ of eigenvalue $\alpha$.
Its nonzero image in $T_e$ belongs to the unique
$\alpha$-eigenline of that cyclic module. In the notation
of the companion identity, there is a nonzero scalar $d$
with
\[
 \pi_T(x)=d\,r(\gamma)v_e,\qquad
 r(z)=q_e(z)/(z-\alpha).
\]
Write $x=d\,r(\gamma)y_e-A$, with $A\in N\otimes_{\mathbb Z_p}L$.
Applying $\gamma-\alpha$ gives
\[
                         d\,C_e=(\gamma-\alpha)A.
\]
Every coefficient of $C_e$ would therefore vanish on
$(1+T_1)(1+T_2)=\alpha$, contradicting the proved cup-pole
lemma. Hence $Z^{[e]}=0$. For a set $E$ of such characters,
$Z\subseteq\bigoplus_{b\notin E}\mathcal E^{[b]}$.
Thus the original map $q_E$ factors through $q$, and
$\ker q\subseteq\ker q_E$. The corresponding field
inclusion is $L_E^{(2)}\subseteq L^{(2)}$.
\end{proof}



\begin{corollary}[Canonical class-two component detection]
\label{thm:class-two-components}
Let $p$ be odd, and let $E$ be a set of even exponents for which
$X^{[e]}$ is finite for every $e\in E$.
Then the pro-$p$ outer-action kernel in $H$ is contained
in $\ker q_E$.  The field
\[
 L_E^{(2)}=\Omega_p^{\ker q_E}
\]
is Galois over $\mathbb Q$ and is contained in $\Omega_p^*$.
Its Galois group over $F_\infty$ is $C_E$, and
\begin{align*}
 C_E^{\rm ab}&=(U/D)\oplus\bigoplus_{e\in E}T_e=B_E,\\
 [C_E,C_E]&=\bigoplus_{e\in E}
       \left(\widehat{\bigwedge}^{\,2}_{\mathbb Z_p}(U/D)\right)^{[e]}.
\end{align*}
The largest subextension of $L_E^{(2)}/F_\infty$ which is
abelian over $F_\infty$ is the field $L_E$ in
Corollary~\ref{thm:abelian}.  For nonempty $E$, the group $C_E$
is nonabelian.
\end{corollary}

\begin{proof}[Proof of Corollary~\ref{thm:class-two-components}]
For every selected $e$, Theorem~\ref{thm:canonical-torsion-quotient}
gives $Z^{[e]}=0$. Hence $q_E$ factors through $q$, and
$L_E^{(2)}\subseteq L^{(2)}\subseteq\Omega_p^*$.
Lemma~\ref{lem:class-two-invariance} gives normality over
$\mathbb Q$. The displayed bracket formula gives the
commutator and abelianization, including the maximal abelian
subfield assertion.

Finally each odd character of $W[1/p]$ is a nonzero rank-one
$R$-module and has infinite dimension over $\mathbb Q_p$.
For any even $e$, choose odd $a,b$ with $a+b=e$.
For $a\ne b$, nonzero vectors in these two summands have
nonzero exterior product.  For $a=b$, choose two
$\mathbb Q_p$-independent vectors in that summand.
Clearing denominators places them in $W$ and preserves their
nonzero wedge.  Thus every selected central bracket
component is nonzero, proving nonabelianness for nonempty $E$.
\end{proof}




\begin{corollary}[Abelian component detection]\label{thm:abelian}
For the same set $E$, put $B_E=W\oplus\bigoplus_{e\in E}T_e$
and let $L_E$ be the fixed field of the kernel of the canonical
map $H\to H^{\rm ab}\to B_E$. Then $L_E/\mathbb Q$ is Galois,
$\operatorname{Gal}(L_E/F_\infty)=B_E$, and
$L_E\subseteq L_E^{(2)}\subseteq\Omega_p^*$.
\end{corollary}
\begin{proof}
The abelianization of $q_E$ is exactly the displayed map to $B_E$.
Thus $\ker q_E$ is contained in its kernel. Galois correspondence
and Corollary~\ref{thm:class-two-components} give both inclusions.
The kernel is stable under the full arithmetic group because the
character projections and $D$ are stable.
\end{proof}

\subsection{Positive cyclotomic eigenvalues}
\label{subsec:positive-eigenvalues}

The arithmetic cohomology also determines which scalar
denominators can vanish at the odd moment weights.
Put $P=\operatorname{Gal}(\Omega_p/F_1)$ and write
$u=\kappa(\gamma)$ for the cyclotomic value of a generator
of $\Gamma$.

\begin{lemma}\label{lem:positive-eigenvalues}
For every integer $m\geq2$,
\[
 \ker(\gamma-u^m:H^{\rm ab}\longrightarrow H^{\rm ab})=0.
\]
Consequently the characteristic polynomial of $\gamma$
on $D\oplus T$, or on any of its rational subquotients,
has nonzero value at $u^m$.
\end{lemma}
\begin{proof}
The freeness theorem for $H$ cited in Section~\ref{sec:abelian}
gives $\operatorname{cd}_pH\leq1$; also
$\operatorname{cd}_p\Gamma=1$.
For $A_n=\mathbb Z/p^n(m)$, Hochschild--Serre for
$1\to H\to P\to\Gamma\to1$ gives the edge isomorphism
\[
 H^2(P,A_n)\simeq
 \frac{\operatorname{Hom}_{\rm cts}(M,A_n)}
      {(\gamma-1)\operatorname{Hom}_{\rm cts}(M,A_n)}.
\]
The action is $(\gamma f)(x)=u^m f(\gamma^{-1}x)$.
Choose a coefficient generator and pair this discrete
module with $M/p^nM$. The adjoint of $\gamma-1$ is
$u^m\gamma^{-1}-1$. Thus its Pontryagin dual is exactly
\begin{equation}\label{eq:positive-eigenvalue-dual}
 H^2(P,A_n)^\vee\simeq
       \ker(\gamma-u^m:M/p^nM\longrightarrow M/p^nM).
\end{equation}
The choice rescales the pairing by a unit and changes
neither this kernel nor its order.

Let $\mathcal G=G_{F_1,S}$, with $S$ the primes above $p$
and the archimedean places, and put
$J=\ker(\mathcal G\to P)$. The coefficients $A_n$ factor
through $P$. A $P$-equivariant homomorphism $f:J\to A_n$
would give the exact extension
\[
 1\longrightarrow\operatorname{im}f\longrightarrow
       \mathcal G/\ker f\longrightarrow P\longrightarrow1.
\]
Its middle group is pro-$p$, so maximality of $P$ forces
$f=0$. Hence $H^1(J,A_n)^P=0$, and the five-term sequence
gives an injection
\[
                    H^2(P,A_n)\lhook\joinrel\longrightarrow
                         H^2(\mathcal G,A_n).
\]

The Chern-class comparison and even $K$-group finiteness
\cite[VI, Theorem~8.2, chapter p.~47]{Weibel} give the
finite group
\[
 H^2(\mathcal G,\mathbb Z_p(m))
     \simeq K_{2m-2}(O_{F_1}[1/p])\otimes\mathbb Z_p.
\]
Here $m\geq2$ is essential. Since $\mathcal G$ has
$p$-cohomological dimension two, continuous
$H^3(\mathcal G,\mathbb Z_p(m))=0$: the finite-coefficient
degree-three groups vanish, and the inverse system of
finite degree-two groups has zero $\varprojlim^1$.
The coefficient exact sequence therefore identifies
\[
 H^2(\mathcal G,A_n)
      =H^2(\mathcal G,\mathbb Z_p(m))/p^n.
\]
Its order, and hence the order in
\eqref{eq:positive-eigenvalue-dual}, is bounded independently of $n$.

A nonzero $u^m$-eigenvector in $M$ can be divided by its
maximal power of $p$ to give one outside $pM$.
Indeed $M$ is a product of copies of $\mathbb Z_p$ and
has no $p$-torsion, since it is the abelianization of
a free pro-$p$ group. Its image modulo $p^n$ would then
have order $p^n$ in the kernel in
\eqref{eq:positive-eigenvalue-dual}, contradicting the
uniform bound. This proves the first assertion.
A rational eigenvector in $D\oplus T$ can have its
denominator cleared in the integral lattice, so none
exists. The characteristic polynomials of rational
subquotients divide that of the whole finite-dimensional
module, proving the last assertion.
\end{proof}

Every nonempty product of positive odd target weights
has eigenvalue $u^{m_1+\cdots+m_d}$ with each $m_i\geq3$.
The lemma therefore shows that all companion denominators
coming from $D\oplus T$ are nonzero on this entire
moment grid. It also gives the exact cohomological reason
for this nonvanishing; the construction by a cofinite
grid in the preceding proof gives the same quotient map.

\section{A canonical detected quotient of class three}
\label{sec:canonical-class-three}

We enlarge the canonical quotient by its central relations.
Let \(q_2=q:H\twoheadrightarrow Q_2=Q\) be the quotient of
Theorem~\ref{thm:canonical-torsion-quotient}. Write
\[
 Q_2=(W\oplus\overline{\mathcal E},\star),\qquad
 \overline{\mathcal E}=\mathcal E/Z,\qquad
 N=\widehat{\bigwedge}^{\,2}_{\mathbf Z_p}W\subseteq\overline{\mathcal E}.
\]
Thus \(W=U/D\), \(D=U_{\Lambda\text{-tors}}\), and
\[
 0\longrightarrow N\longrightarrow\overline{\mathcal E}
 \longrightarrow T/\pi(Z)\longrightarrow0.
\]
Define, using closed commutator subgroups,
\[
 B=\ker(H\to W),\qquad K_2=\ker q_2,\qquad
 J_4=\overline{[H,K_2][B,B]},\qquad
 \mathcal R=K_2/J_4.
\]
Let \(K_3\) be the inverse image of \(\mathcal R_{\Lambda\text{-tor}}\) in \(K_2\), and put \(L^{(2)}=\Omega_p^{K_2}\).

Since \(q_2(B)=\overline{\mathcal E}\) is central,
\([H,B]\subseteq K_2\), and \(\mathcal R\) is central in \(H/J_4\).
The full arithmetic group preserves all these subgroups.
Its action on \(\mathcal R\) factors through \(\Gamma\times\Delta\).

\begin{theorem}[Canonical class-three detection]
\label{thm:canonical-class-three}
For every odd prime $p$, the module
$\mathcal R_{\Lambda\text{-tor}}$ is closed and $K_3$ is
normal in $\operatorname{Gal}(\Omega_p/\mathbb Q)$.
The pro-$p$ outer-action kernel in $H$ is contained in $K_3$.
Consequently
\[
 L^{(2)}\subsetneq L^{(3)}=\Omega_p^{K_3}\subseteq\Omega_p^*,
 \qquad \operatorname{Gal}(L^{(3)}/F_\infty)=Q_3=H/K_3.
\]
The group $Q_3$ is torsionfree and has nilpotency class exactly three.
There is a natural injection
\[
 \operatorname{Lie}^3(W)\lhook\joinrel\longrightarrow
                         \ker(Q_3\to Q_2).
\]
In particular, every nonzero degree-three Lie word in the odd
quotient $W$ survives.
\end{theorem}

\begin{proof}[Proof of Theorem~\ref{thm:canonical-class-three}]
Put \(\mathcal S=[H,B]/J_4\). We first calculate this central module
integrally:
\begin{equation}\label{eq:class3-S}
 \mathcal S\simeq
 \frac{W\widehat\otimes_{\mathbf Z_p}\overline{\mathcal E}}
 {J(\widehat{\bigwedge}^{\,3}W)},\qquad
 J(x\wedge y\wedge z)
 =x\otimes(y\wedge z)+y\otimes(z\wedge x)+z\otimes(x\wedge y).
\end{equation}
The map is the commutator of lifts of \(W=H/B\) and
\(\overline{\mathcal E}=B/K_2\). Changing either lift changes it by
\([B,B]\) or \([H,K_2]\), so it is well-defined.
It is bilinear because its image is central, and Hall--Witt
gives the displayed Jacobi relation.

For the exact presentation, split \(U=D\oplus W\) as compact
\(\mathbf Z_p\)-modules. Choose a convergent basis of \(W\),
finite bases of \(D,T\), and lifts forming a free pro-\(p\)
basis of \(H\). The lifts may be chosen so that \(D\) maps
to \(1\), the \(W\)-basis maps to the odd section of \(Q_2\),
and the \(T\)-basis maps to selected central lifts.
The necessary corrections follow from the surjection
of the derived subgroup onto \(N\).

Split \(\mathcal E=N\oplus\bigoplus_i\mathbf Z_pt_i\), and write a
basis of \(Z\) as
\[
 z_s=n_s+\sum_i a_{si}t_i.
\]
Modulo \(J_4\), the \(D\)-lifts are central,
\([W,t_i]\) and triple \(W\)-commutators are central,
and \([t_i,t_j]\), \([t_i,[W,W]]\), and brackets involving
\(D\) vanish. The surviving central derived generators
have exactly the relations
\[
 \text{Jacobi},\qquad
 [W,z_s]=[W,n_s]+\sum_i a_{si}[W,t_i]=0.
\]
Indeed, collect a free nilpotent group of class three in a
Hall basis. Before the \(z_s\)-relations, its surviving
central module is
\[
 \operatorname{Lie}^3(W)
 \oplus\left(W\widehat\otimes\bigoplus_i\mathbf Z_pt_i\right).
\]
The remaining relations are precisely \(W\widehat\otimes Z\),
which proves \eqref{eq:class3-S}. No coefficient \(a_{si}\)
has been inverted. In particular all integral torsion
arising from \(T/\pi(Z)\) remains in this presentation.
The Hall collection and the leading Magnus terms give
the integral independence of these generators; passing
through finite-coordinate quotients gives the compact
statement. This uses no class-three BCH denominator and
is valid also at \(p=3\).

The abelian image gives the further exact sequence
\begin{equation}\label{eq:class3-R}
 0\longrightarrow \mathcal S\longrightarrow \mathcal R
 \longrightarrow D\oplus\pi(Z)\longrightarrow0.
\end{equation}
To identify its kernel, the derived subgroup of
\(H/[H,B]\) is \(\widehat{\bigwedge}^2W=N\), and the induced
map to \(Q_2\) is injective there. Hence
\(K_2\cap[H,H]=[H,B]\).
The image of \(K_2\) in \(H^{\rm ab}\) is exactly the
kernel of \(M\to Q_2^{\rm ab}\), namely \(D\oplus\pi(Z)\).
This proves \eqref{eq:class3-R}.

Let \(p^b\) kill the finite torsion subgroup of \(T/\pi(Z)\).
We claim
\begin{equation}\label{eq:class3-pbound}
                         p^b\mathcal S[p^\infty]=0.
\end{equation}
Let \(N^{\rm sat}\) be the \(p\)-saturation of \(N\) in
\(\overline{\mathcal E}\). Then \(p^bN^{\rm sat}\subseteq N\), and
\(\overline{\mathcal E}/N^{\rm sat}\) is torsionfree.
The Jacobi image is saturated in \(W\widehat\otimes N\),
because its quotient is the torsionfree module
\(\operatorname{Lie}^3(W)\). Explicitly, in each component
with three distinct basis indices its single relation is
primitive, with coefficients \(1,1,1\) up to sign; repeated
indices give the usual free components.
If \(p^kx\in J\) for \(x\in W\widehat\otimes\overline{\mathcal E}\), its
image in \(W\widehat\otimes(\overline{\mathcal E}/N^{\rm sat})\) vanishes.
Thus \(p^bx\in W\widehat\otimes N\), and saturation there
gives \(p^bx\in J\). This proves \eqref{eq:class3-pbound}.
This bounds the exponent of the entire compact $p$-torsion subgroup.

We next prove that the target length-three map is
injective on \(\mathcal S[1/p]\).
No integral \(\Lambda\)-basis of \(W\) is required.
Choose the established \(\Lambda[1/p]\)-basis of
\(W[1/p]\); finite generation gives a bounded lattice
\(p^a\Lambda^r\subseteq W\subseteq p^{-a}\Lambda^r\).
Every fixed completed tensor degree inherits a common
denominator bound. The integral presentation above used
only a compact \(\mathbf Z_p\)-basis.
Use the cofinite free rational target already chosen for
\(q_2\). The image of \(B\) starts in generator length two,
and the image of \(K_2\) starts in length three.
Therefore \([B,B]\) and \([H,K_2]\) vanish at lengths at
most three, giving a map from \(H/J_4\).
Its restriction to \(\mathcal R\) is additive and has pure length
three.

Over a finite splitting field \(L/\mathbf Q_p\), choose
Jordan blocks for \(\gamma\) on
\(T'=(T/\pi(Z))[1/p]\).
For a block of eigenvalue \(\alpha_i\) and length \(\ell_i\),
choose a lift \(y_i\in\overline{\mathcal E}_L\) of its top vector and put
\[
 C_i=(\gamma-\alpha_i)^{\ell_i}y_i\in N_L.
\]
For each fixed \(\alpha\), the residues of the \(C_i\) with
\(\alpha_i=\alpha\) are linearly independent over \(L\) in
\(N_L/(\gamma-\alpha)N_L\). Otherwise a relation
\(\sum_i d_iC_i=(\gamma-\alpha)A\) would make
\[
 \sum_i d_i(\gamma-\alpha)^{\ell_i-1}y_i-A
\]
a \(\gamma\)-eigenvector in the \(\Lambda\)-torsionfree module
\(\overline{\mathcal E}_L\). Its projection is a linear combination of
independent Jordan socle vectors, so all \(d_i\) must vanish.

In bounded pair coordinates the target kernel of
\((\gamma-\alpha_i)^jy_i\), \(0\le j<\ell_i\), is
\[
 \frac{C_i(X,Y)}
 {((1+X)(1+Y)-\alpha_i)^{\ell_i-j}}.
\]
All pair evaluation denominators were excluded in the
cofinite target chosen for \(Q_2\).
The source \eqref{eq:class3-S}, rationally, is the direct sum
of \(\operatorname{Lie}^3(W)_L\) and the completed tensor
of \(W_L\) with these finitely many lifted Jordan vectors.

Suppose an element has zero target length-three image.
The moment identity, after clearing denominators, gives
an identity in three independent position variables.
Choose a largest nonzero mixed pole of order \(k\) at
the pair \((1,2)\), with root \(\alpha\), and use
\[
 V=(1+T_1)(1+T_2)-\alpha.
\]
Poles at \((2,3)\) are units in the \(V\)-series field:
their constant terms are
\(\alpha(1+T_3)/(1+T_1)-\alpha'\ne0\).
Different roots at \((1,2)\) are units as well.
The coefficient of \(V^{-k}\) is consequently
\[
 \sum_i F_i(T_3)
 C_i\!\left(T_1,\frac{\alpha}{1+T_1}-1\right)=0
\]
as an array over pair colors, for each singleton color.

Constant-field independence suffices here. Choose finitely
many coefficient-extraction functionals in \(T_1\) and
the pair colors giving an invertible matrix over \(L\)
on the finite list of residue vectors.
Apply them coefficientwise in the disjoint variable
\(T_3\). The resulting invertible constant matrix forces
all \(F_i(T_3)=0\), contradicting the chosen pole.
No independence over the pair-variable function field
has been assumed. Thus all mixed terms vanish.
The remaining pure odd triple is detected by the completed
triple moment map.

This proves injection of \(\mathcal S[1/p]\) into rational
independent-position coefficient arrays before evaluation.
On those arrays, a nonzero \(f\in\Lambda\) acts by the
nonzero multiplier
\[
 f((1+T_1)(1+T_2)(1+T_3)-1).
\]
Setting two variables to zero proves nonvanishing.
Hence \(\mathcal S[1/p]\) is \(\Lambda\)-torsionfree.

Choose a characteristic polynomial \(a(z)\) annihilating
\(\gamma\) on the finite free module \(D\oplus\pi(Z)\).
By \eqref{eq:class3-R}, \(a(\gamma)\mathcal R\subseteq \mathcal S\).
If \(r\in \mathcal R_{\rm tor}\), then \(a(\gamma)r\) is torsion
in \(\mathcal S\), so its rational image is zero. Equation
\eqref{eq:class3-pbound} gives \(p^ba(\gamma)r=0\).
Conversely this nonzero \(\Lambda\)-annihilator makes
\(r\) torsion. We have proved the explicit identity
\[
 \mathcal R_{\rm tor}
 =\ker\bigl(p^ba(\gamma):\mathcal R\longrightarrow \mathcal R\bigr).
\]
Thus \(\mathcal R_{\rm tor}\) is closed and arithmetic-stable.
The quotient \(\mathcal R/\mathcal R_{\rm tor}\) is torsionfree, and the
extension
\[
 1\longrightarrow \mathcal R/\mathcal R_{\rm tor}
 \longrightarrow Q_3\longrightarrow Q_2\longrightarrow1
\]
makes \(Q_3\) torsionfree as well.

The pro-$p$ outer kernel lies in \(K_2\).
For its image \(r\in \mathcal R\), the target length-three image
of \(a(\gamma)r\in \mathcal S\) is zero. The rational injection
shows this element is \(p\)-torsion, so
\(p^ba(\gamma)r=0\). Therefore \(r\in \mathcal R_{\rm tor}\), proving
outer-kernel containment in \(K_3\) and the claimed field
inclusions.

Finally the inclusion \(N\hookrightarrow\overline{\mathcal E}\) and
\eqref{eq:class3-S} give
\[
 \operatorname{Lie}^3_{\mathbf Z_p}(W)\hookrightarrow \mathcal S.
\]
Its triple coefficient module is \(\Lambda\)-torsionfree,
so it meets \(\mathcal R_{\rm tor}\) trivially and survives in
\(\ker(Q_3\to Q_2)\).
It is nonzero, since \(W\) has infinite \(\mathbf Z_p\)-rank.
Thus \(Q_3\) has class exactly three and the extension over
\(Q_2\) is strict.
\end{proof}


\subsection{Selected even components in the class-three quotient}
\label{subsec:selected-class-three}

Retain the canonical quotient $Q_3$, put $\bar E=\overline{\mathcal E}$, and write
$R_3=\mathcal R/\mathcal R_{\Lambda\text{-tor}}$ for its central
kernel over $Q_2$.  Let $B_3$ be the image of $B$ in $Q_3$.
It is abelian.  The exact sequence for $\mathcal R$ gives
$\mathcal R^+\simeq\pi(Z)$, which is $\Lambda$-torsion.
Consequently $R_3^+=0$, and projection induces an isomorphism
\[
 B_3^+\xrightarrow{\ \sim\ }\bar E.
\]
Denote its inverse by $s_c$.  Thus
$B_3=s_c(\bar E)\oplus R_3$ as compact abelian groups.
Define
\[
 D_0=\operatorname{im}
   (\mathcal R_{\Lambda\text{-tor}}^-\longrightarrow D).
\]
This records exactly the odd abelian subgroup killed in $Q_3$.

Let $\Sigma$ be a set of even exponents such that $X^{[e]}$ is
finite for every $e\in\Sigma$.  Put
\[
 E_\Sigma=\bigoplus_{e\in\Sigma}\bar E^{[e]},\qquad
 E_{\rm bad}=\bigoplus_{e\notin\Sigma}\bar E^{[e]},
\]
and use the same subscripts for
$N=\widehat\bigwedge^2 W$.  The one-component theorem gives
$Z^{[e]}=0$ for $e\in\Sigma$, and hence
\[
 0\longrightarrow N_\Sigma\longrightarrow E_\Sigma
   \longrightarrow T_\Sigma\longrightarrow0,
 \qquad T_\Sigma=\bigoplus_{e\in\Sigma}T_e.
\]
In particular $T_\Sigma$ is free over $\mathbb Z_p$.
Set
\[
 L_{\rm bad}=[W,s_c(E_{\rm bad})]\subset R_3,
 \qquad
 J_\Sigma=s_c(E_{\rm bad})+L_{\rm bad},
 \qquad Q_{3,\Sigma}=Q_3/J_\Sigma.
\]
The bracket denotes the image of the completed tensor product.
Both summands of $J_\Sigma$ are compact, so this subgroup is
closed.  Since $L_{\rm bad}$ is central, it is exactly the closed
normal subgroup generated by $s_c(E_{\rm bad})$.

\begin{proposition}[Selected class-three quotient]
\label{prop:selected-class-three}
The subgroup $J_\Sigma$ is independent of the chosen complex
conjugation and is stable under the full arithmetic Galois group.
The field defined by $Q_{3,\Sigma}$ is therefore Galois over
$\mathbb Q$ and contained in $\Omega_p^*$.  Its abelianization is
\[
 Q_{3,\Sigma}^{\rm ab}=(U/D_0)\oplus T_\Sigma.
\]
Writing $j$ for the Jacobi map, put
\[
 S_\Sigma=
 \frac{W\widehat\otimes E_\Sigma}
 {\operatorname{pr}_\Sigma j(\widehat\bigwedge^3W)},
 \qquad
 \mathfrak l_{3,\Sigma}=
 \frac{\widehat{\operatorname{Lie}}_3(W)}{[W,N_{\rm bad}]}.
\]
There are exact sequences
\[
 0\longrightarrow S_\Sigma\longrightarrow
 R_3/L_{\rm bad}\longrightarrow D/D_0\longrightarrow0,
 \qquad
 0\longrightarrow \mathfrak l_{3,\Sigma}\longrightarrow S_\Sigma
 \longrightarrow W\widehat\otimes T_\Sigma\longrightarrow0.
\]
The kernel of $Q_{3,\Sigma}\to C_\Sigma$ is the central
group $R_3/L_{\rm bad}$.  The ordinary lower-central graded
groups in degrees two and three are
\[
 \operatorname{gr}_2Q_{3,\Sigma}
   =N_\Sigma\oplus(W\widehat\otimes T_\Sigma),
 \qquad
 \operatorname{gr}_3Q_{3,\Sigma}=\mathfrak l_{3,\Sigma}.
\]
For nonempty $\Sigma$ in the present arithmetic setting,
$\mathfrak l_{3,\Sigma}\ne0$.  Thus this quotient has class exactly three
and strictly enlarges $C_\Sigma$.
\end{proposition}

\begin{proof}
For $q\in Q_3$ and $e\in\bar E$, conjugation on the abelian
group $B_3$ adds the central commutator $[q,s_c(e)]$.
Writing $c'=qcq^{-1}$, direct calculation gives
\[
 s_{c'}(e)=s_c(e)+[q,s_c(e)].
\]
Indeed $c'$ sends $s_c(e)+r$ to
$s_c(e)-r+2[q,s_c(e)]$.  For $e\in E_{\rm bad}$ the displayed
change belongs to $L_{\rm bad}$, so it does not change the normal
subgroup generated by the bad even section.

For an arbitrary arithmetic element $g$, the involutions
$gcg^{-1}$ and $c$ have the same image in the abelian quotient
$\Gamma\times\Delta$.  They are conjugate by an element of $H$:
this follows from conjugacy of the Sylow $2$-subgroups in every
finite quotient of $H\rtimes\langle c\rangle$, followed by
compactness of the sets of conjugators.  Inner $H$ acts trivially
on $\bar E$, so $g$ preserves $E_{\rm bad}$.  Uniqueness of the
even section and the formula above show $gJ_\Sigma g^{-1}
\subseteq J_\Sigma$; applying the same argument to $g^{-1}$
gives equality.  The detected-field assertion follows because
$Q_{3,\Sigma}$ is a quotient of the detected group $Q_3$.

Let $A=W\widehat\otimes\bar E$ and
$J=j(\widehat\bigwedge^3W)$.  The class-three relation module
is $S=A/J$, and its image in $Q_3$ is
$S_3=A/J^{\rm sat}$, where $J^{\rm sat}$ is the $p$-saturation
of $J$ in $A$.  This is because
$S\cap\mathcal R_{\Lambda\text{-tor}}=S[p^\infty]$.
We claim
\[
 \operatorname{pr}_\Sigma(J^{\rm sat})
                 =\operatorname{pr}_\Sigma(J).
\]
If $p^kx\in J$, the $p$-torsionfreeness of
$E_\Sigma/N_\Sigma=T_\Sigma$ first places
$\operatorname{pr}_\Sigma x$ in $W\widehat\otimes N_\Sigma$.
Choose a compact $\mathbb Z_p$-basis of $W$ in
$\Delta$-eigencharacters.  Each unordered triple of distinct
basis elements has one Jacobi row with coefficients $\pm1$ in
its own coordinate block.  Projection onto the selected inner
pair characters retains a subset of those entries.  Every
nonzero projected row is primitive; zero rows give no relation.
The blocks are disjoint, and repeated-index blocks have no
Jacobi row.  Thus the projected Jacobi image is $p$-saturated,
which proves the claim also for the compact product of blocks.
It gives $S_3/L_{\rm bad}=S_\Sigma$ and shows that the kernel
on pure odd triples is exactly $[W,N_{\rm bad}]$.

The canonical relation-module sequence gives
$R_3/S_3=D/D_0$, proving the first exact sequence.  Projection
$E_\Sigma\to T_\Sigma$ proves the second.  The image of
$K_3$ in $H^{\rm ab}$ is $D_0\oplus\pi(Z)$, and the new
normal subgroup adds precisely the omitted even abelian image.
This proves the stated abelianization.  No equality $D_0=D$
has been used.

The derived subgroup of $Q_3$ projects onto $N$ with kernel
$S_3$, since $K_2\cap H'=[H,B]$.  Its $c$-eigenspaces are
therefore $s_c(N)$ and $S_3$.  Quotienting by $J_\Sigma$ gives
$s_c(N_\Sigma)\oplus S_\Sigma$.  The latter summand is central,
so the next lower-central group is precisely $\mathfrak l_{3,\Sigma}$.
The two graded descriptions follow.  Their $c$-splitting is
intrinsic because inner conjugation acts trivially on the
lower-central graded groups.

Finally choose $e\in\Sigma$ and odd characters $a,b$ with
$a+b=e$.  Every odd character of $W$ has positive
$\Lambda$-rank.  Choose distinct compact basis vectors
$x\in W^{[a]}$, $y\in W^{[b]}$, using two independent vectors
in the same eigenspace when $a=b$.  Projection of the chosen compact basis onto $\mathbb Z_p x\oplus\mathbb Z_p y$
extends to a homomorphism of completed free Lie algebras.
Every wedge in $N_{\rm bad}$ maps to zero, since the only
nonzero wedge in the target has character $a+b=e$.
Thus the map factors through $\mathfrak l_{3,\Sigma}$, and sends
$[x,[x,y]]$ to the nonzero Hall word with that same expression
in the free Lie algebra on two generators. This proves strictness.
\end{proof}

For $\Sigma_{30}(p)=\{-2m\bmod(p-1):1\le m\le30\}$,
one may take two independent $x,y\in W^{[-1]}$:
$[x,[x,y]]$ survives because $-2\in\Sigma_{30}(p)$.
This gives an explicit strict enlargement of the selected
class-two field at every odd prime.  The selected quotient's
odd abelian factor is the exact group $U/D_0$ above; its
central extension can retain more of $D$ than the class-two
quotient.

\begin{corollary}[Class-three component detection]
\label{thm:class-three-components}
Let $p$ be odd and let $\Sigma$ be any set of even exponents
such that $X^{[e]}$ is finite for every $e\in\Sigma$.
The quotient $Q_{3,\Sigma}$ is independent of the choice of
complex conjugation. Its defining subgroup is normal over
$\mathbb Q$ and contains the outer-action kernel in $H$.
Consequently its field $L_\Sigma^{(3)}$ satisfies
\[
 L_\Sigma^{(2)}\subseteq L_\Sigma^{(3)}\subseteq\Omega_p^*,
 \qquad \operatorname{Gal}(L_\Sigma^{(3)}/F_\infty)=Q_{3,\Sigma}.
\]
Put $T_\Sigma=\bigoplus_{e\in\Sigma}T_e$.
Its abelianization and ordinary lower-central graded groups are
\begin{align*}
 Q_{3,\Sigma}^{\rm ab}&=(U/D_0)\oplus T_\Sigma,\\
 \operatorname{gr}_2Q_{3,\Sigma}&=N_\Sigma\oplus(W\widehat\otimes T_\Sigma),\\
 \operatorname{gr}_3Q_{3,\Sigma}&=
       \operatorname{Lie}^3(W)/[W,N_{\rm bad}].
\end{align*}
For nonempty $\Sigma$, the last module is nonzero, the first
field inclusion is strict, and $Q_{3,\Sigma}$ has class exactly three.
\end{corollary}
The assertions follow from Proposition~\ref{prop:selected-class-three}.

\subsection{Divisibility on a total-product divisor}
\label{subsec:degree-three-saturation}

We strengthen the coefficient injection by computing
its quotient along a total-product divisor. This controls
the residues of the odd abelian components that first
reappear in the class-three extension.
Take a finite extension $L/\mathbb Q_p$ splitting the
characteristic polynomial on $(T/\pi Z)[1/p]$ and write
$\mathcal S_L=\mathcal S\otimes_{\mathbb Z_p}L$.
Tensor coefficients below have a common bounded
denominator in $L$. Put $z_i=1+X_i$ and
\[
 \mathcal A_3=O_L[[X_1,X_2,X_3]][1/p],\qquad
 \mathcal B_3=\mathcal A_3
       [q(z_1z_2)^{-1},q(z_1z_3)^{-1},q(z_2z_3)^{-1}],
\]
where $q$ is a characteristic polynomial for the finite
rational module $(T/\pi Z)[1/p]$.
The preceding proof constructs an equivariant injection
$\iota_3:\mathcal S_L\to\mathcal B_3^{\oplus d}$,
with $d$ the finite number of ordered color coordinates.

\begin{proposition}\label{prop:degree-three-saturation}
For $\beta\in1+\mathfrak m_L$ and every $k\geq1$,
\[
 \iota_3(\mathcal S_L)\cap
 (z_1z_2z_3-\beta)^k\mathcal B_3^{\oplus d}
       =\iota_3((\gamma-\beta)^k\mathcal S_L).
\]
Thus the induced map on $(\gamma-\beta)$-coinvariants
is injective.
\end{proposition}
\begin{proof}
We first give the exact evaluation map that keeps
divisibility inside the source. For a compact
$O_L$-module $E$ with continuous pro-$p$ action of
$\Gamma$, set $A=\beta\gamma_E^{-1}-1$.
This operator tends to zero in powers on every finite
invariant quotient: modulo the maximal ideal,
$\gamma_E-1$ is nilpotent for the finite $p$-group action.
Consequently the evaluation map is convergent and gives
\begin{equation}\label{eq:total-evaluation}
 \frac{O_L[[X_1]]\widehat\otimes E}
 {(z_1\gamma_E-\beta)(O_L[[X_1]]\widehat\otimes E)}
 \xrightarrow{\ \sim\ }E,\qquad
        \sum_{n\geq0}e_nX_1^n\longmapsto\sum_{n\geq0}A^ne_n.
\end{equation}
Here is the exact kernel calculation. Subtract the value
of the series and divide each $X_1^n-A^n$ by $X_1-A$.
The coefficient of $X_1^j$ in the quotient is the
convergent sum $\sum_{n\geq j+1}A^{n-1-j}e_n$.
This constructs a preimage in the same compact module.
Multiplication by the automorphism $\gamma_E$ identifies
the ideal generated by $X_1-A$ with that in
\eqref{eq:total-evaluation}.
The construction remains valid after inverting $p$.

Over $O_L[[X]][1/p]$, the module $W_L$ is finite free.
The bounded lattice comparison in Section~\ref{sec:abelian}
therefore permits a finite direct sum of
\eqref{eq:total-evaluation} for the first-slot colors.
Apply it both to $N$ and to $\bar E$.
The inclusion $\bar E_L\hookrightarrow N_L[1/q(\gamma)]$
is $e\mapsto q(\gamma)^{-1}q(\gamma)e$ and is compatible
with the evaluation maps. On coefficients the evaluation is
\[
                         z_1=\frac{\beta}{z_2z_3}.
\]
It follows in particular that a class from
$(W\widehat\otimes\bar E)_L$ whose evaluated tensor
lies in the regular $N_L$ coefficient module has a
representative in $(W\widehat\otimes N)_L$.
This is exactly the intersection $\bar E_L\cap N_L=N_L$
inside $N_L[1/q(\gamma)]$, transported by
\eqref{eq:total-evaluation}.

Put $D_\beta=z_1z_2z_3-\beta$.
The quotient $\mathcal A_3/(D_\beta)$ is
$O_L[[X_1,X_2]][1/p]$, by the integral change of
variable $z_3=\beta/(z_1z_2)$. For a root $\alpha$ of $q$,
the three pair divisors reduce, up to units, to
\[
 z_1-\beta/\alpha,\qquad z_2-\beta/\alpha,\qquad
                                  z_1z_2-\alpha.
\]
All roots lie in $1+\mathfrak m_L$ by continuity of the
pro-$p$ action on the finite free lattice.
These are prime factors in the bounded power-series
ring, and the three displayed families have no common
prime factors with one another. For example substituting
$z_1=\beta/\alpha'$ into the third gives a nonzero linear
polynomial in the independent variable $z_2$.
The ring is factorial, being a localization of the
regular local ring $O_L[[X_1,X_2]]$.
The intersection of its localizations in the separate
three factor families is therefore the unlocalized ring.

Represent $s\in\mathcal S_L$ by
$F\in(W\widehat\otimes\bar E)_L$.
Its tensor is alternating in the last two slots and
has poles only along $q(z_2z_3)$.
The bracket map is $F\mapsto F-\tau F$, where
$\tau$ is the cyclic permutation of three slots.
Suppose that $\iota_3(s)$ is divisible by $D_\beta$.
Modulo $D_\beta$, this says $\bar F=\tau\bar F$.
Thus $\bar F$ belongs to all three separate pair-divisor
localizations and is regular by the intersection just proved.
Equation~\eqref{eq:total-evaluation} supplies
$G\in(W\widehat\otimes N)_L$ such that
\[
 F-G\in(\gamma-\beta)(W\widehat\otimes\bar E)_L.
\]
The reduced tensor $\bar G$ is cyclically invariant and
alternating in its last two slots, hence fully alternating.
Apply the alternating idempotent in $L[S_3]$ to $G$.
The total-product divisor is symmetric in all three
positions, so the resulting alternating tensor has the
same reduction. Its denominator $6$ is one fixed rational
scalar and preserves boundedness. Fully alternating
tensors are exactly the images of
$J(\widehat{\bigwedge}^3W)_L$ in
\eqref{eq:class3-S}. Consequently
\[
 F\in J(\widehat{\bigwedge}^3W)_L
             +(\gamma-\beta)(W\widehat\otimes\bar E)_L,
\]
which proves $s\in(\gamma-\beta)\mathcal S_L$.
The converse is equivariance. Cancellation of one
$D_\beta$ in the domain $\mathcal B_3$, followed by
iteration, proves the equality for every $k$.
\end{proof}

For clarity, consider the resulting odd companions.
Put $\overline{\mathcal R}=
(\mathcal R/\mathcal R_{\Lambda\text{-tor}})\otimes L$.
Its quotient by $\mathcal S_L$ is $(D/D_0)\otimes L$.
Choose top Jordan vectors in this finite quotient,
lift them to $y_i\in\overline{\mathcal R}$, and put
$C_i'=(\gamma-\beta_i)^{r_i}y_i\in\mathcal S_L$,
where $r_i$ is the Jordan block length.
For each fixed $\beta$, their residues in
$\mathcal S_L/(\gamma-\beta)\mathcal S_L$ are independent:
a nonzero linear combination divisible by $\gamma-\beta$
would, after subtracting a source preimage, give a nonzero
$\beta$-eigenvector in the torsionfree module
$\overline{\mathcal R}$. Its projection to the independent
socle vectors of the Jordan blocks would be nonzero.
Proposition~\ref{prop:degree-three-saturation} transports
this independence to the three-position coefficient
residues. Thus these odd companions have nonzero
total-product residues in the stated coefficient map.

\section{A canonical detected quotient of class four}
\label{sec:canonical-class-four}

Retain the canonical quotients $Q_2$ and $Q_3$, and their
source kernels $K_2$ and $K_3$. Put $B=\ker(H\to W)$.
Define
\[
 J_4=\overline{[H,K_2][B,B]},\qquad
 J_5=\overline{[H,K_3][B,K_2]},\qquad
 R_4=K_3/J_5,\qquad S_4=J_4/J_5.
\]
Set
\[
 \mathcal R=K_2/J_4,\qquad R_3=K_2/K_3,\qquad
 \bar E=B/K_2,\qquad N=\widehat\bigwedge^2W.
\]
All commutator subgroups and tensor products in this
section are closed and compact, respectively.



\begin{theorem}\label{thm:canonical-class-four}
For every odd prime $p$, the $\Lambda$-torsion submodule
$(R_4)_{\rm tor}$ is closed. Its inverse image
$K_4^{\rm det}$ in $K_3$ is fully Galois-stable and
contains the kernel of the actual outer action. The
canonical quotient
\[
                         Q_4=H/K_4^{\rm det}
\]
is torsionfree of nilpotency class exactly four and
fits into a central extension
\[
 1\longrightarrow R_4/(R_4)_{\rm tor}
 \longrightarrow Q_4\longrightarrow Q_3\longrightarrow1.
\]
The natural map $\operatorname{Lie}^4_{\mathbf Z_p}(W)
\to R_4/(R_4)_{\rm tor}$ is injective. In particular,
\[
 \Omega_p^{K_3}\subsetneq\Omega_p^{K_4^{\rm det}}
                              \subseteq\Omega_p^*.
\]
The quotient is the compact image of $H$ in the full
rational outer image through generator length four.
\end{theorem}

\begin{proof}[Proof of Theorem~\ref{thm:canonical-class-four}]
We first calculate the source module integrally, then
prove its rational coefficient injection, and finally
identify its torsion by one fixed annihilator.

The established filtration
$\mathcal F_1=H,\mathcal F_2=B,\mathcal F_3=K_2,\mathcal F_4=K_3$ satisfies
$[\mathcal F_1,\mathcal F_2]\subset \mathcal F_3$ and
$[\mathcal F_1,\mathcal F_3][\mathcal F_2,\mathcal F_2]\subset \mathcal F_4$.
The three-subgroup identity gives
$[H,[B,B]]\subseteq\overline{[B,K_2]}$.
It follows that $J_5\subseteq J_4$, and that $R_4,S_4$
are central compact $\mathbf Z_p$-modules in $H/J_5$.
Their arithmetic actions factor through
$\Gamma\times\Delta$. The exact quotient is
\begin{equation}\label{eq:class4-old-torsion}
 0\longrightarrow S_4\longrightarrow R_4
     \longrightarrow \mathcal R_{\rm tor}\longrightarrow0.
\end{equation}
Here $\mathcal R_{\rm tor}$ is the entire old torsion module,
including its possibly infinite bounded $p$-torsion.

Let $\nu:\widehat\bigwedge^2W\to\bar E$ and
$\beta:W\widehat\otimes\bar E\to R_3$ be the established
bracket maps. Define
\[
 A_4=(W\widehat\otimes R_3)\oplus
                          \widehat\bigwedge^2\bar E
\]
and let $I_{112}$ be the image of
\begin{align*}
 j_{112}:\widehat\bigwedge^2W\widehat\otimes\bar E
          &\longrightarrow A_4,\\
 (x\wedge y)\otimes e&\longmapsto
 \bigl(x\otimes\beta(y,e)-y\otimes\beta(x,e),
                         -\nu(x\wedge y)\wedge e\bigr).
\end{align*}
The exact integral presentation is
\begin{equation}\label{eq:class4-presentation}
                         S_4=A_4/I_{112}.
\end{equation}
Indeed, the first summand maps by the commutator of an
$H$-lift with a $K_2$-lift. Changing the lifts contributes
$[B,K_2]$ or $[H,K_3]$. The second summand maps by the
commutator of two $B$-lifts; changing a lift contributes
$[B,K_2]$. Centrality gives bilinearity and alternation,
and Hall--Witt gives the displayed relation.

We prove exactness by integral continuous homology.
Write $B_3=B/K_3$. The established class-three structure
gives an abelian group
$B_3=\bar E^+\oplus R_3^-$, with $R_3$ central, and
the $W$-action is
\[
 w\cdot e=e+\beta(w,e),\qquad w\cdot r=r.
\]
The Hopf formula for the free pro-$p$ presentation
$H\twoheadrightarrow Q_3$ identifies
\[
 H_2(Q_3,\mathbf Z_p)
             =(K_3\cap[H,H])/[H,K_3].
\]
Its map to $K_3/J_4=\mathcal R_{\rm tor}$ has image
$\ker(\mathcal R_{\rm tor}\to H^{\rm ab})$. This image is odd:
the kernel lies in the old commutator module $\mathcal S$,
which is odd. Therefore the even part of $H_2(Q_3)$
is represented entirely by $J_4$. The commutator
classes of $[B,K_2]$ are the image of the commuting-pair
map $B_3\widehat\wedge R_3\to H_2(Q_3)$.
Consequently the actual source quotient is
\begin{equation}\label{eq:class4-Hopf}
 S_4=\bigl(H_2(Q_3,\mathbf Z_p)/
             \operatorname{im}(B_3\widehat\wedge R_3)\bigr)^+.
\end{equation}

Apply Hochschild--Serre homology to
$1\to B_3\to Q_3\to W\to1$.
Put $U_4=W\widehat\otimes R_3$,
$V_4=\widehat\bigwedge^2\bar E$ and
$P_4=\widehat\bigwedge^2W\widehat\otimes\bar E$.
Write $\delta:P_4\to U_4$ and $\epsilon:P_4\to V_4$
for the two components of $j_{112}=(\delta,-\epsilon)$.
Use the completed Koszul complex with symmetric
parameters $z_i=w_i^{1/2}-w_i^{-1/2}$ for a compact
$\mathbf Z_p$-basis of $W$. They differ from $w_i-1$
by units and satisfy $c(z_i)=-z_i$. On $B_3$ the
parameter acts as $\beta_i$: if $w_i=1+\beta_i$ on
this module, then $\beta_i^2=0$ and the two half powers
give $z_i=\beta_i$. Thus its chain parity is literally
the exterior-degree parity times the coefficient
parity. The complex gives
\[
 E^2_{1,1}{}^+=U_4/\operatorname{im}\delta,
 \qquad E^2_{2,1}{}^+=\ker\delta.
\]
Indeed, every product of two augmentation operators
on $B_3$ is zero. The even one-chains are the
$W\widehat\otimes R_3$ cycles, and their boundaries
are exactly $\delta(P_4)$. The even two-chains identify
with $P_4$, with differential $\delta$; even three-chain
boundaries vanish because their coefficients lie in
the trivial $W$-module $R_3$.

After quotienting by the image of
$B_3\widehat\wedge R_3$, the even $(0,2)$ term is
$V_4$. Every $W$-coinvariant relation in
$\widehat\bigwedge^2B_3$ contains an $R_3$ factor,
so no further relation is introduced in $V_4$.
The transgression from $\ker\delta$ is $-\epsilon$;
on lifted commuting pairs this is exactly the
Hall--Witt formula already displayed. The only later
possible incoming differential is $d_3$ from a
subquotient of $\widehat\bigwedge^3W$, which is odd;
its even target makes it zero. Likewise the only
possible incoming differential to $E^2_{1,1}{}^+$
has odd source and is zero. Finally the transgression
on $E^2_{2,0}=\widehat\bigwedge^2W$ is $\nu$, whose
even target is $\bar E$ and which is injective.
Thus no $(2,0)$ piece survives.

It follows that \eqref{eq:class4-Hopf} has the exact
two-step filtration
\[
 0\longrightarrow V_4/\epsilon(\ker\delta)
 \longrightarrow S_4
 \longrightarrow U_4/\operatorname{im}\delta
 \longrightarrow0.
\]
The canonical surjection $A_4/I_{112}\to S_4$ has
exactly this filtration by the image of $V_4$.
Its maps on both pieces are the identity, represented
respectively by two $\bar E$ lifts and by a $W$ lift
with an $R_3$ lift. It is therefore an isomorphism.
The completed Koszul calculations follow from finite
coordinate complexes and their compact closed images;
the image of $j_{112}$ is itself compact and closed.
This proves \eqref{eq:class4-presentation} without
assuming primitivity of $N\subset\bar E$ and without
dividing by $3$, so it includes $p=3$.

We now give a uniform integral bound. Choose $b$ such
that $p^b(T/\pi Z)[p^\infty]=0$, as in the class-three
proof, and put
\[
 D_0=\operatorname{im}(\mathcal R_{\rm tor}^-\to D).
\]
Choose $c$ such that $p^c(D/D_0)[p^\infty]=0$.
The established exact sequence is
\[
 0\longrightarrow S_3^{\rm tf}\longrightarrow R_3
 \longrightarrow D/D_0\longrightarrow0,
 \qquad S_3^{\rm tf}=\mathcal S/\mathcal S[p^\infty].
\]
Choose a finite free $\mathbf Z_p$-module $C_2\subset\bar E$
lifting a direct summand complementary to the torsion
of $T/\pi Z$. Then
\[
 E_0=N\oplus C_2\subset\bar E,\qquad p^b\bar E\subset E_0.
\]
The free weighted Hall module
\[
 L_3(W,C_2)=\operatorname{Lie}^3(W)
                      \oplus(W\widehat\otimes C_2)
\]
injects into $S_3^{\rm tf}$, and its cokernel is killed
by $p^b$. To see injectivity, the Jacobi image is already
in $W\widehat\otimes N$; the resulting free Hall module
is torsionfree and meets the removed $p$-torsion trivially.
Choose a finite free $C_3\subset R_3$ lifting a complement
to the torsion of $D/D_0$. Thus
\[
 R_0=L_3(W,C_2)\oplus C_3\subset R_3,
 \qquad p^{b+c}R_3\subset R_0.
\]
These auxiliary lattices need not be $\Lambda$-stable.

Put
\[
 A_0=(W\widehat\otimes R_0)\oplus
                          \widehat\bigwedge^2E_0,
 \qquad d=\max\{b+c,2b\}.
\]
Then $A_0\subset A_4$ and $p^dA_4\subset A_0$.
The tensor inclusions are injective: compact torsionfree
$\mathbf Z_p$-modules are products of copies of
$\mathbf Z_p$, and tensoring in these coordinates
preserves the displayed inclusions. Let $I_0$ be the
Jacobi image with the third input in $E_0$. It lies in
$A_0$, and
\[
 A_0/I_0=L_4(W,C_2,C_3),\qquad p^bI_{112}\subset I_0,
\]
where $L_4$ is degree four of the free graded Lie algebra
on generators of weights $1,2,3$, respectively.
For the first equality, its lower pieces are the free
weighted pieces already displayed. The only new Jacobi
identities in degree four have type $(1,1,2)$, exactly
$I_0$. This is the universal free graded Lie algebra
truncated at degree four; a Hall basis proves its
integral torsionfreeness.

If $a\in A_4$ represents a $p$-torsion element of $S_4$,
then $p^ka\in I_{112}$ for some $k$. Since
$p^da\in A_0$ and $p^{k+b}p^da\in I_0$, saturation
of $I_0$ in $A_0$ implies $p^da\in I_0$. Therefore
\begin{equation}\label{eq:class4-pbound}
                           p^dS_4[p^\infty]=0.
\end{equation}
Furthermore $A_0/I_0\to S_4$ is injective: its kernel
is killed by $p^b$, while $A_0/I_0$ is torsionfree.
In particular
\begin{equation}\label{eq:class4-pure-Hall}
                  \operatorname{Lie}^4(W)\hookrightarrow S_4.
\end{equation}

It remains to prove rational detection of $S_4$.
Work over a finite splitting field $L/\mathbf Q_p$ and
write
\[
 T'=(T/\pi Z)_L,\qquad D'=(D/D_0)_L.
\]
Auxiliary $L$-linear sections give
\[
 \bar E_L=N_L\oplus T',\qquad
 R_{3,L}=L_3(W_L,T')\oplus D'.
\]
Consequently \eqref{eq:class4-presentation} gives
\begin{equation}\label{eq:class4-weighted-free}
                        S_{4,L}=L_4(W_L,T',D'),
\end{equation}
with weights $1,2,3$. These sections are used only for
the proof; the source quotient remains intrinsic.

The target throughout is the free Lie algebra of the
full rational outer image, with its own generator
length. Choose the established finite set of odd
$\Lambda[1/p]$-colors for $W[1/p]$. Coefficient arrays
are bounded power series in independent position
variables $z_i=1+X_i$, with the companion denominators
below inverted.

For a Jordan block of $T'$ with eigenvalue $\alpha$
and length $r$, take a top lift $y\in\bar E_L$ and put
$C=(\gamma-\alpha)^ry\in N_L$. The companions of
distinct blocks at the same root have independent
constant-field residues modulo $\gamma-\alpha$.
Otherwise a linear relation would produce a nonzero
$\Lambda_L$-torsion element in $\bar E_L$, whose image
is a nonzero combination of the independent Jordan
socles. The ordered pair kernels of the successive
Jordan vectors are therefore
\[
              \frac{C(z_1,z_2)}{(z_1z_2-\alpha)^{r-j}},
              \qquad 0\le j<r.
\]
Similarly, for a Jordan block of $D'$ with eigenvalue
$\beta$ and length $s$, a top lift $v\in R_{3,L}$ has
companion $C'=(\gamma-\beta)^sv\in S_{3,L}$.
The $\Lambda_L$-torsionfreeness of $R_{3,L}$ gives
independent constant-field companion residues.
Proposition~\ref{prop:degree-three-saturation} transports
this independence to their actual coefficient residues
on $z_1z_2z_3=\beta$. The ordered triple kernels are
\[
        \frac{C'(z_1,z_2,z_3)}{(z_1z_2z_3-\beta)^{s-j}},
                         \qquad0\le j<s.
\]
The numerator may have internal pair poles, which are
coprime to the displayed triple divisor.
Lemma~\ref{lem:positive-eigenvalues} ensures that these
denominators do not vanish anywhere on the positive
odd-weight moment grid. Thus the formulas apply to
the actual target without discarding any odd generator.

Here is an explicit injectivity argument, which also
avoids any assumption about total degree-four
divisibility. Embed the free weighted Lie algebra
in its free associative tensor algebra and induct on
total degree. Every source word starts with a $W$,
$T'$ or $D'$ token. At the prefix divisor
$z_1z_2z_3-\beta$, only an initial $D'$ token can have
a pole. A later triple token has a different divisor;
all internal pair divisors are also distinct. At the
largest pole order the residue is a finite sum
\[
                   \sum_i\overline C'_i\,G_i,
\]
where the $G_i$ depend only on the disjoint tail
variables. The finite list $\overline C'_i$ is
independent over $L$. Clear its internal pair
denominators and take finitely many color and
coefficient functionals giving an invertible matrix
over $L$. Applying this matrix coefficientwise to the
tail variables forces every $G_i$ to vanish. The
induction hypothesis makes each corresponding source
tail zero. Descend through the finite list of pole
orders and roots to remove all initial $D'$ terms.

Next use the prefix divisor $z_1z_2-\alpha$. Only an
initial $T'$ token can now have this pole, and the same
residue argument removes all such terms. The remaining
terms begin with $W$ and are detected by the ordinary
first-slot coefficient map tensored with the already
injective tail map. The induction starts in degree
one. It proves injectivity in every fixed degree of
this auxiliary free weighted algebra, in particular
in degree four of \eqref{eq:class4-weighted-free}.
The argument extends residue independence only to
disjoint tail variables, never to the function field
of the variables inside the companion.

A zero actual target coefficient first gives zero on
the odd moment grid. Clearing the finitely many
companion denominators and using bounded power-series
uniqueness in each variable gives the coefficient
identity just used. We have proved
\begin{equation}\label{eq:class4-rational-injection}
 S_4[1/p]\hookrightarrow
                     \text{target generator length four}.
\end{equation}
It also follows that $S_4[1/p]$ is $\Lambda$-torsionfree:
a nonzero $f(\gamma)\in\Lambda$ acts on each ordered
coefficient array by the nonzero multiplier
$f(z_1z_2z_3z_4)$ in a domain.

The preceding class-three construction gives
$\mathcal R_{\rm tor}=\ker(p^ba(\gamma):\mathcal R\to\mathcal R)$, where $a$
is the characteristic polynomial on $D\oplus\pi Z$.
By \eqref{eq:class4-old-torsion},
$p^ba(\gamma)R_4\subseteq S_4$. Equations
\eqref{eq:class4-pbound} and
\eqref{eq:class4-rational-injection} now give the exact
closed-submodule formula
\begin{equation}\label{eq:class4-closed-torsion}
 (R_4)_{\rm tor}
       =\ker\bigl(p^{b+d}a(\gamma):R_4\to R_4\bigr).
\end{equation}
Indeed, the $p^ba(\gamma)$-multiple of a torsion element
lies in $S_4$ and is rationally zero, hence is killed
by $p^d$. The converse is immediate from the displayed
nonzero annihilator. The formula proves closedness;
the torsion property and the source construction prove
full arithmetic invariance.

The exact class-three target kernel is $K_3$:
the prior detection argument identifies its possible
central ambiguity as old $\Lambda$-torsion, and
Lemma~\ref{lem:positive-eigenvalues} makes that torsion
zero on every actual positive-weight coordinate.
Consequently $J_5$ has target length at least five,
and the map on $R_4$ takes values in pure length four.
An actual outer-kernel element has image $r\in R_4$
with zero coefficient image. Its multiple
$p^ba(\gamma)r\in S_4$ is zero rationally by
\eqref{eq:class4-rational-injection}, and
\eqref{eq:class4-pbound} gives
$p^{b+d}a(\gamma)r=0$. Formula
\eqref{eq:class4-closed-torsion} proves the asserted
kernel containment.

Conversely the roots of $a$ are eigenvalues of
$D\oplus T$, so Lemma~\ref{lem:positive-eigenvalues}
shows that $a(\gamma)$ is nonzero on every positive
target coordinate of length four. Thus
$(R_4)_{\rm tor}$ itself has zero target image.
This proves the exact compact-image assertion.
The displayed central extension is torsionfree since
$Q_3$ is torsionfree and the central kernel has no
$p$-torsion. The pure Hall submodule
\eqref{eq:class4-pure-Hall} is $\Lambda$-torsionfree by
its four-variable coefficient injection, so it meets
$(R_4)_{\rm tor}$ trivially. Finally
$\gamma_4H\subseteq K_3$ and $\gamma_5H\subseteq J_5$,
while $[x,[x,[x,y]]]$ is a nonzero surviving Hall word
for independent odd coordinates $x,y$. This gives
class exactly four and strictness over $Q_3$.
\end{proof}

The exact abelianization records any even directions
restored at this stage. Since $S_4$ is even, the odd
part of $R_4$ is entirely $\Lambda$-torsion and maps
onto $D_0$ in $H^{\rm ab}$. Put
\[
 Z_0=\operatorname{im}\bigl((R_4)_{\rm tor}^+
                                     \longrightarrow\pi Z\bigr).
\]
Then the abelian image of $K_4^{\rm det}$ is exactly
$D_0\oplus Z_0$, and hence
\[
                       Q_4^{\rm ab}=(U/D_0)\oplus(T/Z_0).
\]
This identity is induced by the actual abelianization
map of the source; no vanishing of $Z_0$ is asserted.

\subsection{Selected components in the class-four quotient}
\label{subsec:selected-class-four}

Let $\Sigma$ be a set of even characters such that
$X^{[e]}$ is finite for each $e\in\Sigma$. In the
canonical quotient $Q_4$, write $B_4$ for the image
of $B=\ker(H\to W)$. Its commutator subgroup is
central, so $B_4$ has class at most two. Since $p$
is odd, its class-two BCH Lie algebra
$\mathfrak b_4=\log B_4$ is defined integrally.
Choose a tame $\Delta$ complement and put
\[
 J_\Sigma^{(4)}=
 \left\langle\!\left\langle
 \exp\!\left(\bigoplus_{\substack{e\text{ even}\\e\notin\Sigma}}
                         \mathfrak b_4^{[e]}\right)
 \right\rangle\!\right\rangle_{Q_4},\qquad
 Q_{4,\Sigma}=Q_4/J_\Sigma^{(4)},
\]
where the normal closure is closed. This definition
includes the omitted even central directions of
relative weight four.

\begin{proposition}\label{prop:selected-class-four}
The subgroup $J_\Sigma^{(4)}$ is independent of the
tame complement and is fully Galois-stable. The
canonical quotient $Q_{4,\Sigma}$ is detected by the
outer action, maps onto $Q_{3,\Sigma}$, and has
\[
             Q_{4,\Sigma}^{\rm ab}=(U/D_0)\oplus T_\Sigma.
\]
If $-2,-4\in\Sigma$, the quotient has nilpotency class
exactly four. In particular this holds for $\Sigma_{30}$,
and its detected field strictly contains the field
of $Q_{3,\Sigma_{30}}$.
\end{proposition}

\begin{proof}
For an arithmetic element $g$, the tame complements
$\Delta$ and $g\Delta g^{-1}$ in the same pro-$p$
extension are conjugate by some $h\in H$.
The element $h^{-1}g$ centralizes $\Delta$ pointwise:
it normalizes it and acts trivially on its image in
the abelian quotient $\operatorname{Gal}(F_\infty/\mathbf Q)$.
It therefore preserves every isotypic submodule of
$\mathfrak b_4$. Conjugation by $h$ preserves the
$Q_4$-normal closure of the displayed generators.
This proves full arithmetic invariance; the same
argument compares two tame complements. The image
in $Q_3$ of the even part of $\mathfrak b_4$ is exactly
$\bar E$, so the induced quotient maps onto the
previous $Q_{3,\Sigma}$. Detection follows from the
surjection $Q_4\twoheadrightarrow Q_{4,\Sigma}$.

The map $\mathfrak b_4\to Q_4^{\rm ab}$ is linear,
because its BCH correction is a commutator.
Its even image is exactly $T/Z_0$, and it commutes
with the integral $\Delta$ idempotents. Thus the
abelian image of $J_\Sigma^{(4)}$ is precisely the
omitted even summand of $T/Z_0$; taking its normal
closure adds no abelian directions. For a selected
character the single-component theorem gives $Z_e=0$,
and hence $(Z_0)_e=0$. The exact abelianization
$Q_4^{\rm ab}=(U/D_0)\oplus(T/Z_0)$ proves the assertion.

For strictness, take distinct positive odd weights
$m,n\geq3$ with $m\equiv n\equiv-1\pmod{p-1}$.
Theorem~\ref{thm:canonical-class-four} embeds $Q_4$
as its compact image in the full rational free Lie
target through generator length four. Choose
$\Delta$-eigen free generators in that target. Such
generators are obtained by applying the appropriate
idempotent to a lift of each linear generator; they
have unchanged linear terms. The resulting substitution
is a filtered automorphism, with its inverse constructed
successively in lengths two, three and four.

There is therefore an exact $\Delta$-equivariant Lie
homomorphism to the free class-four Lie algebra on
$x,y$, sending the generators of weights $m,n$ to
$x,y$ and every other generator to zero. Give $x,y$
the character $\omega^{-1}$. Exponentiation and the
actual target embedding give a group homomorphism
\[
       \Psi:Q_4\longrightarrow
           \exp\bigl(\operatorname{Lie}(x,y)/\gamma_5\bigr).
\]
The image of $B_4$ has zero linear term. Its even
Lie part consequently has only the characters $-2$
in length two and $-4$ in length four. Since these
characters are selected, all omitted even isotypes
map to zero, and $\Psi$ factors through $Q_{4,\Sigma}$.

Choose an integral $w\in W^{[-1]}$ whose moment
coordinates at $m,n$ are nonzero. A rational
$\Lambda$-basis vector, after clearing a bounded
denominator, supplies such a $w$. The linear images
of $w,\gamma w$ have determinant
\[
                   \lambda_m\lambda_n(u^n-u^m)\ne0.
\]
For actual lifts $a,b\in H$, denote these independent
linear images by $A,B$. The length-four image of
$[a,[a,[a,b]]]$ is $[A,[A,[A,B]]]$, which is nonzero
in the free Lie algebra on the independent pair
$A,B$. It survives in $Q_{4,\Sigma}$, proving class
four and strictness over its class-three quotient.
\end{proof}

\begin{corollary}[Class-four component detection]
\label{thm:class-four-components}
Let $p$ be odd and let $\Sigma$ be any set of even exponents
such that $X^{[e]}$ is finite for every $e\in\Sigma$.
The subgroup $J_\Sigma^{(4)}$ is independent of the tame
complement and stable under the full arithmetic Galois group.
The canonical field $L_\Sigma^{(4)}$ defined by $Q_{4,\Sigma}$
is Galois over $\mathbb Q$, and
\[
 L_\Sigma^{(2)}\subseteq L_\Sigma^{(3)}
   \subseteq L_\Sigma^{(4)}\subseteq\Omega_p^*,\qquad
 Q_{4,\Sigma}^{\rm ab}=(U/D_0)\oplus T_\Sigma.
\]
Here $L_\Sigma^{(3)}$ is the class-three field of
Corollary~\ref{thm:class-three-components}.
For $-2,-4\in\Sigma$, the quotient $Q_{4,\Sigma}$ has class
exactly four and $L_\Sigma^{(3)}\subsetneq L_\Sigma^{(4)}$.
\end{corollary}
This follows from Proposition~\ref{prop:selected-class-four}.

\section{Total-product divisibility in the pair-companion algebra}
\label{sec:pair-total-saturation}

This section concerns the precise algebra generated by
$W$ and the even quotient $T'=(T/\pi Z)[1/p]$. Its
source and coefficient maps are defined below. The quotient maps from $H$ to $Q_3$ and $Q_4$ were
constructed in the preceding sections. Here the source
is the displayed free weighted algebra in every degree.

Choose a finite extension $L/\mathbf Q_p$ splitting
the characteristic polynomial of $T'$, and put
$R=O_L[[\gamma-1]][1/p]$. The established rational
module $W_L$ is finite free over $R$. Choose an
$L$-linear section of
\[
 0\longrightarrow N_L\longrightarrow\bar E_L
                  \longrightarrow T'_L\longrightarrow0,
 \qquad N_L=\widehat\bigwedge^2W_L.
\]
Let $A$ be the completed free associative algebra on
$W_L$ in weight one and $T'_L$ in weight two. The
action on $\bar E_L$ defines the Hopf
automorphism
\[
        \gamma t=A_Tt+C_t,\qquad C_t\in N_L,
\]
together with the regular action on $W_L$. All
completed coefficient modules at a fixed degree have
a common bounded denominator in $L$; they are obtained
by taking the indicated compact $O_L$-modules and then
inverting $p$, rather than allowing a separate
denominator for every coordinate.

For a Jordan block with root $\alpha$, length $r$ and
top lift $y$, put $C=(\gamma-\alpha)^ry\in N_L$.
The pair-companion substitution is
\[
 (\gamma-\alpha)^jy\longmapsto
       \frac{C(z_i,z_j)}{(z_iz_j-\alpha)^{r-j}},
                         \qquad0\leq j<r.
\]
The repeated indices on the right denote the two
positions of the token, not its Jordan index.
For a fixed root, the companion residues are
independent over $L$: a relation would give a
$\Lambda_L$-torsion element in $\bar E_L$ with a
nonzero Jordan-socle image in $T'_L$.
Write $q$ for a polynomial annihilating $T'_L$.
In weighted degree $n$ the substitution defines
\[
 \iota_n:A_n\longrightarrow\mathcal B_n^{\oplus c_n},
 \qquad
 \mathcal B_n=O_L[[X_1,\ldots,X_n]][1/p,
                      q(z_iz_j)^{-1}:i<j],\quad z_i=1+X_i,
\]
where the finite array indexes the ordered colors.
It is injective: extract the highest pole of an
initial two-position token, use the independent
companion residues with disjoint tail variables, and
induct on the tail degree; the remaining initial
one-position terms are ordinary coefficient tensors.

\begin{theorem}\label{thm:pair-total-saturation}
For $\beta\in1+\mathfrak m_L$, put
$D_\beta=z_1\cdots z_n-\beta$. For $n=1$ and every
$n\geq3$, and for every $k\geq1$,
\[
 \iota_n(A_n)\cap D_\beta^k\mathcal B_n^{\oplus c_n}
                   =\iota_n((\gamma-\beta)^kA_n).
\]
The same equality holds on the free weighted Lie
subalgebra of $A$.
\end{theorem}

The exclusion of degree two is necessary. At a root
of $q$ its total divisor is already inverted in the
target, whereas the source retains nonzero
coinvariants. The induction below starts with the
separately proved degrees three and four and never
uses that excluded case.

\begin{proof}
Degree one is the usual exact coefficient division.
The initial-token argument below also proves degree
three directly: its total-tail divisor then involves
one variable, is distinct from the pair divisors of
the first-$W$ term, and uses only degree-one source
division. Thus it proves the associative base as
well as the Lie base of
Proposition~\ref{prop:degree-three-saturation}.
We give the additional degree-four argument in
detail, then the exact source induction.

Put $B=A_2$ and $V=W_L\widehat\otimes W_L$. Partition
the words of weight four by their middle cut. There
is an exact $\Gamma$-stable sequence
\begin{equation}\label{eq:pair-four-cut}
 0\longrightarrow B\widehat\otimes_L B
 \longrightarrow A_4
 \longrightarrow W_L\widehat\otimes_L T'_L
                    \widehat\otimes_L W_L\longrightarrow0.
\end{equation}
Its quotient consists of the single token pattern
$W,T',W$. Replacing its middle token by its quadratic
companion lands in the first term, so its induced
action is the diagonal linear action.

Suppose $\iota_4(a)$ vanishes modulo $D_\beta$. On
that hypersurface put $z_4=\beta/(z_1z_2z_3)$.
A middle-pair divisor $z_2z_3-\alpha$ is distinct
from both outer-pair divisor families: their pair
products are $z_1z_2$ and $\beta/(z_1z_2)$.
Consequently a term from $B\widehat\otimes B$ has
no principal part at that middle divisor.

The exact Jordan calculation in the quotient of
\eqref{eq:pair-four-cut} uses $u=z_1z_4$.
For a block of length $r$ with $N=\gamma_T-\alpha$,
the coinvariant relations are
\[
 (u\alpha-\beta)N^jy+uN^{j+1}y=0.
\]
They identify the block cokernel of $uA_T-\beta$
with the singleton coefficient module modulo
$(u\alpha-\beta)^r$, generated by $y$; explicitly
the coefficient vector $F_j(z_1,u)$ maps to
\begin{equation}\label{eq:pair-middle-Jordan}
 \sum_{j=0}^{r-1}F_j(z_1,u)(\beta/u-\alpha)^j
                         \pmod{(\beta/u-\alpha)^r}.
\end{equation}
In the target the middle product is $\beta/u$.
The companion residue depends on the internal
variable $z_2$, whereas the coefficients in
\eqref{eq:pair-middle-Jordan} have the disjoint
singleton variable $z_1$. A finite color and
coefficient minor of the independent residue vectors
detects the largest principal part. Descending through
the pole orders proves that every expression
\eqref{eq:pair-middle-Jordan} vanishes. Division uses
the coordinate $u-\beta/\alpha$ and preserves bounded
coefficients. Thus the quotient image of $a$ is an
$(\gamma-\beta)$-multiple. Lift and subtract
that multiple in $A_4$. We are left with an element
of $B\widehat\otimes B$ having zero total-divisor image.

We compute the required balanced completed tensor
exactly. Let $s=z_1z_2$ and
$\tau_\beta(s)=\beta/s$. The relation
$\gamma_1\gamma_2=\beta$ gives
\begin{equation}\label{eq:pair-balanced-tensor}
 \frac{B\widehat\otimes_L B}
 { (\gamma_1\gamma_2-\beta)(B\widehat\otimes_L B)}
       =B\widehat\otimes_{R,\tau_\beta}B.
\end{equation}
For the generator $\gamma$, this follows by multiplying
the displayed relation by the inverse of the second
$\gamma$. Continuous extension gives balancing for
every bounded series. The involution $\tau_\beta$ is
defined on $R$ because $\beta\in1+\mathfrak m_L$.

The regular pair module $V$ has coordinates
$x=z_1-1$ and $s=z_1z_2$ and is a finite color array of
$O_L[[s-1,x]][1/p]$. Thus
\[
 V\widehat\otimes_{R,\tau_\beta}V
       =O_L[[s-1,x,y]][1/p]^{\oplus c}
\]
with the appropriate finite colors. This is an exact
change of variables; its coefficient map is the
inclusion in the ring with $q(s)$ and $q(\beta/s)$
inverted. In particular its map into
\eqref{eq:pair-balanced-tensor} is injective.
Multiplication by
\[
                           h(s)=q(s)q(\beta/s)
\]
takes \eqref{eq:pair-balanced-tensor} into this regular
submodule. Hence the kernel of its coefficient map
is killed by $h$. It remains to check this kernel
at the finitely many factors of $h$.

Fix $f=s-\alpha$ after splitting these roots in $L$.
Invert the other polynomial factors. The first pair
module has the explicit form
\[
              B=V+\sum_i R f^{-r_i}C_i,
\]
where the sum ranges over the blocks at $\alpha$.
Their residues $C_i\bmod fV$ are independent over $L$.
Choose finitely many internal coefficient and color
coordinates whose matrix on these residues is
invertible. The corresponding matrix $P(s)$ has
determinant $g$ with $g(\alpha)\ne0$. After inverting
this one determinant, the finite-rank coordinate
transformation
\[
 (v_I,v_{\rm rest})\longmapsto
 \bigl(P^{-1}v_I,
               v_{\rm rest}-C_{\rm rest}P^{-1}v_I\bigr)
\]
sends the $C_i$ to distinct coordinate vectors.
Its inverse is the corresponding triangular map.
Both transformations have a fixed common denominator;
their remaining columns are bounded series. The module
$B$ is now a regular coordinate module with finitely
many coordinate lines replaced by $f^{-r_i}R$.

Apply the same transformation to the second factor,
with its twisted action. The balanced completed tensor
is the regular two-internal-variable coordinate module
with selected rows and columns rescaled by
$f^{-r_i-r'_j}$, or by a single one of these factors.
It embeds coordinatewise in the rational coefficient
arrays. This includes the complementary-pair collision:
when the two factors have a pole at the same $f$, their
exponents add. No exclusion of products of eigenvalues
is needed. Every transformation uses finitely many
columns and a fixed denominator, so it acts on the
bounded completion itself. No interchange of arbitrary
products with localization is being presumed.

The kernel therefore vanishes at every factor of $h$.
Since it is killed by $h$, the Bezout identity among
the distinct prime powers makes it zero globally.
This proves injectivity of
\eqref{eq:pair-balanced-tensor} into the specialized
coefficient ring. Returning to
\eqref{eq:pair-four-cut} proves degree-four saturation
for one total divisor. Repetition gives every power.

Now let $n\geq5$. The initial-token decomposition is
the exact $\Gamma$-stable sequence
\begin{equation}\label{eq:pair-initial-token}
 0\longrightarrow P_n=W_L\widehat\otimes A_{n-1}
 \longrightarrow A_n
 \longrightarrow T'_L\widehat\otimes A_{n-2}=Q_n
 \longrightarrow0.
\end{equation}
The right action is diagonal. It has no
$(\gamma-\beta)$-invariants: on each Jordan block
the determinant is a power of
$\alpha\gamma_{A_{n-2}}-\beta$, whose coefficient
action is a nonzero multiplier in a domain.
The adjugate then proves injectivity. Consequently
taking coinvariants in \eqref{eq:pair-initial-token}
is also exact at the left.

There is an exact source evaluation
\begin{equation}\label{eq:pair-source-evaluation}
 (W_L\widehat\otimes A_{n-1})_{\gamma-\beta}
                  \simeq\bigoplus_{\rm colors}A_{n-1},
                 \qquad z_1\longmapsto\beta\gamma_{A_{n-1}}^{-1}.
\end{equation}
Choose a $\Gamma$-stable compact lattice after one
bounded scaling in this fixed degree. The operator
$A=\beta\gamma^{-1}-1$ is topologically nilpotent on
every finite invariant quotient. Evaluation of
$\sum e_jX^j$ at $A$ converges. The identity
\[
 X^j-A^j=(X-A)\sum_{i=0}^{j-1}X^iA^{j-1-i}
\]
gives its exact kernel, with convergent quotient
coefficients. This proves
\eqref{eq:pair-source-evaluation} inside the source,
and bounded scalar extension preserves it.

Suppose $\iota_n(a)$ vanishes on $D_\beta=0$, and put
$s=z_3\cdots z_n$. The pair pole of an initial $T'$
token at $\alpha$ becomes the total-tail divisor
$s-\beta/\alpha$. Since $n-2\geq3$, this divisor is
distinct from every pair divisor in the first-$W$
term and from every proper pair denominator in its
own tail. At its largest pole, the companion residue
is $C(\alpha/z_2,z_2)$, depending only on $z_2$ and
its colors; its coefficient is the tail image on
$s=\beta/\alpha$. The finite constant residue minor
forces every such tail image to vanish. Saturation
in degree $n-2$ therefore puts the source
tail in $(\gamma_A-\beta/\alpha)A_{n-2}$.

The exact Jordan relation in $Q_n$ is
\begin{equation}\label{eq:pair-Jordan-reduction}
 N^jy\otimes(\gamma_A-\beta/\alpha)b
 \equiv-\alpha^{-1}N^{j+1}y\otimes\gamma_A b
                                  \pmod{(\gamma-\beta)Q_n}.
\end{equation}
It follows by expanding
$(\gamma_T\gamma_A-\beta)(N^jy\otimes b)$.
The right side has one lower pole, and is zero at
the end of the Jordan block. Lift each subtraction
to $A_n$; its additional quadratic terms lie in
$P_n$ and have no total-tail pole. Descending through
the finite pole orders and roots shows that the
class of $a$ in $(Q_n)_{\gamma-\beta}$ is zero.

Subtract an $(\gamma-\beta)$-multiple and use
\eqref{eq:pair-source-evaluation}. The remaining
class corresponds to an element of
$\bigoplus A_{n-1}$ whose ordinary coefficient image
in the independent variables $z_2,\ldots,z_n$ is zero.
The already established coefficient injection makes
it zero. This proves saturation at $n$ from saturation
at $n-2$. The bases $3,4$ prove every $n\geq3$,
without using degree two. Cancelling one total
divisor in the coefficient domain and repeating
proves every power $k$.

Finally use the Eulerian primitive projector
\cite[Section~1.1, p.~2]{BandieraSchaetz}
\[
 e=\log^\star(\mathrm{Id})
   =\sum_{j\geq1}\frac{(-1)^{j+1}}j
             (\mathrm{Id}-\eta\epsilon)^{\star j}.
\]
It is a finite sum on every positive weighted degree
of this connected graded cocommutative Hopf algebra,
has image the primitive subspace and is identity
there. Indeed, the coproduct of the identity map is
the convolution product of the maps $a\mapsto a\otimes1$
and $a\mapsto1\otimes a$; these commute by cocommutativity.
Its convolution logarithm therefore satisfies
$\Delta e(a)=e(a)\otimes1+1\otimes e(a)$ for positive
degree $a$. On a primitive element all convolution
powers above the first vanish, so $e(a)=a$.
The finite coefficient formulas extend to the bounded
completion. Naturality makes it commute with the actual
Hopf automorphism $\gamma$, including its quadratic
companion terms. If a primitive $a$ is
$(\gamma-\beta)^k b$ in the associative algebra,
then $a=(\gamma-\beta)^k e(b)$ in the Lie source.
This proves the asserted Lie equality. A projector
depending on ordinary token count would not justify
this step, since that count changes under $\gamma$.
\end{proof}

\section{A canonical detected quotient of class five}
\label{sec:canonical-class-five}

Retain $J_5$, $K_4^{\rm det}$ and the class-four quotient.
Put
\[
 J_6=\overline{[H,K_4^{\rm det}][B,K_3][K_2,K_2]},
 \qquad R_5=K_4^{\rm det}/J_6,
 \qquad S_5=J_5/J_6,
\]
let \(K_5^{\rm det}\) be the inverse image of \((R_5)_{\Lambda\text{-tor}}\) in \(K_4^{\rm det}\), and put \(Q_5=H/K_5^{\rm det}\). For \(2\leq j\leq5\), write \(L^{(j)}\) for the fixed field of the kernel of \(H\to Q_j\). We abbreviate
\[
 E_2=B/K_2=\bar E,
 \qquad E_3=K_2/K_3=R_3,
 \qquad E_4=K_3/K_4^{\rm det}.
\]
These have parities even, odd, even, respectively.
The three-subgroup identity proves $J_6\subseteq J_5$
and centrality of $R_5,S_5$ in $H/J_6$. The exact
quotient is
\begin{equation}\label{eq:class5-old-torsion}
 0\longrightarrow S_5\longrightarrow R_5
                  \longrightarrow(R_4)_{\rm tor}\longrightarrow0.
\end{equation}
It retains all of the preceding bounded $p$-torsion.



\begin{theorem}[Canonical detection through class five]\label{thm:canonical-class-five}
For every odd prime $p$, the $\Lambda$-torsion submodule
$(R_5)_{\rm tor}$ is closed. Its inverse image
$K_5^{\rm det}$ in $K_4^{\rm det}$ is fully
Galois-stable and contains the outer kernel.
The quotient $Q_5=H/K_5^{\rm det}$ is torsionfree of
nilpotency class exactly five and is the compact
image in the full rational outer target through
generator length five. It fits into the central
extension
\[
 1\longrightarrow R_5/(R_5)_{\rm tor}
 \longrightarrow Q_5\longrightarrow Q_4\longrightarrow1.
\]
The natural map
$\operatorname{Lie}^5(W)\to R_5/(R_5)_{\rm tor}$ is injective.
Together with its preceding central quotients it gives
\[
 L^{(2)}\subsetneq L^{(3)}\subsetneq L^{(4)}
                         \subsetneq L^{(5)}\subseteq\Omega_p^*.
\]
\end{theorem}

\begin{proof}[Proof of Theorem~\ref{thm:canonical-class-five}]
The central module is computed using a commutator
surjection that becomes an isomorphism after inverting $p$.
The commutators induce the surjection
\[
 A_5=(W\widehat\otimes E_4)\oplus
                              (E_2\widehat\otimes E_3)
                         \longrightarrow S_5.
\]
The three factors of $J_6$ give precisely the
independence of lifts in its two summands.
Centrality gives bilinearity. The Hall--Witt relations
of types $(1,1,3)$ and $(1,2,2)$ are
\begin{align*}
 j_{113}((x\wedge y)\otimes r)
   &=\bigl(x\otimes[y,r]-y\otimes[x,r],
                                  -\nu(x\wedge y)\otimes r\bigr),\\
 j_{122}(x\otimes(e\wedge f))
   &=\bigl(x\otimes[e,f],
                                  f\otimes[x,e]-e\otimes[x,f]\bigr).
\end{align*}
Here all lower bracket maps are those already defined
in $Q_4$. With the closed relation image, put
\begin{equation}\label{eq:class5-candidate}
 C_5=A_5/\operatorname{im}(j_{113},j_{122}).
\end{equation}
We have a canonical surjection $C_5\twoheadrightarrow S_5$,
and both modules are odd under $c$.

We first prove that this surjection is an isomorphism
after inverting $p$. There is a concrete reason to
make this passage explicitly: the odd Hopf quotient
at this stage can include the earlier bounded module
$S_3[p^\infty]$. It cannot be discarded integrally.
Since $\gamma_5H\subseteq K_4^{\rm det}$, we have
$\gamma_6H\subseteq J_6$. Thus the calculation takes
place in fixed ordinary nilpotency class five.
The rational BCH and inverse-substitution formulas
have only finitely many fixed denominators there.
To justify the bounded completion, use the Magnus
embedding in the completed tensor algebra on
$H^{\rm ab}$. The formulas for $\log(1+X)$, $\exp(X)$
and multiplication through ordinary degree five have
the same finite list of rational coefficients at
every finite rank. One integer clears these
coefficients independently of the number of basis
vectors. Passing to the inverse limit over finite
coordinate sets therefore gives the free rational
nilpotent Lie algebra on the compact module $M$,
with a common bound on every coefficient denominator.

The abelian images of $B,K_2,K_3,K_4^{\rm det}$
are, respectively,
\[
 D\oplus T,\quad D\oplus\pi Z,\quad
 D_0\oplus\pi Z,\quad D_0\oplus Z_0.
\]
Choose rational complements in this finite-dimensional
flag and compatible representatives, assigning weights
\[
 W_L:1,\quad (T/\pi Z)_L:2,\quad (D/D_0)_L:3,
 \quad(\pi Z/Z_0)_L:4,\quad(D_0\oplus Z_0)_L:5.
\]
For $W$ keep a compact $\mathbf Z_p$-basis and lifts.
The remaining choices involve finitely many vectors
and can be made after a fixed bounded scaling.
Substitution by these representatives in the rational
free nilpotent Lie algebra of $H$ is invertible:
its linear term is invertible, and its inverse is
constructed successively through the five ordinary
lower-central steps.

The proved presentations of $Q_2,Q_3,Q_4$ identify
their rational pieces with the free weighted pieces
through weights two, three, four. The map from the
free weighted truncation to $Q_4$, using the compatible
representatives, is therefore an isomorphism on every
graded piece and hence on the finite filtered object.
It identifies the maps to $Q_3,Q_2,W$ with the
corresponding truncations. In the rational free source,
the images of $H,B,K_2,K_3,K_4^{\rm det}$ are exactly
the ideals of weight at least $1,2,3,4,5$.

The image of $J_6$ is consequently all commutators of
weight at least six: every such bracket lies in one
of the three displayed factors. The image of $J_5$
is all commutators of weight at least five. Their
quotient is exactly
\[
 L_5(W_L,T',D',Z'),\qquad
 T'=(T/\pi Z)_L,\quad D'=(D/D_0)_L,\quad
 Z'=(\pi Z/Z_0)_L,
\]
with generator weights $1,2,3,4$.
The weight-five generators themselves do not occur
in this derived quotient. The rational candidate
\eqref{eq:class5-candidate} is the same free weighted
piece: its central bracket types are $(1,4),(2,3)$,
and its new Jacobi types are exactly $(1,1,3),(1,2,2)$.
All lower relations have already been imposed.
The commutator map therefore gives
\begin{equation}\label{eq:class5-rational-source}
              C_5[1/p]\xrightarrow{\ \sim\ }S_5[1/p].
\end{equation}
In particular its integral kernel is $p$-torsion.

We next bound that torsion uniformly. Retain $b,c$
from the class-four proof and put
$d_4=\max\{b+c,2b\}$. Choose $c_4$ killing the finite
torsion subgroup of $\pi Z/Z_0$. The earlier weighted
Hall lattices satisfy
\[
 E_{2,0}=N\oplus C_2\subset E_2,\quad p^bE_2\subset E_{2,0},
\]
\[
 E_{3,0}=L_3(W,C_2)\oplus C_3\subset E_3,
                      \quad p^{b+c}E_3\subset E_{3,0}.
\]
The exact sequence
$0\to S_4/S_4[p^\infty]\to E_4\to\pi Z/Z_0\to0$
and the class-four Hall comparison give a finite
free section $C_4$ and
\[
 E_{4,0}=L_4(W,C_2,C_3)\oplus C_4\subset E_4,
                 \qquad p^{d_4+c_4}E_4\subset E_{4,0}.
\]
These are compact $\mathbf Z_p$-lattices; no integral
$\Lambda$-freeness or $\Lambda$-stability is required.

Set
\[
 A_{5,0}=(W\widehat\otimes E_{4,0})\oplus
                          (E_{2,0}\widehat\otimes E_{3,0}),
 \qquad d_5=\max\{d_4+c_4,2b+c\}.
\]
Then $p^{d_5}A_5\subseteq A_{5,0}$. Let $I_0$ be the
two Jacobi images on these free weighted lattices.
The quotient $A_{5,0}/I_0=L_5(W,C_2,C_3,C_4)$ is a
free Hall module and is $p$-torsionfree. If $I$ is
the full relation image in \eqref{eq:class5-candidate},
then
\[
                              p^{d_4}I\subseteq I_0:
\]
the $(1,1,3)$ relation needs the factor $p^{b+c}$,
and the $(1,2,2)$ relation needs $p^{2b}$.
For $p^ka\in I$, the element $p^{d_5}a$ lies in
$A_{5,0}$ and has a $p$-power multiple in $I_0$.
Saturation of $I_0$ there gives $p^{d_5}a\in I_0$.
Thus $p^{d_5}C_5[p^\infty]=0$, and
$A_{5,0}/I_0\hookrightarrow C_5$ as well: its kernel
is killed by $p^{d_4}$ and is therefore zero.
By \eqref{eq:class5-rational-source}, every $p$-torsion
element of $S_5$ lifts to $p$-torsion in $C_5$.
We have proved
\begin{equation}\label{eq:class5-pbound}
 p^{d_5}S_5[p^\infty]=0,
 \qquad\operatorname{Lie}^5(W)\hookrightarrow S_5.
\end{equation}
This argument keeps the possible integral kernel of
$C_5\to S_5$ instead of asserting that it is zero.

For rational detection, the new finite quotient is
$Z'$ of weight four. Its companions lie in
$S_{4,L}=L_4(W_L,T',D')$. We verify the total-product
divisibility needed for their residues. Let $A^0$
be the pair algebra of
Theorem~\ref{thm:pair-total-saturation}, and let $A$
also contain the cubic generators $D'$. In degree
four there is an exact $\Gamma$-stable sequence
\[
 0\longrightarrow A^0_4\longrightarrow A_4
 \longrightarrow(W_L\widehat\otimes D')
                    \oplus(D'\widehat\otimes W_L)\longrightarrow0.
\]
The quotient action is diagonal, and it has no
$(\gamma-\chi)$-invariants by the regular singleton
factor and its finite Jordan determinant.
On $z_1z_2z_3z_4=\chi$, a prefix triple pole at $\beta$
becomes $z_4-\chi/\beta$; a suffix triple pole becomes
$z_1-\chi/\beta$. These singleton divisors are distinct.
Every pair denominator, including an internal
companion denominator, is a unit at each generic
singleton divisor, as is the other triple divisor.

The Jordan cokernel in a singleton variable $u$ is
that of $uA_D-\chi$, namely the coefficient module
modulo $(u\beta-\chi)^r$ for a block of length $r$.
Its kernel image has the successive principal parts
$C_3/(\chi/u-\beta)^{r-j}$. The already checked total
degree-three residue independence, with a finite
color and coefficient minor in the internal triple
variables, detects every singleton jet. A zero total
image therefore makes the quotient class a
$(\gamma-\chi)$-multiple. Subtract it in the source.
The remainder belongs to $A^0_4$, where
Theorem~\ref{thm:pair-total-saturation} supplies the
division. Adding proper triple denominators does
not change total-divisor membership for this
remainder: each is coprime to the prime total divisor.
Iteration gives every power, and the Eulerian
projector gives the same assertion on $S_{4,L}$.

The $\Lambda_L$-torsionfreeness of $E_{4,L}$ now makes
the $Z'$ Jordan companions independent in
$S_{4,L}/(\gamma-\alpha)S_{4,L}$; the preceding
divisibility transports this to target
residues. Their kernels have a full four-position
denominator and only proper pair or triple poles
in the numerator. To detect the free weighted
degree-five source, extract an initial four-position
pole first, then an initial triple pole, then an
initial pair pole, and finally a singleton. Every
residue coefficient involves only disjoint tail
variables; the finite constant extraction matrices
and induction give an injective associative map.
It restricts to the free Lie piece in
\eqref{eq:class5-rational-source}. This proves
\begin{equation}\label{eq:class5-rational-injection}
 S_5[1/p]\hookrightarrow\text{target generator length five}.
\end{equation}
In particular $S_5[1/p]$ is $\Lambda$-torsionfree,
because a nonzero $f(\gamma)$ acts on its coefficient
arrays by the nonzero multiplier $f(z_1\cdots z_5)$.

The preceding exact formula kills $(R_4)_{\rm tor}$
by $p^{b+d_4}a(\gamma)$. Equations
\eqref{eq:class5-old-torsion},
\eqref{eq:class5-pbound} and
\eqref{eq:class5-rational-injection} give
\begin{equation}\label{eq:class5-closed-torsion}
 (R_5)_{\rm tor}
      =\ker\bigl(p^{b+d_4+d_5}a(\gamma):R_5\to R_5\bigr).
\end{equation}
Indeed a torsion element multiplied by
$p^{b+d_4}a(\gamma)$ lies in the rationally torsionfree
module $S_5$ and hence is killed by $p^{d_5}$.
The converse follows from the nonzero annihilator.
This proves closedness and full arithmetic invariance.

The exact class-four target kernel is $K_4^{\rm det}$.
The three factors of $J_6$ therefore have target
length at least six, so the image of $R_5$ is in pure
length five. For an outer-kernel element,
multiply its class in $R_5$ by $p^{b+d_4}a(\gamma)$.
The result lies in $S_5$ and has zero rational image;
\eqref{eq:class5-pbound} then proves that its class
lies in \eqref{eq:class5-closed-torsion}.
Conversely the roots of $a$ are among the original
abelian torsion eigenvalues. By
Lemma~\ref{lem:positive-eigenvalues}, none occurs on
a positive target coordinate of length five.
Thus the displayed torsion subgroup itself maps to
zero. This proves exactness of the compact-image map.

The central extension is torsionfree since $Q_4$ is
torsionfree and its central kernel has no $p$-torsion.
The pure Hall module in \eqref{eq:class5-pbound}
has an injective five-variable coefficient map,
so meets $(R_5)_{\rm tor}$ trivially. The word
$[x,[x,[x,[x,y]]]]$ is nonzero for independent odd
coordinates. Finally $\gamma_6H\subseteq J_6$,
which proves class exactly five and strictness.
\end{proof}

Put
\[
 D_1=\operatorname{im}\bigl((R_5)_{\rm tor}^-
                                         \longrightarrow D_0\bigr).
\]
Since $S_5$ is odd, the even part of $R_5$ is entirely
$\Lambda$-torsion and has image $Z_0$ in $M$.
The abelianization map therefore gives
\begin{equation}\label{eq:class5-abelianization}
                         Q_5^{\rm ab}=(U/D_1)\oplus(T/Z_0).
\end{equation}

\subsection{Selected components at class five}
\label{subsec:selected-class-five}

Let $B_5$ be the image of $B$ in $Q_5$. It still has
class at most two, since
$[[B,B],B]\subseteq[B,K_3]\subseteq J_6$.
Thus its integral class-two BCH Lie algebra exists
for every odd $p$. For a good even set $\Sigma$, take
the closed $Q_5$-normal subgroup generated by the
exponentials of all omitted even isotypic submodules
of $\operatorname{Lie}(B_5)$, and denote the quotient
by $Q_{5,\Sigma}$. Write $J_\Sigma^{(5)}$ for this closed normal subgroup, and $L_\Sigma^{(5)}$ for the fixed field of $\ker(H\to Q_{5,\Sigma})$.



\begin{corollary}[Class-five component detection]
\label{cor:selected-class-five}
Let $p$ be odd and let $\Sigma$ be any set of even exponents
such that $X^{[e]}$ is finite for every $e\in\Sigma$.
The subgroup $J_\Sigma^{(5)}$ is independent of the tame
complement and stable under the full arithmetic Galois group.
Its canonical field is Galois over $\mathbb Q$, and
\[
 L_\Sigma^{(5)}\subseteq\Omega_p^*,\qquad
 \operatorname{Gal}(L_\Sigma^{(5)}/F_\infty)=Q_{5,\Sigma},\qquad
 Q_{5,\Sigma}^{\rm ab}=(U/D_1)\oplus T_\Sigma.
\]
For $-2,-4\in\Sigma$, the quotient has nilpotency class
exactly five. It maps onto the class-four component
quotient of Proposition~\ref{prop:selected-class-four},
and strictly enlarges its field in this case.
\end{corollary}

\begin{proof}[Proof of Corollary~\ref{cor:selected-class-five}]
Tame complements are $H$-conjugate, and for each
arithmetic element $g$ a suitable $h^{-1}g$ centralizes
the complement pointwise. The same normal-closure
argument as at class four proves independence and
full arithmetic stability. The image of the omitted
even Lie submodules in $B_4$ is exactly the earlier
set of generators, giving the stated surjection.
Their abelian image in \eqref{eq:class5-abelianization}
is precisely the omitted even part of $T/Z_0$.
For selected $e$, the component theorem gives
$(Z_0)_e=0$, proving the displayed abelianization.

Project the free target through length five
onto two distinct odd generators of weights
$m,n\equiv-1\pmod{p-1}$, after the same
$\Delta$-equivariant choice of lifts as before.
The image of $B_5$ has no linear term, so its even
part has only character $-2$ in length two and
character $-4$ in length four. The morphism therefore
kills every omitted even generator and factors
through $Q_{5,\Sigma}$. The two moment coordinates
of $w,\gamma w$ have the nonzero determinant
$\lambda_m\lambda_n(u^n-u^m)$; their fifth Hall
commutator has nonzero length-five term. It survives
in the selected quotient, proving class five and
strictness.
\end{proof}

\section{Divisibility in a weighted companion algebra}
\label{sec:weighted-companion-saturation}

We prove the coefficient divisibility statement used at every finite
stage of the source construction.  Completions and the source division
are included in the statement.

\subsection{The source and its coefficient representation}

Let \(L/\mathbb Q_p\) be finite, with ring of integers \(O\), and put
\[
 R=O[[\gamma-1]][1/p].
\]
Let \(W=R^r\), with its regular \(\Gamma=\langle\gamma\rangle\)-action.
For \(2\leq j\leq d\), let \(V_j\) be a finite-dimensional
\(L\)-space with continuous \(\Gamma\)-action.  A zero space is allowed.
Form the completed free weighted associative algebra
\[
 \mathcal A^{\leq d}
   =\widehat T_L(W,V_2,\ldots,V_d),
 \qquad \operatorname{wt}(W)=1,\quad
 \operatorname{wt}(V_j)=j.
\]
In each fixed degree, the completion is obtained from compact
\(O\)-lattices and then by inverting \(p\).  In particular an entire
coefficient array has one scalar denominator.

Write
\[
 S_j=(\mathcal A^{<j})_j,\qquad
 E_j=(\mathcal A^{\leq j})_j=S_j\oplus V_j
\]
as \(L\)-spaces.  The action on the generators has the form
\begin{equation}\label{eq:wcs-action}
 \gamma v=A_jv+C_j(v),\qquad C_j(v)\in S_j,
\end{equation}
where \(A_j\) is the given finite action.  Thus it defines an actual
weight-preserving algebra automorphism, triangular in generator
weight.  We require that each actual module \(E_j\) is \(R\)-torsion-free.
For the Lie assertion below, the companions \(C_j(v)\) are primitive;
equivalently, \(\gamma\) is a Hopf automorphism for the coproduct that
makes all generators primitive.  These properties will be verified
for the arithmetic source modules when the theorem is applied.

At a fixed degree only finitely many companion coefficients occur.
Choose a free compact lattice in \(W\), and scale the finitely many
generator lattices successively so that the companions for both
\(\gamma\) and \(\gamma^{-1}\) are integral.  This gives a compact
\(\Gamma\)-stable model.  This operation does not assert integral
\(\Lambda\)-freeness of an arithmetic lattice originally giving \(W\).

Set \(z_i=1+X_i\).  The regular coefficient map on \(W\) is the
identity map to the one-variable power-series arrays, with one
coordinate for each of its \(r\) colors.  Products concatenate
ordered variable blocks.  If a weight-\(j\) generator occupies a
contiguous block \(I\), its coefficient is defined by the exact
equation
\begin{equation}\label{eq:wcs-resolvent}
 (z_I I-A_j)\,\iota(v)=\iota(C_j(v)),
 \qquad z_I=\prod_{i\in I}z_i.
\end{equation}
Choose a polynomial \(q_j\) annihilating \(A_j\).  The resulting map in
degree \(n\) is
\[
 \iota_n:(\mathcal A^{\leq d})_n
 \longrightarrow \mathcal B_n^{\,r^n},\qquad
 \mathcal B_n=
 O[[X_1,\ldots,X_n]]
 [1/p,q_j(z_I)^{-1}:|I|=j,\ I\text{ contiguous}].
\]
In this expression a localization uses a common polynomial
denominator for the whole array.  For \(n>d\), all these intervals
are proper.  The action of \(\gamma\) on the target is multiplication
by \(z_1\cdots z_n\).

\begin{theorem}\label{thm:weighted-companion-saturation}
The maps \(\iota_n\) are injective.  For every \(n>d\),
\(\beta\in1+\mathfrak m_L\), and \(k\geq1\), put
\(D_\beta=z_1\cdots z_n-\beta\).  Then
\begin{equation}\label{eq:wcs-saturation}
 \iota_n((\mathcal A^{\leq d})_n)
       \cap D_\beta^k\mathcal B_n^{\,r^n}
 =\iota_n\bigl((\gamma-\beta)^k
                       (\mathcal A^{\leq d})_n\bigr).
\end{equation}
When the companions are primitive, the same equality holds on
the completed free weighted Lie subalgebra.
\end{theorem}

Degree \(n\leq d\) is deliberately absent.  A generator of weight
\(n\) can have its full total divisor inverted in
\eqref{eq:wcs-resolvent}; its source coinvariants need not vanish.

\subsection{Bounded local coordinates and balanced tensors}

\begin{lemma}\label{lem:wcs-local-coordinates}
Each injected fixed-degree source module is locally pro-free
over \(R\), with finitely many coordinate lines rescaled by powers
of a local parameter.  For an extension
\[
 0\longrightarrow S_j\longrightarrow E_j\longrightarrow V_j
       \longrightarrow0
\]
whose old coefficient map is saturated on its total divisor, these
coordinate changes may be chosen simultaneously in the source and
on its actual coefficient columns.  They use finitely many
coefficient functionals and one auxiliary denominator.
\end{lemma}

\begin{proof}
First consider the closed ideal in a fixed-degree source consisting
of words containing at least one \(W\)-letter.  Filter it by the
number of \(W\)-letters.  A companion substitution never decreases
this number.  Its graded pieces are finite sums of a tensor power
of \(W\) and finite modules of zero-\(W\) gap words.

On the first regular \(W\)-factor the exact map
\[
       [g]\otimes m\longmapsto[g]\otimes g^{-1}m
\]
turns the diagonal action into the regular action on that factor.
The map and its inverse are defined on compatible finite
group-algebra quotients and finite invariant coefficient quotients.
Thus each graded piece is compact pro-free over \(O[[\Gamma]]\).

A continuous compact surjection onto such a module has a
continuous module section.  Indeed, its convergent free basis
has lifts tending to zero: choose them through a decreasing
sequence of open submodules, using openness of the compact
surjection.  The series formed from those lifts converges on
every finite quotient and defines the section.  Splitting the
finite filtration proves that the ideal is pro-free.  The
quotient consists of finitely many zero-\(W\) words and is finite
over \(O\).

After enlarging \(L\), fix a root \(\alpha\), and put
\(f=\gamma-\alpha\).  For the top lifts of the Jordan blocks of
the finite quotient write
\[
                    f^{r_i}y_i=C_i
\]
in the pro-free submodule.  The classes of the \(C_i\) modulo
\(f\) are independent over \(L\).  A relation
\(\sum c_iC_i=fs\) would make
\[
 f\left(\sum c_i f^{r_i-1}y_i-s\right)=0.
\]
Torsion-freeness of the full source would force the parenthesized
element to vanish, contradicting its nonzero Jordan-socle image.

Choose finitely many coefficient coordinates giving an invertible
minor on these residues.  Its determinant \(g\) is nonzero at
\(f=0\).  After inverting \(g\), a finite-rank change of coordinates
sends the full \(f\)-dependent columns \(C_i\) to coordinate
vectors.  The local source is obtained by rescaling these finitely
many lines by \(f^{-r_i}\).  Every operation has one common
denominator.

For the stated companion extension, old total saturation injects
\(S_j/fS_j\) into the actual specialized coefficient module.
Multiply the companion columns by a common proper-interval
denominator and choose the same finite minor there.  This gives
bounded source and coefficient projectors
\(\ell_i(C_a)=\delta_{ia}\), as well as their complementary
projector.  Locally,
\begin{equation}\label{eq:wcs-polar-coordinates}
 E_j=S_j^\circ\oplus\bigoplus_i R_f f^{-r_i}C_i,\qquad
 S_j=S_j^\circ\oplus\bigoplus_i R_f C_i .
\end{equation}
Here \(R_f\) indicates the local ring with the one auxiliary
denominator allowed.  The projectors commute with \(f\).  The
regular complement inherits every power of the old saturation.
\end{proof}

An inverse of \(g\) need not have a bounded infinite Taylor
expansion in \(f\).  The proof only uses a fixed finite jet at
each step.  Its coefficients have a common bound, obtained by
inverting a fixed power of \(g(0)\).  This is why the local
coordinate argument remains valid for bounded completions.

\begin{lemma}\label{lem:wcs-balanced}
For two previously constructed degree modules \(E_i,E_j\),
whose smaller alphabets have the total saturation of
Theorem~\ref{thm:weighted-companion-saturation}, the actual
balanced coefficient map is injective:
\[
 (E_i\widehat\otimes E_j)/(\gamma\otimes\gamma-\beta)
   =E_i\widehat\otimes_{R,\tau_\beta}E_j
       \longrightarrow \mathcal B_{i+j}/D_\beta,
 \qquad \tau_\beta(\gamma)=\beta\gamma^{-1}.
\]
All tensor products and source quotients use the compact
models followed by inversion of \(p\).
\end{lemma}

\begin{proof}
The displayed source equality is the quotient by the balancing
relation on pure tensors, followed by its continuous extension.
Choose the total variable \(r\) of the first block; the second
total is then \(\beta/r\).  The two sets of internal variables
remain disjoint.

First take the two old modules \(S_i,S_j\).  At \(r=\alpha\),
their maps modulo \(f=r-\alpha\) are injective by old total
saturation, with the second root transported to \(\beta/\alpha\).
The completed tensor of these injections is injective.
To check this last statement within the bounded convention,
clear the fixed denominators.  A compact torsion-free
\(O\)-module is a product of copies of \(O\); completed tensoring
with such a product gives a product of the original injection.
The target is exactly the bounded coefficient module in the
disjoint internal variables.

In the local coordinates of Lemma~\ref{lem:wcs-local-coordinates},
a nonzero balanced tensor has a least \(f\)-adic order.
Monic division by that power is bounded.  Its nonzero reduction
would have nonzero coefficient by the preceding paragraph.
This proves the old balanced injection, and the same argument
proves its \(f\)-saturation.

Now replace \(S_i,S_j\) by \(E_i,E_j\).  At a root of either
finite companion polynomial, use
\eqref{eq:wcs-polar-coordinates} on both factors.  Their
independent companion columns give finite local coordinate
changes in the source and coefficient modules.  The tensor
coordinates are rescaled by the sums of their two exponents;
unchanged coordinates have exponent zero.  The old injection
therefore remains injective on each rescaled block.

The possible global kernel is killed by the product of the
two finite companion polynomials, the second transported by
\(\tau_\beta\).  Its local vanishing at their finitely many
roots implies global vanishing by Bezout.  These are finite
polynomials and finite coordinate changes; no separate
denominator has been assigned to an infinite coefficient.
\end{proof}

\subsection{Interval poles and the boundary bound}

On \(z_1\cdots z_n=\beta\), two proper contiguous-interval
divisors can coincide only when their intervals and roots
coincide, or when their intervals are complementary prefix
and suffix intervals and their roots multiply to \(\beta\).
Indeed, after fixing the full product and one interval product,
a second monomial is constant exactly when its indicator vector
lies in the span of those two indicator vectors.  For zero-one
indicators this gives the interval, its complement, the empty
set, or the full set.  The complement is contiguous precisely
for a prefix or suffix.

Recursively expanding \eqref{eq:wcs-resolvent} produces finitely
many companion-tree patterns at a fixed weighted degree.
Each finite vertex occupies a contiguous interval and supplies
the denominator of that interval's total.  Ancestor weights
strictly increase.  The leaf coefficients remain completed
bounded \(W\)-arrays; only the finite list of token patterns
is expanded.

\begin{lemma}\label{lem:wcs-boundary-bound}
Put \(n=d+s<2d\), \(1\leq s<d\), and work on the total
hypersurface at
\[
 r=z_1\cdots z_d,\qquad t=z_{d+1}\cdots z_n=\beta/r,
 \qquad f=r-\alpha.
\]
Consider all source terms except those beginning with a new
weight-\(d\) token.  In the local suffix coordinates
\eqref{eq:wcs-polar-coordinates}, their negative principal
parts lie only in the suffix companion columns
\(f^{-q_a}D_a\), and their pole order on that column is at
most \(q_a\).  They have no negative principal part in
the suffix regular complement.
\end{lemma}

\begin{proof}
By the interval calculation, an \(f\)-pole of one of these
terms must come from a vertex occupying the entire complementary
suffix of length \(s\).  No other term has a vertex occupying
the full prefix of length \(d\).  This suffix vertex has
weight \(s\), occurs once in its tree, and has root
\(\beta/\alpha\).

Every ancestor of this vertex is a suffix of length \(j>s\).
On the total hypersurface its total is exactly
\begin{equation}\label{eq:wcs-ancestor}
 z_{n-j+1}\cdots z_n
       =\frac{\beta}{z_1\cdots z_{n-j}}.
\end{equation}
Choose \(z_1,\ldots,z_{d-1}\) as the internal prefix
coordinates and \(r\) as its total.  Since \(n-j<d\),
\eqref{eq:wcs-ancestor} is a function of proper prefix
coordinates, independent of \(f\).  Every coefficient
outside the suffix vertex depends on prefix variables
and the suffix total only.  Its denominators are regular
at \(f\), and its ancestor denominators are the fixed
proper-prefix factors in \eqref{eq:wcs-ancestor}.

The vertex itself supplies an exact suffix column
\((t-\beta/\alpha)^{-q_a}D_a(t)\).  The factor
\(t-\beta/\alpha\) is a unit multiple of \(f\).
Thus the negative principal part is a sum of
\[
 h_a(f,\text{prefix})\,f^{-q_a}D_a(f,\text{suffix}),
 \qquad h_a\text{ regular at }f.
\]
This includes an old token crossing the cut and a new suffix
token whose companion contains the suffix vertex.

There are only finitely many negative orders.  Discard all
\(f\)-regular terms before applying any coefficient
projector.  Their possibly mixed prefix-suffix denominators
therefore never enter the extraction space.  For the remaining
polar terms, take only the finite jets required by their
orders.  A fixed power of each nonzero constant denominator
clears these jets.  The resulting denominators separate into
a common prefix denominator and a common suffix denominator.
The principal-part identity consequently lies in the
completed tensor of the two separate coefficient modules.
The full suffix projectors of
Lemma~\ref{lem:wcs-local-coordinates} apply to that identity
and give the stated bounds.  Using full columns, rather
than just their residues, handles derivatives of unequal
Jordan blocks.
\end{proof}

\subsection{Proof of the saturation theorem}

\begin{proof}[Proof of Theorem~\ref{thm:weighted-companion-saturation}]
We induct on the largest generator weight, retaining the
assertion for every smaller alphabet.  For the regular
alphabet \(W\), change one variable to the full product.
Monic power-series division in that coordinate proves the
statement with a common coefficient bound.

Write \(\mathcal A=\mathcal A^{<d}\),
\(\mathcal B=\mathcal A^{\leq d}\).  Old total saturation
in degree \(d\) and torsion-freeness of \(E_d\) give the
independent full companion residues of Lemma~\ref{lem:wcs-local-coordinates}.  The ordinary
coefficient map of \(\mathcal B\) is injective: take the
largest initial token weight, extract its highest full
prefix pole, and use a finite coefficient minor of those
residues with the disjoint tail variables.  Lower initial
token weights cannot have this full prefix pole.
Descending through its Jordan orders removes that token.
Induction on the tail degree and then the initial weight
finishes.  The initial \(W\)-term uses ordinary completed
tensor injection.  This proves unrestricted injectivity
in every degree while establishing the induction step.

It suffices first to prove the quotient map modulo
\(D_\beta\) is injective in the range \(d<n\leq2d\).

\paragraph{Interior new tokens.}
A new weight-\(d\) token in an interior interval has
nonempty outside blocks of lengths \(a,b\), with
\(a+b=n-d\).  Its full interval has disconnected complement,
so its pole cannot coincide with a distinct interval pole
on the total hypersurface.  The outside source module is
the actual \(\mathcal A_a\widehat\otimes\mathcal A_b\).
Both \(a,b<d\), so
Lemma~\ref{lem:wcs-balanced} injects its quotient at outside
total \(\beta/\alpha\).  The independent internal companion
residues therefore force its outside coefficient to be an
actual source multiple of
\(\gamma_{\rm out}-\beta/\alpha\).

For a Jordan chain \(\gamma y=\alpha y+Ny\), the source
coinvariant relation is exactly
\begin{equation}\label{eq:wcs-Jordan}
 N^k y\otimes(\gamma_{\rm out}-\beta/\alpha)b
 \equiv-\alpha^{-1}N^{k+1}y\otimes\gamma_{\rm out}b
                                      \pmod{\gamma-\beta}.
\end{equation}
At the end of the chain its companion belongs to the old
algebra.  Repeating removes the interior placement by
actual source relations.  The process is finite in the
number of positions and the Jordan lengths.  In particular
it remains valid when \(n=2d\) and \(a+b=d\).

\paragraph{A boundary token when \(d<n<2d\).}
Put \(s=n-d\).  Modulo old prefix terms, the contribution
beginning in \(V_d\) has the actual balanced form
\begin{equation}\label{eq:wcs-new-prefix}
 \sum_i f^{-r_i}C_i(\text{prefix},f)
                       \otimes b_i(\text{suffix},f),
 \qquad b_i\in E_s .
\end{equation}
Use the local suffix decomposition
\[
 b_i=b_i^\circ+\sum_a f^{-q_a}D_a\,b_{ia},
 \qquad b_i^\circ\in S_s^\circ,\quad b_{ia}\in R_f.
\]
Here \(S_s\) is the algebra in degree \(s\) \emph{before}
weight \(s\) was adjoined.  Its total saturation is part of
the retained smaller-alphabet induction, even though
\(s\) is below the current old maximum weight.

Apply the suffix companion-column projectors first.
By Lemma~\ref{lem:wcs-boundary-bound}, every other term
has zero negative principal part in the complementary
coordinates.  A finite prefix minor of the independent
\(C_i\) columns therefore shows
\[
 \iota(b_i^\circ)\in f^{r_i}\mathcal B_s.
 \quad\text{Hence}\quad b_i^\circ\in f^{r_i}S_s^\circ
\]
by old saturation.  On the \(a\)-th polar line, every
other term has order at most \(q_a\), whereas
\eqref{eq:wcs-new-prefix} gives
\(\sum_i f^{-r_i-q_a}C_i b_{ia}\).  Multiply by \(f^{q_a}\)
and use the same prefix minor.  It gives
\(b_{ia}\in f^{r_i}R_f\).  Thus the exact source conclusion is
\begin{equation}\label{eq:wcs-tail-divisibility}
                         b_i\in f^{r_i}E_s .
\end{equation}

Balancing sends \(f=r-\alpha\) to
\(\beta\gamma_s^{-1}-\alpha\), a unit multiple of
\(\gamma_s-\beta/\alpha\).  On this Jordan block the
quotient of the new prefix by the old prefix is exactly
\[
 E_s/(\gamma_s-\beta/\alpha)^{r_i}E_s .
 \]
Equation~\eqref{eq:wcs-tail-divisibility} is zero in this
source quotient.  Its finitely many root-local conclusions
give global zero by Bezout, since the quotient is killed by
the finite transported characteristic polynomial of \(V_d\).
Right exactness of the actual coinvariant sequence now
replaces the prefix contribution by an old source element
plus an actual \((\gamma-\beta)\)-multiple.  The companion
correction has no weight-\(d\) token.

The symmetric argument treats the suffix.  At most one
new token occurs for \(n<2d\).  Having removed its interior
and boundary placements, the remaining source is old;
old degree-\(n\) saturation finishes this range.

\paragraph{The paired boundary case \(n=2d\).}
After removing interior placements, all remaining new
terms lie in the single cut module
\[
                  (E_d\widehat\otimes E_d)_{\gamma=\beta}.
\]
This includes both edge orientations and the possible
two-token word \(V_dV_d\).  Every term outside this cut
has no full prefix-\(d\) or suffix-\(d\) pole.

Use the finite polar coordinates in both factors, with
the second total transported by \(t=\beta/r\).
The old balanced module
\((S_d\widehat\otimes S_d)_{\gamma=\beta}\)
is saturated at \(f=r-\alpha\): modulo \(f\) its map is
the completed tensor of the two old specialized
injections in disjoint internal variables.  Repeated
division gives every power.
The new tensor coordinates are rescaled by
\(f^{-r_i-q_a}\), with exponent zero on each unchanged
line.  Since the total cut coefficient is regular,
its coordinate numerators are divisible by the full
rescaling powers.  It belongs to the old balanced
source locally at every relevant root.  The quotient
is killed by the product of the two companion
polynomials, one transported by \(r\mapsto\beta/r\);
Bezout gives global old-cut membership.  This is an
actual source replacement.  Old saturation finishes
degree \(2d\).

\paragraph{Degrees \(n>2d\).}
Filter by the largest permitted weight \(j\) of the
first token.  Its full prefix pole becomes a total
divisor of the tail of length \(n-j>d\).  No proper
companion interval has that length, and a lower
initial token cannot contain the full prefix.
Thus the highest Jordan residue is separate from the
lower filtration.  Finite coefficient extraction in
the prefix and the already proved lower-degree total
saturation put its tail coefficient in the actual
source image of \(\gamma_{\rm tail}-\beta/\alpha\).
Equation~\eqref{eq:wcs-Jordan} removes it.  This descends
through the first-token weights.

At the regular first token, the exact source evaluation is
\[
 z_1=\beta\gamma_{\rm tail}^{-1}.
\]
For \(X_1=z_1-1\), it sends a series
\(\sum_m X_1^m b_m\) to
\(\sum_m(\beta\gamma_{\rm tail}^{-1}-1)^m b_m\).
The operator in parentheses is topologically nilpotent
on every finite integral quotient, so the series
converges there.  Its kernel is the image of
\(X_1-(\beta\gamma_{\rm tail}^{-1}-1)\), by the exact
identity \(X^m-A^m=(X-A)\sum_{\ell<m}X^{m-1-\ell}A^\ell\)
and continuity.  Unrestricted tail injection finishes
this step.  This proves all \(n>d\).

\paragraph{Powers and primitive elements.}
If \(\iota_n(a)\) is divisible by \(D_\beta^k\), the
first-power assertion gives \(a=(\gamma-\beta)a_1\)
in the actual source.  Since multiplication by
\(D_\beta\) is injective in the coefficient domain,
\(\iota_n(a_1)\) is divisible by \(D_\beta^{k-1}\).
Repeating proves \eqref{eq:wcs-saturation}.  Each fixed
\(n,k\) uses finitely many bounded divisions; the final
source element has a common denominator.

For primitive companions, the Eulerian idempotent
\[
 e=\log^\star(\operatorname{id})
 \]
projects the completed tensor algebra onto its primitive
free Lie algebra; see \cite[\S1.1, p.~2]{BandieraSchaetz}.
In fixed weighted degree this is a finite sum with a
common rational denominator.  Its naturality makes it
commute with the Hopf automorphism \(\gamma\), even when
\(\gamma\) changes ordinary token count.  Apply \(e\)
to the source equality
\(\ell=(\gamma-\beta)^k a\) for primitive \(\ell\).
It gives
\(\ell=(\gamma-\beta)^k e(a)\) in the free weighted
Lie source, which proves the last assertion.
\end{proof}

\subsection{The exact cut relation represented by the proof}

For \(n>d\), the actual source maps
\[
 \bigoplus_{i+j+k=n}\mathcal A_i\widehat\otimes
       \mathcal A_j\widehat\otimes\mathcal A_k
 \xrightarrow{\ \partial\ }
 \bigoplus_{i+j=n}\mathcal A_i\widehat\otimes\mathcal A_j
 \longrightarrow\mathcal A_n\longrightarrow0,
 \qquad
 \partial(a\otimes b\otimes c)=ab\otimes c-a\otimes bc,
\]
are exact.  For each word, its possible cuts form a
nonempty simplex, and these are exactly its cut-identification
relations.  The coefficientwise contraction has coefficients
\(0,1,-1\), so it is bounded on the completions.
Every term and every cycle is \((\gamma-\beta)\)-torsion-free
by unrestricted injection.  The exact sequences of cycles
therefore remain exact on coinvariants.

The proof above shows that every relation among the
specialized coefficient images of all cuts is the image
of this actual three-block boundary.  The boundary
principal-part argument is needed here: intersections
of two specialized cut images can be mediated by a
third cut, and a crossing token can carry a complementary
internal pole.  Equations~\eqref{eq:wcs-ancestor} and
\eqref{eq:wcs-tail-divisibility} give their source
replacement without discarding either phenomenon.

\section{The canonical cyclotomic tower in every finite class}
\label{sec:canonical-all-stages}

Return to the actual cyclotomic source
\[
 H=\operatorname{Gal}(\Omega_p/F_\infty),\qquad
 M=H^{\mathrm{ab}},\qquad
 V=M_{\Lambda\text{-tor}}=D\oplus T,\qquad W=M/V.
\]
The arithmetic results already established give that \(H\) is free
pro-\(p\), \(V\) is finite free over \(\mathbb Z_p\), \(W\) is
torsion-free over \(\Lambda\), and \(W[1/p]\) is finite free over
\(\Lambda[1/p]\).  These use Ferrero--Washington and Iwasawa's
freeness theorem as recorded in \S\ref{sec:abelian}.
Fix the characteristic polynomial
\[
               a(Z)=\det(Z-\gamma\mid V).
\]
Lemma~\ref{lem:positive-eigenvalues} proves
\begin{equation}\label{eq:uniform-positive-exclusion}
                       a(u^m)\ne0\qquad(m\geq2).
\end{equation}

Here is a recursion involving only the source and its arithmetic
action.  Put \(\mathcal F_1=H\), \(\mathcal R_1=M\), and let \(\mathcal F_2\) be the inverse
image of \(V\).  For \(n\geq2\), put
\begin{equation}\label{eq:uniform-recursion}
 J_n=\overline{\prod_{\substack{i+j=n\\i,j\geq1}}[\mathcal F_i,\mathcal F_j]},
 \qquad
 \mathcal R_n=\mathcal F_n/J_{n+1},
 \qquad
 \mathcal F_{n+1}=\operatorname{preim}_{\mathcal F_n}
                  \bigl((\mathcal R_n)_{\Lambda\text{-tor}}\bigr).
\end{equation}
The proof below establishes at each step that these expressions
are defined and that the displayed torsion subgroup is closed.
Set
\[
 Q_n=H/\mathcal F_{n+1},\qquad E_n=\mathcal F_n/\mathcal F_{n+1}.
\]

We prove Theorem~\ref{thm:canonical-all-stages}. In the course of the proof, the map from \((\mathcal R_n)_{\Lambda\text{-tor}}\) to \(M\) will also be shown to have kernel killed by \(p^{h_n}\).

The exact source-intersection question remaining in the theorem is whether
\(\bigcap_n\mathcal F_n=1\).  The source countermodel below shows why the
finite-stage assertions alone do not imply that equality.

\subsection{A uniform source calculation}

We give the integral comparison needed in the induction.
The statement distinguishes the actual commutator surjection
from the rational Hall isomorphism.

Suppose the construction is known through \(Q_{n-1}\).  Then
\([\mathcal F_i,\mathcal F_j]\subseteq \mathcal F_{i+j}\) whenever the right side has
already been constructed.  This follows either from its
established target description or from the corresponding
central filtration.  Thus \(J_{n+1}\subseteq \mathcal F_n\), and \(\mathcal R_n\)
is central in \(H/J_{n+1}\).  There is an actual exact sequence
\begin{equation}\label{eq:uniform-central-sequence}
 0\longrightarrow \mathcal S_n:=J_n/J_{n+1}
 \longrightarrow \mathcal R_n
 \longrightarrow (\mathcal R_{n-1})_{\Lambda\text{-tor}}
 \longrightarrow0.
\end{equation}
Indeed \(\mathcal F_n/J_n\) is exactly the torsion subgroup used at the
preceding step.

\paragraph{Rational Hall coordinates with a common bound.}
Work in \(H/\gamma_{n+1}H\).  A compact
\(\mathbb Z_p\)-basis of \(M\), using a continuous \(\mathbb Z_p\)-module splitting \(M\simeq W\oplus V\), lifts to a convergent
free basis of \(H\), by the Frattini and projectivity argument
used above.  Its free nilpotent Hall coordinates identify
\[
       \log(H/\gamma_{n+1}H)[1/p]
\]
with the completed free rational Lie algebra on \(M[1/p]\),
truncated at ordinary length \(n\).
In fixed length, BCH, logarithm, and exponential use a finite
list of universal rational coefficients.  One integer clears
their denominators independently of the number of basis
coordinates.  Passing through the finite-tail quotients
therefore gives the same common bound on the compact completion.

Let
\[
 A_i=\operatorname{im}(\mathcal F_{i+1}\longrightarrow M)\subseteq V.
\]
Choose a rational basis of \(V[1/p]\) adapted to the finite flag
\[
 V[1/p]=A_1[1/p]\supseteq\cdots\supseteq A_{n-1}[1/p].
\]
Assign weight \(j\) to a complement of \(A_j[1/p]\) in
\(A_{j-1}[1/p]\), and weight \(n\) to the remaining space.
Assign weight one to \(W[1/p]\).
The finite matrices and their inverses have a common denominator.
After scaling, the selected finite basis vectors lift to the
indicated \(\mathcal F_j\).

The literal source maps are obtained as follows.  The image of
a weight-\(j\) lift in \(Q_{j-1}\) has zero abelianization.
Express it in the already established derived Hall coordinates
of weights at most \(j-1\), lift that expression to the free
source, and subtract it.  The corrected generator lies in the
rational kernel of the actual map to \(Q_{j-1}\), with unchanged
abelian image.  Correct the remaining weight-\(n\) vectors
against \(Q_{n-1}\) itself.  There are finitely many such
corrections.  The bounded Hall sections at the preceding
stages and the fixed-class BCH bound give one bound for them.

These substitutions have invertible bounded linear part and
derived corrections.  Their inverse is obtained successively
in ordinary degrees at most \(n\), again with a common bound.
The resulting map to \(Q_{n-1}\) is the literal weighted
truncation map.  To see this, its map on each already
established relative graded piece is the free Hall
isomorphism; a finite descending-filtration argument gives
an isomorphism on the truncation.  The same statement holds
for every earlier quotient.  Thus the rational image of
\(\mathcal F_i\) is exactly the ideal of assigned weights at least
\(i\), as an equality of kernels of the actual maps.

Every nonlinear monomial of assigned weight at least \(n+1\)
is a bracket of monomials whose weights dominate positive
\(i,j\) with \(i+j=n+1\); hence it belongs to \(J_{n+1}\).
The reverse inclusion follows from the assigned weights of
its defining brackets.  Consequently
\begin{equation}\label{eq:uniform-rational-Hall}
 \mathcal S_n[1/p]\simeq
 \operatorname{Lie}^{\mathrm{wt}}_n
      \bigl(W[1/p],C_2,\ldots,C_{n-1}\bigr),
\end{equation}
where the \(C_j\) are the finite rational generator spaces
already recovered in weight \(j\).
The only additional rational generators in \(\mathcal R_n\) are the
remaining abelian vectors assigned weight \(n\).

\paragraph{The integral central presentation and its torsion.}
Form
\[
 \mathcal P_n=
 \bigoplus_{\substack{i<j\\i+j=n}}
       E_i\widehat\otimes_{\mathbb Z_p} E_j
 \ \oplus\
 \begin{cases}
 \widehat{\bigwedge}^{\,2}E_{n/2},&n\text{ even},\\
 0,&n\text{ odd}.
 \end{cases}
\]
Commutators give a continuous surjection
\(\mathcal P_n\to \mathcal S_n\).  Changing a representative introduces
a commutator in \(J_{n+1}\); the resulting central brackets
are \(\mathbb Z_p\)-bilinear.  Quotient by all Jacobi images
for triples of positive weights summing to \(n\), using the
actual lower bracket maps.  Hall--Witt gives
\begin{equation}\label{eq:uniform-Hall-surjection}
       \mathcal C_n:=\mathcal P_n/\mathcal I_n
                         \twoheadrightarrow \mathcal S_n.
\end{equation}
Equation~\eqref{eq:uniform-rational-Hall} makes this a rational
isomorphism.  Its integral kernel may contain inherited
bounded \(p\)-torsion.

Choose bracket-compatible compact free weighted Hall lattices
\(E_{i,0}\subseteq E_i\), \(i<n\), with
\[
 p^{b_i}E_i\subseteq E_{i,0},\qquad b_1=0.
\]
They are generated by the actual compact \(W\)-lattice and
finitely many selected lifts of each recovered \(C_j\).
Define
\begin{equation}\label{eq:uniform-bounds}
 d_n=\max_{i+j=n}(b_i+b_j),\qquad
 j_n=\max_{i+j+k=n}(b_i+b_j+b_k),
\end{equation}
taking \(j_2=0\).  Then
\[
 p^{d_n}\mathcal P_n\subseteq\mathcal P_{n,0},\qquad
 p^{j_n}\mathcal I_n\subseteq\mathcal I_{n,0}.
\]
The free Hall quotient
\(\mathcal P_{n,0}/\mathcal I_{n,0}\) is
\(p\)-torsion-free.  If \(x\) represents torsion in
\(\mathcal C_n\), some \(p^kx\) lies in \(\mathcal I_n\).
The element \(p^{d_n}x\) belongs to \(\mathcal P_{n,0}\),
and a further power \(p^{k+j_n}\) puts it in
\(\mathcal I_{n,0}\).  Saturation of the free Hall relation
module gives
\begin{equation}\label{eq:uniform-S-torsion}
 p^{d_n}\mathcal C_n[p^\infty]=0,\qquad
 p^{d_n}\mathcal S_n[p^\infty]=0.
\end{equation}
The second assertion follows because the kernel of
\eqref{eq:uniform-Hall-surjection} is torsion.
Also \(p^{d_n}\mathcal S_n\) lies in the free Hall image.
The pure module \(\operatorname{Lie}_n(W)\) injects into
the actual \(\mathcal S_n\), since its map is rationally injective
and its source is \(p\)-torsion-free.

\subsection{Induction against the actual outer image}

The rational outer image supplied by Hain--Matsumoto and
Brown is the weight-completed free Lie image on the odd
letters \(s_m\), \(m\geq3\), as used in the preceding
sections \cite{HM,Brown}.  Use its own generator-length
filtration.  After the established normalization,
\(\gamma s_m=u^m s_m\).  Every finite weight piece of
its length-\(n\) quotient is a sum of positive-weight
eigenspaces.

The base \(n=1\) is the already proved moment detection:
the degree-one coefficient map has kernel \(V\), so
\[
 \mathcal F_2=\ker\rho_{\leq1},\qquad Q_1=W.
\]
The torsion formula has \(h_1=0\).

Assume the theorem through \(n-1\).
The source calculation gives
\eqref{eq:uniform-rational-Hall}.  Apply
Theorem~\ref{thm:weighted-companion-saturation} to the
finite alphabet already recovered through weight \(n-1\).
Its unrestricted degree-\(n\) coefficient map is injective,
and its total-divisor map is saturated.

These are the actual normalized target coefficients.
Indeed the action of \(\gamma\) on a finite recovered
generator and its companion gives
\eqref{eq:wcs-resolvent} on every ordered odd-weight tuple.
Every finite generator is a subquotient of \(V[1/p]\);
its companion polynomial therefore has roots among those
of \(a\).  Equation~\eqref{eq:uniform-positive-exclusion}
ensures that none of its denominators vanishes on the
positive-weight evaluation grid.  Clearing the common
denominator reduces vanishing to the bounded multivariable
moment separation already proved in the class-two
comparison.  It follows that
\begin{equation}\label{eq:uniform-S-injection}
 \mathcal S_n[1/p]\longrightarrow
   \operatorname{gr}^{\mathrm{length}}_n
                          \operatorname{Lie}\mathcal U
\end{equation}
is injective.  In particular \(\mathcal S_n[1/p]\) is
\(\Lambda[1/p]\)-torsion-free; its integral
\(\Lambda\)-torsion is precisely its bounded
\(p\)-torsion from \eqref{eq:uniform-S-torsion}.

Put \(h_n=h_{n-1}+d_n\).  The quotient in
\eqref{eq:uniform-central-sequence} is annihilated by
\(p^{h_{n-1}}a(\gamma)\).  This multiple of a torsion
element of \(\mathcal R_n\) consequently lies in the bounded
torsion of \(\mathcal S_n\).  Conversely, an element killed
by a nonzero element of \(\Lambda\) is torsion by
definition.  Hence
\begin{equation}\label{eq:uniform-closed-torsion}
 (\mathcal R_n)_{\Lambda\text{-tor}}
           =\ker\bigl(p^{h_n}a(\gamma):\mathcal R_n\to \mathcal R_n\bigr).
\end{equation}
This proves closedness before taking its quotient.
The action of \(H\) on \(\mathcal R_n\) is trivial, and the
arithmetic action preserves its \(\Lambda\)-torsion.
Thus \(\mathcal F_{n+1}\) is a closed arithmetically stable
normal subgroup.  The quotient
\(E_n=\mathcal R_n/(\mathcal R_n)_{\Lambda\text{-tor}}\) has no
\(\Lambda\)-torsion, including no \(p\)-torsion.

The map \((\mathcal R_n)_{\Lambda\text{-tor}}\to M\) has kernel
killed by \(p^{h_n}\).  An element of that kernel maps,
in \eqref{eq:uniform-central-sequence}, to the preceding
kernel; multiplication by \(p^{h_{n-1}}\) puts it in
\(\mathcal S_n\), and multiplication by \(p^{d_n}\) then kills it.
Its image is therefore a closed finite-rank submodule
\(A_n\subseteq V\), and the new rational generator space is
\[
                         (A_{n-1}/A_n)[1/p].
\]

It remains to verify the exact target kernel.
The actual length-\(n\) map is defined on \(\mathcal R_n\),
because every bracket defining \(J_{n+1}\) has target
length at least \(n+1\).  Its value on a torsion element
is killed by \(p^{h_n}a(\gamma)\).
On every finite geometric weight piece this operator
is invertible by \eqref{eq:uniform-positive-exclusion};
thus torsion maps to zero in the full target.

Conversely, if \(r\in \mathcal R_n\) has zero length-\(n\) image,
then \(p^{h_{n-1}}a(\gamma)r\) lies in \(\mathcal S_n\) and
still has zero image.  By
\eqref{eq:uniform-S-injection} it is \(p\)-torsion,
and \eqref{eq:uniform-S-torsion} kills it.  Thus
\(r\in(\mathcal R_n)_{\Lambda\text{-tor}}\).
We have proved
\[
                    \mathcal F_{n+1}=\ker\rho_{\leq n}.
\]
Compactness gives the asserted isomorphism with the
actual image.

For the next induction step, the rational \(E_n\) is
the primitive companion extension of
\(\mathcal S_n[1/p]\) by the finite space
\((A_{n-1}/A_n)[1/p]\).
Its torsion-freeness and the total saturation of \(\mathcal S_n\)
give independent companion residues.  In the associative
algebra, the Eulerian idempotent decomposes the degree
module as this actual primitive extension plus the
nonprimitive old summand.  Both are torsion-free.
Thus all hypotheses of the weighted companion theorem
are verified for the enlarged alphabet.

We also close the integral lattice induction.
The quotient \(E_n/\operatorname{im}\mathcal S_n\) has bounded
\(p\)-torsion and finite free quotient.  Its bounded
torsion comes from the preceding kernel, already killed
by \(p^{h_{n-1}}\), and from the finite torsion in the
quotient of two finite-rank submodules of \(V\).
Choose \(p^{c_n}\) killing it and lift a basis of its
finite free quotient.  These finitely many lifts,
together with the old free Hall image, define
\(E_{n,0}\).  Since \(p^{d_n}\mathcal S_n\) lies in that Hall
image,
\[
                 p^{b_n}E_n\subseteq E_{n,0},
                 \qquad b_n=d_n+c_n.
\]
This proves the bounds required at the next stage.

Finally, \(Q_n\) is torsion-free because its finite
central filtration has the \(p\)-torsion-free factors
\(E_1,\ldots,E_n\).  It has class at most \(n\).
The integral pure Hall module
\(\operatorname{Lie}_n(W)\subseteq \mathcal S_n\) survives
in \(E_n\): its intersection with the torsion kernel
is zero.  This module is nonzero, since the positive
\(\Lambda\)-rank of \(W\) supplies two independent
\(\mathbb Z_p\)-coordinates and their length-\(n\)
Hall commutator.  Therefore the class is exactly \(n\).

\subsection{The full outer kernel and the inverse limit}

We justify the equality of the rational and pro-\(p\)
outer kernels used in the limit.  For the geometric
free pro-\(p\) group \(\Pi=\langle x,y\rangle\),
the tangential normalization is
\[
 x\longmapsto x^{\chi(\sigma)},\qquad
 y\longmapsto f_\sigma^{-1}y^{\chi(\sigma)}f_\sigma,
 \qquad f_\sigma\in[\Pi,\Pi].
\]
This normalized action on the cyclotomic kernel of \(G_{\mathbb Q}\) is given in
\cite[\S2, (2.1) and the following paragraph, author
p.~7]{NakamuraSchneps2000}. Here \(\chi=1\). The kernel comparison below also proves its descent to \(H\): an element trivial in the pro-\(p\) outer action has identity normalized action, so the kernel of the maximal extension unramified outside \(p\) acts trivially.

The map from \(\Pi\) to its rational prounipotent
completion is injective.  One way to verify this
directly is the Magnus map \(x\mapsto1+X\),
\(y\mapsto1+Y\).  Its reduction modulo \(p\)
detects every finite \(p\)-quotient: substitution
into the augmentation algebra of that quotient
is defined because its augmentation ideal is
nilpotent.  Hence its kernel in the free pro-\(p\)
group is trivial.  Inclusion of integral
coefficients into rational coefficients, followed
by the invertible formal exponential change of
variables, gives the stated rational injection.

Suppose a normalized cyclotomic action
is inner in the rational completion.  A conjugating
element fixing \(\exp X\) is \(\exp(aX)\):
the centralizer of \(X\) in the completed free
associative algebra is \(\mathbb Q_p[[X]]\),
by the following coefficient argument. In a homogeneous series commuting with \(X\), the coefficient of a word \(X^rYv\) equals that of \(YvX^r\) by successively moving the initial \(X\) to the end. Comparing \(XA=AX\) at the word beginning in \(Y\) and ending in \(X\) makes this coefficient zero. Thus every word containing \(Y\) has zero coefficient. Moreover primitivity in \(\mathbb Q_p[[X]]\) is the identity \(f(X_1+X_2)=f(X_1)+f(X_2)\); comparison of powers in characteristic zero leaves only its linear term. The primitive elements are precisely
\(\mathbb Q_pX\).
Use the convention that the inner action is \(z\mapsto g^{-1}zg\). Then \(g=\exp(aX)\), and the formula on \(y\) says that \(f_\sigma\exp(-aX)\) centralizes \(\exp Y\). Hence \(f_\sigma=\exp(bY)\exp(aX)\).
The abelian image of \(f_\sigma\) is zero;
abelianization forces \(a=b=0\).
The Magnus injection forces \(f_\sigma=1\)
in \(\Pi\).  Thus the rationally trivial
outer action is already trivial as a pro-\(p\)
outer action.  The reverse implication is
functorial, proving equality of the kernels.

The same conclusion holds when rational outer triviality is formulated at every finite geometric quotient. In the free Lie algebra truncated at degree \(N\), the centralizer of \(X\) is \(\mathbb Q_pX\) plus the top-degree center, by the same coefficient argument in degrees below \(N\). A conjugator therefore has logarithm \(aX\) modulo that center. The condition \(f_\sigma\in[\Pi,\Pi]\) makes the normalized action have no degree-two term on \(Y\), whereas this conjugator gives \(-a[X,Y]\). For \(N\geq2\), this forces \(a=0\), and the inner action is the identity. Passing through the quotients proves the assertion without choosing compatible conjugators.

The generator-length filtration of the
weight-completed free outer image is separated.
In any fixed geometric weight \(m\), all odd
generators have weight at least three, so a
nonzero word has length at most \(m/3\).
Vanishing at every length therefore means
vanishing at every finite weight.
Together with \(\mathcal F_{n+1}=\ker\rho_{\leq n}\),
this proves the first equality in the theorem.  The second is
the Galois correspondence for the decreasing
closed normal subgroups \(\mathcal F_n\).


\subsection{An exact linear obstruction at the inverse limit}
\label{sec:limit-relation-module}

The remaining intersection can be expressed by the actual
torsion modules and their quotient maps. Put
\[
 K=\bigcap_{n\geq1}\mathcal F_n=\ker\rho_{\mathrm{outer}},\qquad
 T_n=(\mathcal R_n)_{\Lambda\text{-tor}}
       =\mathcal F_{n+1}/J_{n+1},\qquad
 \mathcal L=\varprojlim_nT_n.
\]
The transition \(T_{n+1}\to T_n\) is induced by inclusion
\(\mathcal F_{n+2}\subseteq \mathcal F_{n+1}\), followed by reduction modulo
\(J_{n+1}\).

\begin{proposition}\label{prop:limit-relation-module}
There are canonical topological isomorphisms
\begin{align}
 K/\overline{[H,K]}&\simeq\mathcal L,
 \label{eq:limit-relation-iso}\\
 \mathcal L/(p,\gamma-1)\mathcal L
 &\simeq
 \varprojlim_n\bigl(T_n/(p,\gamma-1)T_n\bigr).
 \label{eq:limit-mod-p-iso}
\end{align}
Both modules on the right detect the triviality of \(K\):
\begin{equation}\label{eq:limit-linear-obstruction}
 K=1\quad\Longleftrightarrow\quad
 \mathcal L=0\quad\Longleftrightarrow\quad
 \varprojlim_n\bigl(T_n/(p,\gamma-1)T_n\bigr)=0.
\end{equation}
Writing \(A_n=\operatorname{im}(T_n\to M)\) and
\(P_n=\ker(T_n\to M)\), there is also an exact sequence
\begin{equation}\label{eq:limit-abelian-torsion}
 0\longrightarrow\varprojlim_nP_n
 \longrightarrow\mathcal L
 \longrightarrow\bigcap_nA_n
 \longrightarrow0,
 \qquad p^{h_n}P_n=0.
\end{equation}
\end{proposition}

\begin{proof}
For \(n\geq2\), the commutator definitions give
\[
 \overline{[H,K]}\subseteq J_n
 \subseteq\overline{[H,\mathcal F_{\lceil n/2\rceil}]}.
\]
In any finite quotient of \(H\), the descending images of
\(\mathcal F_m\) eventually stabilize at the image of \(K\).
Indeed every coset belonging to all these images has
nonempty nested compact intersections with the \(\mathcal F_m\).
Consequently the images of the two commutator bounds
eventually coincide. Passing through all finite quotients
proves
\[
                    \bigcap_nJ_n=\overline{[H,K]}.
\]
A compatible family in \(\varprojlim T_n\) is a nested family
of nonempty compact cosets of \(J_{n+1}\), each contained
in \(\mathcal F_{n+1}\). A point of their intersection belongs to
\(K\). The kernel of the resulting surjection
\(K\to\mathcal L\) is the displayed intersection of the
\(J_n\). Compactness proves
\eqref{eq:limit-relation-iso} as a topological isomorphism.

If \(\mathcal L=0\), then \(K=\overline{[H,K]}\).
Its image \(N\) in any finite \(p\)-group quotient \(G\)
of \(H\) satisfies \(N=[G,N]\). Iteration puts \(N\)
in every lower central term of the nilpotent group \(G\),
so \(N=1\). Every finite image of \(K\) is trivial,
which proves \(K=1\).

Inverse limits of these exact sequences of compact groups
are exact. For the surjectivity assertion, a compatible
family of target points can be lifted over any finite
initial set of indices by lifting first at the largest
index and pushing down. The resulting finite systems
of closed conditions have the finite intersection
property in a product of compact groups.
Apply this observation to the maps
\[
 T_n^2\longrightarrow T_n,\qquad
 (x,y)\longmapsto px+(\gamma-1)y,
\]
their compact images, and their cokernels. It proves
\eqref{eq:limit-mod-p-iso}.

For completeness, compact Nakayama in this setting does
not need finite generation. Open \(\Gamma\)-stable
subgroups form a neighborhood basis in \(\mathcal L\):
an open subgroup has an open stabilizer in the compact
group \(\Gamma\), and one intersects its finitely many
translates. On a finite quotient by such a subgroup,
\((p,\gamma-1)\) acts nilpotently, since \(\Gamma\)
has finite \(p\)-group image. Thus
\(\mathcal L=(p,\gamma-1)\mathcal L\) makes every
finite quotient zero. This proves
\eqref{eq:limit-linear-obstruction}.

Finally the maps \(T_n\to M\) have nested images \(A_n\).
The compact exact sequences
\(0\to P_n\to T_n\to A_n\to0\) give
\eqref{eq:limit-abelian-torsion}. The bound on \(P_n\)
is the finite-stage kernel bound already proved.
Equivalently, nested compact fibers give directly
\(\bigcap A_n=\operatorname{im}(K\to M)\).
\end{proof}

The exponents \(h_n\) are finite at each stage; no uniform
bound in \(n\) is asserted. Thus
\eqref{eq:limit-abelian-torsion} retains the inverse limit
of the inherited bounded torsion as well as the intersection
of the finite-rank lattices. The last module in
\eqref{eq:limit-linear-obstruction}, with its actual
transition maps, is the precise remaining linear
vanishing question.

\paragraph{The homology measured by the torsion limit.}
Put \(G=H/K=\rho_{\mathrm{outer}}(H)\), with its actual compact
quotient topology. The free pro-\(p\) presentation
\(1\to K\to H\to G\to1\) gives the continuous Hopf sequence
\[
0\longrightarrow H_2^{\mathrm{cont}}(G,\mathbb Z_p)
\longrightarrow K/\overline{[H,K]}
\longrightarrow H^{\mathrm{ab}}
\longrightarrow G^{\mathrm{ab}}\longrightarrow0.
\]
For the precise pro-\(p\) morphism, see
\cite[Proposition~2.2, author p.~3, and the Hopf formula on
author p.~4]{Moravec2011}. It also follows directly by
Pontryagin-dualizing the Hochschild--Serre five-term sequence
with discrete coefficients \(\mathbb Q_p/\mathbb Z_p\);
freeness of \(H\) makes its second cohomology zero.
This argument applies to the present countable source and
does not require a finite presentation of \(G\).
The kernel of the middle map is
\((K\cap\overline{[H,H]})/\overline{[H,K]}\).
Comparison with \eqref{eq:limit-abelian-torsion} therefore gives
the actual topological isomorphism
\[
 \varprojlim_nP_n
       \simeq H_2^{\mathrm{cont}}(G,\mathbb Z_p).
\]
The finite-stage bounds on \(P_n\) thus identify an exact
continuous homology term; they do not evaluate its inverse
limit.


\section{Even components in every finite quotient}
\label{sec:selected-all-stages}

This section defines the component quotient directly in the
nilpotent group, so that the construction applies at every
nilpotency class. Continue with the canonical groups
\(Q_n=H/\mathcal F_{n+1}\) of Theorem~\ref{thm:canonical-all-stages},
for \(n\geq2\), and put
\[
 B_n=\operatorname{im}(\ker(H\longrightarrow W)\longrightarrow Q_n),
 \qquad D_n=\operatorname{im}(\mathcal F_{n+1}\longrightarrow U).
\]
A tame complement gives an action of \(\Delta\) on \(Q_n\).
For an exponent \(e\), write
\[
 \mathcal E_{n,e}=
 \{g\in B_n:\delta g\delta^{-1}=g^{\omega(\delta)^e}
                         \text{ for every }\delta\in\Delta\}.
\]
The power is defined in the procyclic subgroup generated by \(g\).
For a set \(\Sigma\) of even exponents, define
\[
 J_{n,\Sigma}=
 \overline{\left\langle\!\left\langle
       \bigcup_{\substack{e\ \mathrm{even}\\e\notin\Sigma}}
                   \mathcal E_{n,e}
                   \right\rangle\!\right\rangle}_{Q_n},
 \qquad Q_{n,\Sigma}=Q_n/J_{n,\Sigma}.
\]
Every element of each displayed set is used in this definition.

\begin{lemma}\label{lem:tame-eigen-element-lift}
Let a finite group \(\Delta\) of order prime to \(p\) act
continuously on a nilpotent pro-\(p\) group \(B\), and let
\(B\twoheadrightarrow V\) be an equivariant homomorphism to
an abelian pro-\(p\) group. For a character
\(\chi:\Delta\to\mathbb Z_p^\times\) and
\(v\in V\) satisfying \(\delta v=\chi(\delta)v\),
there is \(g\in B\) above \(v\) satisfying
\(\delta g=g^{\chi(\delta)}\) for every \(\delta\).
\end{lemma}
\begin{proof}
Lift \(v\) to \(B^{\rm ab}\) and apply the integral projector
\[
 \epsilon_\chi=\frac1{|\Delta|}
               \sum_{\delta\in\Delta}\chi(\delta)^{-1}\delta.
\]
This gives an eigenclass above \(v\). We lift it successively
through the finite lower central series of \(B\).
At one central step, let \(C\) be the central abelian kernel,
and choose a lift \(u\) of the eigen-element already obtained.
Write the central errors additively as
\[
 \delta u\delta^{-1}=k_\delta u^{\chi(\delta)},
 \qquad k_\delta\in C.
\]
Put \(b_\delta=\chi(\delta)^{-1}k_\delta\) and give \(C\)
the action \(\delta\star c=\chi(\delta)^{-1}\delta(c)\).
Multiplication of two automorphisms gives
\[
 b_{\delta\eta}=b_\delta+\delta\star b_\eta.
\]
The finite average \(c=|\Delta|^{-1}\sum_\eta b_\eta\)
therefore satisfies \(b_\delta=c-\delta\star c\).
Since \(c\) is central, replacing \(u\) by \(cu\) changes
its error to
\[
 \delta(c)+k_\delta-\chi(\delta)c=0.
\]
Thus the corrected lift is an eigen-element at this step.
Division by \(|\Delta|\) is an automorphism of every abelian
pro-\(p\) group, including its torsion subgroup. The procedure
is consequently integral at each step. The lower central
series has finite length, so it ends with the required
\(g\in B\).
\end{proof}


\begin{proof}[Proof of Theorem~\ref{thm:selected-all-stages}]
Both tame complements lie in the inverse image of the finite subgroup \(\Delta\subset\Delta\times\Gamma\), whose kernel is \(Q_n\). Profinite Schur--Zassenhaus conjugates them by an element of \(Q_n\). Conjugation sends each full eigen-element set to the corresponding set for the other complement, proving independence of its closed normal closure. For an arithmetic element \(g\), choose \(h\in Q_n\) such that \(h^{-1}g\) normalizes the complement. The quotient \(\Delta\times\Gamma\) is abelian and the complement maps injectively to it, so \(h^{-1}g\) centralizes the complement. It preserves \(B_n\) and every eigen-element set with its character label. Normality handles the factor \(h\), proving full arithmetic stability. The map \(Q_n\to Q_{n-1}\) sends a
\(\Delta\)-eigen-element to a \(\Delta\)-eigen-element of
the same character. It induces a surjection
\(Q_{n,\Sigma}\to Q_{n-1,\Sigma}\).

Let \(V_n=\operatorname{im}(\mathcal F_{n+1}\to M)\). Then
\(Q_n^{\rm ab}=M/V_n\), and the image of \(B_n\) is the
image of \(D\oplus T\). Every bad even eigen-element maps
into its indicated bad even character. Conversely, apply
Lemma~\ref{lem:tame-eigen-element-lift} to the surjection
from \(B_n\) to this abelian image. Every class in a bad
even character has such a lift. Hence the abelian image of
\(J_{n,\Sigma}\) is exactly the sum of the bad even
components of \(T/(V_n)^+\), and contains no odd class.
Since \(\mathcal F_{n+1}\subseteq \mathcal F_3\) for \(n\geq2\),
\[
                   (V_n)^+\subseteq\pi_T(Z).
\]
The proved single-component cup theorem gives
\(Z^{[e]}=0\) for \(e\in\Sigma\). Therefore the selected
even part is exactly \(T_\Sigma\), and the odd part is
\(U/D_n\), proving the displayed abelianization.
Kernel containment follows from the canonical map through
\(Q_n\), and arithmetic stability gives the Galois assertion.

For exact class, choose distinct positive odd weights
\(m,k\geq3\) congruent to \(-1\) modulo \(p-1\).
Project the full free rational outer Lie algebra to the
free Lie algebra on its two letters of weights \(m,k\),
and truncate at generator length \(n\). This gives a
homomorphism from \(Q_n\) to its nilpotent exponential group.
Both letters have tame character \(-1\). Every even Lie
word of length at most \(n\) consequently has one of the
characters in the displayed condition. An element of
\(\mathcal E_{n,e}\) with bad even \(e\) has logarithm in
a missing eigenspace, so its image is zero. The homomorphism
therefore factors through \(Q_{n,\Sigma}\).

Choose \(w\in W^{[-1]}\) for which both linear moments
are nonzero. Such a choice follows from the explicit regular
\(\Lambda\)-coordinate of this character, after clearing
one scalar denominator. The two linear images of \(w\)
and \(\gamma w\) have determinant
\[
             \kappa_m(w)\kappa_k(w)(u^k-u^m)\ne0.
\]
Choose lifts to \(H\). The iterated commutator with initial
word \(\operatorname{ad}(w)^{n-2}[w,\gamma w]\) maps to a
nonzero length-\(n\) Hall word on two independent vectors.
The quotient thus has class at least \(n\); it has class
at most \(n\) as a quotient of \(Q_n\).

For \(n\leq5\), the group \(B_n\) has class at most two.
Its integral BCH logarithm identifies its group
\(\Delta\)-eigen-elements exactly with the Lie
\(\Delta\)-eigenspaces. The defining eigen-element sets are therefore the exponential images of exactly the earlier eigenspaces. Their closed normal closures are equal. This proves the
last identification through the specified maps.
\end{proof}

In particular, the thirty characters
\(\Sigma_{30}=\{-2m:1\leq m\leq30\}\) contain all the
characters required for every \(2\leq n\leq61\).
The exact class assertion gives the strict tower
\[
 L_{2,\Sigma_{30}}\subsetneq L_{3,\Sigma_{30}}
 \subsetneq\cdots\subsetneq L_{61,\Sigma_{30}}
 \subseteq\Omega_p^*.
\]
The arithmetic premise for this application is the posted
\(K\)-group vanishing statement, with the dependency analysis
of Section~\ref{sec:dependencies}. The component theorem
itself was proved for an arbitrary \(\Sigma\) using only
finiteness of its selected Iwasawa components.

\section{The limit of the central quotients}
\label{sec:source-limit}

\subsection{An exact nonzero limit in the source recursion}
\label{subsec:limit-countermodel}

The source recursion has a nontrivial limit even for
a free pro-$p$ group with a regular odd abelian part
and with no positive integral cyclotomic eigenvalue
in its finite abelian torsion. The following explicit
model identifies every subgroup and quotient map.

\begin{proposition}\label{prop:source-limit-countermodel}
For every odd $p$ and every $r,s\geq1$, there is a
free pro-$p$ group $H$ with continuous commuting
actions of $\Gamma\simeq\mathbf Z_p$ and an involution
$c$, such that
\[
 H^{\rm ab}=\Lambda^r\oplus\mathbf Z_p^s,
 \qquad c|_{\Lambda^r}=-1,
 \qquad\Gamma|_{\mathbf Z_p^s}=1.
\]
The signs of $c$ on the finite summand may be prescribed
independently. For the exact source recursion
\[
 \mathcal F_1=H,\qquad \mathcal F_2=\operatorname{preim}(H^{\rm ab}_{\Lambda\text{-tors}}),
 \qquad J_{n+1}=
 \overline{\prod_{i+j=n+1}[\mathcal F_i,\mathcal F_j]},
 \qquad \mathcal R_n=\mathcal F_n/J_{n+1},
\]
with $\mathcal F_{n+1}$ the inverse image of $(\mathcal R_n)_{\rm tor}$,
every torsion submodule is closed and
\[
                          \bigcap_{n\geq1}\mathcal F_n=N\ne1
\]
for an explicitly defined closed normal subgroup $N$.
For an infinite-order $u\in1+p\mathbf Z_p$, the finite
summand has no eigenvalue $u^m$ with $m\geq2$.
\end{proposition}

\begin{proof}
Let $F$ be the free pro-$p$ group on
$a_1,\ldots,a_r,t$, and let
\[
 F\twoheadrightarrow\Gamma,\qquad a_i\longmapsto1,
                         \quad t\longmapsto\gamma.
\]
Put $P=\ker(F\to\Gamma)$. It is free pro-$p$ by
pro-$p$ Nielsen--Schreier. Conjugation by $t$ gives
the stated continuous $\Gamma$-action on $P$.
The completed free-group augmentation resolution,
with differential given by $a_i-1,t-1$, becomes
\[
 \Lambda^{r+1}\longrightarrow\Lambda,\qquad
 (v_1,\ldots,v_r,w)\longmapsto(\gamma-1)w
\]
after applying the quotient to $\Gamma$.
The kernel identifies with $P^{\rm ab}$, by the completed
Schreier, or equivalently Shapiro, identification.
Since $\gamma-1$ is not a zero divisor, its kernel
is exactly $\Lambda^r$. The actual isomorphism sends
\[
       \gamma^j e_i\longmapsto[t^j a_i t^{-j}]
                                  \quad\text{in }P^{\rm ab}.
\]
Thus the module map and its $\Gamma$ action are fixed,
not merely their ranks.

The automorphism of $F$ fixing $t$ and sending every
$a_i$ to $a_i^{-1}$ restricts to an involution $c$
on $P$, commuting with $\Gamma$ and acting as $-1$
on $P^{\rm ab}$. Let $C$ be free pro-$p$ on
$z_1,\ldots,z_s$ and put $H=P\mathbin{*}_{p}C$.
Let $\Gamma$ fix the $z_i$, and prescribe
$c(z_i)=z_i$ or $z_i^{-1}$ independently.
These assignments define the required actions.
One can also define a full tame action by fixing $t$
and using the chosen Teichmuller powers on the $a_i$
and $z_i$.

Let $N$ be the kernel of the retraction
$H\twoheadrightarrow P$ that kills all $z_i$;
equivalently, it is their closed $H$-normal closure.
The quotient imposing centrality of those generators
is exactly
\[
                     H/[H,N]\simeq P\times\mathbf Z_p^s.
\]
In particular $N/[H,N]=\mathbf Z_p^s$ with trivial
$\Gamma$ action. We claim, for every $n\geq2$,
\begin{align*}
 \mathcal F_n&=N\gamma_nP,\\
 J_{n+1}&=[H,N]\gamma_{n+1}P,\\
 \mathcal R_n&\simeq\mathbf Z_p^s\oplus
                              (\gamma_nP/\gamma_{n+1}P),
\end{align*}
where lower-central terms are closed. The first
formula starts with the literal inverse image of
the finite torsion summand. Inductively, the factor
$[H,\mathcal F_n]$ contains $[H,N]$ and
$[P,\gamma_nP]=\gamma_{n+1}P$. Every other factor
of $J_{n+1}$ lies in their product. This proves the
second formula, and the product quotient just
displayed proves the third.

The second summand of $\mathcal R_n$ is the completed free
Lie piece $\operatorname{Lie}_n(\Lambda^r)$.
Its Hall embedding in the completed tensor module
is integral. On each ordered color coordinate that
tensor module is $\mathbf Z_p[[\Gamma^n]]$ and a
nonzero $f(\gamma)$ acts by the nonzero multiplier
$f(z_1\cdots z_n)$ in a domain. Hence the Lie piece
is $\Lambda$-torsionfree. The torsion of $\mathcal R_n$ is
therefore precisely the displayed $\mathbf Z_p^s$,
which is killed by $\gamma-1$ and is closed. Its
inverse image is $N\gamma_{n+1}P$, completing the
induction. The exact finite-stage quotients are
\[
                         H/\mathcal F_{n+1}=P/\gamma_{n+1}P.
\]

Finally $\bigcap_n\gamma_nP=1$ by an exact finite
quotient argument. Every finite $p$-group quotient
of $P$ is nilpotent, so its quotient map kills
$\gamma_nP$ once $n$ exceeds its nilpotency class.
An element in every closed $\gamma_nP$ therefore
lies in every open normal subgroup of $P$, hence is
the identity. The retraction to $P$ now gives
$\bigcap_n\mathcal F_n=N$, and $N\ne1$ since it contains
each nontrivial free generator $z_i$.
The only finite abelian eigenvalue is $1$, which
differs from every $u^m$ for $m\geq2$.
\end{proof}

This is an explicit source model; no arithmetic
realization is asserted. Its quadratic companion
extension is split in the persistent even directions.
It proves that source freeness, the regular odd
module, and the positive-eigenvalue exclusion alone
cannot make the recursion's limit trivial.

% Standalone section; requires amsmath, amssymb, amsthm and the usual
% theorem/proposition/lemma environments.  It uses no private macros.
% Additional bibliography key needed: MS (entry below).
% All arithmetic assertions in this section are unconditional at p=37.

\section{The component quotient and its Galois action}
\label{sec:component-quotient-audit}

Put $F_1=\mathbb Q(\mu_p)$, $F_\infty=\mathbb Q(\mu_{p^\infty})$,
and let $\Omega_p$ be the maximal pro-$p$ extension of $F_1$
unramified outside $p$.  Write
\[
 P=\operatorname{Gal}(\Omega_p/F_1),\qquad
 H=\operatorname{Gal}(\Omega_p/F_\infty),\qquad
 \Gamma=\operatorname{Gal}(F_\infty/F_1),\qquad
 \Lambda=\mathbb Z_p[[\Gamma]].
\]
We use closed commutator and normal subgroups throughout.  Choose a
lift of $\Delta=\operatorname{Gal}(F_1/\mathbb Q)$ by the profinite
Schur--Zassenhaus theorem.  Coprime fixed-point lifting gives a
$\Delta$-fixed element of $P$ mapping to a generator of $\Gamma$;
its procyclic closure is a section.  Thus we may use commuting
lifts and write $P=H\rtimes\Gamma$.  Complex conjugation is denoted
by $c$.  For $M=H^{\rm ab}$ write $U=M^-$ and $T=M^+$.

\subsection{An explicit calculation with arbitrary lifts}

The choice of an abelian component does not specify a subgroup to be
killed in a nonabelian group.  The following calculation gives the
exact maps behind this issue.  Let $F$ be free pro-$p$ on $x,y,t$, where
$p$ is odd, and define
\[
 c(x)=x^{-1},\qquad c(y)=y^{-1},\qquad c(t)=t.
\]
Thus $F^{\rm ab}=\mathbb Z_p x\oplus\mathbb Z_p y\oplus\mathbb Z_p t$
has two negative summands and one positive summand.  With
$[x,y]=x^{-1}y^{-1}xy$, both
\[
 a\longmapsto t^a,\qquad
 a\longmapsto(t[x,y])^a
\]
lift the same homomorphism $\mathbb Z_p\to F^{\rm ab}$.
Their normal closures are the respective kernels of the retractions
to the free pro-$p$ group on $x,y$ which send $t$ to $1$ and
$[x,y]^{-1}$.  These are explicit identifications of both quotient
groups, with their different quotient maps.

Let $E_{ij}$ denote a matrix unit and set, in
$\operatorname{UT}_4(\mathbb F_p)$,
\[
 X=I+E_{12},\qquad Y=I+E_{23}+E_{34},\qquad Z=I-E_{13}.
\]
Matrix multiplication gives
\[
 [X,Y]=I+E_{13},\qquad
 [X^{-1},Y^{-1}]=I+E_{13}-E_{14}.
\]
The homomorphism $x\mapsto X$, $y\mapsto Y$, $t\mapsto Z$
therefore kills $t[x,y]$ and sends $c(t[x,y])$ to $I-E_{14}$.
Consequently
\begin{equation}\label{eq:component-raw-lift}
 c\bigl(\langle t[x,y]\rangle_F^{\rm normal}\bigr)
 \ne\langle t[x,y]\rangle_F^{\rm normal}.
\end{equation}
This is a finite quotient certificate for the failure of invariance.

Exact fixed lifts also leave choices.  In the free group on $x,y$ put
\[
 h=[x,y],\qquad d=h^{-1}c(h),\qquad w=h d^{1/2}.
\]
Unit powers are uniquely defined in a pro-$p$ group.  In particular
$c(d)=d^{-1}$, so $c(w)=w$.  The class of $w$ modulo the third lower
central subgroup is the nonzero class of $[x,y]$.  Its nonzero image
is also visible in $\operatorname{UT}_3(\mathbb F_p)$.
Thus $t$ and $tw$ are exact $c$-fixed lifts of the same abelian
generator with distinct $c$-stable normal closures: the quotient
killing $t$ sends $tw$ to $w\ne1$.

\subsection{The actual arithmetic group at $p=37$}

We first give all arithmetic inputs used below.  Iwasawa's
Theorem~2, together with Ferrero--Washington, makes $H$ a free
pro-$37$ group
\cite[Theorem~2, pp.~273--274]{Iwasawa,FW}.
The maximal extension in that theorem agrees with $\Omega_{37}$:
a finite $37$-extension over $F_\infty$ descends to a finite
cyclotomic layer, and its normal closure over $F_1$ remains a
$37$-extension unramified outside $37$.  The group $H$ is
second countable as a closed subgroup of the finitely generated
pro-$37$ group $P$.

The cyclotomic class-group calculation at $37$ gives
\begin{equation}\label{eq:component-p37-class-input}
 X^+=0,\qquad
 X^-=X^{[5]}\simeq\mathbb Z_{37}\quad
 \text{as a $\mathbb Z_{37}$-module}.
\end{equation}
Here the brackets denote the Teichm\"uller character exponent.
One published verification is
\cite[\S1, items~(2)--(3)]{HHO}: the $37$-class group
of $\mathbb Q(\mu_{37^r})$ is cyclic of order $37^r$ for every
$r\ge1$.  Its character is $5=1-32$ modulo $36$.
The direct limit $A_\infty^-$ under extension of ideals is therefore
$(\mathbb Q_{37}/\mathbb Z_{37})^{[5]}$.
For example, after choosing compatible generators for the
surjective norm maps, the extension maps at all sufficiently high
levels are multiplication by $37$ times a unit, since the composite
of norm and extension is multiplication by $37$.
The exact Kummer duality map
\[
 M^+\simeq(A_\infty^-)^\vee(1)
\]
then gives
\begin{equation}\label{eq:component-p37-even}
 T=T_{32}\simeq\mathbb Z_{37},\qquad T_e=0\quad(e\ne32).
\end{equation}
The character transport is explicit: on the dual functional
$f$ the action is $(\delta f)(b)=\omega(\delta)f(\delta^{-1}b)$,
which sends exponent $5$ to $1-5=32$ modulo $36$.
For the Kummer map see
\cite[Theorem~2.3.18]{SharifiCyclotomicNotes}.

The integral odd module is
\begin{equation}\label{eq:component-p37-odd}
 U\simeq\Lambda^{18}.
\end{equation}
We include the two exceptional character calculations in the proof
of this identity.  For an odd exponent $a\ne1,-1$, apply
\cite[Proposition~2.2.6]{WWEOrd} with
$\theta=\omega^{1-a}$.  The proposition identifies $X_\theta=0$
with integral rank-one $\Lambda$-freeness of
$M^{[a],\#}(1)$.  The inversion and twist transport the ring by
$\gamma\mapsto\kappa(\gamma)\gamma^{-1}$, an automorphism; undoing
them preserves freeness and rank.  This proves the assertion in
sixteen of the eighteen odd characters.  The article's exclusions
of $\theta=1$ and $\theta=\omega^2$ are not used for the other
two characters.

Let $U_{\rm loc}$ be the inverse norm limit of local pro-$37$ units
and $E$ the image of the inverse limit of global units.  The actual
reciprocity maps give
\[
 0\longrightarrow E\longrightarrow U_{\rm loc}
 \longrightarrow M\longrightarrow X\longrightarrow0.
\]
The global odd units are exactly $E^-=\mathbb Z_{37}(1)$.
Indeed $v/c(v)$ is a root of unity for a global unit $v$, and
$2$ is invertible on its pro-$37$ completion.  The root-of-unity
inclusion and the first-moment map in Coleman's sequence are
\begin{equation}\label{eq:component-coleman}
 0\longrightarrow\mathbb Z_{37}(1)\longrightarrow U_{\rm loc}
 \xrightarrow{\operatorname{Col}}
 \mathbb Z_{37}[[\Delta\times\Gamma]]
 \xrightarrow{\lambda\mapsto\int x\,d\lambda(x)}
 \mathbb Z_{37}(1)\longrightarrow0.
\end{equation}
See \cite[Theorem~5.5.31]{SharifiIwa}; the prime-to-$37$
root term in the full sequence disappears on taking pro-$37$
units.  The class modules $X^{[-1]}$ and $X^{[1]}$ vanish, also
directly from \eqref{eq:component-p37-class-input}.  In character
$-1$, both ends of \eqref{eq:component-coleman} vanish and $E^{[-1]}=0$,
so reciprocity identifies $M^{[-1]}$ with $\Lambda$.  In character
$1$, reciprocity and \eqref{eq:component-coleman} identify
\[
 M^{[1]}\simeq U_{\rm loc}^{[1]}/\mathbb Z_{37}(1)
 \simeq\ker\bigl(\Lambda\longrightarrow\mathbb Z_{37}(1)\bigr)
 =(\gamma-u)\Lambda\simeq\Lambda,
 \qquad u=\kappa(\gamma).
\]
The kernel equality is division by the monic linear series
$T-(u-1)$, with $T=\gamma-1$.  Multiplication by $\gamma-u$ is
the displayed isomorphism from $\Lambda$ onto that ideal.  This
finishes the proof of \eqref{eq:component-p37-odd}.

The other input is a cup product at the first finite layer itself.
McCallum--Sharifi's calculation supplies $c$-even Kummer classes
\begin{equation}\label{eq:component-nonzero-cup}
 f,g\in H^1(P,\mathbb F_{37}(1)),\qquad
 f\smile g\ne0\quad\text{in }H^2(P,\mathbb F_{37}(2)).
\end{equation}
Here $P$ is the maximal pro-$37$ Galois group of
$F_1=\mathbb Q(\mu_{37})$ unramified outside $37$.
In their notation one can use the cyclotomic units $\eta_7$
and $\eta_{25}$: Corollary~8.3 has a nonzero coefficient of
$[x_7,x_{25}]$, and their duality formula identifies that
coefficient with this cup product
\cite[Theorem~7.5, p.~20; Corollary~8.3, p.~23]{MS}.
The unit characters are $\omega^{1-7}$ and $\omega^{1-25}$,
which are even.  The cup target has character $\omega^{2-32}$.
The degree-one classes for the full group $G_{F_1,S}$ factor
through its maximal pro-$37$ quotient $P$.  Naturality of cup
products under inflation shows that their nonzero cup for
$G_{F_1,S}$ comes from a nonzero cup for $P$; no surjectivity
assertion for degree-two inflation is required.
The two coefficient modules in \eqref{eq:component-nonzero-cup}
are trivial as $P$-modules because $\mu_{37}\subset F_1$.
Their $c$-actions are respectively $-1$ and $+1$.

\begin{theorem}\label{thm:component-no-galois-p37}
At $p=37$, let $t\in H$ be any lift of a
$\mathbb Z_{37}$-generator of $T=T_{32}$.  Its closed normal
closure $N=\langle t\rangle_H^{\rm normal}$ is not
$\Gamma$-stable.  The field $\Omega_{37}^N$ is therefore not
Galois over $F_1$, and hence is not Galois over $\mathbb Q$.
This applies in particular to every exact $\Delta$-eigenlift of
that generator.
\end{theorem}

\begin{proof}
The image of $t$ in $H/\Phi(H)=M/37M$ is nonzero, since it
generates the direct summand $T/37T$.  It is primitive in the free
pro-$37$ group $H$, so $H/N$ is free pro-$37$.
Here primitivity can be checked directly in countable rank.
Choose a convergent free basis $b_1,b_2,\ldots$ of $H$ and a
coordinate in which $t$ has nonzero Frattini coefficient, calling
it the first coordinate.  Replace $b_1$ by $t$ and leave the
others fixed.  The resulting endomorphism is an isomorphism on
the Frattini quotient, so its closed image is $H$: a proper
closed subgroup lies in an open subgroup of index $37$.
It preserves the normal subgroup $K_n$ generated by the tail
$b_j$, $j>n$, and its map on $H/K_n$ is a surjective
endomorphism of a finitely generated free pro-$37$ group, hence
an automorphism.  The groups $K_n$ have trivial intersection,
because any finite continuous quotient kills all but finitely
many elements of a convergent basis.  Thus the original
endomorphism is injective as well.  This verifies the free
factor containing $t$ and the asserted freeness of $H/N$.

The image of $N$ in $M$ is precisely $T$, whence
\begin{equation}\label{eq:component-quotient-ab}
 (H/N)^{\rm ab}=M/T=U\simeq\Lambda^{18}.
\end{equation}
Suppose for contradiction that $N$ is $\Gamma$-stable.  Then
$Q=P/N$ is defined, with the induced split extension
\[
 1\longrightarrow H/N\longrightarrow Q\longrightarrow\Gamma
 \longrightarrow1.
\]
Take $A=\mathbb F_{37}(2)$, which is a trivial $Q$-module.
Both $H/N$ and $\Gamma$ have $37$-cohomological dimension at
most one.  The Hochschild--Serre spectral sequence therefore
gives the actual isomorphism
\begin{equation}\label{eq:component-h2-zero}
 H^2(Q,A)\simeq
 H^1\bigl(\Gamma,\operatorname{Hom}_{\rm cts}(U,A)\bigr)=0.
\end{equation}
To verify the last equality including its topology, put
$V=\operatorname{Hom}_{\rm cts}(U,A)$, a discrete module.
Procyclic cohomology gives $H^1(\Gamma,V)=V/(\gamma-1)V$.
By \eqref{eq:component-quotient-ab}, the compact dual is
$U/37U\simeq\mathbb F_{37}[[T]]^{18}$.
The dual of $\gamma-1$ on $V$ is multiplication by
$\gamma^{-1}-1=-\gamma^{-1}T$ on this module.  It is injective,
so continuous Pontryagin duality makes $\gamma-1$ surjective
on $V$.  The Tate twist introduces no additional scalar here:
$\kappa(\gamma)^2=1$ in $\mathbb F_{37}^{\times}$.

Now take the two classes in \eqref{eq:component-nonzero-cup}.
Since $P$ acts trivially on $\mathbb F_{37}(1)$, their cocycles
are homomorphisms.  Being $c$-even means, with the coefficient
action included,
\[
 f(chc^{-1})=-f(h),\qquad g(chc^{-1})=-g(h)
 \quad(h\in P).
\]
The restrictions of $f$ and $g$ to $H$ factor through $M$.
The abelian image of $t$ lies in $T=M^+$, so $ctc^{-1}$ and $t$
have the same abelian image.  The displayed identities therefore
give $2f(t)=2g(t)=0$, and hence $f(t)=g(t)=0$.
Each homomorphism kills the whole $H$-normal
closure $N$.  They consequently descend to classes
$\bar f,\bar g\in H^1(Q,\mathbb F_{37}(1))$.
Naturality of cup products gives
\[
 f\smile g=\operatorname{inf}_{Q}^{P}(\bar f\smile\bar g)=0
\]
by \eqref{eq:component-h2-zero}, contradicting
\eqref{eq:component-nonzero-cup}.  Thus $N$ is not
$\Gamma$-stable.  The statements about fixed fields are now
exactly the normality criterion in Galois correspondence.
\end{proof}

Exact $c$-fixed lifts exist: for any lift $b$ of a positive abelian class,
$b(b^{-1}c(b))^{1/2}$ is $c$-fixed with the same abelian image.
The theorem proves a failure of invariance for every primitive
lift, not just the particular lift in
\eqref{eq:component-raw-lift}.  For a component prescription which
kills $T_{32}$, this is the entire nonzero even part at $37$.

\subsection{Ambient saturation and the precise remaining map}

Let $\mathcal G=\operatorname{Gal}(\Omega_p/\mathbb Q)$,
$\pi:H\to M$, and let $A\subset M$ be a closed
$\Delta\times\Gamma$-stable sum of components.  From a set of
lifts with abelian image generating $A$, form its $H$-normal
closure $N$ and then its ambient saturation
\[
 N^{\mathcal G}=
 \overline{\langle gng^{-1}:g\in\mathcal G,\ n\in N\rangle}.
\]
This is a closed $\mathcal G$-normal subgroup of $H$, and
\[
 \pi(N^{\mathcal G})=A,\qquad
 (H/N^{\mathcal G})^{\rm ab}=M/A.
\]
The equality of images follows by taking the closed sum of all
$\mathcal G$-translates of $A$.  Thus saturation repairs the
Galois action while preserving the specified abelian quotient.

\begin{corollary}\label{cor:component-saturation-not-free}
At $p=37$, let $S\subset H$ be any closed $P$-normal subgroup
whose image in $M$ is precisely $T$.  Then
\[
 H^2(P/S,\mathbb F_{37}(2))\ne0,
 \qquad H/S\text{ is not a free pro-}37\text{ group}.
\]
In particular this holds for the full ambient saturation of
the normal closure of any lift of a generator of $T_{32}$.
\end{corollary}

\begin{proof}
The two Kummer homomorphisms in
\eqref{eq:component-nonzero-cup} vanish on $S$, since their
restrictions to $H$ factor through odd functionals on $M$.
They descend to $P/S$, and their descended cup has nonzero
inflation.  It is therefore nonzero.  On the other hand,
$(H/S)^{\rm ab}=M/T=\Lambda^{18}$.
If $H/S$ were free pro-$37$, the Hochschild--Serre calculation
\eqref{eq:component-h2-zero}, with $S$ in place of $N$, would
make this same degree-two cohomology zero.  This contradiction
proves the second assertion.  Ambient saturation has image
$T$ by the preceding image calculation, so it is included.
\end{proof}

Put $K=\operatorname{Gal}(\Omega_p/\Omega_p^*)$ for the actual
outer-action kernel.  Galois correspondence gives the precise
remaining requirement
\[
 \Omega_p^{N^{\mathcal G}}\subseteq\Omega_p^*
 \quad\Longleftrightarrow\quad K\subseteq N^{\mathcal G}.
\]
The automatic intersection construction is
\[
 \Omega_p^{N^{\mathcal G}K}
   =\Omega_p^{N^{\mathcal G}}\cap\Omega_p^*,
 \qquad
 (H/(N^{\mathcal G}K))^{\rm ab}
   =M/(A+\pi(K)).
\]
These formulas describe the exact maps and retain the needed
nonabelian kernel containment as a separate assertion.
The next construction proves that injection on a retained free
factor supplies no such containment, even after full saturation.

\subsection{An operator-group counterexample after full saturation}

This subsection concerns the explicit groups constructed here.
The arithmetic result is Theorem~\ref{thm:component-no-galois-p37}.
Fix $p=37$, an element $\alpha\in1+37\mathbb Z_{37}$, a
cyclic group $\Delta=\langle\delta\rangle$ of order $36$, and
its faithful Teichm\"uller character $\omega$.

We use the following exact eigenprojection.  For a pro-$37$ group
with this $\Delta$-action, put $\beta=\omega(\delta)^e$ and
\[
 \mathcal E_e(v)=
 \left(\prod_{j=0}^{35}\delta^j(v)^{\beta^{-j}}\right)^{1/36},
 \qquad
 v^{(e)}=\lim_{n\to\infty}\mathcal E_e^n(v).
\]
The ordered product uses increasing $j$.
For completeness, if $\delta(v)v^{-\beta}$ belongs to the
$m$th closed lower central subgroup, its conjugates commute
modulo the $(m+1)$st subgroup.  In that quotient the displayed
averaging is the ordinary character idempotent.  It changes $v$
only in degree $m$ and makes the eigen-relation hold in degree
$m+1$.  Repeated averaging therefore converges in every finite
$37$-group quotient, since such a quotient is nilpotent.
It yields $\delta(v^{(e)})=(v^{(e)})^\beta$.
Products, unit powers and limits show that this construction is
functorial for continuous $\Delta$-equivariant homomorphisms.
This is also the eigenprojection of
\cite[Lemma~4]{Sharifi}.

Let $B$ be free pro-$37$ on $\gamma$ and $x_a$ for all odd
residues $a$ modulo $36$.  Map $B$ onto $\Gamma\simeq\mathbb Z_{37}$
by $\gamma\mapsto\gamma$ and $x_a\mapsto1$, and put
$J=\ker(B\to\Gamma)$.  Set
\[
 \delta(\gamma)=\gamma,\qquad
 \delta(x_a)=x_a^{\omega(\delta)^a}.
\]
These formulas define an automorphism of order $36$.
The group $J$ is free pro-$37$, and its abelianization has the
operator-module description
\[
 J^{\rm ab}\simeq\bigoplus_{a\ \mathrm{odd}}\Lambda(\omega^a).
\]
Here are the two maps proving this description.  The normal
generators $x_a$ give a surjection $\Lambda^{18}\to J^{\rm ab}$.
The homomorphism $B\to\Lambda^{18}\rtimes\Gamma$ sending
$x_a$ to the corresponding standard vector gives a map in the
opposite direction.  Their composite fixes all standard vectors,
so the initial surjection is an isomorphism.  This also proves
the stated $\Delta$-action.

Form the free pro-$37$ group $\mathcal F=J*\langle z\rangle$
and extend the commuting operators by
\[
 \delta(z)=z^{\omega(\delta)^{32}},\qquad
 \gamma^b(z)=z^{\alpha^b}\quad(b\in\mathbb Z_{37}).
\]
On $J$, the $\Gamma$-action is the conjugation action from $B$.
Thus
\[
 \mathcal F^{\rm ab}=U_0\oplus T_0,
 \qquad
 U_0=\bigoplus_{a\ \mathrm{odd}}\Lambda(\omega^a),\qquad
 T_0=\Lambda/(\gamma-\alpha)
\]
with $\Delta$-character $\omega^{32}$ on $T_0$.
Choose
\[
 w=[x_3,x_{29}]^{(32)}\in J,
 \qquad t=(zw)^{(32)}\in\mathcal F.
\]
Both $z$ and $t$ are exact $\Delta$-eigenlifts of the same
generator of $T_0$.  The retraction
\[
 \rho:\mathcal F\longrightarrow J,\qquad
 \rho|_J=\operatorname{id},\quad\rho(z)=1
\]
is $\Gamma\times\Delta$-equivariant.  It is injective on the
retained free factor $J$, and $\rho(t)=w$.
Let $\mathcal N$ be the closed normal closure of $t$ in
$\mathcal F\rtimes(\Gamma\times\Delta)$.

In $\operatorname{UT}_3(\mathbb F_{37})$ put
\[
 X=I+E_{12},\qquad Y=I+E_{23},\qquad
 Z=[X,Y]=I+E_{13}.
\]
Map $B$ to this group by
\[
 \gamma\mapsto1,\quad x_3\mapsto X,\quad x_{29}\mapsto Y,
 \quad x_a\mapsto1\quad(a\ne3,29),
\]
and restrict the map to $J$.  Extend it to
$q:\mathcal F\to\operatorname{UT}_3(\mathbb F_{37})$ by
$q(z)=Z^{-1}$.  The map is $\Gamma$-invariant: this is true
on $J$ by its construction, and on $z$ because
$\alpha^b\equiv1\pmod{37}$.  It is $\Delta$-equivariant for
the Heisenberg automorphism that scales $X,Y,Z$ with respective
characters $\omega^3,\omega^{29},\omega^{32}$.
Consequently $\ker q$ is stable under all the ambient operators.

The eigenprojection is already exact on the central element $Z$,
so functoriality gives
\[
 q(w)=Z,\qquad q(t)=(Z^{-1}Z)^{(32)}=1.
\]
Hence $\mathcal N\subseteq\ker q$.  The element $z$ supplies
the required certificate:
\begin{equation}\label{eq:component-saturation-kernel-failure}
 z\in\ker\rho,\qquad q(z)=Z^{-1}\ne1,
 \qquad z\notin\mathcal N.
\end{equation}
Thus $\ker\rho\nsubseteq\mathcal N$ after full ambient
saturation, with all displayed abelian modules, eigenlifts and
operator actions fixed.  The quotient
$(\mathcal F/\mathcal N)^{\rm ab}$ is nevertheless $U_0$,
by the saturation image calculation above.
Equation~\eqref{eq:component-saturation-kernel-failure} also
rules out a factorization of the quotient map through $\rho$:
its evaluation on $z$ would send $1$ to a nonidentity element.

% Suggested additional bibliography entries (do not duplicate a main-file entry):
% \bibitem{MS}
% W. G. McCallum and R. T. Sharifi,
% \emph{A cup product in the Galois cohomology of number fields},
% Duke Math. J. 120 (2003), 269--310.
% Author PDF: https://www.math.ucla.edu/~sharifi/publ/pairing.pdf
% Exact page locators above refer to this author PDF (dated December 5, 2003).
%
% \bibitem{SharifiCyclotomicNotes}
% R. T. Sharifi, \emph{Modular curves and cyclotomic fields},
% Theorem 2.3.18, printed p.19.
% https://www.math.ucla.edu/~sharifi/misc/aws-iwasawa.pdf

% Standalone section: theorem/proof environments and amsmath required.
% Bibliography keys: MagnusKoszul, LalondeRam (entries suggested at end).

\section{The intersection over all lift choices}
\label{sec:all-lift-intersection}

We now determine the exact strength of requiring the field inclusion
for every unrestricted choice of lifts.  No invariance under
$\Delta$ or $\Gamma$ is imposed on the individual normal closures
in this section.

\begin{theorem}\label{thm:all-lift-intersection}
Let $P$ be a free pro-$p$ group on a finite or countably infinite
convergent free basis, of rank at least two, and put
\[
 F=P*F(t_1,\ldots,t_r),\qquad r<\infty.
\]
For $\boldsymbol w=(w_1,\ldots,w_r)\in[P,P]^r$, let
$q_{\boldsymbol w}:F\to P$ fix $P$ and send $t_j$ to $w_j$.
Then
\[
 \bigcap_{\boldsymbol w\in[P,P]^r}\ker q_{\boldsymbol w}=1,
 \qquad
 \ker q_{\boldsymbol w}
   =\overline{\langle t_jw_j^{-1}:1\leq j\leq r\rangle}^{\rm normal}_F.
\]
Each $t_jw_j^{-1}$ has the same image as $t_j$ in $F^{\rm ab}$.
For every nonidentity element the proof supplies a separating
homomorphism to a finite $p$-group which factors through one of
these retractions.
\end{theorem}

\begin{proof}
The assertion for $r=0$ is immediate.  First assume that $P$ has
finite rank $d\geq2$, with generators $x_1,\ldots,x_d$.
Use the Magnus embedding
\[
 \mu:F\hookrightarrow
 1+(X_1,\ldots,X_d,T_1,\ldots,T_r)
 \subset\mathbb F_p\langle\!\langle X_1,\ldots,X_d,T_1,\ldots,T_r
                       \rangle\!\rangle,
\]
where $x_i\mapsto1+X_i$ and $t_j\mapsto1+T_j$.
The completed group-algebra identification and its injectivity
are recalled in \cite[\S3.1, equation~(3.3), p.~8]{MagnusKoszul}.

Fix $h\ne1$.  If the retraction $t_j\mapsto1$ does not kill $h$,
it already separates $h$.  Otherwise every monomial of
$\mu(h)-1$ contains a $T$-letter.  Let $k\geq1$ be the smallest
number of $T$-letters occurring in a monomial with nonzero
coefficient.  Among monomials with exactly $k$ such letters,
let $\ell$ be the smallest number of $X$-letters.
Their sum is a nonzero polynomial $Q$ over $\mathbb F_p$.
Every word of $Q$ has exactly $k$ $T$-letters and $\ell$
$X$-letters.

Choose two of the retained variables and call them $X,Y$,
ordered by $X<Y$.  Choose an integer $D>\ell$ and, for
$1\leq j\leq r$, put
\[
 u_j=X Y^D X Y^{D+j} X Y^{D+r+1-j}.
\]
All these words have the same length
$L=3D+r+4>\ell$.
Each is a Lyndon word.  Indeed a rotation beginning with $Y$
is larger.  A rotation beginning at either later occurrence
of $X$ has an initial $Y$-run longer than $D$, so it is larger
than the rotation beginning at the first $X$.
Let $W_j$ be its integral standard Lyndon bracketing reduced
modulo $p$.  The triangular expansion is
\[
 W_j=u_j+\sum_{v>u_j,\ |v|=L}a_vv.
\]
In particular the coefficient of $u_j$ is one in every
characteristic; see \cite[\S1, equation~(1.2)]{LalondeRam}.
Replacing the Lie brackets by group commutators gives
$w_j\in[P,P]$ whose Magnus expansion begins with $W_j$ in
degree $L$.

We check the word substitution $T_j\mapsto u_j$ precisely.
In any substituted word from $Q$, the occurrences of the
signature
\[
                             X Y^D X
\]
are exactly the beginnings of the inserted markers $u_j$.
Inside a marker its first $Y$-run has length $D$ and its
other two $Y$-runs have lengths strictly larger than $D$.
Between markers the total number of original $X_i$-letters
is $\ell<D$, so those letters cannot form a $Y$-run of
length $D$.  A final $Y$-run of a marker can merge with
following original $Y$-letters, but its length is already
larger than $D$ and remains so.  The initial $X$ of the
next marker prevents merger with its first $Y$-run.
These observations also rule out a signature crossing a
marker boundary.

The signature therefore determines every marker's starting
position.  The length of its middle $Y$-run is $D+j$, so it
determines $j$; the common length $L$ then determines its
ending position, even when its final run merges with
following original letters.  Removing the markers recovers
the original word of $Q$.  Thus distinct words of $Q$ have
distinct substituted words.

Choose the lexicographically least substituted word among
those with a nonzero coefficient in $Q$.
The triangular expansions of the $W_j$ show that it occurs
with that same coefficient in $Q(X_1,\ldots,X_d,W_1,\ldots,W_r)$.
Every other contribution is larger: within one substitution
this follows from the equal lengths of the terms of $W_j$,
and between different words it follows from the choice of
the least substituted word.  Therefore this polynomial is
nonzero.

In the full substituted series, its degree is $kL+\ell$.
A monomial with more than $k$ $T$-letters has degree at least
$(k+1)L>kL+\ell$.  A monomial with exactly $k$ $T$-letters
and more than $\ell$ $X$-letters has larger degree as well.
The higher-degree terms of $\mu(w_j)-1$ also contribute
only in larger degree.  Consequently
\[
 \mu(q_{\boldsymbol w}(h))-1
\]
has the nonzero initial polynomial just exhibited, so
$q_{\boldsymbol w}(h)\ne1$.
Truncating the noncommutative series algebra after degree
$kL+\ell$ gives a finite-dimensional nilpotent augmentation
ideal $J$ over $\mathbb F_p$.
The finite group $1+J$ is a $p$-group and still detects this
nonzero coefficient.  This is the promised finite quotient.

For countably infinite rank, a homomorphism from $F$ to a
finite group kills all but finitely many members of its
convergent free basis.  Choose such a homomorphism detecting
$h$, and retain the finitely many $x_i$ it does not kill,
adding two retained generators when necessary.  The
retraction which kills the other $x_i$ and keeps all $t_j$
sends $h$ to a nonidentity element of a finite-rank free
pro-$p$ group.  Apply the preceding construction there.
Its $w_j$ use only the retained generators, so the finite-rank
retraction commutes with the substitutions.  The resulting
$q_{\boldsymbol w}$ on the original $F$ still detects $h$.

Finally the displayed normal-closure formula follows from
the universal property of the free pro-$p$ product: adjoining
the relations $t_j=w_j$ gives exactly $P$, with the indicated
quotient map.  The abelian image assertion follows from
$w_j\in[P,P]$.
\end{proof}

\begin{corollary}\label{cor:all-lifts-faithfulness}
Let $A$ be a finite-rank direct summand of the abelianization
of a countably based free pro-$p$ group $F$, and suppose the
complementary free rank is at least two.  Fix a
$\mathbb Z_p$-basis of $A$.  The intersection of the closed
normal closures of all tuples lifting this basis is trivial.
For any homomorphism $\rho:F\to G$, therefore,
\[
 \ker\rho\subseteq N\ \text{for every such lift normal closure }N
       \quad\Longleftrightarrow\quad \ker\rho=1.
\]
\end{corollary}

\begin{proof}
A finite tuple lifting the basis of $A$ has independent
Frattini images and extends to a convergent free basis, by
the finite-coordinate and finite-tail argument in
Theorem~\ref{thm:component-no-galois-p37}.
Write the resulting free-product decomposition as in
Theorem~\ref{thm:all-lift-intersection}.  The tuples
$t_jw_j^{-1}$ in that theorem are among the tuples under
consideration.  Their normal closures already have trivial
intersection.  The equivalence is now the literal
intersection of the displayed subgroup containments.
\end{proof}

In the cyclotomic application, the omitted sum of even
components is finite free over $\mathbb Z_p$ and is a direct
summand of $H^{\rm ab}$.  The complementary free rank is
infinite because the odd part has positive $\Lambda$-rank.
Consequently a field inclusion required for every unrestricted
lift choice,
\[
                   \Omega_p^N\subseteq\Omega_p^*,
\]
is, through Galois correspondence and
Corollary~\ref{cor:all-lifts-faithfulness}, equivalent to the
full equality $\Omega_p=\Omega_p^*$.
This statement concerns the quantifier over all lift choices;
it adds no Galois-invariance assertion for an individual $N$.

% Suggested bibliography entries:
% \bibitem{MagnusKoszul}
% J. Minac, F. W. Pasini, C. Quadrelli, N. D. Tan,
% Mild pro-p groups and the Koszulity conjectures,
% Expositiones Math. 40 (2022), 432--455; arXiv:2106.03675v2,
% section 3.1, equation (3.3), printed p.8.
% \bibitem{LalondeRam}
% P. Lalonde and A. Ram, Standard Lyndon bases of Lie algebras
% and enveloping algebras, Trans. Amer. Math. Soc. 347 (1995),
% 1821--1830, section 1, equation (1.2).

% Standalone section. Uses the established p=37 module and cup inputs.
\section{The cyclotomic core of a torsion lift at \(37\)}
\label{sec:p37-core}

At \(p=37\), write
\[
 H=\operatorname{Gal}(\Omega_{37}/F_\infty),\qquad
 P=\operatorname{Gal}(\Omega_{37}/F_1),\qquad
 \Gamma=P/H=\langle\gamma\rangle.
\]
The arithmetic inputs established above give
\[
 H^{\rm ab}=U\oplus T,\qquad U\simeq\Lambda^{18},
 \qquad T=T_{32}\simeq\mathbf Z_{37}.
\]
Write \(\gamma|_T=\alpha\in1+37\mathbf Z_{37}\).
The nonzero cup of \cite[Theorem~7.5, p.~20; Corollary~8.3, p.~23]{MS}, through the proved
companion comparison, implies that the central class-two
extension above \(T_{32}\) has no nonzero \(\Lambda\)-torsion.

\begin{theorem}\label{thm:p37-cyclotomic-core}
Let \(t\in H\) have its entire abelian image equal to a generator
of \(T\), and let \(N_t\) be its closed normal closure in \(H\).
Then
\[
 \operatorname{Core}_{P}(N_t)
   =\bigcap_{\eta\in\Gamma}\eta N_t\eta^{-1}=1.
\]
\end{theorem}

\begin{proof}
Choose lifts \(x_{i,0}\in H\) of an integral \(\Lambda\)-basis
\(u_1,\ldots,u_{18}\) of \(U\), and define
\[
 x_{i,n+1}=\gamma x_{i,n}\gamma^{-1}x_{i,n}^{-1}.
\]
Their abelian images are \((\gamma-1)^nu_i\).
The sequences \(x_{i,n}\) tend to \(1\): their images are
iterated commutators of increasing length in each finite
quotient of the pro-\(37\) group \(P\), and every finite
\(37\)-group is nilpotent.
Together with \(t\), they lift a convergent product
\(\mathbf Z_{37}\)-basis of \(H^{\rm ab}\).

These elements are a free pro-\(37\) basis. Indeed, the
homomorphism from the free pro-\(37\) group on the corresponding
convergent pointed set to \(H\) is an isomorphism on
abelianizations. Its compact image is Frattini-surjective,
hence all of \(H\). Freeness of \(H\) gives a continuous
section; its closed image in the source is again
Frattini-surjective and hence is the entire source.
Thus the homomorphism is an isomorphism.

Let \(P_0\) be the closed free factor on the \(x_{i,n}\).
Then
\[
 H=P_0*F(t),\qquad P_0^{\rm ab}=U.
\]
It is \(\Gamma\)-stable, since
\(\gamma x_{i,n}\gamma^{-1}=x_{i,n+1}x_{i,n}\).
This gives invariance under positive powers; closedness and
\(\gamma^{37^s-1}\to\gamma^{-1}\) give invariance under the
inverse, and then under all of \(\Gamma\).

Let \(q_0:H\to P_0\) fix \(P_0\) and kill \(t\), and for
\(n\ge0\) put
\[
 r_n=\gamma^{-n}|_{P_0}\circ q_0\circ\gamma^n.
\]
It fixes \(P_0\), sends \(t\) to an element \(w_n\in[P_0,P_0]\),
and has kernel \(\gamma^{-n}N_t\gamma^n\).
We show that the intersection of these kernels is trivial.

In the class-two Lie quotient with all brackets involving
\(T\) removed, let \(Y\) be the logarithm of \(t\).
There is a companion identity
\[
 \gamma Y=\alpha Y+C,\qquad
 C\in\widehat{\bigwedge}^{\,2}_{\mathbf Z_{37}}U.
\]
Here \(C\ne0\): otherwise \(Y\) would be a nonzero
\(\Lambda\)-torsion element above the generator of \(T\).
This also follows directly from the nonzero companion
residue, which is unchanged by altering the lift.

Use the integral \(\Lambda\)-basis of \(U\) for completed
tensor coordinates. The action of \(\gamma\) on an ordered
pair is multiplication by \(z=(1+X)(1+Y)\).
Iteration of the companion identity, followed by \(q_0\)
and then \(\gamma^{-n}\), gives the exact quadratic
coefficient of \(w_n\):
\[
 B_{n;a,b}(X,Y)
 =C_{a,b}(X,Y)\sum_{i=0}^{n-1}\alpha^{n-1-i}z^{i-n}
 =\frac{C_{a,b}(X,Y)}{z-\alpha}
       \left(1-\left(\frac{\alpha}{z}\right)^n\right).
                                                        \tag{*}
\]
The sum is a bounded integral series for every \(n\ge0\);
the fraction expression is used in a fraction field.

Suppose \(g\ne1\) is killed by every \(r_n\).
Give the Magnus variables for \(P_0\) assigned degree one
and the variable for \(t\) assigned degree two.
Let \(d\) be the first nonzero assigned degree of \(g-1\).
This degree exists by the finite-tail Magnus argument.
At degree \(d\) there are finitely many token patterns:
singletons of one of eighteen odd colors, and adjacent
pairs occupied by \(t\). The completed coefficients are
bounded series
\[
 A_{\mathcal P}(T_i:i\text{ singleton})
 \in\mathbf Z_{37}[[T_i:i\text{ singleton}]].
\]
Indeed each odd tensor factor is \(\Lambda\), and its
completed \(s\)-fold tensor product over \(\mathbf Z_{37}\)
is \(\mathbf Z_{37}[[T_1,\ldots,T_s]]\).

Every \(r_n\) fixes the odd generators and sends \(t\) to
ordinary degree at least two. Thus higher source degrees,
and higher terms of the \(w_n\), cannot contribute at
target degree \(d\). The degree-\(d\) substitution uses only
\((*)\). It is legitimate on the completed tensors:
finite-tail truncations converge uniformly on compact
\(H\), and the fixed-degree substitutions are continuous.

Choose a nonzero token pattern with the largest number
\(k\) of pairs, and call its pair set \(J\).
For each selected pair choose an ordered color pair on
which \(C\) has a nonzero coefficient. Retain the source
colors at its singleton positions. These choices are
independent because the pairs are disjoint.

For the \(d\) independent target positions put \(Z_i=1+T_i\)
and work in
\[
 K_d=\operatorname{Frac}
       \bigl(\mathbf Z_{37}[[T_1,\ldots,T_d]][1/37]\bigr).
\]
For a pair starting at \(j\), write
\[
 R_j=\frac{C_{a_j,a_{j+1}}(T_j,T_{j+1})}
              {Z_jZ_{j+1}-\alpha},
 \qquad \beta_j=\frac{\alpha}{Z_jZ_{j+1}}.
\]
For every integer \(n\ge0\), the vanishing target coefficient
becomes
\[
 \sum_{\mathcal P}A_{\mathcal P}
   \prod_{j\in J(\mathcal P)}R_j(1-\beta_j^n)=0.
\]
Expanding the products gives a finite sum of exponentials
\[
 \beta_I^n,\qquad
 \beta_I=\prod_{j\in I}\beta_j
    =\frac{\alpha^{|I|}}{\prod_{i\in V(I)}Z_i},
\]
where \(I\) is an adjacent disjoint-pair matching and
\(V(I)\) its covered vertex set.

The \(\beta_I\)'s are distinct. Algebraic independence of
the \(Z_i\)'s makes equality imply the same covered vertex
set. Such a set determines its adjacent matching uniquely:
the leftmost vertex must pair with its right neighbor,
and one then proceeds inductively. The identities at
\(n=0,\ldots,M-1\), where \(M\) is the number of distinct
exponentials, have an invertible Vandermonde matrix over
\(K_d\). Every exponential coefficient is therefore zero.

For \(I=J\), a source pattern can contribute only if its
pair set contains \(J\). A strictly larger set has zero
coefficient by maximality. Equality of pair sets and of
the already selected singleton colors leaves exactly
the chosen pattern, with coefficient
\[
 (-1)^k A_{\mathcal P}
   \prod_{j\in J}
   \frac{C_{a_j,a_{j+1}}(T_j,T_{j+1})}
             {Z_jZ_{j+1}-\alpha}.
\]
All its factors are nonzero in \(K_d\), a contradiction.
If \(k=0\), the retraction \(r_0=q_0\) already detects the
nonzero singleton coefficient.

Hence \(\bigcap_{n\ge0}\ker r_n=1\), proving the theorem.
Closedness and density of integral powers also identify
the integral intersection with the full procyclic one.
\end{proof}

\begin{corollary}\label{cor:one-lift-p37-faithfulness}
For every \(t\) as in Theorem~\ref{thm:p37-cyclotomic-core},
\[
 \Omega_{37}^{N_t}\subseteq\Omega_{37}^*
 \quad\Longleftrightarrow\quad
 \Omega_{37}=\Omega_{37}^*.
\]
\end{corollary}

\begin{proof}
Let \(K\le H\) be the actual outer-action kernel.
It is \(P\)-normal. Thus \(K\subseteq N_t\) implies
\[
 K\subseteq\operatorname{Core}_P(N_t)=1.
\]
The converse is immediate, and the displayed field
equivalence is the exact Galois correspondence.
\end{proof}


\section{The class-group consequence of a single $K$-group}
This section proves the exact arithmetic transport from a $K$-group to one
component in the cyclotomic tower. Put $A_r=\Cl(F_r)\{p\}$ and
$\Lambda=\Z_p[[\Gamma]]$, with $F_r$ and $\Gamma$ as above.
For every character exponent $j$ modulo $p-1$, write
$A_r^{[j]}=\epsilon_j A_r$ and $X^{[j]}=\epsilon_j X$.

\begin{proposition}\label{prop:component-control}
The norm projection induces a $\Delta$-equivariant isomorphism
\[
 X/(\gamma-1)X\ \xrightarrow{\ \sim\ } A_1.
\]
In particular $A_1^{[j]}=0$ is equivalent to $X^{[j]}=0$ for every $j$.
\end{proposition}
\begin{proof}
Let $L_r$ be the Hilbert $p$-class field of $F_r$. Base change preserves
unramified abelian $p$-extensions, so $L_rF_{r+1}\subseteq L_{r+1}$.
The fields $L_r$ are therefore nested. Put $L_\infty=\bigcup_rL_r$ and
$G=\Gal(L_\infty/F_1)$. Class field reciprocity identifies
$\Gal(L_\infty/F_\infty)$ with $X$ and restriction with the norm maps.
Choose a prime of $L_\infty$ above $p$ and its inertia group $I$ in $G$.
At every finite layer the inertia group maps isomorphically to
$\Gal(F_r/F_1)$: $L_r/F_r$ is unramified and $p$ is totally ramified
in $F_r/F_1$. Taking inverse limits gives $I\simeq\Gamma$.
Thus $G=X\rtimes I$, its closed commutator subgroup is $(\gamma-1)X$,
and its abelianization fits into
\[
 0\longrightarrow X/(\gamma-1)X\longrightarrow G^{\rm ab}
   \longrightarrow\Gamma\longrightarrow0,
\]
with the image of $I$ as a section.

The quotient $G^{\rm ab}/I$ describes an abelian pro-$p$ extension
of $F_1$ unramified at $p$; all other primes were unramified already.
Conversely $L_1$ lies in $L_\infty$. Maximality of the Hilbert
$p$-class field therefore identifies this quotient with $A_1$.
The inclusion of $X$ induces the asserted isomorphism. It is the
canonical norm projection, hence commutes with $\Delta$, independently
of the chosen inertia section.

Apply $\epsilon_j$ and topological Nakayama. For completeness, a
nonzero compact $\Lambda$-module has a nonzero finite continuous
$\Lambda$-quotient. Reducing such a quotient modulo $p$ and taking
coinvariants for the finite $p$-group through which $\Gamma$ acts
leaves a nonzero quotient: its group algebra over $\F_p$ is local
with nilpotent augmentation ideal. Thus vanishing modulo
$(p,\gamma-1)$ forces the original compact module to vanish.
This proves the stated equivalence.
\end{proof}
The inertia argument is classical; see, for example,
\cite[\S1.3, pp.~3--4]{Wuthrich}.

\begin{proposition}\label{prop:k-component}
For every odd prime $p$ and every integer $m\geq1$, there is an
isomorphism
\[
 K_{4m}(\Z)/p\ \simeq\ (A_1/pA_1)^{[-2m]}.
\]
The following three groups vanish together:
$K_{4m}(\Z)\{p\}$, $A_1^{[-2m]}$, and $X^{[-2m]}$.
\end{proposition}
\begin{proof}
The odd-primary comparison in \cite[VI, Corollary~8.3]{Weibel}
gives
\[
 K_{4m}(\Z)/p\simeq
 H^2_{\mathrm{\acute et}}(\Z[1/p],\mu_p^{\otimes(2m+1)}).
\]
Over $S'=\operatorname{Spec}\Z[\mu_p,1/p]$, Kummer theory gives
$H^2_{\mathrm{\acute et}}(S',\mu_p)\simeq\operatorname{Pic}(S')/p$.
Indeed the $p$-primary Brauer contribution vanishes: there are no
real places and a single omitted finite prime, so the sum of local
invariants is injective on that possible summand. The omitted prime
is principal, generated by $1-\zeta_p$, whence
$\operatorname{Pic}(S')=\Cl(F_1)$.

Tensor with $\mu_p^{\otimes2m}$ and descend along the finite
\'etale $\Delta$-torsor $S'\to\operatorname{Spec}\Z[1/p]$.
Since $|\Delta|=p-1$ is invertible modulo $p$, descent takes invariants:
\[
 H^2_{\mathrm{\acute et}}(\Z[1/p],\mu_p^{\otimes(2m+1)})
 \simeq\bigl((A_1/pA_1)\otimes\mu_p^{\otimes2m}\bigr)^\Delta.
\]
The twist contributes the character $\omega^{2m}$, so these invariants
select precisely the character $\omega^{-2m}$. Both $K_{4m}(\Z)$
and $A_1$ are finite. Vanishing modulo $p$ therefore means vanishing
of the $p$-primary part, and Proposition~\ref{prop:component-control}
gives the assertion about $X$.
\end{proof}
This is the modern form of Kurihara's comparison
\cite[Corollary~1.5, Proposition~3.7 and Remark~3.11]{Kurihara};
see also \cite[VI, Theorem~10.9]{Weibel}. For the trivial character,
$A_1^{[0]}=0$ because extension after norm to $\Q$ acts as $p-1$
on that component and $\Cl(\Q)=0$.

\section{The posted $K$-group input and its dependencies}
\label{sec:dependencies}

This section states the external vanishing input used in the
application and follows its arithmetic dependencies explicitly.
Mok--Zheng's Theorem~1, p.~2, states
$K_{4m}(\mathbb Z)=0$ for $1\leq m\leq30$ \cite{MZ}.
Their manuscript is posted and unrefereed. The two comparison
maps in Propositions~\ref{prop:k-component} and
\ref{prop:component-control} carry this statement to
$X^{[-2m]}=0$ at every odd prime. For $p\leq61$ these residues
include every nontrivial even character; for $p>61$ they give
thirty distinct characters. The trivial class character is
already zero by the norm map to $\mathbb Q$.

The cases $m=1,2$ are the refereed $K_4$ theorem of Rognes
and the $K_8$ theorem of Dutour Sikiri\'c--Elbaz-Vincent--Kupers--Martinet,
respectively; see \cite[\S2, pp.~2--3]{MZ} for the exact
identifications. In the remaining range, the odd primes below
$2^{31}$ are handled by Hart--Harvey--Ong's Vandiver computation
\cite[\S1, assertion~(2)]{HHO} and the $K$-theory comparison.
The two-primary computation uses Rognes--Weibel and
\cite[VI, Corollary~9.8]{Weibel}. Our Galois application uses
only odd primes.

For odd $p>2^{31}$, Mok--Zheng use a strict Selmer comparison,
an arithmetic upper bound extracted from Alp\"oge's archive,
and exact Hecke determinants. More precisely, put $e=2m$.
Their Proposition~2, pp.~2--3, compares the $p$-primary
$K_{2e}(\mathbb Z)$ obstruction with the strict Selmer group
of $(\mathbb Q_p/\mathbb Z_p)(-e)$, unramified everywhere and
locally trivial at $2$ and $p$. The upper bound used in
\S\S3--5, pp.~3--8, is the pointwise Selmer bound in the
archive, specialized to dimension $2$, the trivial character,
and coefficient ring $\mathbb Z_p$. Thus the dependency is
through that arithmetic statement in the archive. The
numerical Riemann hypothesis for zeta and Dirichlet $L$-functions
does not replace any term in this displayed Selmer comparison.

For $3\leq m\leq26$, the auxiliary fields are
$\mathbb Q(\sqrt{-3})$ and $\mathbb Q(\sqrt{-7})$.
At each Hecke probe, a nonzero Selmer class forces a nonunit
trace difference on one of the admissible classical or
unitary branches. Multiplying the branch factors at each
probe before taking their greatest common divisor is the
logical step that permits the branch to vary with the probe.
The determinant checks in \S\S6--7, pp.~8--11, eliminate
these possibilities. For $27\leq m\leq30$, the argument uses
$\mathbb Q(i)$ and an additional transfer of the archive
construction, described in \S6.5, p.~10, and \S7, p.~11.
That transfer removes one comparison input and replaces six
substantive uses in the source argument. Its public provenance
statement reports that this modified entry point has not
been fully rebuilt.

Sutherland's verification report \cite{Sutherland} records
kernel checks of the original combined archive, including
several independent kernel implementations. Those checks
concern the original archive whose hash is given there.
They do not by themselves check the additional Gaussian
transfer. No new Lean kernel execution is claimed for the
application in this paper.

\subsection{The joint-representation route}

Pintar's posted theorem, p.~1 of \cite{Pintar}, reaches
$m\leq26$ from the arithmetic construction in the same
archive. Its \S2, p.~2, uses the joint-representation
argument: a positive strict Selmer length produces one
admissible integral representation with residual
semisimplification
\[
                 1\oplus\omega^e\oplus\omega^{-e},
\]
valid simultaneously at every Galois element. The logical
map can be stated directly. An admissible local point gives
such an integral representation. If there are no such
representations, the space of local points is empty.
Evaluating the model upper bound at the identity then has
empty covering set and value zero. The Selmer-to-model
inequality forces strict Selmer length zero. This is the
contrapositive used in Pintar's determinant argument.

This route does not use Mok--Zheng's vanishing theorem as
an input. Both dependency paths do use the archive's
arithmetic construction and the $K$-theory comparison.
Pintar's stated range ends at $m=26$, so it gives no
additional route to the final four components. The version
opened here has public identifier 2611.00829v3 and internal
title-page designation ``Version 5, 23 November 2026''.

\subsection{Soul\'e's torsion bound and its exact specialization}

The comparisons above use no general-number-field
torsion-stabilizer bound from Soul\'e's 2003 Lemma~9.
The reviewed joint-representation proof uses the Selmer
and model inequalities just stated; Mok--Zheng and Pintar
do not cite that lemma in their vanishing arguments.
Here is an explicit calculation fixing the scope of the
bound, without identifying unrelated assertions.

Let $F=\mathbb Q(\sqrt5)$, put $t=(\sqrt5-1)/2$, and set
\[
                 A=\begin{pmatrix}0&-1\\1&t\end{pmatrix}
                    \in\operatorname{SL}_2(O_F).
\]
The identities $t^2+t-1=0$ and $A^2-tA+1=0$ give
$A^5=1$, while $A\ne1$. Its order is $5$, exceeding
$1+\max([F:\mathbb Q],2)=3$.
The valid dimension argument is as follows. Regard a
nonidentity order-$p$ endomorphism of $F^N$ as a rational
endomorphism of dimension $[F:\mathbb Q]N$. Its minimal
polynomial has a primitive $p$th root as a zero and hence
contains the irreducible polynomial $\Phi_p$ of degree
$p-1$. Therefore $p\leq[F:\mathbb Q]N+1$. For
$F=\mathbb Q$ this gives $p\leq N+1$, exactly the bound
used in the rational specialization. The displayed
matrix disproves the stated larger-field bound while
the minimal-polynomial argument proves this specialization.



\section{Reflection and a finite list of arithmetic exceptions}
For $1\leq m\leq30$, put $\theta_m=2+2m$ and let
\[
 \mathcal E_m=\{p\text{ prime}:p\geq2^{31},\quad
             p\mid\operatorname{num}(B_{\theta_m})\}.
\]
The nonempty sets are listed in Table~\ref{tab:reflection-exceptions}.

\begin{proposition}\label{prop:reflection-exceptions}
For every $1\leq m\leq30$ and every odd prime
$p\notin\mathcal E_m$, the component $X^{[2+2m]}$ vanishes.
In particular this vanishing holds at every odd prime for
$m=1,\ldots,15$ and for $m=21$.
\end{proposition}
\begin{proof}
We give the reflection morphism explicitly. Put $F=\Q(\mu_p)$,
$A=\Cl(F)\{p\}$, and let $B\subset F^\times$ consist of elements
whose $p$th roots generate everywhere unramified extensions of $F$.
Kummer theory and Artin reciprocity give
\[
 B/F^{\times p}\simeq\operatorname{Hom}(A,\mu_p).
\]
For $b\in B$, every valuation of $b$ is divisible by $p$.
Write $(b)=\mathfrak a^p$. The map $[b]\mapsto[\mathfrak a]$
is well-defined, is $\Delta$-equivariant, and has kernel
$(B\cap\mathcal O_F^\times)/\mathcal O_F^{\times p}$.
For an even $\theta$ with $2\leq\theta\leq p-3$, take
$\chi=\omega^{1-\theta}$. The $\chi$ part of this kernel vanishes:
the odd part of global units modulo $p$ comes from roots of unity,
which have character $\omega$, whereas $\chi\ne\omega$.
Indeed $u/c(u)$ is a root of unity for every unit $u$, since it is
an algebraic integer with all complex conjugates of absolute value one.
Thus the displayed map restricts to an injection
\[
 \operatorname{Hom}(A^{[\theta]},\mu_p)
       \lhook\joinrel\longrightarrow A^{[1-\theta]}[p].
\]
This proves the required direction of Leopoldt reflection, including
its character shift; compare \cite[Theorem~4.4.2 and
Corollary~4.4.4]{SharifiIwa}.

For completeness, the vanishing on the right follows from a precise
Stickelberger annihilator. For $b$ primitive modulo $p$, projection
of $(\sigma_b-b)\Theta$ onto the $\omega^{1-\theta}$ component gives
\[
 (\omega^{1-\theta}(b)-b)B_{1,\omega^{\theta-1}}.
\]
The first factor is a $p$-adic unit, and
$B_{1,\omega^{\theta-1}}\equiv B_\theta/\theta\pmod p$.
Consequently $p\nmid\operatorname{num}(B_\theta)$ implies
$A^{[1-\theta]}=0$ and then $A^{[\theta]}=0$.
These are the Herbrand implication and the Bernoulli congruence in
\cite[Theorem~4.5.7 and Corollary~5.4.15]{SharifiIwa}.

If $p<2^{31}$, the published Vandiver computation
\cite[\S1, assertion~(2)]{HHO} gives $A^+=0$ directly.
For $p\geq2^{31}$, every $\theta_m$ lies between $4$ and $p-3$.
If $p\notin\mathcal E_m$, the preceding injection and annihilator
give $A^{[\theta_m]}=0$. In both cases,
Proposition~\ref{prop:component-control} gives $X^{[\theta_m]}=0$.

The arithmetic table was computed from the defining rational
Bernoulli recurrence. Every numerator factorization was checked
by exact multiplication and every asserted prime was certified
by the fully factored $n-1$ Lucas criterion, recursively down to $2$.
The standard-library verifier and all certificates accompany this paper.
They show $\mathcal E_m=\varnothing$ for $m=1,\ldots,15,21$,
which proves the last assertion.
\end{proof}

\begin{table}[htbp]
\centering
\small
\begin{tabular}{rrl}
$m$ & $\theta_m$ & $\mathcal E_m$ \\
\hline
16&34&$151628697551$\\
17&36&$26315271553053477373$\\
18&38&$154210205991661$\\
19&40&$1897170067619$\\
20&42&$1520097643918070802691$\\
22&46&$67568238839737$\\
23&48&$153289748932447906241$\\
24&50&$47464429777438199$\\
25&52&$4534045619429$\\
26&54&$1120412849144121779$\\
27&56&$19088082706840550550313$\\
28&58&$6235242049,\quad37349583369104129$\\
29&60&$109317926249509865753025015237911$\\
30&62&$28765594733083851481$\\
\hline
\end{tabular}
\caption{The fifteen prime exceptions to the reflection calculation.
The sets for $m=1,\ldots,15$ and $m=21$ are empty.}
\label{tab:reflection-exceptions}
\end{table}

This result is independent of the posted $K$-group vanishing claims.
Combining the proved reflection result with the $K$-group input stated
by Mok--Zheng supplies the pair
$X^{[-2m]}=X^{[2+2m]}=0$ for $1\leq m\leq30$ and
$p\notin\mathcal E_m$. Their input is currently an unrefereed
posted result; its dependency and verification status are recorded
separately in the audit. Pintar's posted argument supplies another
route to the $K$-group input for $1\leq m\leq26$, with shared
archival dependencies. Neither source is needed for
Proposition~\ref{prop:reflection-exceptions} itself.
No vanishing of $X^{[\theta_m]}$ at the primes in
Table~\ref{tab:reflection-exceptions} is established here.

\begin{corollary}[Application of the posted $K$-theory theorem]
\label{cor:k-canonical-detection}
For an odd prime $p$, put
\[
 \Sigma_{30}(p)=\{-2m\bmod(p-1):1\leq m\leq30\}.
\]
With the vanishing theorem of \cite[Theorem~1, p.~2]{MZ} as the
external input described above, Theorem~\ref{thm:selected-all-stages}
gives at every odd prime the canonical Galois subextensions
\[
 L_{2,\Sigma_{30}(p)}\subsetneq L_{3,\Sigma_{30}(p)}
 \subsetneq\cdots\subsetneq L_{61,\Sigma_{30}(p)}
 \subseteq\Omega_p^*.
\]
For every \(2\leq n\leq61\), the group
\(\operatorname{Gal}(L_{n,\Sigma_{30}(p)}/F_\infty)=Q_{n,\Sigma_{30}(p)}\)
has nilpotency class exactly \(n\) and abelianization
\((U/D_n)\oplus T_{\Sigma_{30}(p)}\).
There are two elements whose length-\(n\) commutator is nonzero.
Its class-three quotient has second and third lower-central
graded groups
\[
 N_{\Sigma_{30}(p)}\oplus(W\widehat\otimes T_{\Sigma_{30}(p)}),
 \qquad \operatorname{Lie}^3(W)/[W,N_{\rm bad}],
\]
respectively.
\end{corollary}
\begin{proof}
Proposition~\ref{prop:k-component} transports each stated
$K_{4m}(\mathbb Z)$ vanishing to $X^{[-2m]}=0$.
These are precisely the component hypotheses of
Theorem~\ref{thm:selected-all-stages}. For \(2\leq n\leq61\),
the characters \(-2,-4,\ldots,-2\lfloor n/2\rfloor\) all belong
to \(\Sigma_{30}(p)\), with residues read modulo \(p-1\).
That theorem proves exact class and the strict transitions.
Proposition~\ref{prop:selected-class-three} gives the displayed
lower-central groups of the class-three quotient.
\end{proof}
The Bernoulli prime list in Table~\ref{tab:reflection-exceptions}
records the limits of that reflection calculation. Every prime
in the table is included in the displayed Galois-field application,
which uses the component $X^{[-2m]}$ directly.
The characteristic-zero comparison, the reflection morphism,
and the construction of $C_E$ are independent of the posted
vanishing theorem. The assertion of the corollary uses that
external theorem, with its verification status as described in
Section~\ref{sec:dependencies}. For $m\leq26$, Pintar's posted
argument gives the alternative route with the shared archive
inputs identified there. For $p<2^{31}$ the full equality
$\Omega_p^*=\Omega_p$ already follows from \cite{HHO,Note}.

\section{The first irregular prime}
This section makes the arithmetic indices and a denominator divisor
explicit at $p=37$. The even characters have exponents
\[
 0,2,4,6,8,10,12,14,16,18,20,22,24,26,28,30,32,34.
\]
The Bernoulli recurrence gives the residues
\[
\begin{array}{c|rrrrrrrrr}
 r&2&4&6&8&10&12&14&16&18\\\hline
 B_r\bmod37&31&16&15&16&4&17&32&22&31\\
\end{array}
\]
\[
\begin{array}{c|rrrrrrrr}
 r&20&22&24&26&28&30&32&34\\\hline
 B_r\bmod37&15&15&17&12&29&2&0&2.
\end{array}
\]
In particular $B_{32}=-7709321041217/510$ has $37$-adic valuation one.
The unique irregular index is $32$, and Herbrand--Ribet gives the
odd class character $1-32\equiv5\pmod{36}$.

Here is a small Vandiver certificate. Work modulo $q=149$, put $z=16$,
which has order $37$, and use the corrected cyclotomic-unit test
of \cite[\S4, pp.~4--5]{BH}, with the correction in \cite[\S5]{HHO}:
\[
 V=\prod_{c=1}^{18}
      \bigl(z^{c/2}-z^{-c/2}\bigr)^{r_c}\pmod{149},
 \qquad r_c\in\{0,\ldots,36\},\quad r_c\equiv c^4\pmod{37}.
\]
The exponents $c/2$ are read modulo $37$; the reduced outer
exponents are allowed by the same test. Exact modular
multiplication gives $V=132$ and $V^4=81\ne1$.
Since $32$ is the only irregular index, the criterion proves
$A_1^+=0$. Thus $X^{[e]}=0$ for all eighteen even characters,
and every even characteristic ideal of $X$ has generator $1$.
The four congruences for the Iwasawa invariants in \cite[\S5]{HHO}
are also small here. With
\[
 s=37^{-1}\sum_{a=1}^{18}a^{31},\qquad
 t=37^{-1}\sum_{a=1}^{18}a^{67},
\]
we obtain, modulo $37$,
\[
 2^{32}=7,\quad s=30,\quad t=31,\quad
 (2-32)s=25,\quad(1-32)t=1.
\]
They prove $\lambda_{37}=\nu_{37}=1$. In particular
$X^{[5]}$ is a free $\Z_{37}$-module of rank one, and all other
odd class components vanish.

Choose $\gamma$ with cyclotomic value $u=38$ and write $Z=\gamma-1$.
Let $\vartheta_5(Z)$ be the inverse limit of the negative
Stickelberger projections
\[
 -\frac1{37^{n+1}}
 \sum_{\substack{1\leq a<37^{n+1}\\37\nmid a}}
 a\,\omega(a)^{-5}\gamma_a^{-1},
 \qquad
 \gamma_a\in\Gamma/\Gamma^{37^n}.
\]
This specifies the power series exactly. Its principal ideal is
$\operatorname{char}(X^{[5]})$ by the main conjecture.
Projection at $n=2$ gives
\[
 \vartheta_5(Z)=a_0+a_1Z+a_2Z^2+\cdots,
 \qquad
 \begin{aligned}
 a_0&\equiv14\cdot37+33\cdot37^2&&\pmod{37^3},\\
 a_1&\equiv16+6\cdot37&&\pmod{37^2},\\
 a_2&\equiv29&&\pmod{37}.
 \end{aligned}
\]
The coefficients agree with \cite[\S1.9, p.~8]{Wuthrich} and are
reconstructed by the program cited through the workbench.
The finite-level error is divisible by $(1+Z)^{37^2}-1$, hence
vanishes modulo $37^3$ at every $Z\in37\Z_{37}$.
The displayed precision is therefore sufficient to determine the root
modulo $37^3$. Since $a_1$ is a unit, Weierstrass preparation gives
one root $r\in37\Z_{37}$ and a unit multiple of $Z-r$.
Writing $r=37a+37^2b$, the first digit gives $14+16a\equiv0$, so
$a=13$ and $14+16\cdot13=6\cdot37$. The next digit gives
\[
 6+33+6\cdot13+29\cdot13^2+16b
       \equiv23+16b\equiv0\pmod{37},
\]
so $b=24$. Put $\beta=1+r$ and $\alpha=38/\beta$. Then
\[
 r\equiv13\cdot37+24\cdot37^2,\qquad
 \alpha\equiv1+25\cdot37+20\cdot37^2\pmod{37^3}.
\]
The characteristic polynomials for $\gamma$ on the rational even
parts of $H^{\rm ab}$ are exactly
\[
 q_{32}(z)=z-\alpha,\qquad q_e(z)=1\quad(e\ne32).
\]
Here $\alpha$ is the precisely defined $37$-adic number above,
not a rational approximation substituted for a characteristic root.
The corresponding scalar divisor and its coordinate are
\[
 D(X,Y)=(1+X)(1+Y)-\alpha,\qquad
 Y=\frac{\alpha+S}{1+X}-1,
\]
so $D=S$ and $1/D$ has a simple pole with Laurent coefficient one.
For example the divisor passes through $(0,\alpha-1)$.
The dual cup character is $2-32\equiv6$ and its eigenvalue is
$38^2/\alpha=38\beta$. This scalar calculation specifies the
arithmetic denominator; the survival of a companion numerator
on that divisor is the separate cup-comparison requirement.

\subsection{The exact meaning of vanishing odd torsion}
\label{subsec:odd-torsion-plus-adjoint}

The odd torsion $D$ is related to the plus class module
by the following exact adjoint comparison.

\begin{proposition}\label{prop:D-plus-finiteness}
For every odd prime $p$,
\[
                             D=0\quad\Longleftrightarrow\quad X^+\text{ is finite}.
\]
Under these equivalent assertions the proved
single-component theorem also gives $Z=0$.
\end{proposition}

\begin{proof}
For an odd exponent $a\ne1,-1$, set
$\chi=\omega^{-a}$ and $\theta=\omega^{1-a}$.
The proof of \cite[Proposition~2.2.6, p.~4]{WWEOrd}
gives the actual Poitou--Tate identification
\[
 \operatorname{Ext}^1_\Lambda
       ((M^{[a]})^\#(1),\Lambda)=X^{[1-a]}.
\]
The involution and twist act on the cyclotomic
variable by $\gamma\mapsto u\gamma^{-1}$.
At tame level one, the source assumes $p\geq5$,
$\theta$ primitive and $\theta\ne\omega^2$. These
conditions exclude precisely $a=1,-1$, which we treat below.

Put $R=\Lambda[1/p]$. Restricting extensions along
$D^{[a]}\hookrightarrow M^{[a]}$, after the same
twist and localization, gives an isomorphism
\[
 X^{[1-a]}[1/p]\xrightarrow{\ \sim\ }
 \operatorname{Ext}^1_R
       ((D^{[a]})^\#(1)[1/p],R).
\]
Indeed the torsionfree quotient is projective over
the one-dimensional PID $R$, so the long exact
Ext sequence proves this exact restriction map is
an isomorphism. A nonzero cyclic torsion summand
$R/(f^k)$ has Ext group $R/(f^k)$, by its length-one
resolution. Consequently $D^{[a]}=0$ exactly when
$X^{[1-a]}[1/p]=0$. The first module has no
$\mathbf Z_p$-torsion; Ferrero--Washington makes
the second finitely generated over $\mathbf Z_p$,
so its rational vanishing is precisely finiteness.

For the exceptional components, norm to $\mathbf Q$
and extension back give $A_1^{[0]}=0$, since their
composite is multiplication by the unit $p-1$.
Herbrand gives $A_1^{[1]}=0$ and, for $p>3$,
$A_1^{[-1]}=0$ from $B_2=1/6$; see also
\cite[Lemma~5.2.7, p.~13]{WWEOrd}.
Leopoldt reflection with $\chi=\omega^{-1}$ gives
$r_p(A_1^{[2]})\leq r_p(A_1^{[-1]})=0$.
These are the exact cases of
\cite[Theorems~4.4.2 and~4.5.7]{SharifiIwa}.
Proposition~\ref{prop:component-control} therefore
gives $X^{[0]}=X^{[2]}=X^{[1]}=X^{[-1]}=0$.

The class-field unit sequence and the fact that the
odd global units are exactly $\mathbf Z_p(1)$,
together with Coleman's exact sequence
\cite[Theorem~5.5.31, p.~148]{SharifiIwa}, give the
actual maps
\[
 M^{[1]}\simeq\ker(\Lambda\xrightarrow{\gamma\mapsto u}
                \mathbf Z_p(1))=(\gamma-u)\Lambda,
 \qquad M^{[-1]}\simeq\Lambda\quad(p>3).
\]
Hence both exceptional odd torsion components vanish.
At $p=3$ these odd characters coincide, as do the
even characters $0,2$; the first displayed map and
the trivial-character class calculation cover them.
The correspondence $a\mapsto1-a$ now covers every
odd and every even character, proving the equivalence.
Applying the established finite-$X^{[e]}$ component
theorem to every even $e$ gives $Z=0$.
\end{proof}

The exact adjoint map therefore identifies the assertion
$D=Z=0$ with the finite-$X^+$ hypothesis. The canonical
finite-stage quotients are defined at every odd prime,
including those for which this hypothesis remains undecided.

\section{Integral Soul\'e generation at the first irregular prime}
\label{sec:integral}

This section applies McCallum--Sharifi's integral result to Brown's
freeness theorem. It gives a negative answer to integral generation
by the specified Soul\'e generators, already in one of the
characters singled out by $K$-theory.

Let $\mathfrak g=\bigoplus_{m\geq1}\mathfrak g_m$ be the
integral graded Lie algebra of the outer Galois image, for its
weight filtration, in the convention of
\cite[\S9, pp.~24--28]{MS}. Choose the Soul\'e elements
$\sigma_m$ for odd $m\geq3$ with the normalization used there,
and let $\mathfrak h$ be the graded $\Z_p$-Lie subalgebra
generated by their leading classes.

\begin{proposition}\label{prop:integral37}
At $p=37$, the inclusion $\mathfrak h_{32}\subseteq
\mathfrak g_{32}$ is strict. Its cokernel is a nonzero finite
cyclic $37$-group. The character of this graded piece is
$\omega^{32}=\omega^{-4}$.
\end{proposition}
\begin{proof}
McCallum--Sharifi's Theorem~9.7, p.~28, applies at the first
irregular index $m$ when Vandiver holds, the Deligne--Ihara
map is injective through degree $m$, and one of its cup
coefficients $e_{i,m}$ with $i<m/2$ is nonzero.
At $p=37$ the first irregular index is $32$ and Vandiver
holds, as calculated above. Brown's Theorem~1.1 supplies
the Deligne--Ihara injectivity: the rational free Lie map
is injective, so its integral restriction is injective
because a free $\Z_{37}$-Lie algebra is torsionfree.
The required cup coefficient is explicit in
\cite[Corollary~8.3, p.~23]{MS}: the coefficient of
$[x_7,x_{25}]$ is $-1$ modulo $37$, and $7<16$.
Thus the hypotheses of their Theorem~9.7 all hold.
Its conclusion gives
\[
 \mathfrak h_{32}\varsubsetneq\mathfrak g_{32},
 \qquad
 \mathfrak g_{32}=\Z_{37}\bar x_{32}+\mathfrak h_{32}.
\]
The quotient is consequently nonzero and cyclic.
Hain--Matsumoto's rational generation, with Brown's
freeness, gives
$\mathfrak h_{32}\otimes\Q_{37}
 =\mathfrak g_{32}\otimes\Q_{37}$.
These graded pieces are finite $\Z_{37}$-modules in
each weight, so the quotient is finite. It is therefore
$\Z/37^r$ for some $r\geq1$.
The torus acts in weight $32$ by $\omega^{32}$, and
$32\equiv-4\pmod{36}$.
\end{proof}

This application strengthens the formerly conjectural
injectivity input of the published integral obstruction
by using Brown's theorem; the obstruction and its cup
coefficient are those of McCallum and Sharifi.
The argument proves no particular value of $r$.
In particular the filtered relation involving
$B_{32}\bar x_{32}$ has an error in the third term of an
induced filtration. Identifying that error with the
third term of a Lie filtration generated by the chosen
Soul\'e elements would itself require integral generation.
The displayed cyclic quotient is the conclusion justified
by the cited theorem and rational generation.

% Original section for integration into the main paper.
% Bibliography keys required: IshiiKernel, IshiiLie, BaumardSchneps,
% Rubin1991, GotoIshii2026, SanoVenjakob2026, MilneANT.
\section{The CM elliptic question after Ishii}
\label{sec:cm-audit}

This section identifies the arithmetic and Lie-theoretic statements
needed for the once-punctured CM elliptic curve, and proves the elementary
calculations at its two smallest split primes.
Let \(K\) be an imaginary quadratic field of class number one, let
\(\mathcal O_K\subset\operatorname{End}_K(E)\), and put \(X=E-\{O\}\).
For an ideal \(\mathfrak m\), write \(K(\mathfrak m)\) for the ray class
field of modulus \(\mathfrak m\), with no archimedean modulus. Fix a prime
\(p\geq5\) splitting as \((p)=\mathfrak p\bar{\mathfrak p}\).
Put \(K_\infty=\bigcup_nK(p^n)\), and let \(\Omega\) be the maximal
pro-\(p\) extension of \(K(p)\) unramified outside \(p\).
Ishii's Theorem A, p.~2, equivalently Theorem 2.14, p.~11,
identifies the fixed field of the outer kernel with \(K(E[p])\Omega\)
under three hypotheses: \(p\nmid h(K(p))\), exactly two primes of
\(K(p^2)\) above \(p\), and his Conjecture 2.3
\cite{Ishii}. The third hypothesis is a statement of freeness
for a particular arithmetic graded Lie algebra.

\subsection{The first Pollack relation and the arithmetic morphism}

Ishii's subsequent paper proves generation under the vanishing of
\(H^2_{\mathrm{\acute et}}(\mathcal O_K[1/p],\mathbb Q_p(m_1,m_2))\)
for every admissible pair \((m_1,m_2)\)
\cite[Theorem 1.2, p.~3]{IshiiLie}. Here
\(\mathbb Q_p(m_1,m_2)\) has character
\(\chi_{\mathfrak p}^{m_1}\chi_{\bar{\mathfrak p}}^{m_2}\), with
\(\chi_{\mathfrak p}\) and \(\chi_{\bar{\mathfrak p}}\) the two
characters of the CM Tate module, and admissibility means
\[
 m_1,m_2>0,\qquad m_1\equiv m_2\pmod {w_K},\qquad
 (m_1,m_2)\ne(1,1),
 \quad w_K=|\mathcal O_K^\times|.
\]
The same paper records that the earlier freeness assertion fails
when \(w_K=2\), and retains the two fields \(w_K=4,6\) as open cases
\cite[Former Conjecture 1.8, p.~4, and Remark 1.9, p.~5]{IshiiLie}.
We give the explicit derivation calculation behind this correction.
The relation is Pollack's; its derivations and their relation with
mixed elliptic motives are discussed by Baumard and Schneps
\cite[pp.~2--3]{BaumardSchneps}.

Let \(L=\mathbb Q_p\langle a,b\rangle_{\mathrm{Lie}}\) be the free
Lie algebra on \(a,b\), realized inside the free associative algebra
\(\mathbb Q_p\langle a,b\rangle\), with bracket \([u,v]=uv-vu\).
Set \(A_j=\operatorname{ad}_a^j(b)\). For even \(k\geq2\), define
a derivation by
\[
 \epsilon_k(a)=A_k,\qquad
 \epsilon_k(b)=\sum_{j=0}^{k/2-1}(-1)^j[A_j,A_{k-1-j}].
\]
A derivation is determined by its values on \(a,b\). Its bidegree
is the increase in the numbers of occurrences of \(a,b\); thus
\(\epsilon_k\) has bidegree \((k-1,1)\).

\begin{proposition}[The first Pollack relation]\label{prop:pollack-explicit}
These derivations satisfy
\[
 [\epsilon_4,\epsilon_{10}]-3[\epsilon_6,\epsilon_8]=0.
\]
Both brackets displayed here are nonzero.
\end{proposition}
\begin{proof}
The derivation identity and
\([a,[A_i,A_j]]=[A_{i+1},A_j]+[A_i,A_{j+1}]\)
give a telescoping sum
\(\epsilon_k([a,b])=0\).
For \(n\geq1\), repeated use of the derivation identity gives
\begin{align*}
 \epsilon_k(A_n)
 ={}&\sum_{r=0}^{n-1}\binom{n}{r+1}
              [A_{k+r},A_{n-1-r}]\\
 &+\sum_{j=0}^{k/2-1}(-1)^j
      \sum_{r=0}^{n}\binom nr
              [A_{j+r},A_{k-1-j+n-r}].
\end{align*}
Indeed, differentiating each occurrence of
\(\operatorname{ad}_a\) gives
\(\sum_{i=0}^{n-1}\operatorname{ad}_a^i
[A_k,A_{n-1-i}]\). Expanding this sum by the binomial theorem
and using
\(\sum_{i=r}^{n-1}\binom ir=\binom n{r+1}\)
gives the first line. Differentiating \(b\) gives the second.

Apply this formula to the two brackets and collect the coefficients
of \([A_j,A_{13-j}]\), \(0\leq j\leq6\), using
\([A_i,A_j]=-[A_j,A_i]\). The complete table is
\[
\begin{array}{c|rrrrrrr}
j&0&1&2&3&4&5&6\\ \hline
[\epsilon_4,\epsilon_{10}](a)&0&0&-6&-33&-75&-90&-42\\
{}[\epsilon_6,\epsilon_8](a)&0&0&-2&-11&-25&-30&-14 .
\end{array}
\]
Thus the proposed difference \(D\) annihilates \(a\).
It also annihilates \([a,b]\), since every \(\epsilon_k\) does.
Consequently \([a,D(b)]=0\).

For completeness, a finite associative polynomial commuting with
\(a\) contains only powers of \(a\). After removing those powers,
choose a nonzero word having the least number of initial \(a\)'s.
In \(Pa\) the corresponding word followed by \(a\) has that same
number of initial \(a\)'s; in \(aP\) every word containing \(b\)
has strictly more. This contradicts \(aP=Pa\).
Now \(D(b)\) has three occurrences of \(b\), so \(D(b)=0\).
The associative word \(a^{11}ba^2b\) has coefficient \(6\) in
\([\epsilon_4,\epsilon_{10}](a)\), and coefficient \(2\) in the
second bracket; hence neither bracket is zero.
\end{proof}

We specify how the relation reaches the arithmetic Lie algebra.
Let \(F^nG_K\) be the kernel of the outer action on the quotient
of the geometric pro-\(p\) fundamental group by its \((n+1)\)-st
lower central subgroup, and put
\(\mathfrak g_X=\bigoplus_{n>0}F^nG_K/F^{n+1}G_K\).
Ishii identifies \(\mathfrak g_X\otimes\mathbb Q_p\) with the
graded Lie algebra of the unipotent Galois image
\cite[Lemma 3.6, p.~14]{IshiiLie}. This gives an injection into
the outer derivations of \(L\) preserving the boundary class.
The boundary is \(z=[a,b]\), after the symplectic choice of free
generators. Ishii's Proposition 3.3(4), pp.~12--14, gives the
specific quotient map
\[
 \operatorname{Der}^{*}(L)/\mathbb Q_p\operatorname{ad}_z
       \ \xrightarrow{\ \sim\ }\ \operatorname{OutDer}^{*}(L).
\]
Here positive-degree derivations in \(\operatorname{Der}^{*}(L)\)
annihilate \(z\); the displayed kernel has bidegree \((1,1)\).
Therefore the quotient map is an isomorphism on each of the
bidegrees \((3,1),(5,1),(7,1),(9,1),(12,2)\) used below.

In bidegree \((m,1)\), a derivation annihilating \(z\) is a scalar
multiple of \(\epsilon_{m+1}\). Its value on \(a\) lies in the
one-dimensional Lie space spanned by \(A_{m+1}\): every bracket
with just one \(b\) is, by induction on its length, a multiple of
this iterated bracket. Subtract the corresponding multiple of
\(\epsilon_{m+1}\). The resulting derivation kills \(a\) and \(z\),
so it kills \(b\) by the centralizer argument above. This proves
the asserted one-dimensionality without choosing an arithmetic
normalization.

When \(w_K=2\), all four bidegrees \((3,1),(5,1),(7,1),(9,1)\)
are admissible. Write the images of four proposed nonzero generators
as \(c_m\epsilon_{m+1}\), with \(c_m\in\mathbb Q_p^\times\).
The relation then becomes
\[
 c_5c_7[\sigma_{(3,1)},\sigma_{(9,1)}]
 -3c_3c_9[\sigma_{(5,1)},\sigma_{(7,1)}]=0.
\]
The two commutators are linearly independent in the free Lie algebra
on the proposed generators: in its associative realization they
involve four distinct words. The injective arithmetic derivation
map therefore contradicts the proposed freeness. A vanishing
generator would already contradict that proposal.
For \(K=\mathbb Q(i)\) and \(\mathbb Q(\sqrt{-3})\), the index
\((3,1)\) is excluded by \(m_1\equiv m_2\pmod {w_K}\).
The calculation above gives no freeness proof in these two cases.

\subsection{Complex conjugation in the two-variable module}

Let \(\Gamma=\operatorname{Gal}(K_\infty/K(p))\simeq\mathbb Z_p^2\).
Complex conjugation interchanges its two local unit directions.
Choose generators \(\gamma_1,\gamma_2\) with
\(c\gamma_1c^{-1}=\gamma_2\), and identify
\[
 R=\mathbb Z_p[[\Gamma]]
   =\mathbb Z_p[[S,T]],\qquad S=\gamma_1-1,\quad T=\gamma_2-1 .
\]
For a module with the arithmetic conjugation action, the identity
is \(c(rm)=c(r)c(m)\), where \(c(S)=T\) and \(c(T)=S\).
The following calculation specifies the coefficient ring on which
the fixed subgroup is a module.

\begin{proposition}\label{prop:cm-plus}
Let \(M\) be a finitely generated compact \(R\)-module with a
continuous semilinear involution \(c\). Write \(M^+=\ker(c-1)\).
For \(m\in M^+\),
\[
 Sm\in M^+\quad\Longleftrightarrow\quad(S-T)m=0 .
\]
Consequently \(M^+\) is an \(R\)-submodule exactly when
\((S-T)M^+=0\). It is always a module over
\[
 R^c=\mathbb Z_p[[x,y^2]],\qquad
 x=(S+T)/2,\quad y=(S-T)/2 .
\]
When \(M^+\) is finite, \((S-T)M[1/p]=0\).
\end{proposition}
\begin{proof}
For \(cm=m\), semilinearity gives \(c(Sm)=Tm\). This proves
the first equivalence and the necessity of the second assertion.
Under the stated annihilation, multiplication by both \(S\) and
\(T\) preserves \(M^+\). It follows first for polynomials and then
by compactness and continuity for every element of \(R\).
Since \(p\) is odd, the change of variables \(S=x+y,T=x-y\) is
invertible over \(\mathbb Z_p\). The involution fixes \(x\) and
negates \(y\). A series is fixed exactly when its coefficients of
odd powers of \(y\) vanish. This proves the formula for \(R^c\),
and semilinearity proves its action on \(M^+\).

The two continuous projectors \((1+c)/2\) and \((1-c)/2\) give
the direct sum of the fixed and anti-fixed \(\mathbb Z_p\)-modules.
After inverting \(p\), finiteness kills the first summand.
Thus \(c=-1\) on \(M[1/p]\). Applying semilinearity to \(Sm\)
now gives \(-Sm=-Tm\), which proves the last assertion.
\end{proof}

This proposition identifies a precise arithmetic issue in formulating
the desired finiteness hypothesis. For example, \(R\) with its swapping
involution has \(1\in R^c\) but \(S\notin R^c\).
The module
\[
 M=\bigl(R/(S-T)\bigr)^{\oplus2},\qquad c(m)=-m,
\]
has \(M^+=0\); its action is semilinear because \(S=T\) in the
quotient. Nevertheless \(M[1/p]\) requires two generators at
the maximal ideal \((S,T)\) of \(R[1/p]\), since the quotient by
that ideal has dimension two over \(\mathbb Q_p\).
These are explicit module computations; they establish no occurrence
of either example as an arithmetic class-group module.

A second calculation isolates the information omitted by
characteristic ideals. Put
\[
 M_1=R/(S),\qquad
 M_2=R/(S)\oplus R/(S,T).
\]
Projection onto the first summand has kernel \(R/(S,T)\).
This kernel vanishes at every height-one prime of \(R\):
no such prime contains both \(S\) and \(T\), and either variable
annihilates the kernel. Thus \(M_1\) and \(M_2\) have the same
characteristic ideal \((S)\). After inverting \(p\), localization
at \((S,T)\) gives residue dimensions one and two, respectively;
the second module is noncyclic. Rubin's main conjecture computes
characteristic ideals of the relevant one-prime-ramified modules
in terms of elliptic units and Katz's \(p\)-adic \(L\)-functions
\cite{Rubin1991}. Applying that equality to the sought cyclicity
argument requires arithmetic control of the kernel illustrated
by this calculation.

\subsection{The full elliptic-unit radical for arbitrary split primes}
\label{subsec:cm-elliptic-radical}

This subsection requires only that \(K\) be an imaginary quadratic field
of class number one, that \(p\geq5\) split as
\((p)=\mathfrak p\bar{\mathfrak p}=(\pi)(\bar\pi)\), and that
\(E/K\) have complex multiplication by \(O_K\).
Put \(X=E-\{O\}\), \(L_n=K(p^n)\), and
\(L_\infty=\bigcup_{n\geq1}L_n\).
Let \(\Omega\) be the maximal pro-\(p\) extension of \(L_1\)
unramified outside \(p\), and set
\[
 H=\operatorname{Gal}(\Omega/L_\infty),\qquad A=H^{\mathrm{ab}}.
\]
All fields are taken inside the fixed algebraic closure.
Write \(\rho\) for the actual outer action of \(G_K\) on
\(\pi_1^{(p)}(X_{\bar K})\).
The proof of \cite[Lemma 2.13, p.~10]{Ishii} gives a field
\(\Omega^*\subseteq\Omega\) and the exact identity
\begin{equation}\label{eq:cm-radical-outer-field}
 \bar K^{\ker\rho}=K(E[p])\Omega^*,\qquad
 [K(E[p]):L_1]\ \text{prime to }p.
\end{equation}
The field \(L_\infty\) contains \(\mu_{p^\infty}\).
Also \(K(E[p^\infty])=K(E[p])L_\infty\), so
\(L_\infty\subseteq\Omega^*\) by intersecting with \(\Omega\).

Let \(C(L_n)\) be Rubin's elliptic-unit group, with the convention of
\cite[Definition 2.2, p.~8]{IshiiSoule}: it contains \(\mu(L_n)\)
and the appropriate Galois augmentation differences of theta values.
The convention using suitable twelfth roots gives the same
\(p\)-power radical because \(12\) is prime to \(p\).
Define the full radical field
\begin{equation}\label{eq:cm-full-radical}
 L_{\mathrm{ell}}=
 L_\infty\bigl(c^{1/p^r}:c\in C(L_n),\ n,r\geq1\bigr).
\end{equation}
It does not depend on choices of roots, since
\(\mu_{p^\infty}\subseteq L_\infty\).

\begin{proposition}\label{prop:cm-full-radical}
With the preceding definitions,
\[
             L_{\mathrm{ell}}\subseteq\Omega^*.
\]
Thus the actual outer kernel fixes every \(p\)-power root of every
Rubin elliptic unit in every finite ray layer \(L_n\).
This assertion holds for every split \(p\geq5\) and every such \(K,E\).
\end{proposition}

\begin{proof}
We give finite-level Kummer values throughout the argument.
Choose a \(p\)-compatible CM basis
\(\omega_{1,n}\in E[\mathfrak p^n]\),
\(\omega_{2,n}\in E[\bar{\mathfrak p}^n]\), and compatible
roots of unity \(\zeta_{p^r}\).
For a \(p\)-unit \(u\in L_\infty^\times\), all its \(p\)-power roots
belong to \(\Omega\).
Indeed, \(u\) belongs to a finite layer; after adjoining a sufficiently
large finite ray layer containing \(\mu_{p^r}\), its Kummer extension
has \(p\)-power degree and is unramified outside \(p\).
The normal closure over \(L_1\) is still a finite \(p\)-extension,
by taking all conjugate Kummer generators over a finite ray layer.
It is therefore contained in \(\Omega\).
For \(g\in H\), define \(k_u(g)\in\mathbb Z_p\) by
\[
 \frac{g(u^{1/p^r})}{u^{1/p^r}}
             =\zeta_{p^r}^{\,k_u(g)\bmod p^r}\quad(r\geq1).
\]
This value is independent of compatible root choices.
Moreover \(k_{uv}(g)=k_u(g)+k_v(g)\).
In particular \(k_u(g)=0\) means that \(g\) fixes every
\(p\)-power root of the specific finite-level element \(u\).

Put \(w=|O_K^\times|\) and
\[
 I=\{(m_1,m_2)\in\mathbb Z_{\geq1}^2-\{(1,1)\}:
                      m_1\equiv m_2\pmod w\}.
\]
The elliptic Soul\'e characters \(\kappa_{\boldsymbol m}\),
\(\boldsymbol m\in I\), factor through \(A\)
by \cite[Lemma 4.1, p.~16]{IshiiSoule}.
Their relation with \(\rho\) is given by the following explicit
composition, as in
\cite[Definition 2.5 and Theorem 2.6, p.~8]{Ishii}:
\[
 G_{K(E[p^\infty])}
 \xrightarrow{\rho}F^3\Gamma_{1,1}
 \xrightarrow[\sim]{\text{boundary-preserving lift}}
 F^3\Gamma_{1,1}^{\dagger}
 \longrightarrow\Psi_1^\dagger
 \xrightarrow[\sim]{H}
       \mathbb Z_p[[T_1,T_2]](1).
\]
The middle projection is induced by the maximal metabelian quotient
of the geometric fundamental group.
This composition is \(\alpha_{1,1}\), whose coefficient formula is
\begin{equation}\label{eq:cm-radical-geometric-moments}
 \alpha_{1,1}(g)=
 \sum_{\substack{r,s\geq0\\ r+s\geq2\ {\rm even}}}
 \frac{\kappa_{(r+1,s+1)}(g)}
      {(1-p^{r+s})\,r!\,s!}\,U_1^rU_2^s,\qquad
 U_i=\log(1+T_i).
\end{equation}
Consequently every element of \(\ker\rho\) has
\(\kappa_{\boldsymbol m}=0\) for every \(\boldsymbol m\in I\).
We now prove a stronger intermediate statement: every \(g\in H\)
with all these character values zero fixes \(L_{\mathrm{ell}}\).

Fix a nontrivial ideal \(\mathfrak a\) prime to \(6p\).
Let \(\theta_{\mathfrak a}\) be the rational function of
\cite[Definition 2.1, p.~7]{IshiiSoule}.
For nonzero \(p\)-power torsion \(P\), its value is a \(p\)-unit
in the ray field of the conductor of \(P\).
It is exactly invariant under \(O_K^\times=\operatorname{Aut}(E)\).
The distribution and reciprocity formulas used below are
\cite[Propositions 2.1--2.2, p.~7]{IshiiSoule}.
For
\[
             Y=\mathbb Z_p^2-p\mathbb Z_p^2,
\]
define a measure on each primitive residue cylinder by
\begin{equation}\label{eq:cm-cylinder-kummer}
 \mu_{\mathfrak a,g}\bigl((a,b)+p^n\mathbb Z_p^2\bigr)
   =k_{\theta_{\mathfrak a}(a\omega_{1,n}+b\omega_{2,n})}(g),
 \qquad p\nmid\gcd(a,b).
\end{equation}
These are exact \(\mathbb Z_p\)-values.
The \(p^2\) refinements at level \(n+1\) correspond to all
\([p]\)-preimages of \(a\omega_{1,n}+b\omega_{2,n}\).
The distribution identity
\[
  \prod_{\nu\in E[p]}\theta_{\mathfrak a}(Q+\nu)
                       =\theta_{\mathfrak a}(pQ)
\]
therefore proves finite additivity under refinement.
Every cylinder value lies in \(\mathbb Z_p\), so this defines a bounded
\(\mathbb Z_p\)-valued measure on \(Y\).

The exact finite-level identity
\cite[Lemma 3.1, p.~11]{IshiiSoule} is
\begin{equation}\label{eq:cm-primitive-moment-identity}
 \epsilon_{\boldsymbol m,n}^{
       N\mathfrak a-\chi^{\boldsymbol1-\boldsymbol m}
                                      (\sigma_{\mathfrak a})}
 =\prod_{\substack{0\leq a,b<p^n\\p\nmid\gcd(a,b)}}
 \theta_{\mathfrak a}(a\omega_{1,n}+b\omega_{2,n})^{
                            a^{m_1-1}b^{m_2-1}}
 \quad\text{in }L_\infty^\times/L_\infty^{\times p^n}.
\end{equation}
By definition the \(p^n\)-Kummer character of
\(\epsilon_{\boldsymbol m,n}\) is
\(\kappa_{\boldsymbol m}\bmod p^n\).
This formula applies also to \(m_1=1\) or \(m_2=1\), with
\(0^0=1\); thus it includes the two axes.
Since a polynomial is constant modulo \(p^n\) on the cylinders
of level \(n\), taking Kummer values in
\eqref{eq:cm-primitive-moment-identity} for all \(n\) gives
\begin{equation}\label{eq:cm-primitive-moment-value}
 \int_Y x^r y^s\,d\mu_{\mathfrak a,g}
 =\bigl(N\mathfrak a-\chi^{(-r,-s)}
                          (\sigma_{\mathfrak a})\bigr)
                  \kappa_{(r+1,s+1)}(g)
\end{equation}
whenever \(r,s\geq0\), \(r+s>0\), and \(r\equiv s\pmod w\).
These moments vanish by our choice of \(g\).

For \(\varepsilon\in O_K^\times\), the action on \(Y\) is
\[
 (x,y)\longmapsto
        \bigl(\iota_{\mathfrak p}(\varepsilon)x,
              \iota_{\bar{\mathfrak p}}(\varepsilon)y\bigr).
\]
Here the two unit eigenvalues are inverse roots of unity.
The exact invariance of \(\theta_{\mathfrak a}\) makes
\(\mu_{\mathfrak a,g}\) invariant under this action.
When \(r-s\) is not divisible by \(w\), averaging over this finite
unit group makes the moment of \(x^ry^s\) vanish.
Consequently all nonconstant monomial moments vanish.
The constant is recovered without deleting either axis.
For \(N_j=(p-1)p^j\), the function \(z^{N_j}\) converges uniformly
on \(\mathbb Z_p\) to the indicator of \(\mathbb Z_p^\times\).
Hence
\[
             x^{N_j}+y^{N_j}-x^{N_j}y^{N_j}
                         \longrightarrow1\quad\text{uniformly on }Y.
\]
All three integrals vanish, and so the constant moment vanishes.
Polynomial functions over \(\mathbb Q_p\) are uniformly dense in
continuous \(\mathbb Q_p\)-valued functions on \(Y\):
extend a continuous function by zero to the clopen complement in
\(\mathbb Z_p^2\) and use its two-variable Mahler expansion.
Boundedness now gives \(\mu_{\mathfrak a,g}=0\).
In view of \eqref{eq:cm-cylinder-kummer}, \(g\) fixes every \(p\)-power
root of every \(\theta_{\mathfrak a}(P)\) for
\((\mathfrak a,6p)=1\) and every nonzero \(p\)-power torsion point
\(P\).

We next treat exactly the additional auxiliary ideals allowed on
an axis.  Suppose \((\mathfrak a,6\mathfrak p)=1\), allowing
\(\bar{\mathfrak p}\mid\mathfrak a\).
Choose
\[
              c=1+6pN\mathfrak a,\qquad \mathfrak b=(c).
\]
Then \(\mathfrak b\) is prime to \(6p\mathfrak a\).
For a generator \(\alpha\) of \(\mathfrak a\),
\cite[Lemma 2.4, p.~7]{IshiiSoule} gives the rational-function
identity
\[
 \frac{\theta_{\mathfrak b}(P)^{N\mathfrak a}}
      {\theta_{\mathfrak b}(\alpha P)}
 =
 \frac{\theta_{\mathfrak a}(P)^{c^2}}
      {\theta_{\mathfrak a}(cP)}.
\]
If \(P\) has \(\mathfrak p\)-power order, all four terms are defined.
The preceding paragraph kills the Kummer values of the two
\(\theta_{\mathfrak b}\)-terms, including at \(\alpha P\), and gives
\begin{equation}\label{eq:cm-auxiliary-axis}
 c^2 k_{\theta_{\mathfrak a}(P)}(g)
                      =k_{\theta_{\mathfrak a}(cP)}(g).
\end{equation}

The points \(v_n=\bar\pi^n\omega_{1,n}\) satisfy
\(\pi v_{n+1}=v_n\).
Use them to define a bounded measure on \(\mathbb Z_p^\times\) by
\[
 \nu_{\mathfrak a,g}(x+p^n\mathbb Z_p)
                    =k_{\theta_{\mathfrak a}(xv_n)}(g).
\]
Its refinement identity is the \([\pi]\)-distribution formula;
\((\mathfrak a,\mathfrak p)=1\) is precisely its required
coprimality hypothesis.
Equation \eqref{eq:cm-auxiliary-axis} says
\(\nu_{\mathfrak a,g}(cB)=c^2\nu_{\mathfrak a,g}(B)\)
for every cylinder \(B\).
With convolution defined by
\([t]\nu(B)=\nu(t^{-1}B)\), the resulting equation is exactly
\begin{equation}\label{eq:cm-axis-injective-operator}
                    (c^2-[c^{-1}])\nu_{\mathfrak a,g}=0.
\end{equation}
The measure algebra has the decomposition
\[
 \mathbb Z_p[[\mathbb Z_p^\times]]
   =\prod_{\eta\in\widehat{\mu_{p-1}}}\mathbb Z_p[[T]].
\]
Since \(c\equiv1\pmod p\), in each factor the operator in
\eqref{eq:cm-axis-injective-operator} is
\(c^2-(1+T)^s\), where
\(s=\log(c^{-1})/\log(\gamma)\in\mathbb Z_p\) for a generator
\(\gamma\) of \(1+p\mathbb Z_p\).
Its constant coefficient is \(c^2-1\ne0\).
Each factor is an integral domain, so multiplication by this
operator is injective.
It follows that \(\nu_{\mathfrak a,g}=0\).
Thus \(g\) fixes all roots of all allowed theta values on the
\(\mathfrak p\)-axis.  The same proof with
\(\mathfrak p,\bar{\mathfrak p}\) interchanged handles the other
axis.

Finally, \cite[Proposition 2.9, pp.~9--10]{IshiiSoule} gives all
generators of \(C(L_n)\): roots of unity and the
\(\sigma-1\) differences, \(\sigma\in\operatorname{Gal}(L_n/K)\),
of
\[
 \begin{array}{c|c}
 \text{theta value}&\text{allowed auxiliary ideal}\\ \hline
 \theta_{\mathfrak a}(\omega_{1,n}+\omega_{2,n})
                       &(\mathfrak a,6p)=1\\
 \theta_{\mathfrak a}(\omega_{1,n})
                       &(\mathfrak a,6\mathfrak p)=1\\
 \theta_{\mathfrak a}(\omega_{2,n})
                       &(\mathfrak a,6\bar{\mathfrak p})=1 .
 \end{array}
\]
All conjugate values were included in the preceding cylinder
arguments.
Every root of unity in \(L_\infty\) has all its \(p\)-power roots
in \(L_\infty\): the \(p\)-primary part lies in
\(\mu_{p^\infty}\), and raising to \(p\) is invertible on the
prime-to-\(p\) part.
We have therefore proved
\begin{equation}\label{eq:cm-radical-kernel-inclusion}
 \bigcap_{\boldsymbol m\in I}\ker\kappa_{\boldsymbol m}
        \subseteq
 \ker\!\left(H\longrightarrow
                  \operatorname{Gal}(L_{\mathrm{ell}}/L_\infty)\right).
\end{equation}

In particular every element of the absolute outer kernel fixes
\(L_{\mathrm{ell}}\), by
\eqref{eq:cm-radical-geometric-moments}.
Thus \(L_{\mathrm{ell}}\subseteq K(E[p])\Omega^*\).
On the other hand \(L_{\mathrm{ell}}\subseteq\Omega\) by the
finite Kummer-extension argument at the start of the proof.
Since \(K(E[p])/L_1\) has prime-to-\(p\) degree while
\(\Omega/\Omega^*\) is pro-\(p\),
\[
                 \Omega\cap K(E[p])\Omega^*=\Omega^*.
\]
This exact intersection proves the desired inclusion.
\end{proof}

\begin{corollary}\label{cor:cm-radical-joint-kernel}
Put \(\bar A=\operatorname{Gal}(L_{\mathrm{ell}}/L_\infty)\),
with its canonical quotient map \(A\to\bar A\).
Then
\[
 \bigcap_{\boldsymbol m\in I}\ker
       (\kappa_{\boldsymbol m}:A\longrightarrow
                           \mathbb Z_p(\boldsymbol m))
                         =\ker(A\longrightarrow\bar A).
\]
In particular the geometric elliptic Soul\'e characters factor
through, and jointly separate, this full radical quotient.
\end{corollary}

\begin{proof}
One containment is \eqref{eq:cm-radical-kernel-inclusion}.
For the other, let \(g\in H\) fix \(L_{\mathrm{ell}}\).
For a fixed auxiliary ideal \((\mathfrak a,6p)=1\), the value
\(u_n=\theta_{\mathfrak a}(\omega_{1,n}+\omega_{2,n})\)
is a global unit in \(L_n\).
Writing \(d_n=[L_n:K]\), the identity
\[
 u_n^{d_n}
  =N_{L_n/K}(u_n)
       \prod_{\tau\in\operatorname{Gal}(L_n/K)}
                                (u_n^{\tau-1})^{-1}
                         \in C(L_n)
\]
holds because the norm is a global unit of the imaginary
quadratic field \(K\), hence a root of unity.
Thus \(d_n k_{u_n}(g)=0\) in the torsionfree group
\(\mathbb Z_p\), and \(k_{u_n}(g)=0\).
This applies to all conjugates of \(u_n\).

Let \(\kappa_{\boldsymbol m,\mathfrak a}\) denote the moment
character of the norm-compatible mixed-unit system \((u_n)_n\).
Its explicit reduction modulo \(p^n\) is the Kummer character of
\[
 \prod_{\tau\in\operatorname{Gal}(L_n/K)}
             \tau(u_n)^{\chi^{\boldsymbol m-\boldsymbol1}(\tau)}
\]
by \cite[Lemma 4.2, p.~17]{IshiiSoule}.
Consequently \(\kappa_{\boldsymbol m,\mathfrak a}(g)=0\).
The exact comparison
\cite[Proposition 4.3, pp.~17--18]{IshiiSoule} is
\begin{equation}\label{eq:cm-radical-reverse-scalar}
 w\left(\frac1{1-\iota^{\boldsymbol m-\boldsymbol1}(\pi)}
       +\frac1{1-\iota^{\boldsymbol m-\boldsymbol1}(\bar\pi)}
       -1\right)\kappa_{\boldsymbol m,\mathfrak a}
  =\bigl(N\mathfrak a-
           \chi^{\boldsymbol1-\boldsymbol m}(\sigma_{\mathfrak a})
    \bigr)\kappa_{\boldsymbol m}.
\end{equation}
Here
\(\iota^{(r,s)}(\alpha)=
 \iota_{\mathfrak p}(\alpha)^r
 \iota_{\bar{\mathfrak p}}(\alpha)^s\).
Both denominators are nonzero:
the corresponding algebraic numbers have complex absolute value
\(p^{(m_1+m_2-2)/2}>1\).
Choose \(\mathfrak a=(c)\) with a rational integer \(c>1\)
prime to \(6p\).
The ambiguity of the CM reciprocity action by an element of
\(O_K^\times\) disappears because
\(m_1-m_2\equiv0\pmod w\).
The scalar on the right of
\eqref{eq:cm-radical-reverse-scalar} is therefore exactly
\[
                 c^2-c^{\,2-m_1-m_2}\ne0.
\]
Equation \eqref{eq:cm-radical-reverse-scalar} now implies
\(\kappa_{\boldsymbol m}(g)=0\) for every
\(\boldsymbol m\in I\).
This proves the reverse containment.
\end{proof}

The proposition uses no assertion that an individual character is
nonzero or surjective.  Its conclusion identifies a specific
abelian quotient of \(H\) through an exact radical construction.
The further map from this quotient to a proposed presentation by
free Iwasawa summands is the additional module-theoretic step
discussed in \cite[Remark 4.9.1, p.~24]{Ishii}; no such
presentation is inferred here.

\subsection{Finite conjugation-fixed class modules and source-group freeness}
\label{subsec:cm-finite-plus}

Continue with \(K,p,L_n,L_\infty,\Omega,H,A\) from
Subsection~\ref{subsec:cm-elliptic-radical}.
Complex conjugation is denoted by \(c\).
Put
\[
 \mathcal X_2=\varprojlim_n\operatorname{Cl}(L_n)[p^\infty],
 \qquad
 \mathcal J_2=\varinjlim_n\operatorname{Cl}(L_n)[p^\infty],
\]
where the first transition maps are ideal norms and the second
are extension of ideals.
The finite-level class groups here are the full class groups.
Let \(\mathcal X_2^+=(1+c)\mathcal X_2/2\), considered as a compact
\(\mathbb Z_p\)-module.
This definition does not assume that the fixed subgroup is stable
under the two-variable Iwasawa algebra.
Write
\[
 G=\operatorname{Gal}(L_\infty/K),\quad
 \Gamma=\operatorname{Gal}(L_\infty/L_1)\cong\mathbb Z_p^2,\quad
 L_{\mathrm c}=L_1(\mu_{p^\infty}),\quad
 \Gamma_{\mathrm a}=\operatorname{Gal}(L_\infty/L_{\mathrm c})
                   \cong\mathbb Z_p .
\]
Ray reciprocity identifies
\[
 G=\bigl(\mathbb Z_p^\times\times\mathbb Z_p^\times\bigr)/O_K^\times,
 \qquad
 c(x,y)c^{-1}=(y,x),\qquad
 \chi_{\mathrm{cyc}}(x,y)=xy.
\]
The embedded \(O_K^\times\) has the form
\((\varepsilon,\varepsilon^{-1})\).
Thus \(G/\langle g(g^c)^{-1}:g\in G\rangle\) is precisely the
cyclotomic quotient \(\mathbb Z_p^\times\), and
\(\Gamma_{\mathrm a}\) is its pro-\(p\) kernel.

\begin{lemma}\label{lem:cm-full-class-kummer}
Let
\[
 \mathcal U_\infty=O_{L_\infty}^{\times}
       =\bigcup_n O_{L_n}^{\times},\qquad
 L_{\mathrm u}=L_\infty(u^{1/p^r}:
                         u\in\mathcal U_\infty,\ r\geq1).
\]
There is a \(G\)- and \(c\)-equivariant exact sequence
\begin{equation}\label{eq:cm-full-class-kummer}
 0\longrightarrow
 \mathcal U_\infty\otimes_{\mathbb Z}\mathbb Q_p/\mathbb Z_p
 \longrightarrow
 B:=\operatorname{Hom}_{\mathrm{cont}}(A,\mu_{p^\infty})
 \xrightarrow{\vartheta}\mathcal J_2
 \longrightarrow0 .
\end{equation}
Consequently the restriction map has kernel
\begin{equation}\label{eq:cm-unit-annihilator}
 N_{\mathrm u}:=
 \ker\!\left(A\longrightarrow
            \operatorname{Gal}(L_{\mathrm u}/L_\infty)\right)
 \cong\operatorname{Hom}(\mathcal J_2,\mu_{p^\infty}).
\end{equation}
\end{lemma}

\begin{proof}
We spell out \(\vartheta\); this also proves that the right-hand term
is the full direct-limit class group.
The same construction for a cyclotomic tower is
\cite[Proposition 3.3.9 and Corollary 3.3.14]{SharifiIwasawaNotes}.
Its valuation proof applies to the present tower because it
contains the cyclotomic tower.

Represent a finite-order character in \(B\) by a Kummer class
\(u\in L_\infty^\times/L_\infty^{\times p^r}\).
At every prime outside \(p\), its valuation is divisible by \(p^r\),
because the corresponding Kummer extension is unramified there.
Pass to a finite ray layer containing \(u\) and sufficiently high
in the cyclotomic direction that all valuations at primes over
\(p\) are divisible by \(p^r\) as well.
This is possible since their ramification indices grow without
bound.
At that layer there is a unique fractional ideal \(\mathfrak b\)
with
\[
                         (u)=\mathfrak b^{p^r}.
\]
Set \(\vartheta(u)=[\mathfrak b]\in\mathcal J_2\).
Replacing \(u\) by \(uv^{p^r}\) changes \(\mathfrak b\) by
the principal ideal \((v)\); going to a higher layer extends
\(\mathfrak b\).
Representing the same character at exponent \(p^{r+1}\) replaces
\(u\) by \(u^p\) and leaves \(\mathfrak b\) unchanged.
Thus the map is well defined and equivariant.

Every \(p^r\)-torsion ideal class at a finite layer has a
representative \(\mathfrak b\) with
\(\mathfrak b^{p^r}=(u)\); the associated Kummer extension is
unramified outside \(p\).
This proves surjectivity.
The class \([\mathfrak b]\) dies in the direct limit exactly when,
at a higher finite layer, \(\mathfrak b=(v)\).
Then \(uv^{-p^r}\) is a global unit.
This proves the asserted kernel.
The left Kummer map is injective because a \(p^r\)-th root of a
global unit, whenever it belongs to \(L_\infty\), is again a global
unit.
Taking the Pontryagin dual with values in \(\mu_{p^\infty}\)
gives \eqref{eq:cm-unit-annihilator}.
The Tate twist in this duality is retained: \(c\) acts by \(-1\)
on \(\mu_{p^\infty}\).
\end{proof}

\begin{lemma}\label{lem:cm-anticyclotomic-invariants}
Without a class-group finiteness assumption, one has
\[
 A^{\Gamma_{\mathrm a}}=0,\qquad
 H^2(H,\mathbb Q_p/\mathbb Z_p)=0 .
\]
\end{lemma}

\begin{proof}
Put \(H_{\mathrm c}=\operatorname{Gal}(\Omega/L_{\mathrm c})\).
These are the full maximal pro-\(p\) Galois groups unramified outside
\(p\) of their indicated infinite base fields.
To check the comparison, a finite Galois \(p\)-extension over either
infinite field descends to a finite subfield of that base,
which is a \(p\)-extension of \(L_1\).
After increasing that finite subfield, take the descended
Galois \(p\)-extension and adjoin all its conjugates over \(L_1\);
the resulting normal closure has \(p\)-power degree and is
unramified outside \(p\), and therefore belongs to \(\Omega\).
Thus we have the exact sequence of arithmetic groups
\begin{equation}\label{eq:cm-anticyclotomic-group-extension}
 1\longrightarrow H\longrightarrow H_{\mathrm c}
                 \longrightarrow\Gamma_{\mathrm a}
                 \longrightarrow1 .
\end{equation}

Iwasawa's theorem on weak Leopoldt for a cyclotomic
\(\mathbb Z_p\)-extension gives
\[
 H^2(H_{\mathrm c},\mathbb Q_p/\mathbb Z_p)=0;
\]
see \cite[Remark 4.3, author version p.~26]{VenjakobIwasawa},
which explicitly states the result for every number field and
refers to its cohomological formulation.
This theorem concerns the trivial coefficient
\(\mathbb Q_p/\mathbb Z_p\).
If it is formulated first for the full \(S\)-ramified Galois
group, its inflation to the maximal pro-\(p\) quotient in degree
two is injective.
Indeed, an invariant nonzero homomorphism from the kernel of
the maximal pro-\(p\) quotient to \(\mathbb F_p\) would produce
a larger pro-\(p\) quotient; the same assertion for
\(\mathbb Q_p/\mathbb Z_p\) follows by taking a nonzero
order-\(p\) multiple.
The five-term sequence proves the required injectivity.

Let \(M=H^1(H,\mathbb Q_p/\mathbb Z_p)=A^\vee\).
Since \(\operatorname{cd}_p(\Gamma_{\mathrm a})=1\),
Hochschild--Serre for
\eqref{eq:cm-anticyclotomic-group-extension} gives
\begin{equation}\label{eq:cm-anticyclotomic-HS}
 0\longrightarrow H^1(\Gamma_{\mathrm a},M)
 \longrightarrow H^2(H_{\mathrm c},\mathbb Q_p/\mathbb Z_p)
 \longrightarrow
 H^2(H,\mathbb Q_p/\mathbb Z_p)^{\Gamma_{\mathrm a}}
 \longrightarrow0 .
\end{equation}
For a generator \(\gamma\) of \(\Gamma_{\mathrm a}\),
\[
 H^1(\Gamma_{\mathrm a},M)=M/(\gamma-1)M,
 \qquad
 \bigl(M/(\gamma-1)M\bigr)^\vee=A^{\Gamma_{\mathrm a}}.
\]
The first vanishing follows.
Equation \eqref{eq:cm-anticyclotomic-HS} also makes the
\(\Gamma_{\mathrm a}\)-invariants of the indicated \(H^2(H,-)\)
zero.
A nonzero discrete \(p\)-primary module under a pro-\(p\) group
always has nonzero invariants: choose a nonzero element of order
\(p\), take the finite-dimensional span of its finite orbit,
and use the nonzero fixed subspace of a finite \(p\)-group
acting on a nonzero \(\mathbb F_p\)-space.
Thus \(H^2(H,\mathbb Q_p/\mathbb Z_p)=0\) as well.
No freeness of \(H\) was used.
\end{proof}

\begin{proposition}\label{prop:cm-finite-plus-free}
For the preceding CM ray tower, suppose that the explicitly
defined group \(\mathcal X_2^+\) is finite.
Then
\[
 \mathcal J_2=0,\qquad
 \operatorname{Gal}(L_{\mathrm u}/L_\infty)=A,
 \qquad H\text{ is a free pro-}p\text{ group}.
\]
In particular every finite-layer \(p\)-primary ideal class becomes principal
at some higher ray layer, and the maximal abelian subextension
of \(\Omega/L_\infty\) is generated by \(p\)-power roots of
global units.
\end{proposition}

\begin{proof}
First, every norm
\[
 \operatorname{Cl}(L_{n+1})[p^\infty]
           \longrightarrow\operatorname{Cl}(L_n)[p^\infty]
\]
is surjective.
Ray reciprocity identifies
\(\operatorname{Gal}(L_{n+1}/L_n)\) with the product of the
two principal-unit quotients, each of order \(p\).
The inertia at \(\mathfrak p\) contains the first factor and
the inertia at \(\bar{\mathfrak p}\) contains the second.
Hence \(L_{n+1}/L_n\) has no nontrivial everywhere-unramified
subextension.
By the Artin reciprocity description of the image of the ideal
norm, this is exactly the required surjectivity on class groups.
The same holds for all composite norm maps.
All maps commute with \(c\); averaging by \((1+c)/2\) proves
surjectivity on their fixed subgroups.

It follows that \(\mathcal X_2^+\) surjects onto every
\(\operatorname{Cl}(L_n)[p^\infty]^+\).
Since \(\mathcal X_2^+\) is finite, their orders are eventually
constant, and for all sufficiently large \(m\geq n\) the norm
\[
 N_{m/n}:\operatorname{Cl}(L_m)[p^\infty]^+
                       \xrightarrow{\sim}
                         \operatorname{Cl}(L_n)[p^\infty]^+
\]
is an isomorphism.
Let \(j_{m/n}\) be extension of ideals.
The exact identity
\[
                  N_{m/n}j_{m/n}=[L_m:L_n]
                               =p^{2(m-n)}
\]
holds on \(\operatorname{Cl}(L_n)\).
For fixed sufficiently large \(n\), this multiplication kills
the finite group \(\operatorname{Cl}(L_n)[p^\infty]^+\) when
\(m\) is sufficiently large.
Injectivity of \(N_{m/n}\) on the plus subgroup then gives
\(j_{m/n}=0\) there.
Consequently
\begin{equation}\label{eq:cm-direct-limit-plus-zero}
                         \mathcal J_2^+=0.
\end{equation}
This argument uses the exact norm and extension maps and does not
regard a plus subgroup as a two-variable Iwasawa submodule.

By \eqref{eq:cm-unit-annihilator} and the Tate twist in that
pairing, \eqref{eq:cm-direct-limit-plus-zero} says that \(c\)
acts as \(+1\) on \(N_{\mathrm u}\).
The subgroup \(N_{\mathrm u}\) is stable under \(G\), because
global units and their radical field are stable under \(G\).
For \(g\in G\) and \(a\in N_{\mathrm u}\), one therefore has
\[
                    ga=c(ga)=g^c\,ca=g^c a.
\]
Hence every \(g(g^c)^{-1}\) acts trivially on \(N_{\mathrm u}\).
The explicit ray-reciprocity quotient preceding the lemmas shows
that \(N_{\mathrm u}\subseteq A^{\Gamma_{\mathrm a}}\).
Lemma~\ref{lem:cm-anticyclotomic-invariants} makes this group zero.
Thus \(N_{\mathrm u}=0\).
The perfect Kummer duality
\eqref{eq:cm-unit-annihilator} now also yields
\(\mathcal J_2=0\), and proves the assertion about the maximal
abelian subextension.

It remains to prove freeness.
Equation \eqref{eq:cm-full-class-kummer} and
\(\mathcal J_2=0\) identify
\[
 H^1(H,\mu_{p^\infty})
       =\mathcal U_\infty\otimes_{\mathbb Z}
                                     \mathbb Q_p/\mathbb Z_p.
\]
The group on the right is \(p\)-divisible.
Since \(H\) acts trivially on \(\mu_{p^\infty}\), choosing
compatible roots of unity identifies it with
\(H^1(H,\mathbb Q_p/\mathbb Z_p)\).
The cohomology sequence for
\[
 0\longrightarrow\mathbb F_p
  \longrightarrow\mathbb Q_p/\mathbb Z_p
  \xrightarrow{p}\mathbb Q_p/\mathbb Z_p
  \longrightarrow0
\]
therefore gives \(H^2(H,\mathbb F_p)=0\), using the already
proved vanishing of \(H^2(H,\mathbb Q_p/\mathbb Z_p)\).
The cohomological freeness criterion for pro-\(p\) groups proves
that \(H\) is free.
For the same criterion in this arithmetic setting, see
\cite[Theorem 4.17 and its proof, author version p.~33]
{VenjakobIwasawa}.
\end{proof}

The group whose freeness has just been proved is the
source \(H=\operatorname{Gal}(\Omega/L_\infty)\).
To obtain a statement about an outer quotient, one must still
apply the specific restriction map defined by \(\rho\).
The radical calculation of the preceding subsection supplies
that map on the abelian quotient once the finite-level
elliptic-unit/global-unit index is compared.

\subsubsection{Character components and the exact Kummer twist}

Let
\[
 \Delta=\operatorname{Gal}(L_1/K)
       \cong\bigl((\mathbb F_p^\times)^2\bigr)/O_K^\times .
\]
We regard \(\Delta\) as the canonical prime-to-\(p\) factor of
\(G=\Delta\times\Gamma\).
All its characters take values in \(\mathbb Z_p^\times\), since
its exponent divides \(p-1\).
For such a character \(\chi\), write
\[
 e_\chi=\frac1{|\Delta|}
             \sum_{\delta\in\Delta}\chi(\delta)^{-1}\delta,
 \qquad
 \chi^c(\delta)=\chi(c\delta c^{-1}),\qquad
 \chi^\dagger=\omega_{\mathrm{cyc}}\chi^{-1},
\]
where \(\omega_{\mathrm{cyc}}=\chi_{\mathrm{cyc}}|_\Delta\).
For a \(c\)-stable set \(\mathcal O\) of characters, put
\[
 e_{\mathcal O}=\sum_{\chi\in\mathcal O}e_\chi,\qquad
 \mathcal O^\dagger=\{\chi^\dagger:\chi\in\mathcal O\}.
\]
The dagger operation commutes with \(c\).
In the pair notation
\(\chi_{r,s}(\delta_1,\delta_2)=\delta_1^r\delta_2^s\),
where \(r-s\equiv0\pmod w\), these operations are exactly
\[
                  (r,s)^c=(s,r),\qquad
                  (r,s)^\dagger=(1-r,1-s)
                                \pmod{p-1}.
\]

\begin{proposition}\label{prop:cm-component-unit-detection}
Let \(\mathcal O\) be a \(c\)-stable set of characters of
\(\Delta\), and suppose that
\((e_{\mathcal O}\mathcal X_2)^+\) is finite.
Then
\[
 e_{\mathcal O}\mathcal J_2=0,\qquad
 e_{\mathcal O^\dagger}N_{\mathrm u}=0.
\]
In particular, the canonical abelian field
\[
 L_{\mathcal O}
   =\Omega^{\,\ker(H\longrightarrow A
                       \xrightarrow{\,e_{\mathcal O^\dagger}\,}
                                e_{\mathcal O^\dagger}A)}
\]
satisfies \(L_{\mathcal O}\subseteq L_{\mathrm u}\).
\end{proposition}

\begin{proof}
The idempotent \(e_{\mathcal O}\) commutes with \(c\), with
extension of ideals, and with ideal norm.
The norm maps on
\((e_{\mathcal O}\operatorname{Cl}(L_n)[p^\infty])^+\)
are consequently surjective.
The eventual-isomorphism argument in the preceding proof,
using \(N_{m/n}j_{m/n}=p^{2(m-n)}\), gives
\[
                      (e_{\mathcal O}\mathcal J_2)^+=0.
\]

We record the component map in the Kummer pairing, including
its twist.  For \(a\in N_{\mathrm u}\),
\(b\in\mathcal J_2\), and \(\delta\in\Delta\), equivariance is
\[
              \langle\delta a,\delta b\rangle
                    =\delta\langle a,b\rangle
                    =\omega_{\mathrm{cyc}}(\delta)
                                         \langle a,b\rangle,
\]
where the pairing is written additively.
Thus a \(\psi\)-component of \(N_{\mathrm u}\) pairs with a
\(\chi\)-component of \(\mathcal J_2\) exactly when
\(\psi\chi=\omega_{\mathrm{cyc}}\).
The isomorphism \eqref{eq:cm-unit-annihilator} therefore
restricts to the exact isomorphism
\[
 e_{\mathcal O^\dagger}N_{\mathrm u}
       \cong
 \operatorname{Hom}(e_{\mathcal O}\mathcal J_2,
                                      \mu_{p^\infty}).
\]
Because the source on the right has only the \(c\)-minus part,
the group on the left is \(c\)-fixed.
It is stable under \(\Gamma\), since \(\Delta\) commutes with
\(\Gamma\).
For \(a\) in this subgroup and \(\gamma\in\Gamma\), one has
\(\gamma a=c(\gamma a)=\gamma^c a\).
The elements \(\gamma(\gamma^c)^{-1}\) generate
\(\Gamma_{\mathrm a}\), so this subgroup lies in
\(A^{\Gamma_{\mathrm a}}\).
Lemma~\ref{lem:cm-anticyclotomic-invariants} makes it zero.
Perfect duality then gives \(e_{\mathcal O}\mathcal J_2=0\).

Finally the restriction kernel
\(\ker(H\to\operatorname{Gal}(L_{\mathrm u}/L_\infty))\)
maps into \(N_{\mathrm u}\) in \(A\).
Its \(e_{\mathcal O^\dagger}\)-projection is zero, so it is
contained in the kernel defining \(L_{\mathcal O}\).
The Galois correspondence proves
\(L_{\mathcal O}\subseteq L_{\mathrm u}\).
\end{proof}

\subsection{The exact finite-level elliptic-unit comparison}
\label{subsec:cm-unit-index}

\begin{lemma}\label{lem:cm-unit-radical-comparison}
For every \(n\geq1\), the exact group \(C(L_n)\) in
\cite[Definition 2.2, p.~8]{IshiiSoule} has finite index in
\(O_{L_n}^{\times}\).
Consequently the two radical fields coincide:
\[
                         L_{\mathrm{ell}}=L_{\mathrm u}.
\]
\end{lemma}

\begin{proof}
We prove full rank with the defining theta values.
The formula below identifies their generators with the
Robert--Ramachandra invariants used in the logarithmic calculation.
Let \(F=L_n\), \(D=\operatorname{Gal}(F/K)\), and let
\(\chi\ne1\) be a complex character of \(D\).
Its conductor \(\mathfrak m\mid(p^n)\) is nontrivial, because
\(K\) has class number one.
Put \(F_{\mathfrak m}=K(\mathfrak m)\), and write
\(\chi_{\mathfrak m}\) for the induced primitive character.

Here is the required exact normalization.
For a lattice \(L\), the function
\[
 \varphi(z,L)=
       \bigl(e^{-\eta(z,L)z/2}\sigma(z,L)\bigr)^{12}\Delta(L)
\]
in \cite[(2), p.~255]{Oukhaba2009} is the fundamental theta
function in \cite[Definition 3.1(i), p.~6]{GotoIshii2026}.
The quotient formula in
\cite[Definition 3.1(ii), equation (1), p.~6]{GotoIshii2026}
gives, with precisely the rational function used by Ishii,
\[
 \theta_{\mathfrak a}(z;L)=
              \frac{\varphi(z,L)^{N\mathfrak a}}
                   {\varphi(z,\mathfrak a^{-1}L)}.
\]
Write \(L=\Omega_E O_K\), \(\mathfrak m=(\beta)\), and
\(\tau=\Omega_E/\beta\bmod L\).
Scaling both \(z\) and \(L\) leaves \(\varphi\) unchanged:
the factor \(\sigma^{12}\) scales by the twelfth power,
\(\Delta\) by its inverse, and \(\eta(z,L)z\) is invariant.
Thus \((\tau,L)\) may be replaced by \((1,\mathfrak m)\).

Let \(e_{\mathfrak m}>0\) generate
\(\mathfrak m\cap\mathbb Z\).
Oukhaba's invariants and reciprocity, equations (3)--(4) on
p.~255, state
\[
 \varphi_{\mathfrak m}(\sigma_{\mathfrak a})
      =\varphi(1,\mathfrak a^{-1}\mathfrak m)^{e_{\mathfrak m}},
 \qquad
 \varphi_{\mathfrak m}(1)^{\sigma_{\mathfrak a}}
      =\varphi_{\mathfrak m}(\sigma_{\mathfrak a}).
\]
Combining the displayed formulas gives the exact algebraic identity
\begin{equation}\label{eq:cm-exact-unit-convention}
 \theta_{\mathfrak a}(\tau)^{e_{\mathfrak m}}
       =\varphi_{\mathfrak m}(1)^{
                                  N\mathfrak a-\sigma_{\mathfrak a}}.
\end{equation}
The integer \(e_{\mathfrak m}\), which may be divisible by \(p\),
has been retained.

For \(u\in F^\times\), define
\[
 \ell_\chi^F(u)=
       \sum_{\sigma\in D}\chi(\sigma)^{-1}
                                    \log|\sigma(u)|^2.
\]
Then
\(\ell_\chi^F(u^\delta)=\chi(\delta)\ell_\chi^F(u)\).
For \(u\in F_{\mathfrak m}^\times\), one has
\[
 \ell_\chi^F(u)=[F:F_{\mathfrak m}]
                       \ell_{\chi_{\mathfrak m}}^{F_{\mathfrak m}}(u).
\]
The Kronecker limit formula
\cite[(6), p.~256]{Oukhaba2009}, applied at the primitive
conductor \(\mathfrak m\), gives
\[
 \ell_{\chi_{\mathfrak m}}^{F_{\mathfrak m}}
                         (\varphi_{\mathfrak m}(1))
   =-12r_{\mathfrak m}e_{\mathfrak m}
                          L'(0,\chi_{\mathfrak m}^{-1})\ne0 .
\]
Here \(r_{\mathfrak m}\) counts roots of unity of \(K\)
congruent to \(1\bmod\mathfrak m\).
The simple zero and nonzero derivative for a nontrivial primitive
character are also recorded explicitly in
\cite[\S2.3, p.~6, before (2.3)--(2.4)]{Oukhaba2003}.

Take \(c_0=1+6p\) and \(\mathfrak a=(c_0)\).
This auxiliary ideal is prime to \(6\mathfrak m\), and
\(N\mathfrak a-\chi_{\mathfrak m}(\sigma_{\mathfrak a})\ne0\):
the first term is the real integer \(c_0^2>1\) and the second
is a root of unity.
Equation \eqref{eq:cm-exact-unit-convention} therefore gives
\[
 \ell_\chi^F(\theta_{\mathfrak a}(\tau))
  =-12r_{\mathfrak m}[F:F_{\mathfrak m}]
    \bigl(N\mathfrak a-\chi_{\mathfrak m}(\sigma_{\mathfrak a})\bigr)
                         L'(0,\chi_{\mathfrak m}^{-1})\ne0.
\]
Choose \(\delta\in D\) with \(\chi(\delta)\ne1\).
The element
\[
                       v_\chi=\theta_{\mathfrak a}(\tau)^{\delta-1}
\]
is a generator allowed in Ishii's exact definition of \(C(F)\).
Indeed, \(F_{\mathfrak m}\subseteq F\), so the norm in that
definition is the identity.
The unit-injectivity requirement for \(\mathfrak m\) holds:
the norm of a difference of distinct roots of unity in an
imaginary quadratic field has prime divisors only \(2\) and \(3\),
while \(\mathfrak m\) is supported above \(p\geq5\).
For one-prime conductor the augmentation is still a global unit,
as in that definition; this also follows from the equal valuation
at all primes over that conductor in
\cite[(5), p.~256]{Oukhaba2009}.
Its \(\chi\)-logarithm is
\[
                \ell_\chi^F(v_\chi)
                   =(\chi(\delta)-1)
                      \ell_\chi^F(\theta_{\mathfrak a}(\tau))\ne0 .
\]

Dirichlet's logarithmic embedding of
\(O_F^\times/\mu(F)\) spans the real augmentation space of
\(\mathbb R[D]\).
This follows by choosing the embeddings extending the fixed
complex embedding of \(K\), one from each complex place.
After complexification, each nontrivial character of \(D\)
occurs exactly once.
The image of \(C(F)\) is \(D\)-stable, and the preceding
construction gives its nonzero projection to every such line.
It therefore has the full unit rank.
Since it is a subgroup of the finitely generated unit group
and already contains its torsion, its index is finite.

Finally, for any \(u\in O_{L_n}^\times\) there is an integer
\(d>0\) with \(u^d\in C(L_n)\).
If \(g\in H\) fixes \(L_{\mathrm{ell}}\), its exact Kummer
character satisfies \(d\,k_u(g)=0\) in \(\mathbb Z_p\).
Torsionfreeness gives \(k_u(g)=0\).
Thus \(g\) fixes all \(p\)-power roots of every global unit,
even when \(p\mid d\).
The reverse inclusion of radical fields follows from
\(C(L_n)\subseteq O_{L_n}^\times\), proving
\(L_{\mathrm{ell}}=L_{\mathrm u}\).
\end{proof}

\begin{theorem}\label{thm:cm-finite-plus-abelian}
Let \(\mathcal O\) be a \(c\)-stable set of characters of
\(\Delta=\operatorname{Gal}(K(p)/K)\), and suppose that
\((e_{\mathcal O}\mathcal X_2)^+\) is finite.
With
\(\mathcal O^\dagger=\{\omega_{\mathrm{cyc}}\chi^{-1}:
                                         \chi\in\mathcal O\}\),
the canonical field
\[
 L_{\mathcal O}=
 \Omega^{\,\ker(H\longrightarrow A
                \xrightarrow{\,e_{\mathcal O^\dagger}\,}
                              e_{\mathcal O^\dagger}A)}
\]
is contained in \(\Omega^*\).
Equivalently the image in \(A\) of the outer kernel has
zero projection to \(e_{\mathcal O^\dagger}A\).

For the set of all characters, finiteness of
\(\mathcal X_2^+\) gives
\[
 \Omega^{[H,H]}=L_{\mathrm{ell}}=L_{\mathrm u}
                   \subseteq\Omega^*,\qquad
 \bigcap_{\boldsymbol m\in I}\ker(\kappa_{\boldsymbol m}|_A)=0.
\]
The source \(H\) is then free pro-\(p\), by
Proposition~\ref{prop:cm-finite-plus-free}.
\end{theorem}

\begin{proof}
Proposition~\ref{prop:cm-component-unit-detection} gives
\(L_{\mathcal O}\subseteq L_{\mathrm u}\).
The preceding lemma identifies
\(L_{\mathrm u}=L_{\mathrm{ell}}\), and
Proposition~\ref{prop:cm-full-radical} places this field in
\(\Omega^*\).
These are actual inclusions inside \(\Omega\), so their
composition proves the component assertion.
For all characters, Proposition~\ref{prop:cm-finite-plus-free}
identifies \(L_{\mathrm u}\) with the maximal abelian
subextension \(\Omega^{[H,H]}\).
Corollary~\ref{cor:cm-radical-joint-kernel} then gives the
stated joint-character injectivity.
\end{proof}

\subsection{A source quotient detected by the CM depth-two measure}
\label{subsec:cm-measure-class-two}

Assume \(\mathcal X_2^+\) is finite and retain
\(H=\operatorname{Gal}(\Omega/L_\infty)\) and
\(A=H^{\mathrm{ab}}\).  The proved CM results make \(H\)
free pro-\(p\) and make the joint elliptic Soul\'e map
injective on \(A\).

There is a second useful measure realization of this map.
Put \(V=T_p(E)\simeq\mathbb Z_p^2\), written additively.
By \cite[Lemma 2.4, Definition 2.5 and Theorem 2.6,
pp.~7--8]{Ishii}, the canonical boundary-preserving
lift of the outer action gives
\[
 \alpha:A\longrightarrow
       \mathbb Z_p[[V]](1).
\]
Here \(\mathbb Z_p[[V]]\) denotes bounded measures on \(V\).
If \(U_i=\log(1+T_i)\), its exponential transform is
\[
 f_h(U_1,U_2)
 =\sum_{\substack{r,s\geq0\\r+s\geq2\ {\rm even}}}
 \frac{\kappa_{(r+1,s+1)}(h)}
      {(1-p^{r+s})r!s!}\,U_1^rU_2^s .
\]
The coefficients outside the admissible CM weight set
vanish.  The already proved joint-kernel equality makes
\(\alpha\) injective.

Define two automorphisms of \(V\times V\) by
\[
 t_1(x,y)=(x,x+y),\qquad t_2(x,y)=(x+y,y),
\]
and define the bounded-measure operator
\[
 \mathcal J=(t_1)_*+(t_2)_*-\operatorname{id}.
\]
The one-dimensional restriction of this operator occurs in
\cite[equations (29) and (32), p.~233]{LuqueNovelliThibon}.
To specify the comparison, send their coefficient
\(\phi_{n+1}\) to \(U^n/n!\), so that
\(\Phi(t)\) becomes \(\exp(tU)\), the exponential transform
of the point measure \(\delta_t\). Send an ordinary bracket
\([f,g]\) to \(f(U)g(V)-g(U)f(V)\). Their equation (32)
then becomes
\[
 \delta_s\wedge\delta_{s+t}
 +\delta_{s+t}\wedge\delta_t-\delta_s\wedge\delta_t
                 =\mathcal J(\delta_s\wedge\delta_t).
\]
For every \(v\in V\), the map \(i_v(t)=tv\) preserves
addition and gives
\((i_v\times i_v)_*\mathcal J_1
 =\mathcal J(i_v\times i_v)_*\).
The source calculation below proves the displacement formula
for the CM arithmetic action on the two-dimensional Tate module.

Use the alternating tensor inclusion to define
\[
 J_A:\widehat{\bigwedge}^2 A
 \longrightarrow\mathbb Z_p[[V\times V]](2),\qquad
 a\wedge b\longmapsto
 \mathcal J\bigl(\alpha(a)\otimes\alpha(b)
                         -\alpha(b)\otimes\alpha(a)\bigr).
\]
Let \(N_2^{\rm CM}\) be the inverse image of
\(\ker J_A\) under the canonical map
\(\gamma_2H\to\gamma_2H/\gamma_3H
 \simeq\widehat{\bigwedge}^2A\).

\begin{proposition}\label{prop:cm-decomposable-commutators}
For \(a,b\in A\) whose images in \(A[1/p]\) are linearly
independent, \(J_A(a\wedge b)\ne0\). Thus every such pair
has a nonzero commutator in \(Q_2^{\rm CM}\).
\end{proposition}
\begin{proof}
Write \(f=\alpha(a)\), \(g=\alpha(b)\) in the exponential
coordinates above. Both have zero constant and linear terms.
The exponential transform is injective: the Amice transform
identifies bounded measures with integral power series, and
the substitution \(T_i=\exp(U_i)-1\) has the inverse
\(U_i=\log(1+T_i)\) over \(\mathbb Q_p\).
The transform of \(J_A(a\wedge b)\) is
\begin{align*}
 &f(u)g(u+v)-g(u)f(u+v)\\
 &\quad+f(u+v)g(v)-g(u+v)f(v)-f(u)g(v)+g(u)f(v).
\end{align*}
Its coefficient linear in \(v_i\) is
\(f(u)\partial_i g(u)-g(u)\partial_i f(u)\), since every
term containing \(f(v)\) or \(g(v)\) has order at least two.
If the displayed transform vanishes, both derivatives of
\(g/f\) vanish in \(\operatorname{Frac}(\mathbb Q_p[[U_1,U_2]])\).
Embed this field into \(\mathbb Q_p((U_1))((U_2))\).
Every nonzero power series is invertible there after extracting
its first nonzero \(U_2\)-coefficient; the embedding commutes
with formal derivatives. The \(U_2\)-derivative forces the
ratio to belong to \(\mathbb Q_p((U_1))\), and its
\(U_1\)-derivative forces it to belong to \(\mathbb Q_p\),
by coefficient comparison in characteristic zero. Thus
\(g=cf\). Injectivity of \(\alpha[1/p]\) gives \(b=ca\),
contradicting independence. The commutator assertion follows
from the quotient theorem below.
\end{proof}

\begin{theorem}\label{thm:cm-measure-class-two}
The subgroup \(N_2^{\rm CM}\) is closed and stable under
the arithmetic action, and contains the outer
kernel.  The quotient
\[
                       Q_2^{\rm CM}=H/N_2^{\rm CM}
\]
is torsion-free of class at most two, has abelianization
\(A\), and has central commutator module
\(\operatorname{im}J_A\).
Its class is exactly two for every class-number-one
CM field in the stated range.
\end{theorem}

\begin{proof}
Let \(\Pi\) be the geometric pro-\(p\) fundamental group
of \(E-\{O\}\), with the generators \(x_1,x_2\)
of \cite[Lemma 2.1]{Ishii}. Put \(C=[\Pi,\Pi]\), \(C_r=\gamma_rC\), and
\(M=C/C_2=\mathbb Z_p[[V]]\bar z\), with
\(z=[x_2,x_1]\) as in \cite[Lemma 2.1]{Ishii}.
The closed subgroup \(C\) of the free pro-\(p\) group
\(\Pi\) is free pro-\(p\).  The integral class-two Hall
basis, first on finite subsets of a convergent free
basis and then by inverse limits, gives
\[
 C_2/C_3=\widehat{\bigwedge}^{\,2}M,
 \qquad \gamma_2H/\gamma_3H=
            \widehat{\bigwedge}^{\,2}A.
\]
The exterior square identifies with the antisymmetric
submodule of the completed tensor square by
\(a\wedge b\mapsto a\otimes b-b\otimes a\).
Since \(p\) is odd, half the exterior projection is
its inverse on this submodule.  Thus
\(C_2/C_3\) is a closed submodule of
\(\mathbb Z_p[[V\times V]](2)\).

For \(v\in V\), write \(e_v=v\bar z\), and let
\(T_v\) denote translation on \(M\) and simultaneous
translation on its tensor square.  Let \(f=s(h)\)
be the canonical special lift of \(h\in H\), and set
\(a=\alpha(h)\bar z\).  This lift fixes \(z\) exactly:
its exponent on the inertia generator is its
abelian determinant, which is one.  For \(q\in\Pi\)
with abelian image \(v\), put
\(k_f(q)=f(q)q^{-1}\in C\).
The crossed-homomorphism identity and Ishii's defining
formula on \(x_1,x_2\) give
\[
                 k_f(q)\bmod C_2=(T_v-1)a.
\]
This follows first on words and then by continuity.
Because \(C\) is normally generated by \(z\), the identity
\[
 f(qzq^{-1})=k_f(q)(qzq^{-1})k_f(q)^{-1}
\]
shows that \(f\) acts trivially on \(M\).  Its first
nontrivial action on \(C\) is the continuous map
\[
 \delta_a:M\longrightarrow C_2/C_3,
 \qquad \delta_a(e_v)=[T_va-a,e_v].
\]
In particular, it depends only on \(\alpha(h)\), and
\begin{equation}\label{eq:cm-delta-translation}
 \delta_a(T_vb)=T_v\delta_a(b)+[T_va-a,T_vb].
\end{equation}
The formula follows on point measures by substitution
and extends to all bounded measures by continuity.

Take another lift \(g=s(k)\), with
\(b=\alpha(k)\bar z\).  Their metabelian actions
commute.  In the class-two group \(C/C_3\), the BCH
law involves only the invertible denominator \(2\).
Use
\(k_{fg}(q)=f(k_g(q))k_f(q)\) and its reversed-order
counterpart.  Their unknown central displacements
cancel, and the surviving terms are
\begin{align*}
 k_{[f,g]}(q)\bmod C_3
  ={}&\delta_a((T_v-1)b)-\delta_b((T_v-1)a)\\
     &+[(T_v-1)b,(T_v-1)a].
\end{align*}
Indeed a special automorphism with all displacements
in \(C_2\) acts trivially on \(C/C_3\), by the
conjugate-of-\(z\) identity.  Hence replacing \(q\)
by \(gf(q)\) when passing from \(fg,gf\) to
\([f,g]\) does not change this displayed displacement.
Writing \(A_0=T_va,B_0=T_vb\), the bracket terms
introduced by \eqref{eq:cm-delta-translation} simplify as
\[
 [A_0-a,B_0]-[B_0-b,A_0]+[B_0-b,A_0-a]
                         =[A_0,B_0]-[a,b].
\]
We therefore obtain the exact identity
\[
 k_{[f,g]}(q)\bmod C_3
       =(T_v-1)(\delta_ab-\delta_ba+[a,b]).
\]
From the point-measure formula for \(\delta\),
\begin{align*}
 \delta_ab-\delta_ba+[a,b]
 &=\iint\bigl([e_x,e_{x+y}]+[e_{x+y},e_y]
                         -[e_x,e_y]\bigr)\,da(x)\,db(y)\\
 &=\mathcal J(a\wedge b).
\end{align*}
This proves the integral group comparison
\begin{equation}\label{eq:cm-integral-J-comparison}
 k_{[s(h),s(k)]}(q)\bmod C_3
      =(T_v-1)J_A(h^{\rm ab}\wedge k^{\rm ab}).
\end{equation}

For completeness, let \(\mathcal D_r\) be the special
automorphisms acting trivially on \(V\) whose
displacements all lie in \(C_r\).
An element of \(\mathcal D_1\) acts trivially on
\(C_2/C_3\), whereas an element of \(\mathcal D_2\)
acts trivially on \(C/C_3\).  Comparing their product
displacements and using \([C,C_2]\subset C_3\) proves
\([\mathcal D_1,\mathcal D_2]\subset\mathcal D_3\).
Consequently \(s(\gamma_3H)\subset\mathcal D_3\),
and additivity of displacement on \(\mathcal D_2\)
extends \eqref{eq:cm-integral-J-comparison} continuously
to every \(h\in\gamma_2H\):
\[
 k_{s(h)}(q)\bmod C_3=(T_v-1)J_A(h\bmod\gamma_3H).
\]
For \(q=x_1\), the operator \(T_v-1\) is multiplication
by
\[
 (1+T_1)(1+S_1)-1=T_1+S_1+T_1S_1
\]
in the power-series domain
\(\mathbb Z_p[[T_1,T_2,S_1,S_2]]\), so it is injective.
An outer-kernel element belongs to \(\gamma_2H\) by
abelian detection and has \(s(h)=1\), since the
canonical special lift is inverse to the outer
projection in \(F^3\).  The last identity therefore
places it in \(N_2^{\rm CM}\).

Here is also the exact arithmetic compatibility.
An arithmetic outer automorphism has a representative
preserving \(\langle z\rangle\); any two such
representatives differ by a power of
\(\operatorname{inn}(z)\), which commutes with every
\(s(h)\).  Uniqueness in \(F^3\) makes the special lift
equivariant.  On \(M\), conjugation is the linear
Tate action with determinant twist.  Both shears
commute with simultaneous linear changes on \(V\),
so \(J_A\) is equivariant with the displayed second
twist.  Thus its inverse-image kernel is stable
under the arithmetic action.  It is closed, and its
image is compact.  The resulting exact central extension
\[
 1\longrightarrow\operatorname{im}J_A
 \longrightarrow Q_2^{\rm CM}\longrightarrow A
 \longrightarrow1
\]
has commutator subgroup exactly \(\operatorname{im}J_A\).
Both end modules are \(p\)-torsion-free, so the group
is torsion-free and has the claimed abelianization.

It remains to prove that the commutator module is nonzero.
Let \(I_K\) denote the admissible weights, and put
\[
 \mathcal S=\{\boldsymbol m\in I_K:
  1\leq m_i\leq4,\ (5,5)-\boldsymbol m\in I_K\}.
\]
These are exactly the moments that can contribute
to the bidegree-\((5,5)\) bracket coefficient.
The ray-class description and two CM Tate characters
in \cite[pp.~4--5]{Ishii} show that
\(\Gamma=\operatorname{Gal}(L_\infty/L_1)\) has open
rank-two principal-unit image. Write
\(\chi_{\boldsymbol m}=\chi_1^{m_1}\chi_2^{m_2}\)
for the corresponding product of the two Tate characters. Choose \(\gamma\)
so that \(\chi_{\boldsymbol m}(\gamma)\), for
\(\boldsymbol m\in\mathcal S\), are pairwise distinct.
This is possible because after the \(p\)-adic logarithm
the finitely many equality loci are proper linear
subspaces of an open \(\mathbb Z_p^2\).
For \(\boldsymbol m_0=(2,2),(3,3)\), define the integral
finite-difference operator
\[
 P_{\boldsymbol m_0}(X)=
 \prod_{\substack{\boldsymbol m\in\mathcal S\\
                   \boldsymbol m\ne\boldsymbol m_0}}
       (X-\chi_{\boldsymbol m}(\gamma)).
\]
Both \(\kappa_{22}\) and \(\kappa_{33}\) are nonzero
by \cite[Theorem 2.8(2) and the following Soul\'e-finiteness
paragraph, p.~9]{Ishii}.
Choose \(a_0,b_0\in A\) on which these respective
characters are nonzero, and set
\(a=P_{22}(\gamma)a_0,b=P_{33}(\gamma)b_0\).
Equivariance gives
\[
 \kappa_{\boldsymbol m}(P(\gamma)d)
   =P(\chi_{\boldsymbol m}(\gamma))\kappa_{\boldsymbol m}(d).
\]
Thus all unwanted moments in \(\mathcal S\) vanish,
while the two desired moments remain nonzero.

In exponential moment variables \(U=(x,y),V=(z,t)\),
\(\mathcal JF=F(U+V,V)+F(U,U+V)-F(U,V)\).
Direct multiplication gives
\[
 J(XY,X^2Y^2)
       =(xy-zt)(xt+zt+yz)(xt+xy+yz)\ne0.
\]
The bidegree-\((3,3)\) polynomial component of
\(J_A(a\wedge b)\), corresponding to bracket weight
\((5,5)\), is this polynomial multiplied by
\[
 \frac{\kappa_{22}(a)\kappa_{33}(b)}
      {(1-p^2)(1-p^4)(2!)^2}\ne0.
\]
The definition of \(\mathcal S\) excludes every other
contribution.  Hence \(\operatorname{im}J_A\ne0\),
proving exact class two for every \(K\) in the statement.
\end{proof}

\subsection{Centralizers and free class-two source subgroups}
\label{subsec:cm-class-two-subgroups}

Retain the hypotheses and notation of
Theorem~\ref{thm:cm-measure-class-two}, and write
\[
 Q=Q_2^{\rm CM},\qquad Z=\operatorname{im}J_A,
 \qquad \pi:Q\longrightarrow A.
\]
The following statements use the proved integral map from the
actual source group to this quotient.

\begin{proposition}\label{prop:cm-class-two-centralizers}
The centre of \(Q\) is exactly \(Z\).  For every \(q\in Q\)
with \(a=\pi(q)\ne0\),
\[
 C_Q(q)=\pi^{-1}\bigl(A\cap\mathbb Q_p a\bigr).
\]
The intersection is taken inside \(A[1/p]\), is a free
\(\mathbb Z_p\)-module of rank one, and the displayed
centralizer is, noncanonically, \(Z\times\mathbb Z_p\).
In particular, every abelian subgroup of \(Q\) has image
of \(\mathbb Z_p\)-rank at most one in \(A\).
\end{proposition}

\begin{proof}
The commutator of two elements with abelian images \(a,b\)
is \(J_A(a\wedge b)\).  Proposition~\ref{prop:cm-decomposable-commutators}, together
with injectivity of \(\alpha[1/p]\), gives the exact criterion
\[
 J_A(a\wedge b)=0\quad\Longleftrightarrow\quad
 b\in\mathbb Q_pa\qquad(a\ne0).
\]
This proves the asserted centralizer formula.

Here is the integral topology of the line intersection.
Choose a continuous \(\mathbb Z_p\)-linear functional
\(\ell:A\to\mathbb Z_p\) with \(\ell(a)\ne0\).
Such a functional exists because the compact torsion-free
abelian pro-\(p\) group \(A\) is a product of copies of
\(\mathbb Z_p\).  If \(\lambda a\in A\), then
\(\lambda\ell(a)\in\mathbb Z_p\), which bounds the
denominator of \(\lambda\).  The allowed scalars form a
nonzero closed \(\mathbb Z_p\)-submodule of this bounded
rank-one lattice.  They are consequently \(p^e\mathbb Z_p\)
for some integer \(e\).  Thus \(A\cap\mathbb Q_pa\)
has a generator \(a_0\).  Choose \(q_0\in Q\) above
\(a_0\).  The continuous powers of \(q_0\) identify
its closed procyclic group with \(\mathbb Z_p\), and
\((z,t)\mapsto zq_0^t\) is an isomorphism from
\(Z\times\mathbb Z_p\) onto the displayed centralizer.

The module \(A[1/p]\) has dimension at least two: for
example, the moment-isolated elements constructed below
have distinct nonzero initial homogeneous terms.
Every nonzero \(a\) therefore has a \(\mathbb Q_p\)-independent
\(b\), so an element above \(a\) is not central.
Hence the centre is precisely \(Z\).  The assertion about
abelian subgroups follows from the same centralizer formula.
\end{proof}

\begin{corollary}\label{cor:cm-three-generator-class-two}
Let \(r\le3\), and let \(h_1,\ldots,h_r\in H\) have
\(\mathbb Q_p\)-linearly independent images in \(A[1/p]\).
For the closed subgroup \(U\) they generate, the map
\(U\to Q\) has kernel exactly \(\gamma_3U\), and
\(U/\gamma_3U\) is the free class-two pro-\(p\) group
on these \(r\) generators.
\end{corollary}

\begin{proof}
Let \(B\) be the \(\mathbb Z_p\)-span of their abelian
images.  Every bivector in the exterior square of a
vector space of dimension at most three over
\(\mathbb Q_p\) is decomposable.  The preceding
derivative argument thus makes
\(J_A:\bigwedge^2B\to Z\) injective.
Map the free pro-\(p\) group \(P_r\) to \(H\) using
the chosen elements.  The composite to \(Q\) kills
\(\gamma_3P_r\); its maps on the abelian quotient and
on \(\gamma_2P_r/\gamma_3P_r\simeq\bigwedge^2\mathbb Z_p^r\)
are injective.  Its kernel is therefore exactly
\(\gamma_3P_r\).  The image of that subgroup in \(U\)
is \(\gamma_3U\), which proves both conclusions.
\end{proof}

\begin{proposition}\label{prop:cm-countable-class-two}
There is a closed subgroup \(U\subset H\), generated
by a sequence tending to one, for which
\[
 U\cap N_2^{\rm CM}=\gamma_3U,
 \qquad U/\gamma_3U\simeq P/\gamma_3P,
\]
where \(P\) is the free pro-\(p\) group on a countably
infinite convergent sequence of generators.
Thus the complete free class-two quotient of this
actual source subgroup is detected.  Every finite
initial subsequence gives the corresponding conclusion
in finite rank.
\end{proposition}

\begin{proof}
For \(i\ge1\), put \(k_i=2^i\) and
\(d_i=2k_i-2\).  The character
\(\kappa_{(k_i,k_i)}\) is nonzero by
\cite[Theorem 2.8(2) and the Soul\'e-finiteness paragraph,
p.~9]{Ishii}; this is the ordinary \(k_i\)-th
Tate twist, with \(k_i\ge2\).

Let
\[
 \mathcal S_i=\{\boldsymbol m\in I:
                              |\boldsymbol m|\le2k_i\}.
\]
Put \(\Gamma_2=\operatorname{Gal}(L_\infty/L_1)\), and choose \(\gamma_i\in\Gamma_2\) for which the finitely
many characters \(\chi_{\boldsymbol m}(\gamma_i)\),
\(\boldsymbol m\in\mathcal S_i\), are pairwise distinct.
The open rank-two principal-unit image of the two CM
Tate characters, established by the ray-class morphism
in \cite[pp.~4--5]{Ishii}, permits this choice:
the forbidden loci after taking logarithms are a finite
union of proper linear subspaces.  Put
\[
 P_i(X)=\prod_{\substack{\boldsymbol m\in\mathcal S_i\\
                         \boldsymbol m\ne(k_i,k_i)}}
                 (X-\chi_{\boldsymbol m}(\gamma_i)).
\]
Choose \(b_i^0\in A\) with
\(\kappa_{(k_i,k_i)}(b_i^0)\ne0\), and set
\(b_i=P_i(\gamma_i)b_i^0\).  Equivariance shows that
all its moments in \(\mathcal S_i\), except the desired
one, vanish, and that the desired one remains nonzero.

Choose a lift \(\widetilde h_i\in H\) of \(b_i\), and
put \(h_i=\widetilde h_i^{p^i}\), with abelian image
\(a_i=p^ib_i\).  This sequence tends to one: in every
finite quotient of \(H\), a sufficiently large power
\(p^i\) kills every element.  The exponential transform
\(f_i\) of \(\alpha(a_i)\) has the form
\[
 f_i(U_1,U_2)=c_i(U_1U_2)^{k_i-1}
                    +\text{terms of total degree greater than }d_i,
 \qquad c_i\in\mathbb Q_p^\times.
\]
In particular, the \(d_i\) are strictly increasing.

Let \(P\) be free on the convergent sequence
\(x_i\to1\).  Its abelianization is
\(\prod_{i\ge1}\mathbb Z_p x_i\), and the integral
class-two Hall calculation gives
\[
 \gamma_2P/\gamma_3P
       =\prod_{i<j}\mathbb Z_p[x_i,x_j].
\]
These assertions follow by taking inverse limits of
the corresponding finite-generator, finite-quotient
Hall presentations.  The continuous map \(x_i\mapsto h_i\)
induces an injective abelian map: in a nonzero completed
linear combination of the \(f_i\), the least index
with nonzero coefficient supplies a unique nonzero
initial degree.

The initial homogeneous term of \(J(f_i,f_j)\) is
\[
 c_ic_j J\bigl((U_1U_2)^{k_i-1},
              (U_1U_2)^{k_j-1}\bigr),
\]
of degree
\[
 d_i+d_j=2^{i+1}+2^{j+1}-4.
\]
It is nonzero by the derivative computation above:
the two displayed monomials have nonconstant ratio.
These degrees are pairwise distinct for \(i<j\), by
uniqueness of binary expansion.  A nonzero completed
alternating combination therefore again has a unique
least initial degree, which cannot cancel.
For each total degree there are only finitely many
pairs to consider, so this argument applies to the
completed product, not only to finite combinations.
Consequently the composite
\(P\to H\to Q\) induces an injection on both its
abelian quotient and its class-two central quotient.
Its kernel is exactly \(\gamma_3P\).

Let \(U\) be the image of \(P\) in \(H\).  It is compact
and hence closed.  The image of \(\gamma_3P\) is
\(\gamma_3U\), by the commutator definition and
compactness.  The exact kernel just calculated gives
\(U\cap N_2^{\rm CM}=\gamma_3U\) and the asserted
isomorphism.  No injectivity assertion for the full
map \(P\to H\) is needed for these conclusions.
\end{proof}

The subgroups in Proposition~\ref{prop:cm-countable-class-two}
depend on the moment-isolating choices.  The surrounding
quotient \(Q_2^{\rm CM}\) remains the canonical arithmetic
quotient of Theorem~\ref{thm:cm-measure-class-two}.


\subsection{A mixed-unit measure and the exact lifting obstruction}
\label{subsec:cm-nonabelian-gate}

Keep the notation of Subsection~\ref{subsec:cm-elliptic-radical}, and put
\[
 N_{\mathrm{ell}}=\ker\!\left(A\longrightarrow
                  \operatorname{Gal}(L_{\mathrm{ell}}/L_\infty)\right),
 \qquad \overline A=A/N_{\mathrm{ell}},\qquad
 \Lambda_2=\mathbb Z_p[[\Gamma]].
\]
The following construction uses only the mixed torsion points whose
two coordinates are units.  The passage from its moments to the
full radical is justified by an explicit scalar formula, including
the two boundary families of weights.

\begin{lemma}\label{lem:cm-mixed-measure-embedding}
Fix an integer \(c_0>1\) prime to \(6p\), put
\(\mathfrak a=(c_0)\), and set
\[
 u_n=\theta_{\mathfrak a}(\omega_{1,n}+\omega_{2,n}).
\]
There is a continuous injective \(G\)-equivariant homomorphism
\[
 \overline\Phi_{c_0}:\overline A
            \lhook\joinrel\longrightarrow \mathbb Z_p[[G]](1).
\]
Here the target action is
\(\gamma\cdot\mu=\chi_{\mathrm{cyc}}(\gamma)[\gamma]\mu\);
translation is defined by
\([\gamma]\mu(B)=\mu(\gamma^{-1}B)\).
In particular \(\overline A\) is torsion-free over \(\Lambda_2\).
\end{lemma}

\begin{proof}
Let \(G_n=\operatorname{Gal}(L_n/K)\).  The exact norm relation
of \cite[Proposition 2.5, p.~8]{IshiiSoule} gives
\(N_{L_{n+1}/L_n}(u_{n+1})=u_n\).  Its root-of-unity factors are
one because no nonidentity unit of \(O_K\) is congruent to one
modulo \(p\), for \(p\geq5\).
For \(h\in H\), define a bounded measure by the finite-level
Kummer values
\begin{equation}\label{eq:cm-mixed-measure}
 \Phi_{c_0}(h)(\sigma\operatorname{Gal}(L_\infty/L_n))
                      =k_{\sigma(u_n)}(h)\qquad(\sigma\in G_n).
\end{equation}
The norm identity says exactly that the value of a cylinder is
the sum of the values of its refinements.  Thus these values give
an element of \(\mathbb Z_p[[G]]\).  Additivity of Kummer characters
makes \(\Phi_{c_0}\) a continuous homomorphism on \(A\).
For a lift of \(\gamma\in G\), the defining root equation gives
\[
 k_u(\gamma h\gamma^{-1})
   =\chi_{\mathrm{cyc}}(\gamma)k_{\gamma^{-1}u}(h).
\]
This proves the stated action and its independence of the lift,
since inner conjugation acts trivially on \(A\).

We identify its kernel using the moment formula for all weights.  By the proof of
Corollary~\ref{cor:cm-radical-joint-kernel}, a positive integral
power of each \(u_n\) belongs to \(C(L_n)\).  Consequently
\(N_{\mathrm{ell}}\subseteq\ker\Phi_{c_0}\).
Conversely, suppose \(\Phi_{c_0}(h)=0\).
For every \(\boldsymbol m=(m_1,m_2)\in I\), the exact moment
formula of \cite[Lemma 4.2, p.~17]{IshiiSoule} gives
\[
 \kappa_{\boldsymbol m,\mathfrak a}(h)
   =\int_G\chi^{\boldsymbol m-\boldsymbol1}(\sigma)\,
                                 d\Phi_{c_0}(h)(\sigma)=0.
\]
In this source \(\kappa_{\boldsymbol m,\mathfrak a}\) is defined
from this same mixed norm system for every \(\boldsymbol m\in I\).
When one of \(m_1,m_2\) equals one, the proof of
\cite[Proposition 4.3, pp.~17--18]{IshiiSoule} takes the norm
to the corresponding one-prime ray tower and includes its Euler
factor.  Its conclusion still reads
\begin{equation}\label{eq:cm-mixed-all-weight-scalar}
 w_K\,\frac{1-p^{r+s}}{(1-A_{\boldsymbol m})(1-B_{\boldsymbol m})}
       \kappa_{\boldsymbol m,\mathfrak a}
    =(c_0^2-c_0^{-r-s})\kappa_{\boldsymbol m},
 \quad r=m_1-1,\quad s=m_2-1,
\end{equation}
where
\[
 A_{\boldsymbol m}=i^{(r,s)}(\pi),\qquad
 B_{\boldsymbol m}=i^{(r,s)}(\bar\pi),\qquad
 A_{\boldsymbol m}B_{\boldsymbol m}=p^{r+s}.
\]
These algebraic numbers have complex absolute value
\(p^{(r+s)/2}>1\), because \(r,s\geq0\) and \(r+s>0\).
Both denominators and the numerator on the left are therefore
nonzero, including \(r=0\) and \(s=0\).
The right scalar is also nonzero.  The unit ambiguity in the
Artin action disappears since \(r-s\equiv0\pmod{w_K}\).
Thus \(\kappa_{\boldsymbol m}(h)=0\) for every
\(\boldsymbol m\in I\).
Corollary~\ref{cor:cm-radical-joint-kernel} now gives
\(h\bmod[H,H]\in N_{\mathrm{ell}}\).
This proves the exact equality
\(\ker\Phi_{c_0}=N_{\mathrm{ell}}\).

Finally \(G/\Gamma\) is finite of order prime to \(p\), so the
untwisted \(\mathbb Z_p[[G]]\) is finite free over \(\Lambda_2\).
The cyclotomic twist changes scalars by the continuous ring
automorphism \(\gamma\mapsto\chi_{\mathrm{cyc}}(\gamma)\gamma\)
of \(\Lambda_2\).  Its target remains torsion-free, and the
displayed injection proves the assertion.
\end{proof}

\begin{lemma}\label{lem:cm-exact-eigenlift-obstruction}
Let \(\gamma\neq1\) belong to \(\Gamma\), let
\(\lambda\in\mathbb Z_p^\times\), and let \(t\) be a lift of
\(\gamma\) to \(\operatorname{Gal}(\Omega/K)\).
If an element \(h\in H\) satisfies
\[
                         tht^{-1}=h^\lambda ,
\]
then every elliptic Soul\'e character vanishes on \(h\).
If \(\mathcal X_2^+\) is finite, then \(h\in[H,H]\).
\end{lemma}

\begin{proof}
Write \(a\) for the image of \(h\) in \(\overline A\).
The asserted group equation gives the exact module equation
\[
                              (\gamma-\lambda)a=0.
\]
The element \(\gamma-\lambda\) is nonzero in
\(\Lambda_2=\mathbb Z_p[[\mathbb Z_p^2]]\), an integral domain.
Lemma~\ref{lem:cm-mixed-measure-embedding} therefore gives \(a=0\).
The joint-kernel description implies the first assertion.
Under finite \(\mathcal X_2^+\),
Theorem~\ref{thm:cm-finite-plus-abelian} gives
\(N_{\mathrm{ell}}=0\), hence \(h\) has zero image in \(A\).
\end{proof}

This gives an obstruction to lifting homogeneous
Galois eigenclasses by group elements satisfying the corresponding
eigen-equations.  Rational eigenspace decompositions in a graded
quotient do not provide such equations in \(H\).

\begin{proposition}\label{prop:cm-no-normalized-cyclic}
Assume that \(\mathcal X_2^+\) is finite.  Put
\(G_1=\operatorname{Gal}(\Omega/L_1)\), so that
\[
             1\longrightarrow H\longrightarrow G_1
               \longrightarrow\Gamma\longrightarrow1 .
\]
If \(t\in G_1\) has nonidentity image in \(\Gamma\), then the
equation
\[
                 tht^{-1}=h^\lambda,\qquad
                    h\in H,\quad\lambda\in\mathbb Z_p^\times
\]
has only the solution \(h=1\).
Consequently, for every nontrivial closed procyclic subgroup
\(C\subset H\), the image of \(N_{G_1}(C)\) in \(\Gamma\) is
trivial.
\end{proposition}

\begin{proof}
The finite-plus theorem gives that \(H\) is free pro-\(p\) and
that \(A=H^{\mathrm{ab}}\) embeds equivariantly in
\[
 P=\mathbb Z_p[[G]](1).
\]
As a \(\Lambda_2\)-module, \(P\) is a finite direct sum of
copies of the regular module, with its cyclotomic twist.

We first justify taking completed tensor powers of this
injection.  Every compact torsion-free \(\mathbb Z_p\)-module
is topologically a product of copies of \(\mathbb Z_p\).
Indeed, its Pontryagin dual is a divisible \(p\)-primary group,
and hence a direct sum of copies of
\(\mathbb Q_p/\mathbb Z_p\).  For such a product \(M\), and any
compact \(\mathbb Z_p\)-module \(V\), the definition by finite
continuous quotients gives
\[
 \left(\prod_{i\in I}\mathbb Z_p\right)
       \widehat\otimes_{\mathbb Z_p}V
             \simeq\prod_{i\in I}V.
\]
This functor preserves injections.  Applying this assertion
one tensor factor at a time gives, for every \(d\geq1\), an
injection
\[
 A^{\widehat\otimes d}\lhook\joinrel\longrightarrow
 P^{\widehat\otimes d}
     \simeq\mathbb Z_p[[G^d]](d).
\]

The action of \(\Gamma\) on the last module is diagonal
translation with the scalar twist
\(\chi_{\mathrm{cyc}}^d\).  The continuous group isomorphism
\[
 \Gamma\times\Gamma^{d-1}\longrightarrow\Gamma^d,\qquad
 (g,h_2,\ldots,h_d)\longmapsto
                (g,gh_2,\ldots,gh_d)
\]
identifies diagonal translation with translation in the first
factor.  After choosing the finitely many tame cosets, its
completed group algebra is therefore a finite direct sum of
power-series modules
\[
             \Lambda_2[[T_1,\ldots,T_{2d-2}]].
\]
It is torsion-free over \(\Lambda_2\).  The scalar twist changes
the coefficient action by the ring automorphism
\(\gamma\mapsto\chi_{\mathrm{cyc}}(\gamma)^d\gamma\), which
preserves that property.  Thus every completed tensor power
of \(A\) is \(\Lambda_2\)-torsion-free.

Write \(\gamma_dH\) for the closed lower central series.
For a free pro-\(p\) group, the canonical map
\[
 \gamma_dH/\gamma_{d+1}H
        \longrightarrow A^{\widehat\otimes d},
 \qquad
 [h_1,\ldots,h_d]\longmapsto
          [\bar h_1,\ldots,\bar h_d]_{\mathrm{tensor}},
\]
identifies its source with the completed degree-\(d\) free
Lie submodule.  This assertion is integral: for a finite free
basis, the Hall commutators give the free
\(\mathbb Z_p\)-basis, and their Magnus expansions have the
corresponding leading Lie tensors.  For a general free
pro-\(p\) basis converging to one, projection onto finite
subsets and reduction modulo \(p^n\) give the completed
assertion.  These projections also prove injectivity, since a
nonzero coefficient is seen in a finite set of basis
coordinates and modulo some \(p^n\).  The construction is
functorial in automorphisms.  In particular, the conjugation
action on this graded piece is the action induced by that on
\(A\), and so it factors through \(\Gamma\).
It follows that every
\(\gamma_dH/\gamma_{d+1}H\) is
\(\Lambda_2\)-torsion-free.

Every pro-\(p\) group is separated for its closed lower central
series: any nonidentity element survives in a finite \(p\)-group,
which is nilpotent.  Therefore, if \(h\neq1\), it has a nonzero
image \(a\) in some
\(\gamma_dH/\gamma_{d+1}H\).  If \(t\) has image
\(\gamma\neq1\), the proposed equation gives
\[
                         (\gamma-\lambda)a=0 .
\]
The power series \(\gamma-\lambda\) is nonzero in
\(\Lambda_2\), contradicting the torsion-freeness just proved.
Thus \(h=1\).

Finally \(H\) is torsion-free, so a nontrivial closed
procyclic \(C\) is isomorphic to \(\mathbb Z_p\).
Conjugation by a normalizer element sends a topological
generator \(h\) to \(h^\lambda\) for a unit \(\lambda\).
The first assertion forces the image of that element in
\(\Gamma\) to be one.  This proves the precise normalizer
claim.
\end{proof}

The proposition gives a further constraint on a proposed
arithmetic use of the geometric recognition theorem.  A
nontrivial procyclic subgroup selected inside \(H\), including
one selected inside the arithmetic outer kernel, cannot
acquire an open normalizer image in
\(\operatorname{Gal}(L_\infty/L_1)\): its normalizer image
there is exactly the identity subgroup.  This conclusion
concerns the displayed arithmetic extension and follows from
its actual conjugation action on every lower-central quotient.

\paragraph{The remaining correction equation.}
Let \(\widetilde F^\bullet H\) be the weighted free filtration
attached to a proposed free set of arithmetic lifts, with its
filtered outer map to \(F^\bullet\rho(H)\).
A homogeneous derivation relation, such as the relation in
bidegree \((12,2)\) proved above, produces a word
\(r\in\widetilde F^{14}H\) satisfying
\[
 \rho(r)\in F^{15}\rho(H).
\]
To cancel that outer image without changing its initial
free class, one must solve the stronger, explicit equation
\[
 \rho(r)=\rho(s)\quad\hbox{for some }s\in\widetilde F^{15}H.
\]
Its solution gives \(rs^{-1}\in\ker\rho\).
The first inclusion alone does not supply \(s\):
it concerns the target filtration, whereas the displayed equation
asks for membership in the image of a specified source subgroup.
Rational surjectivity of successive graded maps gives no
uniform integral bound on the denominators in successive
corrections.  Clearing a growing denominator by successive
\(p\)-powers can erase the proposed initial class.
The argument in \cite[proof of Theorem 2.14, p.~27]{Ishii} uses graded
freeness precisely to identify these two filtrations.
The source-group freeness established here does not itself
solve the displayed correction equation.

\subsection{The exact mixed-chart image of the geometric measure}
\label{subsec:cm-measure-image}

Write \(w_K=|O_K^\times|\), put \(U=(\mathbb Z_p^\times)^2\subset V\), and let
\(j:U\hookrightarrow V\) be its inclusion.
The ray-class identification is
\(G=U/O_K^\times\).  For the auxiliary ideal
\(\mathfrak a=(c)\), let \(\Phi_c\) denote the map of Lemma~\ref{lem:cm-mixed-measure-embedding} composed with \(A\to\overline A\), and let \(\mu_c\) be its normalized
\(O_K^\times\)-invariant lift of the mixed-unit
measure \(\Phi_c\) to \(U\): its mass on each finite
unit orbit is the quotient mass divided by \(w_K\).
This is integral since \(p\nmid w_K\).
Let \(P,Q\) act on \(V\) by multiplication by
\(\pi,\bar\pi\), and put
\[
 \mathcal L=(1-P_*)(1-Q_*),\qquad
 D_c=c^2-[c^{-1}]_*.
\]
Here all starred maps denote pushforward of bounded
measures, and \(PQ=p\,\operatorname{id}_V\).

\begin{lemma}\label{lem:cm-measure-euler-image}
The actual geometric measure satisfies
\[
 w_Kj_*\mu_c=D_c\mathcal L\alpha,\qquad
 w_Kj_*\mu_c=D_cj_*(\alpha|_U),\qquad
 \mathcal L\alpha=j_*(\alpha|_U).
\]
Its restriction \(\alpha|_U\) has kernel
\(N_{\mathrm{ell}}\), and uniquely determines
\(\alpha\) with its zero-total-mass normalization.
Under finite \(\mathcal X_2^+\), this restriction
embeds \(A\) as a full-rank closed submodule of
\(\mathbb Z_p[[U]]^{O_K^\times}(1)\).
After identifying this target with
\(\mathbb Z_p[[G]](1)\) and absorbing the twist,
the image is an ideal containing a nonzerodivisor.
\end{lemma}

\begin{proof}
The \((r,s)\)-moment of \(\mu_c\) is
\(\kappa_{(r+1,s+1),\mathfrak a}\).  Rearranging
\cite[Proposition 4.3, pp.~17--18]{IshiiSoule}
and using \cite[Theorem 2.6, p.~8]{Ishii}
gives exactly the first identity on every admissible
positive-degree monomial.  Inadmissible moments
vanish by \(O_K^\times\)-invariance.  The constant
moment of \(\mu_c\) is zero: it is the Kummer value
of \(N_{L_n/K}(u_n)\), a global unit of \(K\), hence
a root of unity with all \(p\)-power roots in
\(L_\infty\).  The constant moment of \(\alpha\)
is zero by its source formula, and the odd moments
vanish by the action of \(-1\).
Polynomial moments determine bounded measures,
so the first identity holds integrally.

Neither \(P(V)\) nor \(Q(V)\) meets \(U\).
Restriction to \(U\), which commutes with \(D_c\),
therefore gives the second identity.
The \((r,s)\)-moment multiplier of \(D_c\) is
\(c^2-c^{-r-s}\ne0\), including degree zero.
Thus \(D_c\) is injective, and cancellation gives
the third identity.

The kernel of \(\mathcal L\) on all bounded
measures is precisely \(\mathbb Z_p\delta_0\).
Indeed its positive-degree moment multipliers
are
\[
 (1-i^{(r,s)}(\pi))(1-i^{(r,s)}(\bar\pi)),
\]
and both algebraic numbers occurring here have
complex absolute value \(p^{(r+s)/2}>1\).
A measure killed by \(\mathcal L\) consequently
has all positive-degree moments zero and is its
total mass times \(\delta_0\).
This proves uniqueness with zero total mass.
The second identity and injectivity of \(D_c\)
identify \(\ker(\alpha|_U)\) with
\(\ker\Phi_c=N_{\mathrm{ell}}\).

Under finite \(\mathcal X_2^+\), the latter kernel
vanishes.  The exact rank formula
\cite[Theorem 4.4, author version p.~27]{VenjakobIwasawa},
together with weak Leopoldt, gives
\[
 \operatorname{rank}_{\Lambda_2}A
     =r_2(L_1)=[L_1:K]=\frac{(p-1)^2}{w_K}.
\]
The maximal pro-\(p\) group comparison used here
is the \((p,S)\)-closed comparison at the start of
\cite[Section 4, author version p.~26]{VenjakobIwasawa}.
The unit-chart target has this same rank.
Equivariance makes the image an ideal after the
cyclotomic twist is absorbed.  The tame
decomposition of \(\mathbb Z_p[[G]]\) is a product
of copies of \(\Lambda_2\), and full rank makes
each ideal component nonzero.  A choice of a
nonzero element in every component gives the
claimed nonzerodivisor.
\end{proof}

These identities retain the radial and axis constraints
on the actual image.  In particular, for \(i,j\geq0\),
\[
 \alpha|_{P^iQ^jU}=(P^iQ^j)_*(\alpha|_U).
\]
This follows by restricting the Euler identity
successively to the disjoint interior shells.
No infinite geometric-series expansion at the axes
is required.  In the class-two question, the remaining
map is still
\[
 \mathcal J\circ(\alpha\widehat\otimes\alpha)
              \bigm|_{\widehat{\bigwedge}^{\,2}A}.
\]
The separate-factor operators \(\mathcal L\otimes1\)
and \(1\otimes\mathcal L\) do not commute with its
addition shears.  The displayed ideal description
therefore retains, rather than removes, this exact
tensor-image constraint.

\subsection{The finite elliptic Euler complexes and the actual kernel}
\label{subsec:cm-elliptic-euler-kernel}

There is a finite-level description of the precise arithmetic
image entering \(J_A\).  Retain the finite-\(\mathcal X_2^+\)
hypothesis and write
\[
 V_n=E[p^n],\qquad
 \alpha_n:A\longrightarrow\mathbb Z_p[V_n](1),\qquad
 A_n=\operatorname{im}\alpha_n.
\]
The transition on \(V_n\) is multiplication by \(p\).
Each \(A_n\) is its actual compact, finite free
\(\mathbb Z_p\)-lattice.

Use Goncharov's elliptic units \(\theta_E(x)\), and let
\(C_1(n)\) be their \(\mathbb Q_p\)-span for
\(0\ne x\in V_n\), inside \(L_n^\times\otimes\mathbb Q_p\).
Here \(B_2(F)\) denotes Goncharov's rational Bloch group.
It is generated by symbols \(\{z\}\), \(z\in\mathbb P^1(F)\),
with differential \(\delta\{z\}=(1-z)\wedge z\), set equal to
zero at \(0,1,\infty\). Its relations identify the two
specializations of \(\sum_i\{f_i\}\) along a connected curve
whenever \(\sum_i(1-f_i)\wedge f_i=0\) in its function field;
see \cite[Section 2.1, p.~10]{GoncharovEuler}.
For number fields this agrees rationally with the five-term
presentation through the canonical map stated there.
Let \(C_2(n)\) be the \(\mathbb Q_p\)-span of his
double elliptic units \(\theta_E(a,b,c)\), with
\(a,b,c\in V_n\) and \(a+b+c=0\), in
\(B_2(L_n)\otimes\mathbb Q_p\).
The ray field and these theta-value spans are described in
\cite[Definition 5.2, Proposition 5.3 and Definition 5.4,
p.~27]{GoncharovEuler}.
Its differential is
\begin{equation}\label{eq:cm-elliptic-triangle-delta}
 \delta\theta_E(a,b,c)
   =\theta_E(a)\wedge\theta_E(b)
    +\theta_E(b)\wedge\theta_E(c)
    +\theta_E(c)\wedge\theta_E(a),
\end{equation}
by \cite[Theorem 2.10(a), equation (38), p.~15]{GoncharovEuler}.
The differential of a triple with a zero entry vanishes;
no value of \(\theta_E(0)\) is needed in its image.
Put
\[
 \mathcal R_n^{\rm ell}
   =\operatorname{coker}\bigl(C_2(n)
                          \xrightarrow{\delta}\bigwedge^2C_1(n)\bigr),
 \qquad
 \mathcal D_n=\ker\bigl(J_n:\bigwedge^2A_n
                   \longrightarrow\mathbb Z_p[V_n\times V_n](2)\bigr),
\]
where \(J_n\) is the same pair of finite additive shears
as in Theorem~\ref{thm:cm-measure-class-two}.

\begin{theorem}\label{thm:cm-elliptic-euler-kernel}
There are natural equivariant identifications
\[
 A_n[1/p]=\operatorname{Hom}_{\mathbb Q_p}
                         (C_1(n),\mathbb Q_p(1)),
 \qquad
 \mathcal D_n[1/p]=\operatorname{Hom}_{\mathbb Q_p}
                         (\mathcal R_n^{\rm ell},\mathbb Q_p(2)).
\]
The transition \(\mathcal D_{n+1}\to\mathcal D_n\)
is induced by the actual pushforward
\(A_{n+1}\to A_n\).  Rationally, it is dual to the
natural map \(\mathcal R_n^{\rm ell}\to
\mathcal R_{n+1}^{\rm ell}\) induced by field extension
and the elliptic-unit distributions.  Finally,
\begin{equation}\label{eq:cm-J-kernel-elliptic-limit}
                       \ker J_A=\varprojlim_n\mathcal D_n.
\end{equation}
Here the inverse limit retains the specified integral
lattices \(\mathcal D_n\).
\end{theorem}

\begin{proof}
We first check the normalization linking the two primary
sources.  Scale the period lattice to
\(L=\mathbb Z\tau+\mathbb Z\), and write
\(z=a\tau+b\), \(q=e^{2\pi i\tau}\),
\(\xi=e^{2\pi iz}\).  Put
\(\eta_1=\eta(\tau)\), \(\eta_2=\eta(1)\).
The sigma product and the Legendre relation give
\[
 \sigma(z,L)=\frac{e^{\eta_2z^2/2}}{2\pi i}
       \xi^{-1/2}(\xi-1)
       \prod_{m\ge1}
       \frac{(1-q^m\xi)(1-q^m\xi^{-1})}{(1-q^m)^2},
 \qquad \eta_2\tau-\eta_1=2\pi i.
\]
The product identity follows from the zeros and the two
quasiperiods of \(\sigma\), with \(\sigma'(0)=1\)
fixing its constant.  Substitution in Nakamura's
\(\Theta(z,L)=\Delta(L)e^{-6\eta(z)z}\sigma(z,L)^{12}\),
using \(\eta(z)=a\eta_1+b\eta_2\), yields
\[
 \Theta(z,L)=q^{6a^2-6a+1}e^{12\pi ib(a-1)}
    (1-\xi)^{12}
    \prod_{m\ge1}(1-q^m\xi)^{12}(1-q^m\xi^{-1})^{12}.
\]
This is precisely the twelfth power of
\cite[equation (19), p.~7]{GoncharovEuler}.
Homogeneity removes the scaling of the lattice.
The root-of-unity ambiguity from torsion representatives
has zero Kummer character on \(H\).  Thus
\(\kappa_{\Theta(x)}=12\kappa_{\theta_E(x)}\);
the factor \(12\) is invertible because \(p\ge5\).

Nakamura's formulas are stronger than a moment identity.
By \cite[equation (3.5.1), p.~209, and
equations (4.11.3)--(4.11.5), p.~220]{Nakamura1995},
the actual finite coefficients satisfy
\begin{equation}\label{eq:cm-alpha-theta-finite}
 \alpha_n(h)(x)=12\kappa_{\theta_E(x)}(h)\quad(x\ne0),
 \qquad
 \alpha_n(h)(0)=-n\kappa_{p^{12}}(h).
\end{equation}
The latter formula also uses
\cite[Proposition (3.8), p.~210]{Nakamura1995}.
Ishii's special-lift measure is this measure, with basis
compatibility recorded in
\cite[Remark 2.6.1(2), p.~8]{Ishii}.
CM-unit invariance and the ray-field description in
\cite[p.~27]{GoncharovEuler} place these rational theta
classes over \(L_n\).

We verify that their restriction to \(H\) loses no
nonzero rational class.  Set
\(G_n=\operatorname{Gal}(\Omega/L_n)\) and
\(\Gamma_n=\operatorname{Gal}(L_\infty/L_n)\).
Finite-field Kummer theory injects the \(p\)-unit
space of \(L_n\) into \(H^1(G_n,\mathbb Q_p(1))\):
the relevant Kummer extensions are pro-\(p\) and
unramified outside \(p\), and hence lie in \(\Omega\).
The restriction kernel to \(H\) is the image of
\(H^1(\Gamma_n,\mathbb Q_p(1))\), which is zero.
Indeed, choose \(\gamma_0\) with
\(\chi_{\rm cyc}(\gamma_0)\ne1\).  Since \(\Gamma_n\)
is abelian, every cocycle \(f\) satisfies
\[
 (\chi_{\rm cyc}(\gamma_0)-1)f(\gamma)
       =(\chi_{\rm cyc}(\gamma)-1)f(\gamma_0),
\]
so it is a coboundary over \(\mathbb Q_p\).
This proves the required Kummer injectivity, including
the possible infinitely \(p\)-divisible kernel in the
base tower.

The nonzero coordinate functionals in
\eqref{eq:cm-alpha-theta-finite} span precisely
\(C_1(n)\); the origin coordinate is their negative
sum, since \(\alpha\) has total mass zero.
A proper rational image of \(A_n[1/p]\) in the dual
of \(C_1(n)\) would have a nonzero annihilator there,
contradicting the Kummer injectivity just proved.
This proves the first identification, while retaining
\(A_n\) as its actual integral lattice.

For \(\eta\in\bigwedge^2A_n[1/p]\), let
\(B_\eta\) denote the alternating pairing on \(C_1(n)\)
under that identification.  The inverse-image formula
for the shears is
\[
 (J_n\eta)(x,y)=\eta(x,y-x)+\eta(x-y,y)-\eta(x,y).
\]
When \(x,y,x-y\ne0\), evenness
\(\theta_E(-u)=\theta_E(u)\) and
\eqref{eq:cm-alpha-theta-finite} give
\begin{equation}\label{eq:cm-J-elliptic-triangle}
 (J_n\eta)(x,y)
       =144 B_\eta\bigl(\delta\theta_E(x,y-x,-y)\bigr).
\end{equation}
When \(x=0\), \(y=0\), or \(x=y\), the coefficient
vanishes by alternation and evenness, independently
of the origin coefficient.  Every nondegenerate
zero-sum triple occurs in
\eqref{eq:cm-J-elliptic-triangle}.
Thus the rational kernel is exactly the dual of
\(\mathcal R_n^{\rm ell}\), with the second Tate
twist.  Its integral kernel \(\mathcal D_n\) is
closed and \(p\)-saturated because its target is
finite free over \(\mathbb Z_p\).

Here are the transition maps, including the origin.
For \(x\ne0\), the isogeny distribution gives, in
additive notation for units tensored with \(\mathbb Q_p\),
\[
 \theta_E(x)=\sum_{[p]y=x}\theta_E(y).
\]
This is the ordinary inclusion of the theta class at
the same torsion point into the next field.  Hence
pushforward \(A_{n+1}\to A_n\) is dual to
\(C_1(n)\hookrightarrow C_1(n+1)\).
At zero the exact calculation is
\[
 \sum_{[p]y=0}\alpha_{n+1}(h)(y)
   =-(n+1)\kappa_{p^{12}}(h)
         +\sum_{0\ne y\in E[p]}\kappa_{\Theta(y)}(h)
   =-n\kappa_{p^{12}}(h),
\]
using \cite[Proposition (2.2)(3), p.~204]{Nakamura1995}.
For double units, independence of the torsion exponent
and the sum over preimage triples are
\cite[Lemma 2.12, p.~17, and equation (41), p.~16]{GoncharovEuler}.
Their differential is the exterior square of the
unit inclusion: two chosen preimages determine the
third by the zero-sum condition.  This proves the
commuting diagram of complexes and the asserted
natural map on cokernels.

Finally, compactness and injectivity of \(\alpha\)
identify \(A\) with \(\varprojlim A_n\).
The coefficient quotients modulo \(p^r\), for all
\(n,r\), form a cofinal system of neighborhoods.
Completed tensor and exterior squares therefore give
\(\widehat{\bigwedge}^2A=\varprojlim\bigwedge^2A_n\).
The shear maps commute with these projections,
so taking their kernels proves
\eqref{eq:cm-J-kernel-elliptic-limit}.
\end{proof}

The transition just used has an explicit comparison
with arithmetic norms.  On a torsion point whose two
Tate coordinates are units,
\(\operatorname{Gal}(L_{n+1}/L_n)\) acts by
\((1+p^nu,1+p^nv)\), \(u,v\in\mathbb F_p\), and
its orbit is the entire \(p^2\)-element \([p]\)-fiber.
The distribution product is then the field norm.
If one coordinate is divisible by \(p\) and the other
is a unit, the arithmetic orbit has \(p\) elements,
each repeated \(p\) times in the norm, whereas the
isogeny fiber has \(p^2\) elements.  The all-chart
transition in the theorem is the isogeny distribution,
with the displayed correction at zero.

For the Gaussian and Eisenstein fields,
\cite[Theorem 5.16, p.~34]{GoncharovEuler} supplies a
surjective morphism from the Bianchi modular complex
at the ideal \((p^n)\) to the complex with target
\(\bigwedge^2\widehat C_1(n)\), where
\(\widehat C_1(n)=C_1(n)\oplus\mathbb Q_p\theta_E(0)\).
Compose its degree-two map with the exterior square
of the projection \(\widehat C_1(n)\to C_1(n)\)
that sends the formal \(\theta_E(0)\) to zero.
The differentials of degenerate triangles vanish,
so this is a surjective morphism to the exact Euler
complex used above.  One can equally send the formal
origin to its actual normalized class \(-n[p]\).
At a prime ideal \(\mathfrak P\),
\cite[Theorem 1.2, p.~9]{GoncharovEuler} identifies
\[
 \operatorname{coker}\bigl(C_2(\mathfrak P)
       \longrightarrow\bigwedge^2C_1^{\rm un}(\mathfrak P)\bigr)
       =H^2_{\rm cusp}(\Gamma_1(\mathfrak P),\mathbb Q),
\]
where \(C_1^{\rm un}\) is the global-unit space.
Since the differential lands in that smaller exterior
square, this gives a canonical subspace of the full
unit cokernel after extension of scalars.  The
global-unit qualification is part of this precise map.

The remaining persistence issue has a concrete form.
Injectivity of the natural map on cokernels requires
\[
 \delta C_2(n+1)\cap\bigwedge^2C_1(n)=\delta C_2(n).
\]
The inclusions of unit spaces alone do not establish
this equality.  Moreover,
\eqref{eq:cm-J-kernel-elliptic-limit} retains the
integral lattices, so rational finite-level classes
must have compatible bounded duals to contribute to
the actual kernel.  The theorem supplies this exact
finite-level and inverse-limit comparison without
inferring such persistence from a dimension count.

\subsection{The centralizer pullback and its kernel}
\label{subsec:cm-actual-centralizer}

Let \(F\) be the geometric free pro-\(p\) group of the
once-punctured CM elliptic curve. The prime-to-\(p\) degree
of \(K(E[p])/L_1\) gives
\(K(E[p])\cap\Omega=L_1\), and restriction identifies
\(\operatorname{Gal}(K(E[p])\Omega/K(E[p]))\) with
\(G_1=\operatorname{Gal}(\Omega/L_1)\).
Ishii's kernel-field inclusion therefore defines the map
\[
 \rho:G_1\longrightarrow\operatorname{Out}(F)
\]
used here.
Put \(O=\rho(G_1)\), \(K=\ker\rho\), and let
\(\widetilde O\) be the inverse image of \(O\) in
\(\operatorname{Aut}(F)\).  The following pullback records the
arithmetic kernel without making a choice of automorphism lifts:
\[
 \mathcal P=\widetilde O\times_O G_1
    =\{(\alpha,g):[\alpha]=\rho(g)\}.
\]
Its geometric subgroup is
\(F\longrightarrow\mathcal P,\ f\longmapsto(\operatorname{inn}(f),1)\).
This is injective because \(F\) has trivial center.

\begin{proposition}\label{prop:cm-actual-centralizer-pullback}
There are exact identifications
\[
 C_{\mathcal P}(F)=\{(1,k):k\in K\}\simeq K,
 \qquad
 \mathcal P|_K=F\times K.
\]
For every closed subgroup \(A\subset F\),
\[
 N_{\mathcal P}(A)=N_{\widetilde O}(A)\times_O G_1.
\]
In particular, if \(V_A\subset O\) is the image of
\(N_{\widetilde O}(A)\), then the arithmetic normalizer image
is precisely \(\rho^{-1}(V_A)\).  It is open in \(G_1\) exactly
when \(V_A\) is open in \(O\).
\end{proposition}

\begin{proof}
Conjugation by \((\alpha,g)\) on the geometric subgroup is
the automorphism \(\alpha\).  It centralizes all of \(F\)
exactly when \(\alpha=1\), which forces \(g\in K\).
Above \(K\), every \(\alpha\) is inner.  The unique element
of \(F\) inducing it gives the displayed direct product,
with multiplication inherited from the pullback.

The same conjugation formula shows that an element normalizes
\(A\) exactly when \(\alpha(A)=A\).  This proves the
normalizer identity and its projection formula.  Finally
the continuous surjection \(G_1\to O\) has the quotient
topology, so a subgroup of \(O\) is open exactly when
its inverse image is open.
\end{proof}

Thus the geometric open-normalizer test factors through the
outer image \(O\).  To prove arithmetic faithfulness
from this pullback one must prove the further assertion
\(C_{\mathcal P}(F)=1\), or equivalently \(K=1\).
The proposition gives the exact map at that step.

There is also an arithmetic restriction on possible mixed
procyclic subgroups of the pullback.  Assume
\(\mathcal X_2^+\) is finite.  If \(A\subset F\times K\)
is closed procyclic and its projection to \(K\) is nontrivial,
then its normalizer in \(\mathcal P\) has trivial image
in \(\Gamma\).  Indeed the projection of any normalizer
element normalizes the nontrivial procyclic image of \(A\)
inside \(K\subset H\), so
Proposition~\ref{prop:cm-no-normalized-cyclic} applies.
Consequently a procyclic subgroup of \(F\times K\)
with open \(\Gamma\)-image normalizer must lie in \(F\).
This conclusion uses the centralizer pullback and
the arithmetic conjugation modules; it leaves the possible
nontrivial, infinitely generated kernel \(K\) as the
remaining centralizer question.

For completeness, the asserted infinite generation also
follows from those same modules.  If a nontrivial finitely
generated closed subgroup \(J\subset H\) were normalized
by a lift of \(\gamma\neq1\), choose the first nonzero
image of \(J\) in \(\gamma_dH/\gamma_{d+1}H\).
It is a finitely generated torsion-free
\(\mathbb Z_p\)-module preserved by \(\gamma\).
Its characteristic polynomial would be a nonzero
element of \(\Lambda_2\) annihilating that image,
contrary to the proved torsion-freeness of the
lower-central module.  Apply this to the normal subgroup
\(K\) to obtain the stated conclusion whenever \(K\neq1\).


\subsection{The examples over \(\mathbb Q(i)\) and
\(\mathbb Q(\sqrt{-3})\)}

We verify both arithmetic hypotheses of Ishii's theorem by
finite calculations, leaving its freeness hypothesis explicit.
The cases are those recorded in
\cite[Example 1.2, p.~3]{Ishii}.

\begin{proposition}\label{prop:cm-small-primes}
For \((K,p)=(\mathbb Q(i),5)\) and
\((\mathbb Q(\sqrt{-3}),7)\), the ray class field \(K(p)\) has
class number prime to \(p\), and \(K(p^2)\) has exactly two primes
above \(p\). For the reductions of
\(E:y^2=x^3-x\) at \(5\) and \(E:y^2=x^3+1\) at \(7\),
the Frobenius polynomials are respectively
\[
 U^2+2U+5,\qquad U^2+4U+7 .
\]
\end{proposition}
\begin{proof}
For a class-number-one imaginary quadratic field, the ray class
sequence gives
\[
 \operatorname{Gal}(K(p^n)/K)
  \simeq(\mathcal O_K/p^n)^\times/\mathcal O_K^\times .
\]
The units inject modulo \(p\), and the Chinese remainder map at
the two split primes gives degree
\((p^{n-1}(p-1))^2/w_K\).
At \(n=1\) this is \(4\) and \(6\), respectively.
The cyclotomic extension \(K(\mu_p)/K\) has degree \(p-1\) and
conductor dividing \(p\); hence
\[
 \mathbb Q(i)(5)=\mathbb Q(\zeta_{20}),\qquad
 \mathbb Q(\sqrt{-3})(7)=\mathbb Q(\zeta_{21}) .
\]
Equivalently, at level one the norm
\((u,v)\mapsto uv\) identifies
\((\mathbb F_p^\times)^2/\mathcal O_K^\times\)
with \(\mathbb F_p^\times\), since \(w_K=p-1\) in these two cases.

For \(K=\mathbb Q(i)\), take \(\pi=2+i\), \(\bar\pi=2-i\).
The isomorphism \(\mathcal O_K/\bar\pi^2\simeq\mathbb Z/25\)
sends \(i\) to \(7\). It sends \(\pi\) to \(9\), and
\(9^4\equiv11\not\equiv1\pmod {25}\).
The quotient
\((\mathcal O_K/\bar\pi^2)^\times/\mathcal O_K^\times\)
has order \(5\); the class of \(\pi\) generates it.
For \(K=\mathbb Q(\sqrt{-3})\), put
\(\omega=(-1+\sqrt{-3})/2\), \(\pi=3+\omega\), and
\(\bar\pi=2-\omega\). The quotient by \(\bar\pi^2\)
sends \(\omega\) to \(30\pmod {49}\), so \(\pi\) maps to \(33\).
Here \(33^6\equiv8\not\equiv1\pmod {49}\), and the corresponding
unit quotient has order \(7\).
In each case inertia at \(\mathfrak p\) supplies the first local
unit factor; the Frobenius at \(\mathfrak p\) supplies the generator
just calculated of the remaining quotient. The decomposition group
is therefore the full ray class group. Each of the two primes of
\(K\) above \(p\) has one extension to \(K(p^2)\).
This is the explicit reciprocity criterion in
\cite[Remark 2.14.1, p.~11]{Ishii}.

We next verify the prime-to-\(p\) class numbers. Minkowski's bound
provides representatives for every ideal class of norm at most
\[
 B_L=(4/\pi)^{r_2}\frac{[L:\mathbb Q]!}
 {[L:\mathbb Q]^{[L:\mathbb Q]}}\sqrt{|d_L|}
\]
\cite[Theorem 4.3, p.~70]{MilneANT}.
For \(L=\mathbb Q(\zeta_{20})\), its degree is \(8\),
\(|d_L|=2^8 5^6\), and \(\sqrt{|d_L|}=2000\).
The elementary estimate \(\pi>3\) gives
\(B_L<4375/288<16\).
The residue degree at \(2\) is \(4\), so the prime norms above
\(2\) equal \(16\). At \(3,7,11,13\) the residue degrees are
\(4,4,2,4\); their prime norms exceed the bound.
Only the primes above \(5\), of norm \(5\), can generate the
class group. The two primes above \(5\) have ramification index
\(4\) over \(\mathbb Q(i)\), and their fourth powers are the
principal ideals \((2+i)\) and \((2-i)\). The class group is
therefore generated by elements killed by \(4\).
Its order is prime to \(5\).

For \(L=\mathbb Q(\zeta_{21})\), its degree is \(12\),
\(|d_L|=3^6 7^{10}\), and \(\sqrt{|d_L|}=453789\).
Using \(\pi>31/10\) gives
\[
 B_L<
 \frac{8088368750000}{71887798161}<113.
\]
At \(2\) the residue degree is \(6\), giving norm \(64\).
At \(3\) the residue degree is \(6\), giving norm \(729\).
The primes above \(7\) have norm \(7\). Among unramified primes
with residue degree one below \(113\), the congruence
\(q\equiv1\pmod {21}\) leaves only \(q=43\).
For residue degree at least two, \(q^2<113\) restricts \(q\)
to \(2,3,5,7\), and \(5\) has residue degree \(6\).
Thus the primes of possible norm below the bound lie over
\(2,7,43\).

Let \(\xi=\zeta_7\) and
\(\alpha=\xi+\xi^2+\xi^4\). The identity
\(\alpha^2+\alpha+2=0\) shows that
\(|N_{L/\mathbb Q}(\alpha)|=64\).
Thus \((\alpha)\) is one of the primes above \(2\), and conjugation
shows that all of them are principal.
Let \(\zeta=\zeta_{21}\) and \(\beta=1+\zeta-\zeta^3\).
Its norm is \(43\). One direct computation is
\[
 \Phi_{21}(u)\equiv3u^2+4u+4\pmod {u^3-u-1},
 \qquad
 \operatorname{Res}(u^3-u-1,3u^2+4u+4)=43 .
\]
Consequently the primes above \(43\) are also principal.
The primes above \(7\) have sixth powers \((3+\omega)\) and
\((2-\omega)\). The class group is generated by elements killed
by \(6\); its order is prime to \(7\).

Finally, for \(y^2=x^3-x\) over \(\mathbb F_5\), the numbers of
solutions in \(y\) for \(x=0,1,2,3,4\) are
\(1,1,2,2,1\). Adding the point at infinity gives \(8\) points
and trace \(5+1-8=-2\).
For \(y^2=x^3+1\) over \(\mathbb F_7\), the corresponding numbers
are \(2,2,2,1,2,1,1\). There are \(12\) points and trace
\(7+1-12=-4\). These prove the two displayed polynomials for
the Frobenius endomorphisms of the reductions.
\end{proof}

\subsection{Locality in the two examples, with the unit calculation made explicit}

\begin{proposition}\label{prop:cm-locality}
For the two pairs of Proposition~\ref{prop:cm-small-primes},
put \(L=K(\mu_p)=K(p)\), and fix a prime \(\mathfrak p\) above \(p\).
The decomposition homomorphism is an isomorphism
\[
 G_{\mathbb Q_p(\mu_p)}^{(p)}
    \xrightarrow{\ \sim\ }\operatorname{Gal}(\Omega/L).
\]
\end{proposition}
\begin{proof}
Put \(\Delta=\operatorname{Gal}(L/K)\), and identify its
cyclotomic character with \(\omega\). Let
\(V=\mathcal O_L[1/p]^\times/p\). Kummer theory and
\(p\nmid h(L)\) identify \(V\) with \(H^1(G,\mu_p)\),
where \(G=\operatorname{Gal}(\Omega/L)\).
The two localization maps
\[
 V\longrightarrow
 L_{\mathfrak p}^{\times}/p\ \oplus\
 L_{\bar{\mathfrak p}}^{\times}/p
\]
are jointly injective. Indeed a unit in their joint
kernel would give a degree-\(p\) Kummer extension that
is split at the two primes above \(p\) and unramified
elsewhere, contradicting \(p\nmid h(L)\).

Complex conjugation exchanges the two primes.
Consequently localization at \(\mathfrak p\) alone is
injective on the image of
\(\mathcal O_{L^+}[1/p]^\times/p\).
The extension map from this real unit group is injective,
since norm followed by extension is multiplication by two.
Its dimension is \(p-1\), and its \(\Delta\)-representation
is the regular representation. In contrast \(V\) has
dimension \(p+1\) and representation
\[
 V\simeq\mathbb F_p[\Delta]\oplus\mathbb F_p
                      \oplus\mathbb F_p(1).
\]
The first two terms describe the free \(S\)-units, and the
last the \(p\)-power roots of unity. Hence the kernel of
localization at \(\mathfrak p\), being \(\Delta\)-stable,
can occur only in the characters \(1\) and \(\omega\).

The trivial-character component is generated by the
two units \(\pi,\bar\pi\) from \(K\). At \(\mathfrak p\),
\(\pi\) has valuation one and \(\bar\pi\) has nonzero
principal-unit class in \(\mathbb Q_p^\times/
(\mathbb Q_p^\times)^p\). The latter assertion is exactly
the conjugate of the congruence modulo \(p^2\) proved in
Proposition~\ref{prop:cm-small-primes}. The two local
classes are independent. They remain so after adjoining
\(\mu_p\), since that local extension has degree \(p-1\)
and its norm proves injectivity on classes modulo \(p\).
Thus the trivial-character part of the kernel is zero.

The \(\omega\)-component of \(V\) corresponds, after
Kummer untwisting, to \(H^1(G,\mathbb F_p)^\Delta\).
This has dimension two and is the inflation of
\(H^1(\Gamma,\mathbb F_p)\), where
\(\Gamma=\operatorname{Gal}(\bigcup_nK(p^n)/L)\).
The decomposition group in \(\Gamma\) is all of \(\Gamma\):
its image modulo \(p\) is already full by the calculation
at modulus \(p^2\), and a closed subgroup of \(\mathbb Z_p^2\)
surjecting modulo \(p\) is the whole group. Therefore
restriction to the local Galois group is injective on
this last character component as well.
We have proved that the whole localization map on \(V\)
is injective.

The local-to-global homomorphism \(Q\to G\), with
\(Q=G_{\mathbb Q_p(\mu_p)}^{(p)}\), is consequently
surjective: a proper closed image in a pro-\(p\) group
would be annihilated by a nonzero homomorphism to
\(\mathbb F_p\). Both \(H^1(G,\mathbb F_p)\) and
\(H^1(Q,\mathbb F_p)\) have dimension \(p+1\), so the
induced injection is an isomorphism.

Restriction \(H^2(G,\mu_p)\to H^2(Q,\mu_p)\) is
injective. The global group injects into the
\(p\)-torsion of the \(S\)-integer Brauer group, identified
by local invariants with
\(\{(a,b)\in\mathbb F_p^2:a+b=0\}\); localization at
\(\mathfrak p\) is projection to the first coordinate
and is injective on this group. These maps commute
with restriction in Galois cohomology.

For \(N=\ker(Q\to G)\), the five-term sequence now
gives \(H^1(N,\mathbb F_p)^G=0\).
A nontrivial pro-\(p\) group \(N\) has a nonzero
homomorphism to \(\mathbb F_p\). Its finite orbit
under the continuous pro-\(p\) action spans a nonzero
finite-dimensional representation of a finite
\(p\)-group, which has a nonzero fixed vector.
Thus \(H^1(N,\mathbb F_p)^G\ne0\) whenever \(N\ne1\).
The contradiction proves \(N=1\).
\end{proof}

The real \(S\)-unit dimension and the extra trivial-character
component in this proof repair the corresponding dimension
step in the proof of \cite[Theorem A.1, pp.~28--29]{Ishii}.
For these examples the numerical dimensions are \(4,6\) at
\(p=5\), and \(6,8\) at \(p=7\), for the real and complex
groups respectively.

For completeness, this also repairs the locality input to
nonvanishing of all the admissible elliptic Soul\'e characters
in these examples. Here \(w_K=p-1\), so an admissible pair
satisfies \(m_1\equiv m_2\pmod{p-1}\). The local character is
\(\chi_{\mathrm{cyc}}^{m_1}\eta^{m_2-m_1}\), where \(\eta\)
is the unramified CM unit-root character. Its Frobenius
value is an algebraic number of complex absolute value
\(\sqrt p\), and hence is not a root of unity.
Local Tate duality gives
\[
 H^2(\mathbb Q_p(\mu_p),\mathbb Q_p(\boldsymbol m))=0
 \qquad(\boldsymbol m\in I_K):
\]
the dual representation has an invariant only when inertia
forces \(m_1=1\) and the unramified character then forces
\(m_2=1\), the excluded pair. Locality and prime-to-\(p\)
descent give global \(H^2\) finiteness. The nonvanishing
assertion follows from \cite[Theorem 2.8(2), p.~9]{Ishii}.

\subsection{An unconditional abelian quotient for the two CM curves}

Let \(\Omega^*\subset\Omega\) be the pro-\(p\) field constructed
in the proof of Ishii's Lemma 2.13, p.~10. Its proved field
identity is
\[
 \overline K^{\ker\rho_{X,p}}=K(E[p])\Omega^* .
\]
The finite extension \(K(E[p])/K(p)\) has degree prime to \(p\),
and its factor in the outer image is central. The action on
the CM Tate module gives \(L_\infty=\bigcup_nK(p^n)\subset\Omega^*\).
Thus the kernel of \(G=\operatorname{Gal}(\Omega/K(p))
\to\operatorname{Gal}(\Omega^*/K(p))\) is contained in
\(H=\operatorname{Gal}(\Omega/L_\infty)\).

\begin{proposition}\label{prop:cm-abelian-detection}
For \(E:y^2=x^3-x\) over \(K=\mathbb Q(i)\), \(p=5\), and
for \(E:y^2=x^3+1\) over \(K=\mathbb Q(\sqrt{-3})\), \(p=7\),
the joint elliptic Soul\'e map is injective:
\[
 A=H^{\mathrm{ab}}
   \longrightarrow\prod_{\boldsymbol m\in I_K}
                 \mathbb Z_p(\boldsymbol m),
 \qquad a\longmapsto(\kappa_{\boldsymbol m}(a))_{\boldsymbol m}.
\]
In particular,
\[
 \operatorname{Gal}(\Omega/\Omega^*)\subset
       \overline{[H,H]},
 \qquad
 \Omega^{\overline{[H,H]}}\subset\Omega^* .
\]
\end{proposition}
\begin{proof}
The arithmetic hypotheses used in this proof have been
established in Proposition~\ref{prop:cm-small-primes}.
Theorem 2.8(2), p.~9, of \cite{Ishii} states that
\(\kappa_{\boldsymbol m}\) factors through \(A\).
Lemma 4.1, p.~17, gives its nonvanishing for every
\(\boldsymbol m\in I_K\).

Write \(\Lambda_2=\mathbb Z_p[[S,T]]\) for the group ring
of \(\operatorname{Gal}(L_\infty/K(p))\). The construction
in Lemma 4.3 and Proposition 4.4, pp.~19--20, gives an
equivariant surjection from one rank-one module for each
tame character onto its component of \(A\). The ordinary
source components are copies of \(\Lambda_2\); the exceptional
component is \(\operatorname{Ann}_{\Lambda_2}\mathbb Z_p(1)\).
Their total rank is \([K(p):K]\), which equals the rank of \(A\).
Since the source is torsion-free, the surjection has zero
kernel. Consequently every tame-character component of \(A\)
is torsion-free of rank one over \(\Lambda_2\).

We spell out the rank-one separation argument, including the
exceptional ideal. Fix a tame component \(M\), choose
\(0\ne v\in M\), and use its rank-one fraction space to choose
\(0\ne b\in\Lambda_2\) with
\[
 bM\subset\Lambda_2v .
\]
This is possible by clearing the denominators of finitely
many \(\Lambda_2\)-generators. For \(a\in M\), write
\(ba=f_av\), \(f_a\in\Lambda_2\).

Choose the two local generators so that their nontrivial
CM-character values are \(u=1+p\). For the fixed tame
component, the pairs
\[
 m_1\equiv r_1\pmod{p-1},\qquad
 m_2\equiv r_2\pmod{p-1}
\]
with both coordinates sufficiently large form a Cartesian
grid in \(I_K\). Set \(t_i=u^{m_i}-1\).
Equivariance gives
\[
 b(t_1,t_2)\kappa_{\boldsymbol m}(a)
       =f_a(t_1,t_2)\kappa_{\boldsymbol m}(v).
\]
At a grid point where \(b(t_1,t_2)\ne0\),
\(\kappa_{\boldsymbol m}(v)\ne0\): otherwise the same relation
for every generator of \(M\) would force the nonzero character
\(\kappa_{\boldsymbol m}\) to vanish identically.
Thus, when all characters kill \(a\), \(f_a\) vanishes at
every grid point at which \(b\) is nonzero. The product
\(bf_a\) vanishes at every grid point, including those where
\(b=0\).

A bounded one-variable power series has only finitely many
zeros in \(p\mathbb Z_p\), unless it is zero, by Weierstrass
preparation. Applying this fact first with \(t_2\) fixed,
then to the coefficients as \(t_2\) varies, proves that a
bounded two-variable series vanishing on this Cartesian
grid is zero. Hence \(bf_a=0\), so \(f_a=0\), and
torsion-freeness gives \(a=0\). Applying this argument to
each tame component proves the injectivity.

The elliptic Soul\'e characters are obtained from the
metabelian geometric outer action, so they kill the actual
outer kernel. Its image in \(A\) is therefore zero by the
injectivity just proved. This is the first containment.
The field containment follows by Galois correspondence.
\end{proof}

This proof uses the explicit module structure and the
nonzero characters. It does not use the freeness assertion
in Ishii's Conjecture 2.3 or the asserted converse in his
Proposition 2.11. The remaining subgroup
\(\overline{[H,H]}\) is explicit in the conclusion.

\subsection{The nonzero two-variable transgression}

\begin{proposition}\label{prop:cm-nonsplit}
For each of the two pairs in Proposition~\ref{prop:cm-small-primes},
put \(L=K(p)=K(\mu_p)\), \(L_\infty=\bigcup_nK(p^n)\),
\(G=\operatorname{Gal}(\Omega/L)\), and
\(H=\operatorname{Gal}(\Omega/L_\infty)\).
The exact sequence
\[
 1\longrightarrow H\longrightarrow G\longrightarrow
 \Gamma=\operatorname{Gal}(L_\infty/L)\simeq\mathbb Z_p^2
 \longrightarrow1
\]
has no continuous section. In fact, every pair of lifts to \(G\)
of a \(\mathbb Z_p\)-basis of \(\Gamma\) has a commutator detected
by a \(\Gamma\)-invariant homomorphism \(H\to\mathbb F_p\).
\end{proposition}
\begin{proof}
Let \(\Delta=\operatorname{Gal}(L/K)\), and write \(\omega\)
for its action on \(\mu_p\). Here \(\Delta\) has order \(p-1\)
and \(\omega\ne1\). Its action on \(\Gamma\) is trivial, since
\(L_\infty/K\) is abelian. A prime-to-\(p\) complement supplies
the \(\Delta\)-action on \(G\), and all invariant calculations
are independent of that choice.

Kummer theory and the already proved \(p\nmid h(L)\) give
\[
 H^1(G,\mu_p)\simeq\mathcal O_L[1/p]^\times/p .
\]
There are precisely two primes above \(p\), both fixed by
\(\Delta\). The rational \(S\)-unit representation is
\(\mathbb Q[\Delta]\oplus\mathbb Q\): the archimedean places
give the regular permutation representation, the two finite
places give two trivial representations, and the product
formula removes one trivial representation. Therefore the
nontrivial character \(\omega\) occurs once in its free part.
It occurs once more in the torsion subgroup \(\mu_p\).
Since \(p\nmid|\Delta|\), the integral character projectors
preserve these multiplicities upon reduction modulo \(p\).
Untwisting gives
\[
 \dim_{\mathbb F_p}H^1(G,\mathbb F_p)^\Delta=2.
\]

The comparison map from maximal pro-\(p\) group cohomology
to \'{e}tale cohomology is injective in degree two.
One can verify the needed injectivity by the five-term
sequence: a nonzero invariant homomorphism from the kernel
of the maximal pro-\(p\) quotient to a finite \(p\)-group
would produce a larger pro-\(p\) quotient. Thus
\[
 H^2(G,\mu_p)\hookrightarrow
 H^2_{\mathrm{\acute et}}(\mathcal O_L[1/p],\mu_p)
 \simeq\operatorname{Br}(\mathcal O_L[1/p])[p].
\]
The last isomorphism uses the vanishing of the class-group
term. Local Brauer invariants identify the last group with
\[
 \{(a,b)\in\mathbb F_p^2:a+b=0\}.
\]
The action of \(\Delta\) on these invariants is trivial,
because it fixes each of the two primes. On removing the
coefficient twist, it follows that
\[
 H^2(G,\mathbb F_p)^\Delta=0.
\]

The inflation map from \(H^1(\Gamma,\mathbb F_p)\) is
injective and both it and \(H^1(G,\mathbb F_p)^\Delta\)
have dimension two. It is consequently an isomorphism.
Taking \(\Delta\)-invariants in the Hochschild--Serre
five-term sequence now gives an isomorphism
\[
 d_2:\operatorname{Hom}(H,\mathbb F_p)^{\Gamma,\Delta}
       \xrightarrow{\ \sim\ }
       H^2(\Gamma,\mathbb F_p)\simeq\mathbb F_p .
\]
Choose \(\varphi\) with \(d_2\varphi\ne0\). The quotient
of \(G\) by \(\ker\varphi\) is the associated central
extension of \(\Gamma\) by \(\mathbb F_p\).
For coordinates \((a,b)\) on \(\Gamma\), the cocycle
\[
 ((a,b),(a',b'))\longmapsto ab'\pmod p
\]
represents a generator of \(H^2(\Gamma,\mathbb F_p)\).
Its alternating part has value one on a basis. Hence
\(\varphi([\widetilde\gamma_1,\widetilde\gamma_2])\ne0\)
for every pair of lifts of a basis. In particular the
lifts cannot commute, and a section cannot exist.
\end{proof}

This is the specific obstruction to using two commuting
arithmetic lifts in these examples. It agrees with Ishii's
construction, in which \([\gamma_1,\gamma_2]\) supplies the
trivial tame-character component of the abelianized kernel
\cite[\S4.1--4.2]{Ishii}. It is distinct from the
semilinear fixed-part calculation only in the explicit
sense established by the two displayed maps: the latter
calculates the coefficient-ring action, while \(d_2\)
calculates the extension class.

\subsection{The cyclotomic descent map}

Continue with either of the two explicit pairs, and put
\[
 L_\infty=K(p^\infty),\qquad
 L_{\mathrm c}=K(\mu_{p^\infty}),\qquad
 \Gamma_{\mathrm a}=\operatorname{Gal}(L_\infty/L_{\mathrm c}).
\]
The Weil-pairing identity \(\chi_1\chi_2=\chi_{\mathrm{cyc}}\)
identifies \(\Gamma_{\mathrm a}\) with the subgroup generated
by \(\gamma_1\gamma_2^{-1}\). Thus the actual group-ring quotient is
\[
 \mathbb Z_p[[S,T]]\longrightarrow\mathbb Z_p[[Z]],
 \qquad S\longmapsto Z,\quad T\longmapsto Z,
 \qquad \ker=(S-T).
\]
Complex conjugation exchanges \(\gamma_1,\gamma_2\), inverts
\(\Gamma_{\mathrm a}\), and acts trivially on the cyclotomic
quotient.

Put \(H_{\mathrm c}=\operatorname{Gal}(\Omega/L_{\mathrm c})\)
and \(\mathfrak X_{\mathrm c}=H_{\mathrm c}^{\mathrm{ab}}\).
The exact sequence
\(1\to H\to H_{\mathrm c}\to\Gamma_{\mathrm a}\to1\)
splits, since \(\Gamma_{\mathrm a}\simeq\mathbb Z_p\) is
projective in the category of pro-\(p\) groups. Abelianizing
this semidirect product gives the canonical exact sequence
\[
 0\longrightarrow A_{\Gamma_{\mathrm a}}
  \longrightarrow\mathfrak X_{\mathrm c}
  \longrightarrow\Gamma_{\mathrm a}\longrightarrow0.
\]
Taking fixed parts for complex conjugation is exact, because
two is invertible, and \(\Gamma_{\mathrm a}^+=0\). Hence
\[
 \mathfrak X_{\mathrm c}^+
 \simeq\bigl(A/(S-T)A\bigr)^+
 \simeq A^+/(S-T)A^- .
\]
For the last equality, every fixed quotient class has a
fixed lift by averaging. Moreover, if \((S-T)a\) is fixed,
it equals \((S-T)(a-ca)/2\). This proves the asserted kernel
without assuming that \(A^+\) is a module over the whole
two-variable ring.

Let \(X_{\mathrm c}\) be the unramified abelian pro-\(p\)
module over \(L_{\mathrm c}\), and let \(I_{\mathrm c}\)
be the closed subgroup of \(\mathfrak X_{\mathrm c}\)
generated by inertia at the primes above \(p\).
The exact ramification quotient is
\[
 0\longrightarrow I_{\mathrm c}^+
 \longrightarrow\mathfrak X_{\mathrm c}^+
 \longrightarrow X_{\mathrm c}^+\longrightarrow0.
\]
Thus the class-group fixed part is reached through both
the displayed anti-cyclotomic coinvariants and this
explicit inertia quotient.

\begin{proposition}\label{prop:cm-unramified-zero}
For the two explicit examples, \(X_{\mathrm c}^+=0\).
In addition, the full unramified abelian pro-\(p\)
module over \(L_\infty\) is zero.
\end{proposition}
\begin{proof}
Write \(L_r=K(\mu_{p^r})\) and \(L_r^+\) for its real
subfield. Extension of ideals and norm identify the
\(p\)-primary class group of \(L_r^+\) with the fixed
part of that of \(L_r\): their composites are multiplication
by two, and extension followed by norm on a fixed class
is also multiplication by two. The initial real class
group therefore has order prime to \(p\).

Each extension \(L_{r+1}^+/L_r^+\) is cyclic of degree
\(p\) and has exactly one ramified prime, the unique
prime above \(p\), which is totally ramified.
Indeed \(p\) splits into two primes in \(K\), and complex
conjugation exchanges the two primes in \(L_r\);
passing to the real field identifies them. The local
cyclotomic tower is totally ramified.

Here is the class-field induction. In a cyclic degree-\(p\)
extension \(F'/F\) with just one totally ramified prime,
suppose the \(p\)-class group of \(F\) is zero.
The Hilbert \(p\)-class field \(B\) of \(F'\) is Galois
over \(F\), with a finite \(p\)-group \(P\). Its inertia
subgroup at that prime has order \(p\), since \(B/F'\)
is unramified. The normal closure of this inertia is
all of \(P\): a nontrivial unramified quotient would
have a cyclic quotient of order \(p\), contrary to
the assumption on \(F\). Thus \(P\) is normally
generated by one element, so \(P/\Phi(P)\) is cyclic,
and the Burnside basis theorem makes \(P\) cyclic.
Its inertia is then \(P\), which has order \(p\);
therefore \(B=F'\). Induction proves that every
\(p\)-class group of \(L_r^+\) is zero. Passing to
norm limits gives \(X_{\mathrm c}^+=0\).

For the second assertion, use the proved locality
isomorphism of Proposition~\ref{prop:cm-locality}.
Let \(I_{\mathfrak p}\subset G\) be local inertia.
Then \(G/I_{\mathfrak p}\simeq\mathbb Z_p\), the
maximal pro-\(p\) unramified local quotient.
Under \(G\to\Gamma=\operatorname{Gal}(L_\infty/L)\),
the inertia image is exactly the first principal-unit
factor \(\mathbb Z_p\times0\). This is the ray-class
reciprocity map; the quotient is the second
\(\mathbb Z_p\)-factor. Since decomposition is all
of \(\Gamma\), the induced map
\[
 G/I_{\mathfrak p}\longrightarrow
 \Gamma/\operatorname{im}(I_{\mathfrak p})
\]
is a surjection between two copies of \(\mathbb Z_p\),
and hence an isomorphism. Its kernel calculation
gives \(H\subset I_{\mathfrak p}\).

Every finite unramified pro-\(p\) extension of
\(L_\infty\) descends to a finite layer as an extension
unramified outside \(p\); its normal closure over \(L\)
is still a \(p\)-extension and lies in \(\Omega\).
In \(\Omega/L_\infty\), the inertia subgroup at
\(\mathfrak p\) is \(H\cap I_{\mathfrak p}=H\).
An unramified quotient must kill it and is therefore
trivial. This proves the assertion about the full
unramified abelian module.
\end{proof}

\subsection{The two-variable cup edge}

The extension class affects the precise cohomological map.
For a finite coefficient module \(B\) on which \(H\) acts
trivially, the Hochschild--Serre sequence for
\(1\to H\to G\to\Gamma\to1\), using
\(\operatorname{cd}_p(H)=1\) and
\(\operatorname{cd}_p(\Gamma)=2\), gives
\[
 0\longrightarrow
 \operatorname{coker}\!\left(
  H^1(H,B)^\Gamma\xrightarrow{d_2}H^2(\Gamma,B)\right)
 \longrightarrow H^2(G,B)
 \longrightarrow H^1(\Gamma,H^1(H,B))
 \longrightarrow0 .
\]
The first arrow is inflation. The differential \(d_2\)
is the extension transgression calculated in
Proposition~\ref{prop:cm-nonsplit}.
For a \(\Gamma\)-module \(M\), put \(D_i=\gamma_i-1\).
The two-variable Koszul complex computes the right
side as
\[
 H^1(\Gamma,M)=
 \frac{\{(m_1,m_2):D_2m_1=D_1m_2\}}
      {\{(D_1m,D_2m):m\in M\}} .
\]
For \(B=\mathbb Z/p^n(2)\), complex conjugation acts on
this complex by \(c\) in degree zero, by
\((m_1,m_2)\mapsto(cm_2,cm_1)\) in degree one,
and by \(-c\) in degree two. The sign in degree two
comes from exchanging the two generators.
In particular the fixed part of the inflated term is
zero, whereas the fixed part of the edge is represented
by the equation
\[
 D_2m=D_1cm
\]
modulo \(D_1M^+\). These explicit maps identify what
must replace the procyclic coinvariant calculation
in a two-variable cup argument.

\subsection{An explicit nonzero commutator in finite geometric quotients}
\label{subsec:cm-low-class-two}

This computation specifies a small nonabelian quotient of the
outer image.  Its definition does not assert detection of the full
class-two quotient of the arithmetic source.
Let \(\Pi=\pi_1^{(p)}((E-\{O\})_{\bar K})\), with descending central
series \(\Pi(n)\), and put
\[
 P=\Pi/\Pi(13),\qquad
 P_r=P/\overline{\langle x^{p^r}:x\in P\rangle}\quad(r\geq1).
\]
These are characteristic quotients of \(\Pi\). The successive
lower-central quotients of \(P_r\) are finitely generated abelian
groups of exponent dividing \(p^r\), so every \(P_r\) is a finite \(p\)-group.
The field identity~\eqref{eq:cm-radical-outer-field} gives the
actual induced outer map \(\rho_H:H\longrightarrow\operatorname{Out}(\Pi)\):
indeed, restriction identifies
\(\operatorname{Gal}(K(E[p])\Omega/K(E[p]))\) with
\(\operatorname{Gal}(\Omega/L_1)\), because the intersection is
\(L_1\), and the outer action factors through its quotient
\(\operatorname{Gal}(\Omega^*/L_1)\).
Define the finite image
\[
 Q_r=\operatorname{im}\bigl(H\xrightarrow{\rho_H}
                              \operatorname{Out}(P_r)\bigr).
\]

\begin{proposition}\label{prop:cm-small-nonabelian-image}
For every class-number-one imaginary quadratic \(K\), every
split \(p\geq5\), and every \(E/K\) as above, all \(Q_r\) have
nilpotency class at most two.  They are nonabelian for all
sufficiently large \(r\).
\end{proposition}

\begin{proof}
We first give the rational Lie calculation.
Let \(L\) be the free Lie algebra over \(\mathbb Q\) on \(a,b\).
Let \(D_{rs}\) denote its space of derivations with
\[
 D(a)\in L_{r+1,s},\qquad D(b)\in L_{r,s+1},\qquad
                         D([a,b])=0.
\]
For a Lyndon word \(w\), write \([w]\) for its standard bracket:
split off its longest proper Lyndon suffix and bracket
recursively, with \(a<b\).
The following coefficient vectors define \(d\in D_{22}\) and
\(e,f\in D_{33}\).  The entries of each row are coefficients in
the displayed ordered Lyndon basis:
\[
\begin{array}{c|c|r@{\,,\,}r@{\,,\,}r@{\,,\,}r@{\,,\,}r}
 &\text{ordered basis}&\multicolumn{5}{c}{\text{coefficients}}\\ \hline
d(a)&[aaabb],[aabab]&\multicolumn{5}{c}{(-2,-1)}\\
d(b)&[aabbb],[ababb]&\multicolumn{5}{c}{(2,3)}\\
e(a)&[aaaabbb],[aaababb],[aaabbab],[aabaabb],[aababab]
       &1&-5&-6&1&-2\\
e(b)&[aaabbbb],[aababbb],[aabbabb],[aabbbab],[abababb]
       &-1&4&6&8&0\\
f(a)&[aaaabbb],[aaababb],[aaabbab],[aabaabb],[aababab]
       &-3&-9&-6&-3&-2\\
f(b)&[aaabbbb],[aababbb],[aabbabb],[aabbbab],[abababb]
       &3&12&6&0&8
\end{array}
\]
Direct expansion gives
\(d([a,b])=e([a,b])=f([a,b])=0\).
The boundary map
\[
 L_{r+1,s}\oplus L_{r,s+1}\longrightarrow L_{r+1,s+1},
 \qquad (u,v)\longmapsto[u,b]+[a,v],
\]
is surjective: a free Lie algebra generated by \(a,b\) is the
sum of \([a,L]\) and \([b,L]\) in every degree greater than one.
The multigraded Witt formula gives
\[
 \dim L_{3,2}=\dim L_{2,3}=2,\quad \dim L_{3,3}=3,
 \qquad
 \dim L_{4,3}=\dim L_{3,4}=5,\quad \dim L_{4,4}=8.
\]
Thus \(\dim D_{22}=1\) and \(\dim D_{33}=2\).
The displayed nonzero \(d\) spans \(D_{22}\);
the distinct last two \(b\)-coefficient vectors show that \(e,f\)
are independent and hence form a basis of \(D_{33}\).

To verify injectivity of the bracket map, only two associative
word coefficients are needed.  Expansion by the derivation rule
gives
\[
\begin{array}{c|rr}
 \text{coefficient in the image of }a &[d,e]&[d,f]\\ \hline
 aaaaabbabbb&-2&6\\
 aaaababbabb&-22&18
\end{array}
\qquad
\det\begin{pmatrix}-2&6\\-22&18\end{pmatrix}=96.
\]
Consequently \([d,\cdot]:D_{33}\longrightarrow D_{55}\) is
injective over \(\mathbb Q_p\), for every \(p\).
All equalities here are integral free-associative-algebra
identities.  The accompanying exact program
\texttt{check\_cm\_low\_class\_two.py} reconstructs the Lyndon
bases, checks the boundary equations, and records the complete
coefficient certificate.

We now identify the arithmetic classes to which this calculation
applies.  By \cite[Theorem 2.8 and Lemma 2.10, p.~9]{Ishii},
\(\kappa_{(2,2)}\) and \(\kappa_{(3,3)}\) are nonzero on the
weight-four and weight-six graded quotients, respectively.
Their nonvanishing uses the proven finiteness of the
second cohomology groups with Tate twists two and three.
The CM character decomposition of the rational graded pieces
is semisimple by \cite[Proposition 2.2]{Ishii}.
It therefore supplies nonzero rational classes of pure
bidegrees \((2,2)\) and \((3,3)\).
Clearing their denominators gives integral graded classes,
represented by actual elements \(h_4\in F^4H\) and
\(h_6\in F^6H\).
These elements are not asserted to satisfy exact eigen-equations
under conjugation.

The symplectic derivation comparison
\cite[Proposition 3.3(4), pp.~12--14]{IshiiLie} identifies these
two outer graded classes with nonzero elements of \(D_{22}\)
and \(D_{33}\).  Its inner kernel is generated by
\(\operatorname{ad}[a,b]\) in bidegree \((1,1)\);
it has zero component in bidegrees \((2,2),(3,3),(5,5)\).
The calculation therefore proves that
\([h_4,h_6]\) has nonzero outer class in weight ten.
In particular its image in \(\operatorname{Out}(P)\) is nontrivial.

By \cite[Proposition 2.2(1)]{Ishii}, the weight-one,
weight-two, and weight-three arithmetic graded pieces vanish.
Thus \(\rho_H(H)\subseteq F^4\operatorname{Out}(\Pi)\).
Centrality of the weight filtration gives
\[
 \rho_H(\gamma_3H)\subseteq F^{12}\operatorname{Out}(\Pi)
       =\ker\bigl(\operatorname{Out}(\Pi)\longrightarrow
                              \operatorname{Out}(\Pi/\Pi(13))\bigr).
\]
Every \(Q_r\) consequently has class at most two.

Finally, the natural map
\(\operatorname{Out}(P)\longrightarrow
 \varprojlim_r\operatorname{Out}(P_r)\) is injective.
To see this, suppose an automorphism is inner modulo every
power subgroup defining \(P_r\).  For each \(r\), its possible
conjugating elements form a nonempty closed subset of the
compact group \(P\), and these subsets are nested.
Their intersection contains a conjugator on \(P\), since
the intersection of those power subgroups is trivial.
The latter assertion follows because every finite quotient of
the pro-\(p\) group \(P\) has finite \(p\)-power exponent.
Hence the nontrivial outer image of \([h_4,h_6]\) survives
in some \(\operatorname{Out}(P_r)\) and in every later one.
This proves the proposition.
\end{proof}

The proposition gives an explicit nonzero central bracket in
the actual finite outer images.  Its source is the displayed
homomorphism \(H\to Q_r\); it does not identify the kernel of
the different homomorphism \(H\to H/\gamma_3H\).


\subsection{What the geometric common kernel gives}
\label{subsec:cm-geometric-common-kernel}

Write \(\Pi=\widehat{\langle x,y\rangle}\) for the geometric
fundamental group of a once-punctured elliptic curve and
\[
 R_{\mathrm{geom}}=
 \bigcap_\ell\ker(\Pi\longrightarrow\Pi^{(\ell)}).
\]
For cusp inertia \(I=\overline{\langle[x,y]\rangle}\), one has
\(I\cap R_{\mathrm{geom}}=1\). Indeed, for every \(n\geq1\)
the homomorphism to \(\operatorname{UT}_3(\mathbb Z/\ell^n)\)
sending \(x\) to \(1+E_{12}\) and \(y\) to \(1+E_{23}\)
sends \([x,y]\) to \(1+E_{13}\), of order \(\ell^n\).
The map \(\widehat{\mathbb Z}\to I\) is consequently injective:
each nonzero component \(\alpha_\ell\) is detected modulo some
\(\ell^n\). This also proves the asserted intersection, for every
conjugate of \(I\).

The arithmetic outer kernel occurs in the quotient row of the
fundamental exact sequence. More precisely, for
\[
 1\longrightarrow\Pi^{(p)}\longrightarrow
 \pi_1(X)^{(p,\mathrm{geom})}\longrightarrow G_K\longrightarrow1
\]
and \(N=\ker(\rho_{X,p})\), the inverse image of \(N\) equals
\(\Pi^{(p)}C\), where \(C\) is the centralizer of \(\Pi^{(p)}\)
in the middle group. To check this equality, take an element
above \(N\). Its action on \(\Pi^{(p)}\) is conjugation by some
\(u\in\Pi^{(p)}\); multiplication by \(u^{-1}\) places it in \(C\).
The converse follows by reversing this operation.
Since the center of the rank-two free pro-\(p\) group is trivial,
\(C\to N\) is an isomorphism. Thus Theorem~\ref{thm:geometric-procyclic-recognition} identifies the
procyclic subgroups of \(R_{\mathrm{geom}}\) with open arithmetic
normalizer image, while the displayed isomorphism places \(N\)
in the centralizer \(C\). A proof using that
recognition assertion to annihilate \(N\) still requires a map
from a nontrivial element of \(C\) to a subgroup to which the
recognition assertion applies. No such map is established here.

The limitation of that particular input also has an explicit
group-theoretic test. Start with an action \(G\to\operatorname{Aut}(\Pi)\)
and its extension \(\Pi\rtimes G\), and choose a nonzero pro-\(\ell\)
group \(D\). Pull back by the projection
\(q:G\times D\to G\). The pulled-back extension is
\((\Pi\rtimes G)\times D\). For every subgroup \(A\subset\Pi\),
its normalizer in this extension is
\(N_{\Pi\rtimes G}(A)\times D\). Its image in \(G\times D\)
is therefore the inverse image under \(q\) of the old normalizer
image. Such an inverse image is open exactly when the original
image is open: one implication is continuity, and the converse
follows because \(q\) is an open quotient map. Every recognition
assertion expressed by these normalizer images is consequently
preserved. The new outer action is \(\rho\circ q\), with kernel
\(\ker(\rho)\times D\). Thus this recognition input alone leaves
an arbitrary factor \(D\) in the arithmetic outer kernel.
This calculation concerns the stated group extension; it claims
no realization of the pullback as a number-field Galois action.

\subsection{The remaining elliptic gate}

For the two examples just computed, Ishii's generation theorem
applies: its cohomological input follows from the two arithmetic
hypotheses by \cite[Lemma 1.5, p.~3]{IshiiLie}.
The source group is free by Proposition~\ref{prop:cm-finite-plus-free}
and the vanishing proved above for these examples. Freeness of the
arithmetic graded Lie algebra in Ishii's proposed presentation remains
open for these two fields.
For the other class-number-one CM fields, the explicit Pollack
relation above proves that the proposed free presentation must
be replaced by a presentation incorporating relations.
The weighted-completion map is a surjection onto the unipotent
outer image; its injectivity is precisely an additional assertion,
as the known Pollack relation demonstrates in the indicated
bidegree. This gives a specific morphism to test, rather than
an inference from the existence of elliptic polylogarithms.

Goto and Ishii's more recent Coleman--Ihara formula
\cite[Theorem 1.1 and Corollary 5.5]{GotoIshii2026}
expresses restrictions of elliptic Soul\'e characters to local
reciprocity images through two-variable Coates--Wiles
homomorphisms and Katz \(p\)-adic \(L\)-values. Its domain is the
norm limit of semilocal principal units, and its target is a
one-dimensional character. Sano and Venjakob's extended
interpolation formula
\cite[Theorem 1.2.1, p.~3, and Theorem 2.4.1, p.~8]{SanoVenjakob2026}
uses \(V_\chi=\chi_{\mathrm{cyc}}\chi^{-1}\) with
\(\chi=\psi^k\bar\psi^j\). The Weil-pairing identity
\(\psi\bar\psi=\chi_{\mathrm{cyc}}\) identifies the exponents of
\(V_\chi\) with \((1-k,1-j)\). Hence the positive target
\((m_1,m_2)\) has \(k=1-m_1\leq0\), \(j=1-m_2\leq0\).
Its comparison is the noncritical one in
\cite[Corollary 2.4.7, pp.~9--10]{SanoVenjakob2026};
the source does not assert nonvanishing of every such regulator.
Theorem~\ref{thm:cm-finite-plus-abelian} gives the abelian component
fields through the exact Kummer pairing. Full nonabelian faithfulness
remains open, including for the two examples. A homogeneous Pollack
relation gives an element whose outer image lies deeper in the
weight filtration. An element of the group kernel requires a correction
inside the correspondingly deeper filtration of the source. The required
strictness of the filtered outer map is not proved here: rational
generation of each graded piece allows denominators that can grow
without bound during successive corrections.

\section{The degree-seven geometric detector}
\label{sec:geometric-detector}

We give the tensor calculation used in the geometric normalizer
argument, and identify the published representation-theoretic input.
Let $V_g$ be the rational symplectic homology of an oriented genus-$g$
surface with one boundary component. Write $\mathfrak h_{g,1}(k)$
for the degree-$k$ derivations of the free Lie algebra on $V_g$
which annihilate its symplectic boundary element, and let
$J_{g,k}$ be the rational image of the degree-$k$ Johnson homomorphism.
The symplectic identification of a derivation $D$ with its tensor is
\[
 T_D=\sum_{i=1}^g\bigl(x_i\otimes D(y_i)-y_i\otimes D(x_i)\bigr).
\]
Thus $T_D\in V_g^{\otimes(k+2)}$.
For a matching $C$ of positions $1,\ldots,8$, define
\[
 K_C(u_1\otimes\cdots\otimes u_9)
  =\left(\prod_{(i,j)\in C}\langle u_i,u_j\rangle\right)u_9.
\]
Each pairing is ordered with $i<j$.
Morita--Sakasai--Suzuki's five-term detector
\cite[Proposition~4.1, p.~11]{MSS2026} is
\begin{align*}
 R={}&-2K_{(12)(35)(47)(68)}
       +2K_{(12)(36)(47)(58)}\\
    &+2K_{(12)(36)(48)(57)}
       +2K_{(12)(37)(46)(58)}
       +K_{(13)(26)(47)(58)}.
\end{align*}

\begin{proposition}\label{prop:johnson-seven-detector}
For every $g\geq3$, $R(J_{g,7})=0$. The same identity holds
after scalar extension to every $\mathbb Q_\ell$.
\end{proposition}
\begin{proof}
We reconstruct the vanishing from six explicit brackets and the
older stable multiplicity bound in
\cite[proof of Theorem~7.5(iv), author pp.~29--30]{MSS2015}.
That bound says that the standard-representation isotypical
part of $J_{G,7}$ has multiplicity at most six for sufficiently
large $G$. The source obtains it from multiplicities $10,3$
in $\mathfrak h_{G,1}(7)$ and its boundary ideal, Enomoto--Satoh
ranks $3,2$, and the zero intersection in its
Proposition~3.5(ii), author pp.~16--17.

Take $G$ in that stable range, with $G\geq5$, and put
$e_0=x_1,e_1=y_1,\ldots,e_8=x_5,e_9=y_5$.
For an ordered triple $t=(i,j,k)$ set
\[
 \delta_t(z)=\langle z,e_i\rangle[e_j,e_k]
       +\langle z,e_j\rangle[e_k,e_i]
       +\langle z,e_k\rangle[e_i,e_j].
\]
The identification $\bigwedge^3V_G=J_{G,1}$ places these
vectors in the degree-one Johnson image. Use the following
six sequences of triples:
\[
\begin{array}{c|lllllll}
1&(3,7,5)&(1,2,0)&(9,8,6)&(6,8,5)&(2,4,7)&(9,0,3)&(4,6,7)\\
2&(9,5,0)&(6,4,7)&(6,2,0)&(6,5,8)&(7,3,8)&(1,4,3)&(9,2,7)\\
3&(1,3,4)&(4,2,0)&(5,7,0)&(9,2,7)&(5,6,8)&(5,3,0)&(1,4,6)\\
4&(0,1,4)&(3,5,0)&(9,0,5)&(2,3,9)&(4,6,8)&(8,1,0)&(2,7,1)\\
5&(3,7,0)&(4,9,6)&(3,5,2)&(7,3,8)&(6,7,2)&(0,8,9)&(1,2,6)\\
6&(5,4,0)&(0,8,1)&(5,2,3)&(7,4,3)&(0,4,1)&(6,5,2)&(6,7,9)
\end{array}
\]
For row $r$, define $Z_{r,1}=\delta_{t_1}$ and recursively
$Z_{r,j}=[Z_{r,j-1},\delta_{t_j}]$; put $Z_r=Z_{r,7}$.
The Johnson image is a Lie subalgebra, so $Z_r\in J_{G,7}$.
Every tensor term of $Z_r$ has torus weight $x_1$.

In order, take the six contraction matchings
\[
\begin{array}{lll}
(12)(34)(57)(68),&(12)(34)(58)(67),&(12)(35)(47)(68),\\
(12)(36)(47)(58),&(12)(36)(48)(57),&(12)(37)(46)(58).
\end{array}
\]
Their values on $T_{Z_r}$ are multiples of $x_1$, with coefficient
matrix
\[
 A=\begin{pmatrix}
18&12&-2&15&-7&-10\\
0&3&0&-4&4&0\\
15&-24&3&2&-5&6\\
-17&14&-10&11&-11&-4\\
-3&-10&2&-15&3&16\\
2&-16&-1&2&-3&0
\end{pmatrix},\qquad \det A=-203784.
\]
The five-term evaluations give $R(T_{Z_r})=0$ for all six rows.
All these integer entries are obtained by the following finite
recursion, which also specifies the calculation completely.
Represent a word $w=e_{i_1}\cdots e_{i_m}$ in the free associative
algebra, use $[v,w]=vw-wv$, and extend a derivation by
\[
 D(w)=\sum_{j=1}^m e_{i_1}\cdots e_{i_{j-1}}
                    D(e_{i_j})e_{i_{j+1}}\cdots e_{i_m}.
\]
Apply this identity to $[D,E](e_i)=D(E(e_i))-E(D(e_i))$ at each
of the six bracket steps. Form $T_{Z_r}$ by the displayed tensor
identification and apply the four scalar pairings for each matching.
The nonzero tensor supports in the six rows have respectively
$22896,43902,57276,33768,26082,27990$ words. The recurrence gives
$Z_r(\sum_i[x_i,y_i])=0$, the stated weights, the matrix $A$ and
the six zero $R$-values. The program \texttt{check\_johnson\_detector.py}
in the \href{\workbenchurl}{workbench} executes this recursion over
$\mathbb Z$ and records the full verification.

Let $e_V$ be the symplectic-equivariant projection onto the
standard-representation isotypical summand. It preserves $J_{G,7}$.
The contractions and $R$ have target $V_G$, so they kill every other
irreducible summand. The nonzero determinant therefore makes
$e_VZ_1,\ldots,e_VZ_6$ independent in its weight-$x_1$ space.
The published upper bound makes them a basis of that space.
Their symplectic orbits span the whole standard isotypical part.
Equivariance and the six zero values now show that $R$ kills this
part, and its target kills all remaining isotypical parts.
Thus $R(J_{G,7})=0$.

For any given $g\geq3$, choose such a $G\geq g$ and extend a
boundary mapping class by the identity on the added handles.
Its Johnson derivation extends by zero on those handles.
For the symplectic inclusion $i:V_g\hookrightarrow V_G$, the
literal contraction formulas give
\[
 R_G(i^{\otimes9}T)=iR_g(T).
\]
The left side vanishes on the extended Johnson image and $i$ is
injective, proving the assertion at genus $g$. Every step is rational;
scalar extension proves the statement over $\mathbb Q_\ell$, including
primes dividing the displayed integer determinant.
\end{proof}

\subsection{The line-tangency calculation}
\label{subsec:geometric-polynomial}

We now compute the invariant obstruction to a degree-six
arithmetic derivation preserving one nonzero homology line.
All vector spaces and maps in this subsection are over a
field of characteristic zero. The formulas are rational,
so they apply over every $\mathbb Q_\ell$.
Write $V=V_g$, $\mathbb L(V)$ for its free Lie algebra, and
$e_{\operatorname{Sp}}$ for the canonical projection onto
the symplectic-invariant part of a finite-dimensional
rational symplectic representation.

\begin{lemma}\label{lem:geometric-line-tangency}
Let $g\geq10$, let $0\ne a\in V$, and let
$D\in\mathfrak h_{g,1}(6)$ satisfy
\[
 D(a)\in[\mathbb L_6(V),a],\qquad [D,J_{g,1}]\subseteq J_{g,7}.
\]
Then $e_{\operatorname{Sp}}D\in J_{g,6}$.
Consequently every symplectically trivial quotient of a
symplectic submodule containing $D$ and annihilating its
intersection with $J_{g,6}$ sends $D$ to zero.
\end{lemma}

\begin{proof}
Choose $b$ with $\langle a,b\rangle=1$ and put
$W=\langle a,b\rangle^\perp$, $s=g-1$.
The reductive group $S=\operatorname{Sp}(W)\times\mathbb G_m$
acts by $t,t^{-1}$ on $a,b$. Both displayed conditions are
$S$-stable linear conditions, so its invariant projector
$e_S$ preserves them. Moreover
$e_{\operatorname{Sp}}e_S=e_{\operatorname{Sp}}$.
We calculate on the space of $S$-invariant line-tangent
derivations.

\paragraph{The complete invariant input.}
Regard the tensor $T_D$ as a word of length eight.
Let $\rho$ rotate its first letter to the last position,
let $S_8=1+\rho+\cdots+\rho^7$, and let
\[
 \mathcal D_n(v_1\cdots v_n)
       =[\cdots[[v_1,v_2],v_3],\ldots,v_n]
\]
be the left-normed Dynkin map. The conditions that $T_D$
be a special derivation are exactly
\[
 (1-\rho)T_D=0,\qquad
 (1\otimes\mathcal D_7)T_D=7T_D.
\]
Indeed the first equality is the kernel of the bracket
$V\otimes\mathbb L_7(V)\to\mathbb L_8(V)$, and the
second says that the last seven letters lie in the free
Lie algebra. The identity $\mathcal D_n^2=n\mathcal D_n$
gives the latter equivalence.

An $S$-invariant word contains the same number $m$ of
$a$'s and $b$'s, where $0\leq m\leq4$.
For $m>0$, its $b$-initial part is
$-b\otimes D(a)=-b\otimes[A,a]$ for a unique
$A\in\mathbb L_6(V)$ with $m-1$ copies each of $a,b$.
Uniqueness follows because the centralizer of $a$ in
degree six is zero. Summing the $b$-initial part over the
eight rotations counts every word $m$ times. Hence
\[
 T_D=\Psi_m(A):=-\frac1m S_8(b\otimes[A,a]).
\]
Conversely this formula gives the required derivation
exactly when
\begin{equation}\label{eq:line-input-constraints}
 (1\otimes\mathcal D_7)\Psi_m(A)=7\Psi_m(A),\qquad
 \operatorname{pr}_b\Psi_m(A)=-b\otimes[A,a].
\end{equation}
For $m=0$, every cyclic invariant tensor satisfying the
Dynkin condition has $D(a)=0$ and is allowed.

Here is a finite construction of every input in these
spaces. Place the indicated copies of $a,b$ in the six
slots for $A$, and match the remaining slots in all ways.
Each ordered pair represents
$\omega_W=\sum_{i=2}^g(x_i\otimes y_i-y_i\otimes x_i)$.
Apply $\mathcal D_6$, choose the independent columns of
this integer matrix, and impose
\eqref{eq:line-input-constraints} on their coefficients.
For $m=0$, start instead with all matchings of eight slots,
apply $S_8$, choose independent columns, and impose the
Dynkin condition. The symplectic invariant tensor theorem
\cite[Proposition~2.1, author p.~8]{MSS2015}
makes these lists spanning. They are linearly independent
as raw decorated diagrams in the range in use, since
$s\geq9$ and there are at most four pairs.
The exact row reductions have the following sizes:
\[
\begin{array}{c|rrrr}
m&\text{raw diagrams}&\text{parameter rank}&
       \text{constraint rank}&\text{kernel dimension}\\\hline
1&15&3&0&3\\
2&90&15&12&3\\
3&90&16&15&1\\
4&20&3&3&0\\
0&105&17&12&5
\end{array}
\]
Thus the full invariant line-tangent space has dimension
twelve. All entries of the reduction are given by the
two displayed word operations; no geometric relations
are imposed in this calculation.

\paragraph{The seven necessary scalar equations.}
Take two further symplectic pairs $(c,d),(e,f)$ in $W$.
For three vectors use the degree-one derivation
\[
 \delta_{uvw}(z)=\langle z,u\rangle[v,w]
          +\langle z,v\rangle[w,u]
          +\langle z,w\rangle[u,v].
\]
Its tensor is $-\operatorname{Alt}(u,v,w)$.
For a degree-seven special derivation $E$, write
$r_t(E)=\langle R(T_E),t\rangle$, with $R$ as in
Proposition~\ref{prop:johnson-seven-detector}.
The Enomoto--Satoh map first contracts the first two
slots of $T_E$ and takes the remaining seven slots
modulo cyclic rotation. It annihilates $J_{g,7}$
\cite[diagram (2), p.~640; \S7.1, p.~656]{EnomotoSatoh2014}.
For a matching $C$ of positions $1,\ldots,6$, define
on seven-letter cyclic words the scalar functional
\[
 \lambda_{C,t}(v_1\cdots v_7)
  =\sum_{j=0}^6
       \left(\prod_{(i,k)\in C}
          \langle v_{i+j},v_{k+j}\rangle\right)
                 \langle v_{7+j},t\rangle,
\]
where subscripts are read modulo seven. This is
well-defined on the cyclic quotient. Let $c_{C,t}(E)$
be its value after the Enomoto--Satoh contraction.
Put $C_1=(12)(35)(46)$ and $C_2=(12)(36)(45)$.
Define the seven-component column $M_7(D)$, in order, by
\[
\begin{pmatrix}
r_d([D,\delta_{abc}])\\
r_b([D,\delta_{acd}])\\
r_f([D,\delta_{cde}])\\
c_{C_1,d}([D,\delta_{abc}])\\
c_{C_2,d}([D,\delta_{abc}])\\
c_{C_1,b}([D,\delta_{acd}])\\
c_{C_1,f}([D,\delta_{cde}])
\end{pmatrix}.
\]
Every $\delta$ here belongs to $J_{g,1}=\bigwedge^3V$.
The second hypothesis of the lemma and the two detector
vanishings therefore give $M_7(e_SD)=0$.

\paragraph{The invariant quotient in degree six.}
Let $K(D)$ be the five scalar contractions of $T_D$
with the following ordered matchings:
\[
\begin{gathered}
(12)(34)(56)(78),\quad(12)(34)(57)(68),\quad
(12)(34)(58)(67),\\
(12)(36)(47)(58),\quad(12)(38)(46)(57).
\end{gathered}
\]
Morita--Sakasai--Suzuki's invariant calculation
\cite[Proposition~7.2, author pp.~26--27]{MSS2015}
says that $K$ restricts to an isomorphism
$\mathfrak h_{g,1}(6)^{\operatorname{Sp}}\to\mathbb Q^5$.
Their contraction convention differs by
$\operatorname{diag}(1,-1,1,-1,-1)$ from this ordered
pairing convention. Their two Johnson vectors have
coordinates $(0,0,0,-2,2)$ and $(4,-1,-4,4,-4)$
in their convention
\cite[proof of Theorem~7.5(i), author pp.~28--29]{MSS2015}.
Thus in the present convention their span is the kernel of
\[
 L=\begin{pmatrix}-1&4&0&0&0\\1&0&1&0&0\\0&0&0&1&1\end{pmatrix}.
\]
This identifies $K(J_{g,6}^{\operatorname{Sp}})=\ker L$.
Since each contraction is symplectically invariant,
$K(D)=K(e_{\operatorname{Sp}}D)$.

\paragraph{The polynomial identity.}
On all twelve inputs above, the literal contraction maps
satisfy
\begin{equation}\label{eq:geometric-polynomial-identity}
 N_0(s)M_7(s)=d_0(s)LK(s),\qquad
 d_0(s)=9(2s+3)(4s^2+6s-1).
\end{equation}
Here the entries of the $3$ by $7$ matrix $N_0$ are as follows:
\begin{align*}
(N_0)_{11}&=-2s(2s+3)(128s^4+528s^3+776s^2+465s+101),\\
(N_0)_{12}&=-4s(2s+3)(16s^4-48s^3-305s^2-402s-152),\\
(N_0)_{13}&=2s(s-1)(2s+3)(128s^3+464s^2+550s+199),\\
(N_0)_{14}&=-3s(576s^4+2720s^3+4684s^2+3546s+1029),\\
(N_0)_{15}&=-6s(s+2)(4s+3)(4s+5)(8s+7),\\
(N_0)_{16}&=3s(2s+1)(2s+3)(80s^2+240s+187),\\
(N_0)_{17}&=24s(s-1)(2s+3)(4s^2+6s-1).
\end{align*}
\begin{align*}
(N_0)_{21}&=2s(2s+1)^2(2s+3)(8s^2+25s+20),\\
(N_0)_{22}&=2s(2s+1)(2s+3)(4s^3-14s^2-70s-61),\\
(N_0)_{23}&=-4s(s-1)(2s+1)(2s+3)(8s^2+25s+20),\\
(N_0)_{24}&=3s(2s+3)(72s^3+232s^2+224s+69),\\
(N_0)_{25}&=6s(s+2)(2s+1)(2s+3)(8s+7),\\
(N_0)_{26}&=-3s(2s+3)(40s^3+140s^2+146s+37),\\
(N_0)_{27}&=-6s(s-1)(2s+3)(4s^2+6s-1).
\end{align*}
\begin{align*}
(N_0)_{31}&=-6s(2s+3)(16s^3+66s^2+80s+27),\\
(N_0)_{32}&=-6s(2s+3)(4s^3-12s^2-82s-81),\\
(N_0)_{33}&=6s(s-1)(2s+3)(16s^2+60s+53),\\
(N_0)_{34}&=-18s(s+2)(36s^2+102s+67),\\
(N_0)_{35}&=-36s(s+2)^2(8s+7),\\
(N_0)_{36}&=18s(s+2)(2s+3)(10s+13),\\
(N_0)_{37}&=0.
\end{align*}

We specify how the polynomial equality is evaluated.
Expand every Lie bracket by $uv-vu$, and every composition
of derivations by substitution in each letter. In the
symplectic tensor identification a substitution contracts
the output's first slot with the substituted slot. The
formula $T_{\delta_{uvw}}=-\operatorname{Alt}(u,v,w)$
fixes its sign. Thus all columns of $M_7$ and $K$ are
finite sums of fully paired decorated diagrams.
For each of the at most four copies of $\omega_W$, choose
one of its two orientations, recording its sign.
Each scalar pairing is zero unless the two letters have
opposite types $x,y$ and the same handle index; a pairing
in the order $y,x$ contributes a further minus sign.
Propagate these handle equalities. A $W$-index cannot
equal the handle of $a,b$, and the fixed handles of
$(a,b),(c,d),(e,f)$ are distinct. Each remaining free
handle has exactly $s$ choices, independently of all
other free handles. Its contribution is therefore $s$.
Summing the signed terms gives a polynomial in $s$,
including the signs and factors from closed contraction
cycles. Distinct cycles are allowed to choose the same
handle; no injectivity condition is inserted.

The row reductions in the table and this orientation
sum compute every entry of the two sides of
\eqref{eq:geometric-polynomial-identity} by rational
arithmetic. Subtraction gives three zero polynomials
for each of the twelve columns. The five pure-pair
columns of $K$ have determinant, up to the orientation
of the chosen kernel basis,
\[
 3538944s^5(s-2)(s-1)^4(s+1)^3
             (2s-1)(2s+1)^5(2s+3).
\]
This also checks their independence in the range used.
For a reproducible choice of columns, order the positions
of the $a$'s and $b$'s lexicographically and enumerate
matchings by pairing the first unused slot with each
later slot in increasing order. Use pivot columns in
reduced row echelon form for each parameter basis and
the associated free-coordinate basis for each kernel;
clear denominators and make the first nonzero coefficient
positive. The program \texttt{verify\_single\_line.py},
cited through the \href{\workbenchurl}{workbench},
performs exactly these word operations and orientation
sums, and gives all thirty-six zero differences.

Since $s\geq9$, $d_0(s)$ is nonzero. For example
$d_0(3+t)=9(2t+9)(4t^2+30t+53)>0$ for $t\geq0$.
The equality $M_7(e_SD)=0$ and
\eqref{eq:geometric-polynomial-identity} imply
$LK(e_SD)=0$. Hence
$K(e_{\operatorname{Sp}}D)\in K(J_{g,6}^{\operatorname{Sp}})$.
Injectivity of $K$ on symplectic invariants proves
$e_{\operatorname{Sp}}D\in J_{g,6}$.
Every symplectically trivial quotient has
$q(D)=q(e_{\operatorname{Sp}}D)$, giving the last assertion.
\end{proof}

\subsection{The geometric recognition maps}
\label{subsec:geometric-recognition-maps}

This subsection composes the line-tangency lemma with the
covering and monodromy maps for the punctured elliptic curve.  We arrange that the
detecting compactification has genus at least ten, so the published
stable-range inputs apply directly.
Let \(X_0=E_0-\{O\}\), where \(E_0:y^2=x^3-x\) over \(\mathbb Q\),
and choose a nonzero rational tangent vector at \(O\).
Write
\[
 F=\pi_1(X_{0,\overline{\mathbb Q}})
       =\widehat{\langle x,y\rangle},\qquad
 \Pi=F\rtimes s_O(G_{\mathbb Q}).
\]
The cusp inertia \(I_O\) is generated by \([x,y]\), up to inversion.
All subgroups and commutator closures below are closed.

\paragraph{A fixed geometric cover in the stable range.}
Consider the explicit surjection
\[
 F\longrightarrow \operatorname{UT}_3(\mathbb F_3),\qquad
 x\longmapsto1+E_{12},\qquad y\longmapsto1+E_{23}.
\]
The two matrices and their commutator generate the group of order
\(27\), and the cusp generator maps to \(1+E_{13}\), of order three.
Let \(F_{10}\) be its kernel and \(Y_{10}\to X_{0,\overline{\mathbb Q}}\)
the corresponding connected finite \'etale cover.
Its compactification \(C_{10}\) has nine points over \(O\),
each of ramification index three.  Riemann--Hurwitz gives
\[
 2g(C_{10})-2=27(2g(E_0)-2)+9(3-1)=18.
\]
Thus \(g(C_{10})=10\).
The cover and its morphism descend to a finite number field,
since they are given by finitely many algebraic equations.
Only this actual geometric covering subgroup will be used.

\begin{theorem}\label{thm:geometric-procyclic-recognition}
Every nontrivial closed procyclic subgroup of \(F\) whose
arithmetic normalizer has open image in \(G_{\mathbb Q(i)}\)
is contained in a conjugate of \(I_O\).
The same conclusion holds after every finite arithmetic base
extension.  Conversely, every closed subgroup of a cusp inertia
group has an arithmetic normalizer with open image.
\end{theorem}

\begin{proof}
Let \(A=\overline{\langle z\rangle}\neq1\) and let \(N\) be its
arithmetic normalizer.  Its image is open in \(G_{\mathbb Q}\),
because \(G_{\mathbb Q(i)}\) is open.
Suppose that \(z\) is outside the union \(\mathcal Z\) of the
conjugates of cusp inertia.
Nakamura's full-class proper-homology detector
\cite[Lemma 2.1.4, pp.~96--97]{Nakamura1994}, applied before any
pro-prime quotient, gives a prime \(\ell\) and an open
\(H\subset F\) containing \(z\) such that
\begin{equation}\label{eq:geo-proper-detector}
 z\notin[H,H]H^\ell\langle\mathcal Z\cap H\rangle.
\end{equation}
Let \(C_H\) compactify the finite geometric cover defined by \(H\).
Killing the boundary inertia is the actual filling homomorphism
\(H\twoheadrightarrow\pi_1(C_H)\), so the quotient in
\eqref{eq:geo-proper-detector} is
\(H_1(C_H,\mathbb F_\ell)\).
Consequently the specified element \(z\) has a nonzero class
in the free \(\mathbb Z_\ell\)-module \(H_1(C_H,\mathbb Z_\ell)\).

The homomorphism to this module kills the primary factors
of \(A\) at primes different from \(\ell\).
If \(u\) is the \(\ell\)-primary component of \(z\), then
\(A_\ell=\overline{\langle u\rangle}\cong\mathbb Z_\ell\)
and \([u]\neq0\) in \(H_1(C_H,\mathbb Q_\ell)\).
This primary factor is characteristic in \(A\), hence is
normalized by \(N\).

Form the full arithmetic core
\begin{equation}\label{eq:geo-arithmetic-core}
 S=\bigcap_{\gamma\in\Pi}\gamma(H\cap F_{10})\gamma^{-1}.
\end{equation}
It is open in \(F\): a finitely generated profinite group
has only finitely many subgroups of any specified finite index,
as follows from their finite permutation representations.
Thus the intersection in \eqref{eq:geo-arithmetic-core} is finite.
It is normal in \(\Pi\) and contained in \(H\cap F_{10}\).
For some \(r\geq0\), the element \(u'=u^{\ell^r}\) generates
\(A_\ell\cap S\).
The inclusions of covering groups and the associated maps of
proper compactifications give the exact pushforward
\begin{equation}\label{eq:geo-pushforward-class}
 H_1(C_S,\mathbb Q_\ell)\longrightarrow H_1(C_H,\mathbb Q_\ell),
 \qquad [u']\longmapsto\ell^r[u]\neq0.
\end{equation}
The map \(C_S\to C_{10}\) is finite and nonconstant.
Riemann--Hurwitz therefore gives \(g(C_S)\geq10\).

The arithmetic subgroup
\[
                  \Pi_Y=S\rtimes s_O(G_{\mathbb Q})
\]
defines an actual geometrically connected cover \(Y/\mathbb Q\)
of \(X_0\).  At the chosen branch over \(O\), its decomposition
subgroup is exactly
\[
 \Pi_Y\cap\bigl(I_O\rtimes s_O(G_{\mathbb Q})\bigr)
               =(S\cap I_O)\rtimes s_O(G_{\mathbb Q}).
\]
Its residue Galois image is all of \(G_{\mathbb Q}\).
Hence that branch is a rational point \(y\) on the smooth
proper compactification \(C=\overline Y\).
This is an explicit descent through an arithmetic covering
subgroup; it does not use a field-of-moduli assertion.

The open immersion
\[
                  Y\longrightarrow X=C-\{y\}
\]
fills the other punctures.  It induces a surjection of geometric
fundamental groups and the identity on the arithmetic quotient.
On proper homology it is the identity.
Consequently the image of \(u'\) in
\[
                         P=\pi_1(X_{\overline{\mathbb Q}})^{(\ell)}
\]
has a nonzero class \(a\) in
\(P^{\mathrm{ab}}\otimes\mathbb Q_\ell
 =H_1(C,\mathbb Q_\ell)=:V\).
In particular the image of \(A_\ell\cap S\) is still
isomorphic to \(\mathbb Z_\ell\): a nonzero closed kernel
would give a finite image, whose rational abelian class is zero.
The image of \(N\cap\Pi_Y\) is a closed normalizer with the
same open arithmetic image.
Choose a finite number field \(L\) such that this image
contains \(G_L\).

\paragraph{The actual degree-six arithmetic quotient.}
For the same curve \(X\), Nakamura
\cite[Corollary 4.2 and the table in 4.3, p.~115;
Lemma 4.5, p.~116; (4.7.1)--(4.7.2), p.~117]{Nakamura1996}
provides the following exact maps.
The universal outer arithmetic degree-six image \(U_{g,6}\)
has a quotient
\[
                  q_{g,6}:U_{g,6}\twoheadrightarrow Q_{g,6},
 \qquad \dim_{\mathbb Q_\ell}Q_{g,6}=1.
\]
The action on the quotient is the cube of the symplectic
similitude character.
The actual degree-six graded Galois image of \(X/L\)
surjects onto this quotient.
The degree-six dimension follows from the rank sandwich
whose two endpoints are both one; the source also states
finite-base invariance immediately before Lemma 4.5.
After clearing a rational denominator, there is therefore
an actual
\[
 \sigma\in G_L(6),\qquad D_\sigma\in U_{g,6},\qquad
                              q_{g,6}(D_\sigma)\neq0.
\]
Here \(D_\sigma\) is the canonical special representative
of its outer graded class.

Lift this same \(\sigma\) to the normalizer of the image
of \(A_\ell\cap S\).  Since \(\sigma\) acts trivially on \(V\),
its normalizer exponent \(\lambda\in\mathbb Z_\ell^\times\)
satisfies \(\lambda a=a\), hence \(\lambda=1\).
Write \(U=\log u'=a+\text{terms of degree at least two}\)
in the rational completed free Lie algebra.
The chosen normalizing automorphism and the canonical
representative differ by an inner automorphism, so
\[
                 \phi_\sigma(U)=\exp(\operatorname{ad}Z)U.
\]
Its degree-one conjugator term is a multiple of \(a\).
Multiplying the conjugator by the inverse of the appropriate
rational power of \(\exp U\) removes that term and preserves
its action on \(U\).
The degree-two through degree-five terms then centralize \(a\)
and vanish, because the free-Lie centralizer of a nonzero
free generator is its one-dimensional span.
Taking degree seven gives the exact equality
\begin{equation}\label{eq:geo-normalizer-tangency}
                         D_\sigma(a)=[Z_6,a].
\end{equation}
Using the full \(U\), rather than only its linear term,
is necessary for this conjugator correction.

\paragraph{The geometric image in the certificate.}
We verify the monodromy comparison needed to apply the
polynomial identity to this \(D_\sigma\).
Use the moduli stack with one marked point and a nonzero
tangent vector.  Its geometric mapping class group fixes
the boundary word and has Zariski closure \(R_g\) in
the automorphisms of the rational punctured Malcev group.
At each finite nilpotent level, the actual profinite
geometric image has the same Zariski closure after
extension to \(\mathbb Q_\ell\): the discrete image is
dense and the defining polynomials are continuous.

Hain's Torelli-to-relative-completion surjection
\cite[Theorem 3.4, p.~605]{Hain1997} identifies the
prounipotent radical mapping onto \((R_g)_u\).
His weight and boundary-action comparison
\cite[Theorem 4.10, p.~611; Remark 4.11, p.~613;
\S14.4, p.~643]{Hain1997} is a morphism of mixed
Hodge structures and is strict.
Consequently its graded image is generated by its
degree-one image.  Enomoto--Satoh
\cite[Theorem 3.5 and Proposition 3.6, p.~640]{EnomotoSatoh2014}
identifies that generated image with the rational Johnson
image \(J_m\).  In particular,
\[
              \operatorname{gr}_{-m}\operatorname{Lie}(R_g)_u=J_m,
              \qquad J_1=\bigwedge^3V.
\]
This proves the equality for the image filtration; density
alone is not used to identify individual filtration kernels.
In degree six, the special-to-outer comparison is injective:
an inner derivation preserving the symplectic boundary class
has generator in its free-Lie centralizer, which is the
one-dimensional span of that degree-two class.
Thus the geometric kernel in Nakamura's exact sequence is
the same \(J_6\), and \(q_{g,6}\) kills the actual Johnson
brackets used in the certificate.

Arithmetic monodromy normalizes the geometric algebraic image.
For an outer degree-six class its tangential representative
can have a lower inner term only from the degree-two
boundary class.  Subtracting this term by
\(\exp(-c\,\operatorname{ad}\log(\partial))\) stays inside
the geometric algebraic image: the boundary Dehn twist
generates its one-parameter Zariski closure.
The centralizer of the boundary class removes the
remaining lower inner terms.
The resulting canonical special representative therefore
still normalizes \(R_g\) to the required order.
Taking the leading term and using strictness gives
\[
                         [D_\sigma,J_1]\subseteq J_7.
\]

Thus the Enomoto--Satoh cyclic contraction and the additional
five-term contraction verified by the six Johnson vectors
both vanish on these brackets.  Their seven scalar rows
are precisely the necessary rows of the polynomial certificate.
All stable-range requirements hold because \(g\geq10\).
Choose \(\langle a,b\rangle=1\), set
\(W=\langle a,b\rangle^\perp\), and let
\(\operatorname{Sp}(W)\times\mathbb G_m\) act by its standard
action on \(W\) and by \(t,t^{-1}\) on \(a,b\).
Its reductive invariant projector preserves both the actual universal
image and the tangency condition
\eqref{eq:geo-normalizer-tangency}; it preserves \(q_{g,6}\),
because that quotient is symplectically trivial.
Lemma~\ref{lem:geometric-line-tangency} therefore proves
\[
 D\in U_{g,6},\quad D(a)\in[\mathbb L_6(V),a]
                       \quad\Longrightarrow\quad q_{g,6}(D)=0.
\]
Applied to \(D_\sigma\), this contradicts its nonzero quotient.
Hence \(z\in\mathcal Z\), and its procyclic closure lies
in that same cusp inertia group.

The argument depends only on an open arithmetic image.
Starting after a finite arithmetic base extension gives
the same geometric group and the same inertia conjugates;
the normalizer image remains open in \(G_{\mathbb Q}\).
Finally, a cusp decomposition group over a field defining
the cusp acts on its procyclic inertia by the cyclotomic
character.  Every closed subgroup of a procyclic group
is characteristic.  That decomposition group therefore
normalizes it and has open arithmetic image, proving
the converse.
\end{proof}

\begin{corollary}\label{cor:geometric-common-kernel-recognition}
Put
\[
 R_{\mathrm{geom}}=
 \bigcap_\ell\ker\!\left(F\longrightarrow F^{(\ell)}\right).
\]
No nontrivial closed procyclic subgroup of \(R_{\mathrm{geom}}\)
has an arithmetic normalizer with open image.
\end{corollary}

\begin{proof}
The intersection of each cusp inertia group with
\(R_{\mathrm{geom}}\) is trivial by an explicit finite quotient.
For every \(n\geq1\), the substitutions
\[
 x\longmapsto1+E_{12},\qquad y\longmapsto1+E_{23}
       \quad\text{in }\operatorname{UT}_3(\mathbb Z/\ell^n\mathbb Z)
\]
send the cusp commutator to \(1+E_{13}\), of exact order
\(\ell^n\).
Thus if a cusp element has exponent
\(\alpha=(\alpha_\ell)_\ell\in\widehat{\mathbb Z}\),
its pro-\(\ell\) image is nontrivial whenever
\(\alpha_\ell\neq0\).
Every nonzero exponent has such a component.
The equality persists under conjugation, since
\(R_{\mathrm{geom}}\) is normal.
The theorem now gives the corollary.
\end{proof}

The arithmetic outer kernel is related to the centralizer in the
fundamental exact sequence by the isomorphism proved in
Subsection~\ref{subsec:cm-geometric-common-kernel}.
The exact pullback calculation there identifies the extra
morphism needed to apply this corollary to that kernel.


\section{What remains open}

For the cyclotomic question, Theorem~\ref{thm:canonical-all-stages} and Proposition~\ref{prop:limit-relation-module} give the exact remaining vanishing problem
\[
 \ker\rho_{\mathrm{outer}}=1
 \quad\Longleftrightarrow\quad
 \varprojlim_n
 \frac{(\mathcal R_n)_{\Lambda\text{-tor}}}
      {(p,\gamma-1)(\mathcal R_n)_{\Lambda\text{-tor}}}=0.
\]
The maps are induced by the inclusions of the recursively defined source subgroups. The kernel of the map from their integral relation-module limit to $M$ is the continuous homology group of the compact outer image, as identified above. The separate finite-stage $p$-bounds leave these transition maps unevaluated. Proposition~\ref{prop:source-limit-countermodel} computes a nonzero limit in a free pro-$p$ model with the regular odd module and the same positive-eigenvalue exclusion. Triviality for the arithmetic source at every prime remains open.

Theorem~\ref{thm:selected-all-stages} constructs the detected component fields at every finite stage. For the lift-defined quotient in the original component question, the calculation at $37$ proves failure of Galois invariance for every pure even lift, and its cyclotomic core calculation identifies the field-containment claim with full faithfulness at that prime. Proposition~\ref{prop:D-plus-finiteness} identifies the vanishing of the odd $\Lambda$-torsion exactly with finiteness of $X^+$. The general single-component hypotheses have been used in the proved canonical quotient through their cup comparison.

For CM curves under finite $\mathcal X_2^+$, the next class-two question is the kernel of
\[
 J_A=\mathcal J\circ(\alpha\widehat\otimes\alpha)
           \bigm|_{\widehat{\bigwedge}^{\,2}A}.
\]
Its kernel has no nonzero decomposable tensor, and it meets the exterior square of every subspace of dimension at most three trivially. The countably generated subgroups constructed above also have their complete class-two quotient detected. Lemma~\ref{lem:cm-measure-euler-image} identifies the mixed-chart image and its exact Euler relation; Theorem~\ref{thm:cm-elliptic-euler-kernel} further identifies
\(\ker J_A=\varprojlim_n\mathcal D_n\), where
\(\mathcal D_n[1/p]=\operatorname{Hom}_{\mathbb Q_p}
(\mathcal R_n^{\rm ell},\mathbb Q_p(2))\).
The cokernel transitions are induced by the exact isogeny distributions.
The remaining issue is the persistence of these triangle-relation
cokernels together with boundedness of compatible integral duals. Full nonabelian faithfulness further requires control in every depth. The Pollack relation gives a kernel in the associated graded derivation map for $w_K=2$; a group-kernel element requires the explicit deeper-source correction described in Section~\ref{sec:cm-audit}. The geometric normalizer theorem reaches the arithmetic outer kernel through the centralizer pullback proved there, whose remaining centralizer is precisely that kernel.

\section*{Method and records}
The mathematics and text were produced by GPT-6 Astra (ChatGPT Pro mode). Independent checks in the accompanying records are re-proofs by the same model; effort settings were inherited and are not externally specified. No human has checked this text. The accompanying workbench is prepared for \url{\workbenchurl}; its permanent record address is \url{\recordurl}, pending the assignment of a paper-specific record page. This text is dedicated to the public domain under CC0 1.0.

\begin{thebibliography}{99}
\bibitem{AI}G.~Anderson and Y.~Ihara, Pro-$l$ branched coverings of $\mathbb P^1$ and higher circular $l$-units, Ann. of Math. \textbf{128} (1988), 271--293. \url{https://doi.org/10.2307/1971443}.
\bibitem{Brown}F.~Brown, Mixed Tate motives over $\Z$, Ann. of Math. \textbf{175} (2012), 949--976, Theorem~1.1. \url{https://arxiv.org/abs/1102.1312}.
\bibitem{Sharifi}R.~Sharifi, Relationships between conjectures on the structure of pro-$p$ Galois groups unramified outside $p$, Proc. Sympos. Pure Math. \textbf{70} (2002), 275--284. Locators refer to arXiv:math/0104116v2, Theorem~1, p.~2. \url{https://arxiv.org/abs/math/0104116v2}.
\bibitem{Note}The question of Anderson and Ihara on the pro-$p$ Galois action on the projective line minus three points, at every prime satisfying Vandiver's conjecture, Hexagon 2611.00939v2, 24 November 2026, Theorem~1.1(a). Unrefereed. \url{https://hexagonmath.org/2611.00939}.
\bibitem{Ishii}S.~Ishii, On the kernels of the pro-$p$ outer Galois representations associated to once-punctured CM elliptic curves, arXiv:2312.04196v3, 6 January 2026, Theorem~A, p.~2, and Theorem~2.14, p.~11. \url{https://arxiv.org/abs/2312.04196v3}.
\bibitem{FW}
B.~Ferrero and L.~C.~Washington,
\emph{The Iwasawa invariant $\mu_p$ vanishes for abelian number fields},
Ann. of Math. (2) \textbf{109} (1979), 377--395.
\url{https://doi.org/10.2307/1971116}.

\bibitem{FK}T.~Fukaya and K.~Kato,
\emph{On conjectures of Sharifi}, author preprint,
15 August 2012, 121 pages.
\url{https://www.math.ucla.edu/~sharifi/sharificonj.pdf}.
Locators in this author version:
\S\S6.1.4--6.1.8, pp.~76--77;
\S7.1.3, pp.~84--85;
Conjecture~7.1.8 and \S\S7.1.9--7.1.10, p.~86;
Theorem~7.2.8 and \S\S7.2.10--7.2.12, p.~89;
Lemma~10.2.2 and \S10.2.3, p.~107.

\bibitem{HM}R.~Hain and M.~Matsumoto, \emph{Weighted completion of Galois groups
and Galois actions on the fundamental group of $\mathbb P^1-\{0,1,\infty\}$},
arXiv:math/0006158v2, 11 August 2001.
\url{https://arxiv.org/pdf/math/0006158}.

\bibitem{HHO}W. B. Hart, D. Harvey and W. Ong,
Irregular primes to two billion,
\emph{Math. Comp.} 86 (2017), 3031--3049.
\url{https://doi.org/10.1090/mcom/3211}.
Author preprint:
\url{https://arxiv.org/abs/1605.02398}.
Section 1, item (2), verifies Vandiver for every odd prime
\(\ell<2^{31}\).

\bibitem{Iwasawa}K.~Iwasawa, \emph{Riemann--Hurwitz formula and $p$-adic Galois
representations for number fields},
Tohoku Math. J. \textbf{33} (1981), 263--288.
\url{https://www.jstage.jst.go.jp/article/tmj1949/33/2/33_2_263/_pdf/-char/en}.

\bibitem{Kurihara}M.~Kurihara, \emph{Some remarks on conjectures about cyclotomic fields and
$K$-groups of $\mathbb Z$}, Compositio Math. \textbf{81} (1992), 223--236.
Cited through \cite[Chapter~VI, Theorem~10.9]{Weibel}.

\bibitem{MilneEtale}J.~S.~Milne,
\emph{Lectures on \textup{\'Etale} Cohomology},
version~2.21, 22 March 2013, 202 pages.
\url{https://www.jmilne.org/math/CourseNotes/LEC.pdf}.
Theorem~14.9, Remark~14.10, Example~14.11 and Proposition~14.13,
printed pp.~96--98.

\bibitem{Ohta2020}M.~Ohta, \emph{$\mu$-type subgroups of $J_1(N)$
and application to cyclotomic fields}, J.~Math.~Soc.~Japan
\textbf{72} (2020), 333--412.
\url{https://doi.org/10.2969/jmsj/78327832}.
Part~I, Proposition~(7.2.4)(iv), p.~375, and pp.~378--381;
Part~II, (6.1.26), p.~408, and Theorem~(6.2.2), p.~409.

\bibitem{OhtaDensity}M.~Ohta, \emph{Zariski density of CM points
on products of modular curves}, J.~Math.~Soc.~Japan
\textbf{78} (2026), 1159--1202.
\url{https://doi.org/10.2969/jmsj/95469546}.
Author manuscript: \url{https://math.or.jp/publication/JMSJ/pdf/JMSJ9546.pdf}.
Theorem~(2.3.4), author pp.~15--16.

\bibitem{SharifiProgram}R.~Sharifi, \emph{Through the lens of the
Eisenstein ideal}, author preprint, retrieved 25 November 2026.
\url{https://math.ucla.edu/~sharifi/publ/program.pdf}.
Arithmetic action convention, p.~3; Ohta's isomorphism and
Theorem~8, pp.~13--14.

\bibitem{SharifiFK}R.~Sharifi, \emph{An extension of the
Fukaya--Kato method}, J.~Eur.~Math.~Soc. \textbf{28} (2026),
2877--2943. \url{https://doi.org/10.4171/JEMS/1519}.
Theorem~5.21.

\bibitem{MSS2015}S.~Morita, T.~Sakasai and M.~Suzuki,
\emph{Structure of symplectic invariant Lie subalgebras of
symplectic derivation Lie algebras}, Adv. Math. \textbf{282}
(2015), 291--334. \url{https://doi.org/10.1016/j.aim.2015.06.017}.
Author version: \url{https://arxiv.org/abs/1404.3351v2}.

\bibitem{MSS2026}S.~Morita, T.~Sakasai and M.~Suzuki,
\emph{The first Galois obstruction in the Johnson cokernel},
arXiv:2608.17673v1, 18 August 2026.
\url{https://arxiv.org/abs/2608.17673}.

\bibitem{MZ}C.~P.~Mok and H.~Zheng, \emph{Vanishing of $K_{4m}(\mathbb Z)$ through degree $120$},
Hexagon 2611.00998, version~1, dated 23 November 2026 and posted 24 November 2026, 12 pages.
\url{https://hexagonmath.org/2611.00998}.
Theorem~1, p.~2; Section~2, pp.~2--3.

\bibitem{SharifiIwa}
R.~T.~Sharifi, \emph{Iwasawa Theory}, author lecture notes, 184 pp.,
retrieved 23 November 2026.
\url{https://www.math.ucla.edu/~sharifi/notes/iwasawa.pdf}.
Locators: Propositions 3.4.1 and 3.4.7, Corollary 3.4.8, pp.~90--92;
Theorem 5.5.31, p.~148; Proposition 6.1.1, pp.~151--152.

\bibitem{WWEOrd}P.~Wake and C.~Wang-Erickson,
\emph{Ordinary pseudorepresentations and modular forms},
Proc. Amer. Math. Soc. Ser.~B \textbf{4} (2017), 53--71.
\url{https://doi.org/10.1090/bproc/29}.
Final author preprint arXiv:1510.01661v4, 30 January 2017:
\url{https://arxiv.org/pdf/1510.01661v4}.
Locators in that version: \S2.1, pp.~2--3;
Proposition~2.2.6, p.~4; Theorem~4.2.8, pp.~10--11;
Corollaries~5.1.1--5.1.2, p.~11; Lemma~5.2.7, p.~13.

\bibitem{WWEPseudo}P.~Wake and C.~Wang-Erickson,
\emph{Pseudo-modularity and Iwasawa theory},
Amer. J. Math. \textbf{140} (2018), no.~4, 977--1040.
\url{https://doi.org/10.1353/ajm.2018.0022}.
Author version: arXiv:1505.05128v3,
7 September 2017, 54 pages:
\url{https://arxiv.org/pdf/1505.05128v3}.
Locators in that version: Corollary~2.1.3 and Conjecture~2.1.5, p.~8;
Corollaries~8.1.2--8.1.3, pp.~48--49;
Theorem~8.2.3, pp.~50--51.

\bibitem{Weibel}C.~A.~Weibel, \emph{The $K$-book: an introduction to algebraic $K$-theory},
Graduate Studies in Mathematics \textbf{145}, American Mathematical Society, 2013.
Chapter~VI, Theorem~8.2, Corollary~8.3 and Theorem~10.9, read in the author's online version of that chapter,
pp.~47--48 and 65: \url{https://sites.math.rutgers.edu/~weibel/Kbook/Kbook.VI.pdf}.

\bibitem{McCallum}W.~McCallum, Greenberg's conjecture and units in multiple $\Z_p$-extensions, Amer. J. Math. \textbf{123} (2001), 909--930. Author version, Theorem~3, p.~3. \url{https://archive.math.arizona.edu/wmc/Research/GreenbergConjecture.pdf}.
\bibitem{Wuthrich}C.~Wuthrich, Overview of some Iwasawa theory, Heidelberg lectures, 2012, \S1.3, pp.~3--4, and \S1.9. \url{https://www.maths.nottingham.ac.uk/plp/pmzcw/download/heidelberg.pdf}.
\bibitem{IshiiSoule}
S. Ishii, \emph{On Kummer characters arising from Galois actions on
the pro-\(p\) fundamental groups of once-punctured CM elliptic curves},
Research in Number Theory \textbf{11} (2025), article 38.
DOI: 10.1007/s40993-025-00622-0.
The version audited here is arXiv:2312.04175v2, 4 September 2024;
page locators refer to that version.


\bibitem{IshiiLie}S.~Ishii, On graded Lie algebras associated to once-punctured elliptic curves with complex multiplication, arXiv:2602.00615v2. Theorem~1.2, p.~3; Former Conjecture~1.8, p.~4; Remark~1.9, p.~5; Proposition~3.3 and Lemma~3.6, pp.~12--14. \url{https://arxiv.org/abs/2602.00615v2}.
\bibitem{BaumardSchneps}S.~Baumard and L.~Schneps, On the derivation representation of the fundamental Lie algebra of mixed elliptic motives, author version, pp.~2--3. \url{https://webusers.imj-prg.fr/~leila.schneps/BSPollack.pdf}.
\bibitem{Rubin1991}K.~Rubin, The ``main conjectures'' of Iwasawa theory for imaginary quadratic fields, Invent. Math. \textbf{103} (1991), 25--68.
\bibitem{GotoIshii2026}Y.~Goto and S.~Ishii, CM elliptic analogue of the Coleman--Ihara formula, arXiv:2610.21052v1 (2026). \url{https://arxiv.org/abs/2610.21052v1}.
\bibitem{SanoVenjakob2026}T.~Sano and O.~Venjakob, On extended interpolation properties of the Katz $p$-adic $L$-function, arXiv:2610.36607v1 (2026), Theorem~2.4.1 and Corollaries~2.4.3--2.4.7, pp.~8--10. \url{https://arxiv.org/abs/2610.36607v1}.
\bibitem{MilneANT}J.~S.~Milne, Algebraic Number Theory, course notes, ray class fields and the Minkowski bound. \url{https://www.jmilne.org/math/CourseNotes/ANT.pdf}.
\bibitem{BH}J.~Buhler and D.~Harvey, Irregular primes to 163 million, Math. Comp. \textbf{80} (2011), 2435--2444. \url{https://arxiv.org/abs/0912.2121}.
\bibitem{MS}W.~McCallum and R.~Sharifi, A cup product in the Galois cohomology of number fields, Duke Math. J. \textbf{120} (2003), 269--310. Author version dated 5 December 2003: Theorem~7.5, p.~20; Corollary~8.3, p.~23; Theorem~9.7, p.~28. \url{https://www.math.ucla.edu/~sharifi/publ/pairing.pdf}.
\bibitem{SharifiCyclotomicNotes}R.~Sharifi, Modular curves and cyclotomic fields, lecture notes, Theorem~2.3.18, p.~19. \url{https://www.math.ucla.edu/~sharifi/misc/aws-iwasawa.pdf}.
\bibitem{Pintar}L.~Pintar, \emph{Vanishing of $K_{4k}(\mathbb Z)$ through degree $104$ from Alp\"oge's Step~17}, Hexagon 2611.00829v3, internal Version~5, 23 November 2026. Unrefereed. \url{https://hexagonmath.org/2611.00829v3}.
\bibitem{Sutherland}A.~Sutherland, \emph{Verification files for the Lean proof of the Riemann Hypothesis}, public verification report, updated 24 November 2026. \url{https://math.mit.edu/~drew/grh/}.
\bibitem{MagnusKoszul}J.~Min\'a\v c, F.~W.~Pasini, C.~Quadrelli and N.~D.~T\^an, \emph{Mild pro-$p$ groups and the Koszulity conjectures}, Expositiones Math. \textbf{40} (2022), 432--455. Author version arXiv:2106.03675v2, \S3.1, equation~(3.3), p.~8. \url{https://arxiv.org/abs/2106.03675v2}.
\bibitem{LalondeRam}P.~Lalonde and A.~Ram, \emph{Standard Lyndon bases of Lie algebras and enveloping algebras}, Trans. Amer. Math. Soc. \textbf{347} (1995), 1821--1830, \S1, equation~(1.2).
\bibitem{VenjakobIwasawa}
O. Venjakob, \emph{On the Iwasawa theory of \(p\)-adic Lie extensions},
Compositio Math. \textbf{138} (2003), 1--54.
DOI: 10.1023/A:1025413030203.
The author version used for the numbered locators is
\url{https://www.mathi.uni-heidelberg.de/~otmar/papervenjakob/iwasawa.pdf}.

\bibitem{SharifiIwasawaNotes}
R. Sharifi, \emph{Iwasawa Theory}, lecture notes, Chapter 3,
Proposition 3.3.9 and Corollary 3.3.14,
\url{https://www.math.ucla.edu/~sharifi/notes/iwasawa-ch03.html},
accessed 25 November 2026.

\bibitem{Oukhaba2009}
H. Oukhaba, \emph{The index of elliptic units in
\(\mathbb Z_p\)-extensions, II},
Tohoku Math. J. \textbf{61} (2009), 253--265.
\url{https://www.jstage.jst.go.jp/article/tmj/61/2/61_253/_pdf}.

\bibitem{Oukhaba2003}
H. Oukhaba, \emph{Index formulas for ramified elliptic units},
Compositio Math. \textbf{137} (2003), 1--22.

\bibitem{Nakamura1994}
H. Nakamura, \emph{Galois rigidity of pure sphere braid groups
and profinite calculus}, J. Math. Sci. Univ. Tokyo
\textbf{1} (1994), 71--136.

\bibitem{Hain1997}
R. Hain, \emph{Infinitesimal presentations of the Torelli groups},
J. Amer. Math. Soc. \textbf{10} (1997), no.~3, 597--651.
DOI: 10.1090/S0894-0347-97-00235-X.

\url{https://nsalter.science.nd.edu/teaching/torelli/hain.pdf}.

\bibitem{EnomotoSatoh2014}
N. Enomoto and T. Satoh,
\emph{New series in the Johnson cokernels of the mapping
class groups of surfaces}, Algebr. Geom. Topol.
\textbf{14} (2014), no.~2, 627--669.
DOI: 10.2140/agt.2014.14.627.
\url{https://msp.org/agt/2014/14-2/agt-v14-n2-p01-s.pdf}.

\bibitem{Nakamura1996}
H.~Nakamura, \emph{Coupling of universal monodromy representations
of Galois--Teichm\"uller modular groups}, Math. Ann.
\textbf{304} (1996), 99--119.
\url{https://resolver.sub.uni-goettingen.de/purl?PPN235181684_0304}.

\bibitem{BandieraSchaetz}R.~Bandiera and F.~Sch\"atz,
\emph{Eulerian idempotent, pre-Lie logarithm and combinatorics of trees},
arXiv:1702.08907v1 (2017), Section~1.1, p.~2.
\url{https://arxiv.org/abs/1702.08907v1}.
\bibitem{NakamuraSchneps2000}H.~Nakamura and L.~Schneps, \emph{On a subgroup of the Grothendieck--Teichm\"uller group acting on the tower of profinite Teichm\"uller modular groups}, Invent. Math. \textbf{141} (2000), 503--560. Author manuscript, Section~2, equation~(2.1), p.~7. \url{https://webusers.imj-prg.fr/~leila.schneps/newGT.pdf}.
\bibitem{LuqueNovelliThibon}
J.-G.~Luque, J.-C.~Novelli and J.-Y.~Thibon,
\emph{Period polynomials and Ihara brackets},
J. Lie Theory \textbf{17} (2007), 229--239,
equations (29) and (32), p.~233.
\url{https://www.heldermann-verlag.de/jlt/jlt17/luquela2e.pdf}.
\bibitem{Moravec2011}
P.~Moravec, \emph{On the Schur multipliers of finite $p$-groups of given coclass},
Israel J. Math. \textbf{185} (2011), 189--205.
Author version, Proposition~2.2, p.~3, and the Hopf formula, p.~4.
\url{https://users.fmf.uni-lj.si/moravec/Papers/H2cc_revised3.pdf}.
\bibitem{Nakamura1995}
H.~Nakamura, \emph{On exterior Galois representations associated with open elliptic curves},
J. Math. Sci. Univ. Tokyo \textbf{2} (1995), 197--231.
\url{https://www.ms.u-tokyo.ac.jp/journal/pdf/jms020107.pdf}.
\bibitem{GoncharovEuler}
A.~B.~Goncharov, \emph{Euler complexes and geometry of modular varieties},
arXiv:math/0510310v5, 9 June 2007.
\url{https://arxiv.org/abs/math/0510310v5}.
Locators refer to this author version: Theorem~2.10(a), equation~(38),
p.~15; Lemma~2.12, p.~17; Theorems~1.2 and~5.16, pp.~9 and~34.
\end{thebibliography}
\end{document}
